\documentclass[10pt,a4paper]{article}
\usepackage[T1]{fontenc}\usepackage[utf8]{inputenc}
\usepackage{lmodern,amsmath,amssymb,amsthm,mathtools,geometry,hyperref,longtable,mathrsfs,booktabs,array,calc,newunicodechar}
\providecommand{\tightlist}{\setlength{\itemsep}{0pt}\setlength{\parskip}{0pt}}
\newunicodechar{ℝ}{\ensuremath{\mathbb R}}
\newcounter{none}
\geometry{margin=24mm}\allowdisplaybreaks[2]\setlength{\emergencystretch}{3em}
\hypersetup{hypertexnames=false,colorlinks=true,linkcolor=blue,urlcolor=blue}
\makeatletter
\renewcommand*\l@subsection{\@dottedtocline{2}{1.5em}{3.5em}}
\makeatother
\newcommand{\R}{\mathbb R}\newcommand{\T}{\mathbb T}\newcommand{\Z}{\mathbb Z}
\newcommand{\curl}{\operatorname{curl}}\newcommand{\Div}{\operatorname{div}}
\newcommand{\supp}{\operatorname{supp}}\newcommand{\tr}{\operatorname{tr}}
\newtheorem{theorem}{Theorem}[section]\newtheorem{proposition}[theorem]{Proposition}
\newtheorem{lemma}[theorem]{Lemma}\newtheorem{corollary}[theorem]{Corollary}
\newtheorem{definition}[theorem]{Definition}\newtheorem{remark}[theorem]{Remark}
\title{Finite-time Navier--Stokes construction:\\reconstruction and validation in progress}
\author{}\date{8 September 2026}
\newcommand{\D}{\mathcal D}
\newcommand{\AX}{\mathcal A_X}
\newcommand{\hh}{\mathcal H}
\newcommand{\res}{\mathcal R}
\newcommand{\diag}{\operatorname{diag}}
\newcommand{\dist}{\operatorname{dist}}
\newcommand{\J}{\mathcal J_K}
\newcommand{\E}{\mathcal E_K}
\newcommand{\N}{\mathcal N_K}
\begin{document}\maketitle
\begin{abstract}
This reader records independent derivations of the released classical three-dimensional construction, retaining its coordinates, signs, constants, domains and displayed remainder terms. The source is the 165-page paper \href{https://cdn.openai.com/pdf/32d9f210-8b73-45e0-91bc-82a30aef8a9a/navier-stokes.pdf}{FINITE TIME BLOWUP FOR NAVIER--STOKES}, SHA-256 \texttt{8c8a94ad9ac824c8b605b9827cadf7beaca48bd10b380de3cfc872a2c37afa81}. The reader remains under independent validation; imported dependencies and the Section 9 assembly are recorded in their ledgers. No prize claim is made.
\end{abstract}\tableofcontents
\section{Scope and provenance}
The target is $\partial_tu+(u\cdot\nabla)u-\nu\Delta u+\nabla p=f$, $\nabla\cdot u=0$ on $\mathbb R^3$ and, after the proved periodization map, $\mathbb T^3$. The source construction uses zero initial velocity, compact support, a smooth compactly supported force, finite energy, and velocity divergence at the terminal time. Every map is written on the full triple $(u,p,f)$.

The included bodies are the independently written axis/profile, base/heat, oscillation/curl, mean-correction, finite-stage correction-cycle, and terminal localization/comparison derivations. They are assembled with their exact definitions and dependencies. The corrected finite-stage calculation retains the full phase vector and the longitudinal part of the complete curl amplitude. A newer 166-page source revision has also been obtained. The component derivations remain explicitly pinned to the 165-page source hash above; later source audits identify their own capture. A completed revision audit records exact page/glyph correspondences and citation changes, with its finite comparison scope stated in the accompanying JSON.
\section{Validation ledger}
The static audit of the pinned formal repository found 580 local Navier--Stokes modules and matched 26 of 26 official statement declarations. That audit did not run Lean and does not certify a kernel proof. Actual import compilation is a separate ongoing check. The accompanying JSON contains the component proof source and check programs. The source-hash comparisons, exact algebra checks, numerical evaluations, and structural text checks are identified separately in their receipts; none alone certifies an analytic estimate or the full theorem. The terminal chapter was read against source pages 117--126.
\par
The continuation chapters contain the complete additional derivations named in
their headings. They retain the separately identified 165-page and 166-page
source captures. Their inclusion records those calculations; the independent
whole-construction and formal checks remain in progress.

% --- included proof body: profiles_body.tex ---
% Independent Section 4 / Appendix B derivation. No preamble.
% Requires standard amsmath/amssymb; no custom macros.
% Source body: core_exact.md
\section{Exact leading-field maps and radial primitives}

This is an independent derivation of the mathematical identities used on printed pp. 24-31 and 41-45 of \emph{Finite Time Blowup for Navier-Stokes}, source SHA-256 \texttt{8c8a94ad}\allowbreak\texttt{9ac824c8}\allowbreak\texttt{b605b982}\allowbreak\texttt{7cadf7be}\allowbreak\texttt{aca48bd1}\allowbreak\texttt{0b380de3}\allowbreak\texttt{cfc872a2}\allowbreak\texttt{c37afa81}. The symbols below retain their source values and orientations. Equations are mathematical identities; the accompanying arguments are independently written. The construction of the complete outer profile and the wave cone is recorded separately in the dependency ledger, not assumed to have been validated by this file.

The completed source construction chooses and fixes \(0<h<1/100\). That restriction is retained throughout its use here. The coordinate lemma itself has the larger validity range \(0<h<1/2\), which is stated below solely to record its precise domain; it does not enlarge the construction's chosen parameter range.

\subsection{1. Coordinates and their exact Jacobian}

Take right-handed cylindrical coordinates
\[
x=(r\cos\theta,r\sin\theta,z),\quad e_\theta=(-\sin\theta,\cos\theta,0),
\quad \theta\in\mathbb R/(2\pi\mathbb Z),\quad s=r^2/2.
\]
Put
\[
\begin{aligned}
&\tau=1-t>0,\quad A=\tfrac12+h,\quad D=\tfrac12-h,\quad 0<h<\tfrac12,\\
&z=q^D\eta,\quad \tau=q(1-\eta^2),\quad X=s/q,\\
&d=1-\eta^2,\quad L=1-2h\eta^2.
\end{aligned}
\]
For fixed \(\tau>0,z\in\mathbb R\), eliminating \(\eta\) gives
\(g(q)=q-z^2q^{2h}=\tau\). At \(q_*=|z|^{1/D}\), \(g(q_*)=0\); for \(z=0\) take \(q_*=0\). For \(q>q_*\), \(z^2q^{2h-1}<1\), and
\[
g'(q)=1-2hz^2q^{2h-1}\ge1-2h>0.
\]
Moreover \(g(q)\to\infty\). Thus exactly one \(q>q_*\) exists, \(|\eta|<1\), and the implicit function theorem makes \(q,\eta\) smooth for \(t<1\). The values \(\eta=\pm1\) used for profiles are one-sided parameter endpoints; they are not additional points with \(\tau>0\).

Differentiation at fixed physical coordinates gives
\[
q_t=-L^{-1},\quad \eta_t=\frac{D\eta}{qL},\quad X_t=\frac X{qL},
\qquad q_z=\frac{2\eta q^{1-D}}L,
\quad \eta_z=\frac d{q^DL},\quad X_z=-\frac{2\eta X}{q^DL}.
\]
The formula for \(\eta_z\) uses the exact equality \(L-2D\eta^2=d\). Consequently, for every real \(b\) and smooth profile \(f\),
\[
\partial_t(q^bf)=q^{b-1}L^{-1}(-bf+D\eta f_\eta+Xf_X),
\]
\[
\partial_z(q^bf)=q^{b-D}L^{-1}(2b\eta f+df_\eta-2\eta Xf_X).
\]
No constant or power of \(q\) is discarded in these maps.

\subsection{2. Cartesian axis extension and incompressibility}

Let \(C>1\), \(F=\phi/C\), \(E=\sqrt{2X}F\), and prescribe
\[
u_\theta=q^{-A}E,\quad u_z=q^{-A}U,\quad ru_r=V_0=Xv_0,
\quad p=q^{-2A}\Pi.
\]
For \(F,U,v_0,\Pi\) smooth on \([0,R]\times[-1,1]\), the Cartesian components are exactly
\[
u_1=\frac{v_0}{2q}x_1-q^{-A-1/2}Fx_2,
\quad u_2=\frac{v_0}{2q}x_2+q^{-A-1/2}Fx_1,
\quad u_3=q^{-A}U.
\]
Since \(X=(x_1^2+x_2^2)/(2q)\), these formulas prove smoothness across \(r=0\) for each \(t<1\). Smoothness is imposed on \(F\), not on \(E\) as a function of \(X\). This explicitly proves the map from scalar profile regularity to Cartesian velocity regularity.

Incompressibility is
\(\partial_s(ru_r)+\partial_zu_z=0\). Because \(A+D=1\), its exact scalar equation is
\[
(V_0)_X=L^{-1}(2A\eta U-dU_\eta+2\eta XU_X).
\]
Integrate with \(V_0(0,\eta)=0\). Writing \(M(X,\eta)=\int_0^XU(x,\eta)\,dx\), the result is
\[
V_0=\frac{2\eta XU-2D\eta M-dM_\eta}{L}.
\]
The extension of \(M/X\) at zero is smooth because
\(M/X=\int_0^1U(vX,\eta)\,dv\); differentiating this integral proves every mixed derivative assertion. Thus the formula also proves smoothness of \(V_0/X\).

The radial derivative of pressure and centrifugal term have the common factor \(q^{-2A-1/2}\), with respective coefficients \(\sqrt{2X}\Pi_X\) and \(E^2/\sqrt{2X}\). Their exact equality is
\[
\Pi_X=\frac{E^2}{2X}=F^2.
\]
The other radial terms are still present in the full Navier-Stokes residual. Radial time differentiation, both material-transport contributions and radial viscosity have scale \(q^{-3/2}=q^{2h}q^{-2A-1/2}\). Axial viscosity of the radial velocity has scale \(q^{-1/2-2D}=q^{-3/2+2h}=q^{4h}q^{-2A-1/2}\). Axial viscosity has scale \(q^{-A-2D}=q^{2h}q^{-A-1}\) in the tangential equations. These equalities identify every remaining type of term without setting any of them to zero.

\subsection{3. Exact transport and stress integration}

All residual identities here have viscosity one, as defined on source p. 6. Define
\[
H=\sqrt{2X}E=2XF,\quad l=X\partial_X\log H,
\quad W=1-\frac{2D\eta M+dM_\eta}{X},\quad H_c=D\eta+dU.
\]
Applying the preceding derivative maps and \(u_r=V_0/r\) gives, for arbitrary \(b,f\),
\[
(\partial_t+u_r\partial_r+u_z\partial_z)(q^bf)
=\frac{q^{b-1}}L\{W Xf_X+H_cf_\eta-b(1-2\eta U)f\}.
\]
Indeed the coefficient of \(Xf_X\) before substituting \(V_0\) is \(1+LV_0/X-2\eta U\), exactly \(W\). Angular momentum is \(ru_\theta=q^{-h}H\). Hence the angular material derivative equals \(-q^{-h-1}HS_q/L\), where
\[
S_q=-Wl-h(1-2\eta U)-H_c(\log E)_\eta.
\]
For the axial velocity, add \(\partial_zp\). Since \(-2A-D=-A-1\), the result is \(-q^{-A-1}S_n/L\), with
\[
S_n=-WXU_X-A(1-2\eta U)U-H_cU_\eta
-d\Pi_\eta+4A\eta\Pi+2\eta X\Pi_X.
\]

Let \(D_X=X\partial_X\). Solve
\[
D_XQ_s+(1+l)Q_s=S_q,\qquad D_XN_s+N_s=S_n.
\]
Their integrating factors are \(XH\) and \(X\). For \(\phi>0\) smooth at the axis, \(H=2X\phi/C\), and any nonzero integration constant produces an \(X^{-2}\) singularity in \(Q_s\) or an \(X^{-1}\) singularity in \(N_s\). Therefore the unique regular choices are
\[
Q_s=\frac{\int_0^XHS_q\,dx}{XH},\qquad
N_s=\frac1X\int_0^XS_n\,dx.
\]
In particular,
\[
Q_s=\frac1{F(X,\eta)}\int_0^1vF(vX,\eta)S_q(vX,\eta)\,dv,
\quad N_s=\int_0^1S_n(vX,\eta)\,dv,
\]
which prove the scalar smooth extensions and give \(Q_s(0)=S_q(0)/2\), \(N_s(0)=S_n(0)\).

Define
\[
\begin{aligned}
&p_s=\left(\frac{XQ_s}L,\frac{XN_s}{LE}\right),\quad a=1-2D_X\log E=2-2l,\\
&b_s=\frac{2D_XU}{E},\quad s_{\rm shear}=(a,-b_s),\quad T_0=F(p_s-s_{\rm shear}).
\end{aligned}
\]
Here \(s_{\rm shear}\) is only a typographical disambiguation from the unchanged coordinate \(s=r^2/2\); it is exactly the manuscript's vector \(s\). The stress components remain ordered \((r\theta,rz)\).

For any profile \(G\), at fixed \(t,z\),
\[
(\partial_r+k/r)(q^{-A-1/2}G)
=q^{-A-1}\left(\sqrt{2X}\,G_X+\frac k{\sqrt{2X}}G\right).
\]
Substituting \(G=FXQ_s/L=HQ_s/(2L)\), \(k=2\), produces \(q^{-A-1}ES_q/L\). Substituting \(G=FXN_s/(LE)=\sqrt{X/2}N_s/L\), \(k=1\), produces \(q^{-A-1}S_n/L\). The remaining stress terms are exactly
\[
\partial_ru_\theta-u_\theta/r=-q^{-A-1/2}Fa,
\qquad \partial_ru_z=q^{-A-1/2}Fb_s.
\]
The differential identities
\[
(\partial_r+2/r)(\partial_ru_\theta-u_\theta/r)
=\partial_{rr}u_\theta+r^{-1}\partial_ru_\theta-r^{-2}u_\theta,
\]
\[
(\partial_r+1/r)\partial_ru_z=\partial_{rr}u_z+r^{-1}\partial_ru_z
\]
now prove, including both signs,
\[
R^{(0)}_\theta=-(\partial_r+2/r)(q^{-A-1/2}T_{0,\theta}),
\quad R^{(0)}_z=-(\partial_r+1/r)(q^{-A-1/2}T_{0,z}),
\]
where \(R^{(0)}_j=\mathcal R(u,p)\cdot e_j+\partial_{zz}u_j\). Thus only axial viscosity has been removed from these exact identities.

Zero stress is equivalent to
\(Q_s=-2L\phi_X/\phi\) and \(N_s=-2LU_X\). Applying the two radial operators and canceling the two opposite quadratic logarithmic-derivative terms gives exactly
\[
-2L\frac{X\phi_{XX}+2\phi_X}{\phi}=S_q,
\qquad -2L(XU_{XX}+U_X)=S_n.
\]
Conversely, these equations and regularity imply the same \(Q_s,N_s\), by uniqueness of the regular radial primitives, hence zero stress. The first displayed fraction was checked directly against PDF p. 27.

\subsection{4. The five moments and their exact matching map}

Retain all five quantities
\[
\begin{aligned}
&M=\int_0^XU\,dx,\quad I=\int_0^XH\,dx,\quad J=\int_0^XUH\,dx,\\
&S=\int_0^X(U^2-E^2/2)\,dx,\quad C_p=\int_0^X\frac{E^2}{2x}\,dx,\quad \Pi=\Pi_{\rm ax}+C_p.
\end{aligned}
\]
Since \((XW)_X=1-2D\eta U-dU_\eta\), integration by parts in \(-WXH_X\) gives \(-XWH\) and the remaining integrand
\((1-h)H-D\eta H_\eta-d(UH)_\eta+2(h-D)\eta UH\). Hence
\[
Q_s=-W+\frac{(1-h)I-D\eta I_\eta-dJ_\eta+2(h-D)\eta J}{XH}.
\]
For the axial primitive, integrate \(-WXU_X\) in the same way. Its velocity terms become
\(D(U-\eta U_\eta)+4h\eta U^2-d(U^2)_\eta\). Pressure contributes
\[
X(4A\eta\Pi-d\Pi_\eta)-2h\eta\int_0^XE^2\,dx
+\frac d2\partial_\eta\int_0^XE^2\,dx,
\]
because \(\int_0^X\Pi\,dx=X\Pi-\frac12\int_0^XE^2\,dx\). Thus
\[
N_s=-WU+\frac{D(M-\eta M_\eta)+4h\eta S-dS_\eta}{X}
+4A\eta\Pi-d\Pi_\eta.
\]
The independent exact check differentiates both right-hand primitives, retains \(\Pi_X=H^2/(4X^2)\), and obtains precisely \(HS_q\) and \(S_n\).

If two pairs \((U_i,E_i)\) agree for \(X\ge X_h>0\) and their five moment values at \(X_h\) agree as functions of \(\eta\), all five moment differences are identically zero beyond \(X_h\): their radial derivatives are zero and their initial values are zero. All parameter derivatives also agree. With the same \(\Pi_{\rm ax}\), the explicit formulas then give identical \(\Pi,V_0,Q_s,N_s,p_s\). Agreement of profiles on an interval gives agreement of radial derivatives, hence of \(a,b_s,T_0\). This proves the complete morphism from matching data to the exterior field and stress; pressure matching alone would not suffice.

On a compact \(X\)-rectangle bounded away from zero and with \(E_i\ge e_{\min}>0\), the displayed formulas are compositions of addition, multiplication and reciprocal on denominators \(X,L,E,H\) bounded away from zero. For example
\(H_2^{-1}-H_1^{-1}=-(H_2-H_1)/(H_1H_2)\). The product rule through \(k\) parameter derivatives bounds the output difference in \((Q_s,N_s,p_s)\) by the input differences in \((U,E,M,I,J,S,C_p)\) through order \(k+1\). This uses no radial derivative of the difference. The one extra parameter derivative is necessary because \(W\) contains \(M_\eta\) and the primitive formulas contain moment derivatives.

Under the exact coordinate substitution \(X=\rho x\), use
\[
\widehat U(x,\eta)=U(\rho x,\eta),\quad\widehat E(x,\eta)=E(\rho x,\eta),
\quad H=\rho^{1/2}\widehat H,
\]
\[
(M,I,J,S,C_p)=(\rho\widehat M,\rho^{3/2}\widehat I,
\rho^{3/2}\widehat J,\rho\widehat S,\widehat C_p),
\quad\widehat Q_s=Q_s,\quad\widehat N_s=N_s,\quad\widehat p_s=\rho^{-1}p_s.
\]
Direct substitution proves that precisely the same rational formulas hold on the fixed \(x\)-rectangle. This is an invertible change of variables retaining all factors, not a replacement of the original moment data.

\subsection{5. Cone equivalence and finite-dimensional solve}

For \(a>0\), define
\[
t_s=-b_s/a,\quad v_s=a(1+t_s^2),\quad P_c=p_{s,1}+t_sp_{s,2},\quad
J_c=p_{s,2}-t_sp_{s,1}.
\]
The roots in \(v\) of \(2(P_c-v)^2-(v-2)J_c^2\) are
\[
v_\pm=P_c+J_c^2/4\pm |J_c|\sqrt{(P_c-2)/2+J_c^2/16}.
\]
For \(P_c>2\), evaluate the polynomial at \(2\) and \(P_c\): it is positive at \(2\), nonpositive at \(P_c\), and has its vertex at or above \(P_c\). Consequently \(2<v_-\le P_c\). This includes \(J_c=0\), where the two roots are \(P_c\). Therefore for \(v_s>2\), the admissible inequality \(v_s<v_-\) is exactly
\[
P_c>v_s,\qquad (v_s-2)J_c^2<2(P_c-v_s)^2.
\]
In particular \(v_s\le2\), \(P_c>2\) always implies the strict relaxed inequality \(v_s<v_-\).

For \(p_{s,2}=wp_{s,1}\), put \(c=1-b_sw/a\), \(j=w+b_s/a\), \(v=a+b_s^2/a\). Direct expansion gives
\[
(v-2)j^2-2c^2=(1+b_s^2/a^2)((a-2)w^2+2b_sw+b_s^2/a-2).
\]
On the compact set specified in Lemma 4.5, choose actual minima \(c_K>0\), \(\gamma_K=\min(2c^2-(v-2)j^2)>0\), and \(B_K\ge\max\{1,v,cv\}\). For
\[
p_{s,1}\ge P_K=\max\{(B_K+2)/c_K,8B_K/\gamma_K\},
\]
one has \(P_c>\max\{2,v\}\) and
\[
\frac{2(P_c-v)^2-(v-2)J_c^2}{p_{s,1}^2}
\ge\gamma_K-4B_K/p_{s,1}\ge\gamma_K/2.
\]
These are exactly the source constants, with the endpoint \(v=2\) covered by the preceding root calculation.

Stress coordinates satisfy, by direct scalar products,
\[
T_{0,\theta}+t_sT_{0,z}=F(P_c-v_s),\quad
T_{0,z}-t_sT_{0,\theta}=FJ_c.
\]
Thus a positive margin for the two stress inequalities proves nonvanishing and controls direction independently of magnitude.

For completeness, the moment matrix invertibility used later is proved as follows. Distinct real exponents \(\alpha_i\) form a Chebyshev family on \((0,\infty)\). For \(m\) exponents, divide a nonzero combination by its smallest power and differentiate. If the original combination had \(m\) distinct positive zeros, Rolle's theorem would give \(m-1\) zeros of a combination of \(m-1\) distinct powers, contradicting the induction starting with one nonzero power. Therefore the determinant \(\det[x_j^{\alpha_i}]\) never vanishes for ordered \(x_1<\cdots<x_m\), and its sign is constant on that connected set. Multilinearity expresses the bump moment determinant as the integral of this determinant against \(\prod_j\beta_j(x_j)\ge0\). Ordered disjoint nontrivial supports make the integral nonzero. Compact smooth parameter families then have bounded inverse matrices and bounded fixed derivatives by the adjugate formula.

For the exact quadratic moment equation \(Bc+Q(c,c)=d\), multiplication by \(B^{-1}\) has norm \(\nu_k\), and \(Q\) has bilinear norm \(q_k\) on \(C^k\). On the ball of radius \(r_k=2\nu_k\|d\|_{C^k}\), the map \(c\mapsto B^{-1}(d-Q(c,c))\) has image norm at most \(r_k/2+\nu_kq_kr_k^2\) and Lipschitz constant at most \(2\nu_kq_kr_k\). The source condition \(8\nu_k^2q_k\|d\|_{C^k}\le1\) gives invariance and Lipschitz constant at most \(1/2\). Iteration from zero in \(C^0\) and \(C^k\) gives the same fixed point. Its pointwise derivative matrix is \(B+Q(c,\cdot)+Q(\cdot,c)\); the corresponding \(B^{-1}\)-perturbation has norm at most \(1/2\), proving invertibility. The finite-dimensional implicit function theorem therefore gives all higher fixed parameter derivatives without imposing infinitely many smallness conditions.

\subsection{6. Direct heat-exterior verification}

The function named \(\mathcal H\) here is exactly the source's \(H(Z)\) in (4.29), distinguished typographically from angular momentum \(H(X,\eta)\):
\[
\mathcal H(Z)=\frac1{\Gamma(1+h)}\int_0^\infty e^{-v}v^h(1+Zv)^{-h}\,dv,
\qquad Z\ge0.
\]
For each \(k\), its \(k\)-th derivative is the integral with factor \((-1)^k(h)_kv^k(1+Zv)^{-h-k}\). The integrable majorant is \((h)_ke^{-v}v^{h+k}\), proving smooth one-sided derivatives at \(Z=0\) and positivity. Differentiating
\(e^{-v}v^{h+1}(1+Zv)^{-h-1}\) with respect to \(v\) shows
\[
Z^2\mathcal H''+((2+2h)Z+1)\mathcal H'+h(1+h)\mathcal H=0.
\]
In fact the numerator of the ODE integrand after taking the common denominator \((1+Zv)^{h+2}\) is \(h(h+1-v-Zv^2)\), exactly \(h\) times that derivative numerator; both boundary values vanish. This is a full integration-by-parts proof.

For \(K=c_\infty s^{-A}\mathcal H(2\tau/s)\), let \(Z=2\tau/s\). Then
\[
\partial_tK=-2c_\infty s^{-A-1}\mathcal H',\quad
\partial_sK=c_\infty s^{-A-1}(-A\mathcal H-Z\mathcal H'),
\]
\[
(\partial_{rr}+r^{-1}\partial_r-r^{-2})K
=c_\infty s^{-A-1}\{(2A^2-\tfrac12)\mathcal H
+(4A+2)Z\mathcal H'+2Z^2\mathcal H''\}.
\]
Using \(2A^2-1/2=2h(1+h)\) and the preceding ODE proves the exact radial heat equation. Also
\[
A\mathcal H+Z\mathcal H'
=\frac1{\Gamma(1+h)}\int_0^\infty e^{-v}v^h
\frac{A+Zv/2}{(1+Zv)^{h+1}}\,dv>0,
\]
so \(K_r=rK_s<0\) for \(r>0\). Finally
\(q^{-A}c_\infty X^{-A}\mathcal H(2d/X)=c_\infty s^{-A}\mathcal H(2\tau/s)\) exactly; the exterior has no residual hidden by a coordinate substitution.


% Source body: axis_exact.md
\section{Analytic axis construction and endpoint inequalities}

This independent proof calculation covers source pp. 144-150 (B.1-B.4 through Proposition B.3). It uses the exact pressure function \(\Pi_0\) prepared in Appendix A. The separate outer-profile audit must establish that preparation; the present file proves the axis operations on that specific datum and records the dependency instead of pretending to have proved Appendix A in this lane.

\subsection{1. Axis data and the complementary parameter regions}

Keep \(A=1/2+h\), \(D=1/2-h\), \(d=1-\eta^2\), \(L=1-2h\eta^2\), and the already fixed \(0<h\le10^{-2}\). The imported datum is holomorphic near \([-1,1]\), real and even on that interval, and has
\[
\Pi_0\le-\frac52P_*^2(1+\eta^2)^{-2},\qquad
\eta\Pi_0'(\eta)>0\quad(\eta\ne0).
\]
Choose \(0<j_0\le0.05\) and set exactly
\[
U_*=4\eta+j_0,\quad H_*=D\eta+dU_*,
\quad W_*=1-dU_*'-2D\eta U_*,
\]
\[
Z_*=-A(1-2\eta U_*)U_*-H_*U_*'-d\Pi_0'+4A\eta\Pi_0.
\]
For \(-1<\eta<1\),
\[
\frac{H_*}{d}=\frac{D\eta}{1-\eta^2}+4\eta+j_0,
\quad \partial_\eta(H_*/d)
=D\frac{1+\eta^2}{(1-\eta^2)^2}+4>0.
\]
Its endpoint limits are \(-\infty,+\infty\), so it has exactly one zero \(\eta_0\). Its value at zero is \(j_0>0\), hence \(\eta_0<0\). The root equation yields
\[
|\eta_0|=\frac{j_0}{4+D/(1-\eta_0^2)}.
\]
First this is at most \(j_0/4\le1/80\). Substitution back gives the useful explicit bounds \(j_0/5<|\eta_0|<j_0/4\). Further,
\[
-W_*=3-8h\eta^2+(1-2h)j_0\eta
\ge3-0.08-0.05=2.87>2.8.
\]
At the zero, \(H_*U_*'=0\) and \(-d\Pi_0'\ge0\). The term \(4A\eta_0\Pi_0\) is positive. Using \((1+\eta_0^2)^{-2}\ge1/4\), it is at least \((A/2)j_0P_*^2\). The remaining negative term has magnitude bounded by a fixed multiple of \(j_0\), since \(U_*=-D\eta_0/d\) at the zero. Therefore the already large fixed \(P_*\) gives \(Z_*(\eta_0)>0\), quantitatively at least a positive constant times \(j_0P_*^2\).

Continuity gives an open neighborhood \(V\) of \(\eta_0\) on which \(Z_*\) is bounded below by a positive number. Choose \(\delta_*>0\) smaller than half that lower bound. The compact set \(K_\delta=\{\eta:|Z_*(\eta)|\le\delta_*\}\) avoids \(V\), hence \(|H_*|\ge m>0\) there (if this set is empty, the following restriction is vacuous). Choose \(\sigma_*>0\) with \(\sigma_*^2<m^2/99\). Then
\[
\chi=\frac{H_*^2}{H_*^2+\sigma_*^2}>0.99\quad\hbox{on }K_\delta.
\]
Thus \(\chi\le0.99\) implies \(|Z_*|>\delta_*\), exactly the later endpoint alternative. Strict margins allow the same choices on a slightly larger compact interval \(I\supset[-1,1]\). There is a bounded simply connected complex neighborhood \(\Omega\) of \(I\) on whose closure \(L\) and \(H_*^2+\sigma_*^2\) are nonzero and \(\Pi_0\) is holomorphic: take a sufficiently thin neighborhood of the compact real interval, avoiding the closed zero sets.

Define
\[
\zeta_*=-\frac{LH_*}{H_*^2+\sigma_*^2},\qquad \xi_0=\Lambda\zeta_*,
\quad\phi_*(\eta)=\exp\left(\Lambda\int_0^\eta\zeta_*(w)\,dw\right).
\]
The primitive is single-valued on \(\Omega\), and \(\phi_*>0\) on \(I\). The exact identity needed below is \(H_*\zeta_*/L=-\chi\). After choosing \(\Lambda\), impose
\[
C\ge\sup_{\eta\in\Omega}|\phi_*(\eta)|.
\]
This supremum is finite after choosing \(\Omega\) with compact closure inside the holomorphic domain. Then \(g=\phi_*/C\) has a complex bound independent of both later large parameters. A real-interval bound alone would not establish the required parameter derivative bounds.

\subsection{2. The coefficient space, with all derivative losses accounted for}

Set \(Y=\Lambda X\). For coefficient sequences \(F=\sum_{\alpha\ge0}F_\alpha(\eta)Y^\alpha\), use precisely
\[
a_{\alpha\beta}=
\frac{20^{-\alpha}\rho^{-\beta}\beta!\binom{\alpha+\beta}{\beta}}
{(\alpha+1)^2(\beta+1)^2},
\qquad
\|F\|_\rho=\sup_{\alpha,\beta\ge0}\sup_{\eta\in I}
\frac{|\partial_\eta^\beta F_\alpha(\eta)|}{a_{\alpha\beta}}.
\]
Completeness follows coefficient by coefficient: a norm-Cauchy sequence converges uniformly for every derivative, and the fundamental theorem of calculus identifies each derivative limit with the derivative of its predecessor. The supremum bound passes to the limit, so this is a Banach space.

Let \(C_{\rm sq}=8\sum_{n\ge1}n^{-2}\). For \(0\le i\le N/2\), \((N+1)^2/(N-i+1)^2\le4\), while the other half has the analogous bound after \(i\mapsto N-i\). Therefore
\[
\sum_{i=0}^N\frac{(N+1)^2}{(i+1)^2(N-i+1)^2}\le C_{\rm sq}.
\]
In a coefficient product and its \(\beta\)-th derivative, the Leibniz factor \(\binom\beta k\) cancels the factorials in the two input weights. The remaining binomial quotient is at most one because
\[
\binom{i+k}{k}\binom{\alpha-i+\beta-k}{\beta-k}
\le\binom{\alpha+\beta}{\beta}.
\]
It counts subsets selecting specified numbers from the two fixed blocks. Applying the preceding convolution bound once in the radial and once in the parameter index proves
\(\|FG\|_\rho\le C_{\rm sq}^2\|F\|_\rho\|G\|_\rho\).

Define \(\mathcal I F(Y)=\int_0^YF(v)\,dv\), \(\mathcal A F=Y^{-1}\mathcal I F\), and let \(\mathcal J_\nu F\), \(\nu=1,2\), be the regular solution with zero value at \(Y=0\) of \(YG_{YY}+\nu G_Y=F\). Comparing coefficients gives
\[
(\mathcal J_\nu F)_{\alpha+1}=\frac{F_\alpha}{(\alpha+1)(\alpha+\nu)},
\quad (\mathcal J_\nu F)_0=0.
\]
The exact weight ratios are
\[
\frac{a_{\alpha\beta}}{a_{\alpha+1,\beta}}
=20\frac{(\alpha+2)^2}{(\alpha+1)^2}
\frac{\alpha+1}{\alpha+\beta+1}\le80,
\]
\[
\frac{a_{i,k+1}}{a_{i+1,k}}
=\frac{20}{\rho}(i+1)\frac{(i+2)^2}{(i+1)^2}
\frac{(k+1)^2}{(k+2)^2}\le\frac{80}{\rho}(i+1).
\]
These prove boundedness of multiplication by \(Y\), \(\mathcal I\), \(\mathcal J_\nu\), averaging (which divides coefficient \(\alpha\) by \(\alpha+1\)), and \(\partial_\eta\mathcal I\). The last assertion uses its integration divisor \(i+1\) to cancel the last factor in the second ratio.

For \(\mathcal J_\nu[(\partial_\eta F)(Y\partial_YG)]\), consider output degree \(N=\alpha+1\), split the input radial degrees as \(i\) and \(\alpha-i\), and allocate the added degree to the differentiated \(F\)-factor. Put \(I_1=i+1\), \(J_1=\alpha-i\). The weight comparison and radial inverse give, after dividing by the output weight, at most
\[
\frac{80}{\rho}\|F\|_\rho\|G\|_\rho
\sum_{i=0}^\alpha\sum_{k=0}^\beta
\frac{I_1J_1}{N(\alpha+\nu)}
\frac{(N+1)^2}{(I_1+1)^2(J_1+1)^2}
\frac{(\beta+1)^2}{(k+1)^2(\beta-k+1)^2}.
\]
The first fraction is at most one, including the zero endpoints; the two sums are bounded by \(C_{\rm sq}\). If the logarithmic derivative is absent, replace \(J_1\) by one, still giving a fraction at most one. If \(F\) is averaged, its divisor \(i+1\) cancels \(I_1\). If the parameter derivative is absent, use the first weight ratio instead. Inserting further undifferentiated factors adds the same convolution bounds with unchanged degree allocation. This proves every nonlinear operator estimate invoked below, including the case where both differentiated factors are copies of the same unknown.

A holomorphic coefficient \(b(\eta)\) bounded on a larger neighborhood belongs as degree zero to this space: Cauchy's estimate gives \(|b^{(\beta)}|\le M\beta!r^{-\beta}\), and choosing \(\rho<r\) bounds \((\beta+1)^2(\rho/r)^\beta\). This applies uniformly to \(g\), since its complex supremum is at most one.

\subsection{3. Exact rescaling and contraction}

Write the original fields as
\[
\phi=\phi_*\Phi,\quad U=U_*+\Lambda^{-1}u,
\quad \Pi=\Pi_0+\Lambda^{-1}p,
\quad p=\mathcal I(g^2\Phi^2),\quad \Phi(0)=1,\quad u(0)=0.
\]
Thus \(\Pi_X=g^2\Phi^2=\phi^2/C^2\) exactly. Set
\[
B=-2D\eta\mathcal A u-d\partial_\eta\mathcal A u,
\quad W=W_*+\Lambda^{-1}B,\quad H_c=H_*+\Lambda^{-1}du.
\]
Substitute these into the two stress-free equations in \texttt{core\_exact.md}. The unchanged equations are
\[
2(Y\Phi_{YY}+2\Phi_Y)=-\chi\Phi+\Lambda^{-1}R_1,
\qquad 2(Yu_{YY}+u_Y)=-Z_*/L+\Lambda^{-1}R_2,
\]
where
\[
LR_1=[W+h(1-2\eta U)+du\zeta_*]\Phi+WY\Phi_Y+H_c\Phi_\eta,
\]
\[
LR_2=[A(1-4\eta U_*)+dU_*']u-2A\eta\Lambda^{-1}u^2
+WYu_Y+H_*u_\eta+d\Lambda^{-1}uu_\eta
-4A\eta p+d\partial_\eta p-2\eta Yp_Y.
\]
Both identities have been checked as exact rational expressions by \texttt{exact\_checks.py}, before any estimate. In particular the pressure terms retain all three signs and coefficients. The angular \(-\chi\Phi\) comes from \(H_*\zeta_*/L=-\chi\); this term is order one and must be inverted, not treated as small.

The only products with both an unintegrated parameter derivative and a logarithmic radial derivative are \((\partial_\eta\mathcal A u)Y\Phi_Y\) and \((\partial_\eta\mathcal A u)Yu_Y\). They have the precise structure proved above. The term \(\partial_\eta p=\partial_\eta\mathcal I(g^2\Phi^2)\) is bounded, and \(Yp_Y=Yg^2\Phi^2\). Every other monomial is controlled by multiplication and a radial inverse, with at most one of the two derivative types. Subtracting two monomials by replacing their factors one at a time proves the corresponding local Lipschitz estimates. All constants are finite on a fixed ball, uniformly for \(\Lambda\ge1\) and \(C\) satisfying the complex bound.

Let \(M_\chi\) be the bounded multiplier norm of \(\chi\), and set \(T=\mathcal J_2\chi/2\), multiplication preceding integration. A sequence vanishing below degree \(b\) obeys
\[
\|\mathcal J_2F\|_\rho\le\frac{80}{(b+1)(b+2)}\|F\|_\rho.
\]
Multiplication by \(\chi\) preserves that vanishing order and \(\mathcal J_2\) increases it by one. Therefore
\[
\|T^k\|\le\frac{(40M_\chi)^k}{k!(k+1)!}.
\]
The absolutely convergent series \(\sum_{k\ge0}(-T)^k\) is a two-sided inverse of \(1+T\): multiply finite partial sums by \(1+T\) and use \(\|T^{k+1}\|\to0\). No smallness of \(\chi\) is required.

The center of the contraction is
\[
\Phi_0=f_0(Y\chi),\quad u_0=-\frac{YZ_*}{2L},
\qquad f_0(z)=\sum_{\alpha\ge0}\frac{(-z/2)^\alpha}{\alpha!(\alpha+1)!}.
\]
Its recurrence verifies \(2(zf_0''+2f_0')=-f_0\), \(f_0(0)=1\), so \(\Phi_0=(1+T)^{-1}1\). On the closed radius-one ball about \((\Phi_0,u_0)\), let \(M\) bound the pair
\(((1+T)^{-1}\mathcal J_2R_1/2,\mathcal J_1R_2/2)\), and let \(K\) be its Lipschitz constant. These are the finite constants supplied by the proved estimates. Choosing
\(\Lambda\ge\max\{1,2M,2K\}\) makes the map
\[
(\Phi,u)\longmapsto
\left(\Phi_0+\frac{(1+T)^{-1}\mathcal J_2R_1}{2\Lambda},
u_0+\frac{\mathcal J_1R_2}{2\Lambda}\right)
\]
invariant and a contraction with factor at most \(1/2\). Its unique fixed point has norm distance at most \(M/\Lambda\). This explicitly supplies the existential large-parameter threshold from finite operator constants.

To recover analytic functions, for \(R<20\) use
\[
\sum_{\alpha\ge0}\binom{\alpha+\beta}{\beta}(R/20)^\alpha
=(1-R/20)^{-\beta-1}.
\]
The coefficient norm therefore bounds parameter derivatives on \(|Y|\le R\) by a constant times \(\beta![\rho(1-R/20)]^{-\beta}\). It gives analytic continuation in a common positive parameter radius. Choose \(4.1<R_1<R_2<20\); Cauchy's formula on the larger radial disk and a slightly smaller parameter neighborhood gives, for every fixed \(r,s\),
\[
\sup_{[0,4.1]\times I}
\left|\partial_Y^r\partial_\eta^s(\Phi-f_0(Y\chi),u+YZ_*/(2L))\right|
\le C_{r,s}/\Lambda.
\]
The same parameter domain serves every fixed radial derivative because those derivatives use radial Cauchy estimates on the fixed larger disk. For original radial derivatives the factor is exactly \(\partial_X^r=\Lambda^r\partial_Y^r\), giving \(C_{r,s}\Lambda^{r-1}\). Real coefficients and uniqueness imply real profiles on the real rectangle.

\subsection{4. Positivity and endpoint alternative with the original constants}

For \(0\le z\le4.1\), put \(t=z/2\le2.05\). The magnitude ratio of consecutive terms of \(f_0\) is \(t/[(n+1)(n+2)]\), so the alternating tail after the cubic has decreasing magnitudes. Consequently
\[
f_0(z)\ge1-t/2+t^2/12-t^3/144
\ge305719/1152000>0.265.
\]
The cubic derivative is \(-1/2+t/6-t^2/48<0\): its discriminant is negative and its leading coefficient is negative. Its minimum is thus the displayed exact endpoint value. Grouping the alternating series after the constant term gives \(f_0\le1\). The proved \(O(\Lambda^{-1})\) bound gives \(\Phi\ge c_0>0\) when \(\Lambda\) is sufficiently large. In particular division by \(\phi\) in the original equation is justified. Uniform complex derivative bounds then give a smaller common complex domain with no zeros; derivatives of \(\log\Phi\) are bounded there.

Differentiate \(\chi=H_*^2/(H_*^2+\sigma_*^2)\):
\[
|H_*\chi'|=2|H_*'|\frac{H_*^2\sigma_*^2}{(H_*^2+\sigma_*^2)^2}
\le2\|H_*'\|_\infty\chi.
\]
Since \(f_0\) is positive and has bounded fixed derivatives on the stated compact range,
\[
|D_X\log\Phi|+|H_*\partial_\eta\log\Phi|
\le C\chi+C/\Lambda.
\]
Here \(D_X=Y\partial_Y\) and no division by \(\chi\) is used. Also
\[
-H_c\xi_0=\Lambda L\chi-du\zeta_*,\qquad
|du\zeta_*|\le C_{\sigma_*}\sqrt\chi,
\]
because \(|H_*|/(H_*^2+\sigma_*^2)\le\sigma_*^{-1}\sqrt\chi\). Thus the full source has the lower bound
\[
S_q\ge\Lambda L\chi-W_*-10h-C\chi-C_{\sigma_*}\sqrt\chi-C/\Lambda.
\]
For any fixed small \(\varepsilon>0\), Young's inequality in the exact form
\(K\sqrt\chi\le\varepsilon\Lambda L\chi+K^2/(4\varepsilon\Lambda L)\)
absorbs the square-root term. Increase \(\Lambda\) so the linear \(C\chi\) term uses the remaining part of a total \(0.05\Lambda L\chi\) allowance. Since \(-W_*>2.8\), \(10h\le0.1\), and the remaining constant error can be made at most \(0.2\), this gives
\[
S_q\ge2.5+0.95\Lambda L\chi.
\]

The integral formula
\[
Q_s(X,\eta)=\frac{\int_0^Xx\Phi(\Lambda x,\eta)S_q(x,\eta)\,dx}
{X^2\Phi(\Lambda X,\eta)}
\]
proves \(p_1/X=Q_s/L\ge c_1>0\). Stress freedom gives the exact identities
\[
p_1=\frac{XQ_s}L=a=-2Y\Phi_Y/\Phi,
\quad n_s=N_s/L=-2U_X=-2u_Y=Z_*/L+O(\Lambda^{-1}),
\quad p_2=Xn_s/E.
\]
All fixed parameter derivatives of \(\Phi,\log\Phi,u,p_1,n_s\) are bounded uniformly in subsequent \(C\), with \(\Lambda\) fixed; the radial scaling factor has already been accounted for.

At \(X_0=4/\Lambda\), first take \(\chi>0.99\), so \(t=2\chi\in(1.98,2]\). The series for \(f_0(z)+zf_0'(z)\) has terms \((-t)^n/(n!)^2\). Its alternating remainder beyond degree four is negative, giving
\[
f_0+zf_0'\le1-t+t^2/4-t^3/36+t^4/576.
\]
The derivative is \(-1+t/2-t^2/12+t^3/144\); on \([1.98,2]\) its first two terms are nonpositive and its last two have sum \(-t^2(12-t)/144<0\). Thus the polynomial decreases, and its value at \(1.98\) is exactly \(-75535511/400000000<-0.188<-0.18\). Since \(f_0\le1\),
\[
-2zf_0'/f_0=2-2(f_0+zf_0')/f_0>2.36.
\]
Taking \(\Lambda\) large gives \(p_1>2.3\).

On \(\chi\le0.99\), \(|Z_*|>\delta_*\). Since \(L\le1\) on \([-1,1]\), the preceding error estimate gives \(|n_s|\ge\delta_*/2\). At the same endpoint, exactly
\[
|p_2|=C\sqrt{2/\Lambda}\,\frac{|n_s|}{\phi_*\Phi}.
\]
For fixed \(\Lambda\), \(\phi_*\Phi\) has a finite positive upper bound independent of subsequent \(C\); \(p_1\) has a finite upper bound independent of \(C\). Thus choosing \(C\) large makes \(p_2^2/p_1>2.3\) uniformly on this complementary compact set. The two alternatives yield
\[
p_1+p_2^2/p_1>2+0.2
\quad\text{at }X_0=4/\Lambda,
\]
retaining the source's allowed \(c_{\rm ex}=0.2\). The proof uses distinct regions only after constructing their exact covering through \(\chi\) and \(Z_*\); no missing case at \(H_*=0\) remains.


% Source body: continuation_exact.md
\section{Continuation, exact moment repair, and assembly of the annular profile}

This file gives the independent continuation derivation for source pp. 150-157 and checks the assembly on pp. 36-45. Original source parameters remain \(0<h<1/100\), \(X_0=X_a=4/\Lambda\), \(X_b^{\rm ref}=100\), \(X_i=110\). The notation \(X_b^{\rm ref}\) distinguishes the temporary number 100 in Appendix B from the final outer annular edge \(X_b=e^3X_{\rm tail}\) in Section 4; neither numerical value or role is changed. Refer to \texttt{core\_exact.md} and \texttt{axis\_exact.md} for all field definitions and proved identities.

\subsection{1. The reference continuation and its estimates}

The smooth step is exactly
\[
\sigma(y)=\frac{e^{-1/y^2}}{e^{-1/y^2}+e^{-1/(1-y)^2}}\quad(0<y<1),
\qquad \sigma=0\ (y\le0),\quad\sigma=1\ (y\ge1).
\]
It is increasing: the ratio of its second denominator term to its first has logarithmic derivative \(-2/(1-y)^3-2/y^3<0\). It is smooth and flat at both endpoints by the corresponding exponential formulas. Choose fixed \(\bar t>0\) with \(4e^{2\bar t}<4.1\). Set \(y=\log(X/X_0)\). For \(0<t_1\le\bar t\), retain the analytic axis profile through \(y=t_1\), then prescribe on \((t_1,2t_1)\)
\[
\partial_y\log\phi_r=(1-\sigma((y-t_1)/t_1))D_X\log\phi_{\rm nat},
\quad \partial_yU_r=(1-\sigma((y-t_1)/t_1))D_XU_{\rm nat}.
\]
After \(2t_1\) keep \(\phi_r,U_r\) constant in \(y\). All endpoint jets match because the step is flat. Pressure and the integrated coefficients continue to be computed from the axis by the original integral definitions.

The axis approximation bounds \(U_{\rm nat}-U_*\) and its fixed parameter derivatives by \(O(\Lambda^{-1})\); its cutoff variation integrates over an interval of length \(t_1\) and is \(O(t_1/\Lambda)\). Thus \(U_r-U_*\) and its radial average remain small uniformly through \(X_i\): explicitly
\[
\|\mathcal A_XU_r-U_*\|_{C^k_\eta}
\le\sup_{0\le X\le X_i}\|U_r-U_*\|_{C^k_\eta}.
\]
No factor equal to the long logarithmic interval appears. Likewise integrate the bounded natural logarithmic slope to bound \(\log\phi_r-\log\phi_*\). Its leading parameter derivative remains \(\xi_0\), with the \(O(\chi)+O(\Lambda^{-1})\) weighted-gradient errors proved for the natural solution and an additional error tending to zero with \(t_1\) after \(\Lambda\) is fixed. In particular \(CE_r=\sqrt{2X}\phi_r\), its reciprocal on \([X_0,X_i]\), and all their fixed parameter derivatives have finite bounds independent of subsequent large \(C\) and sufficiently small \(t_1\).

Pressure satisfies the exact identity and estimate
\[
\Pi_r-\Pi_0=C^{-2}\int_0^X\phi_r(x,\eta)^2\,dx,
\qquad \sup_{X\le X_i}\|\Pi_r-\Pi_0\|_{C^k_\eta}\le K_kC^{-2}.
\]
The axis factor makes this integral finite at zero. The natural shear has \(a=p_1>0\); cutting its nonpositive \(D_X\log\phi\) to zero proves \(l_r\le1\). In the angular source the unchanged large term is \(-H_*\xi_0=\Lambda L\chi\). Replacing \(H_*\) by \(H_{c,r}\) costs \(K_{\sigma_*}\sqrt\chi+o_{t_1}(1)\); the other weighted gradients cost \(K\chi+K/\Lambda\). Young's inequality as in \texttt{axis\_exact.md}, now allocating a total \(0.06\Lambda L\chi\), and the strict \(-W_*>2.8\) margin give
\[
S_{q,r}\ge0.94\Lambda L\chi+2.4.
\]
Here one first enlarges \(\Lambda\), then \(C\), and then decreases \(t_1\); all residual constant errors have a strictly positive allowance. The axial source is
\[
S_{n,r}=Z_*+O(\Lambda^{-1})+o_{t_1}(1)+O(C^{-2})
\]
uniformly on the finite interval. Indeed all changes of \(U\) and of its parameter derivative are small, the logarithmic radial derivative is \(O(\Lambda^{-1})\) on the cutoff and zero later, and pressure is controlled by the preceding exact integral.

The defining primitive equations imply
\[
D_Xp_{1,r}=XS_{q,r}/L-l_rp_{1,r},\qquad
D_Xn_{s,r}+n_{s,r}=S_{n,r}/L.
\]
Positive regular initial data and \(l_r\le1\) imply \(p_{1,r}>0\). Multiplying the first inequality by the integrating factor \(e^y\), keeping the two positive source terms separately, gives
\[
p_{1,r}(y)\ge p_{1,r}(0)e^{-y}
+1.88\chi(e^y-e^{-y}),
\]
\[
p_{1,r}(X)\ge p_{1,r}(X_0)X_0/X
+1.2(X-X_0^2/X).
\]
For the first, \(X=X_0e^y=(4/\Lambda)e^y\) supplies the exact factor \(0.94\Lambda X_0/2=1.88\). For the second, use \(L\le1\), integrate \(2.4X\) against the same factor, and convert back to \(X\). On \(\chi>0.99\), the first lower bound remains above 2.3: replace its initial value by 2.3; its derivative is positive for \(y\ge0\), since \(1.88\chi>2.3/2\). On \(\chi\le0.99\), \(|Z_*|>\delta_*\), and the exact primitive for \(n_{s,r}\) retains its sign and a fixed positive lower bound when the displayed errors are small. Because \(p_{2,r}=Xn_{s,r}/E_r\), increasing \(C\) then makes \(p_{2,r}^2/p_{1,r}\) large uniformly. This proves
\[
v_r:=p_{1,r}+p_{2,r}^2/p_{1,r}>2+c
\]
on the entire reference interval, for a fixed \(c>0\). The second comparison yields \(p_{1,r}>3\) on \([100,110]\), since \(X_0\le4\). Upper bounds for \(p_{1,r},n_{s,r}\) and their fixed parameter derivatives follow directly from their integral representations and the finite coefficient bounds. They are independent of subsequent \(C\) and of the small reference widths.

\subsection{2. Activation of stress, with division by the flat factor proved}

Fix a sufficiently large finite \(C\). Choose \(t_*>0\) so the preceding bounds hold for every \(0<t_1\le t_*\). On this family, \(v_r\) has a finite common upper bound \(V_{\max}\); it may depend on \(C\), which has already been fixed. Choose \(0<\kappa_0<1/2\) later small enough relative to this bound. On \(0<y<t_1\) put
\[
e_a=(1-\kappa_0)\sigma(y/t_1),\quad\kappa=1-e_a,
\quad a=\kappa p_{1,r},\quad D_XU=-\kappa Xn_{s,r}/2,
\quad \partial_y\log\phi=-\kappa p_{1,r}/2.
\]
The natural reference is stress-free here. Therefore the exact differences obey
\[
\partial_y(\log\phi-\log\phi_r)=e_ap_{1,r}/2,
\quad \partial_y(U-U_r)=e_aXn_{s,r}/2.
\]
For a fixed parameter derivative, the integral from zero is bounded by a constant times \(ye_a(y)\), since \(e_a\) is increasing. The constants are uniform in smaller \(t_1,\kappa_0\) after \(\Lambda,C\) are fixed. The five moment differences are integrals of these value differences, with fixed smooth weights on the activation interval; their parameter derivatives have the same bound. The proved moment map in \texttt{core\_exact.md}, using one further parameter derivative, then gives
\[
p_s-p_{s,r}=O(ye_a),\qquad E_r/E=1+O(ye_a).
\]
The actual shear vector is exactly
\[
s_{\rm shear}=\kappa(p_{1,r},p_{2,r}E_r/E).
\]
Because \(\kappa\) cancels before forming its slope, there is no inverse power of \(\kappa_0\) in
\[
t_s=(p_{2,r}/p_{1,r})(E_r/E),\quad
v_s=\kappa v_r+O(ye_a),\quad P_c=v_r+O(ye_a),\quad J_c=O(ye_a).
\]
Consequently \(P_c-v_s=e_av_r+O(ye_a)\ge ce_a\) for small fixed \(t_1\), and \(P_c>2+c'\). The ratio
\((v_s-2)_+J_c^2/(P_c-v_s)^2\) is at most a fixed constant times \(y^2\), hence below 2 for small \(t_1\). The exact cone equivalence proves admissibility wherever \(v_s>2\) and the relaxed condition everywhere else. Since \(v_s(0)=v_r(0)>2+c\), continuity on the compact parameter interval supplies a first collar of positive width with \(v_s>2+c/2\).

The stronger endpoint assertion requires more than an \(O(ye_a)\) estimate. Here is the exact flat-integral calculation. For \(c>0\) and smooth \(b\), substitute \(u=y/\sqrt{1+y^2v}\) to obtain
\[
\int_0^y e^{-c/u^2}u^{-j}b(u,\eta)\,du
=\tfrac12e^{-c/y^2}y^{3-j}
\int_0^\infty e^{-cv}(1+y^2v)^{(j-3)/2}
b\left(\frac y{\sqrt{1+y^2v}},\eta\right)\,dv.
\]
For each fixed derivative the final integrand is bounded by \(e^{-cv}\) times a polynomial in \(v\), uniformly on the compact parameter set; denominators \(1+y^2v\) are at least one. Differentiation under the integral proves a smooth coefficient through \(y=0\), with limiting value \(b(0,\eta)/c\). This proves the needed result from Appendix A.9 directly.

Near zero, the specific step gives \(e_a=e^{-t_1^2/y^2}g_a(y)\), where
\[
g_a(y)=\frac{1-\kappa_0}{e^{-t_1^2/y^2}+e^{-1/(1-y/t_1)^2}}
\]
extends smoothly with \(g_a(0)=(1-\kappa_0)e>0\). Taking \(j=0\) in the flat-integral formula proves
\[
\log\phi-\log\phi_r,\ U-U_r\ \in e_ay^3C^\infty.
\]
This class is preserved by smooth compositions, products with smooth functions, and division by nonvanishing fields: for example \((e^{e_ay^3b}-1)/(e_ay^3)=b\int_0^1e^{ve_ay^3b}\,dv\) is smooth. Moment and pressure differences acquire a second integration and belong to \(e_ay^6C^\infty\). Substitution into the exact moment formulas therefore shows
\[
T_0=e_aB_0(y,\eta),\qquad
B_0\in C^\infty,\qquad
B_0(0,\eta)=F(X_0,\eta)p_{s,r}(X_0,\eta)\ne0.
\]
At the endpoint the latter vector is a positive multiple of the limiting shear. Its component along \((1,t_s)\) is positive, and that along \((-t_s,1)\) is zero. Both cone-direction margins consequently stay positive on a smaller collar. Every derivative of the exponential adds only finitely many powers of \(y^{-1}\), so for every fixed mixed derivative
\[
|T_0|\ge c e^{-t_1^2/y^2},\qquad
|\partial^\alpha T_0|\le C_\alpha e^{-t_1^2/y^2}y^{-N_\alpha}.
\]
Since \(X_0>0\), converting \(y\)-derivatives to \(X\)-derivatives changes only fixed smooth coefficients. Thus the vector direction is smooth through the zero-stress edge, not merely convergent along individual paths.

\subsection{3. Continuation to 110 and preservation of parameter analyticity}

After activation retain the same shear prescription with \(\kappa=\kappa_0\) through the reference cutoff and then through \(X=100\). The natural reference slopes are nonzero only in its initial interval of length \(2t_1\). Integrating the difference on that interval gives \(O(t_1)\). Beyond it, the reference slopes vanish, and the new slopes have size at most \(\kappa_0\) times the bounded functions \(p_{1,r}/2\), \(Xn_{s,r}/2\). Integration on the fixed interval gives \(O(\kappa_0)\). Thus field values, logarithms and the required parameter derivatives differ from reference by \(O(t_1+\kappa_0)\). Their integrated coefficients obey the same estimate with finite constants depending on the already fixed \(C,\Lambda\).

Choose \(\kappa_0\) small enough that \(\kappa_0V_{\max}<1/4\) and that its comparison errors use at most one quarter of the minimum \(v_r-2\). Then choose \(t_1\) small for the other error quarters and for the activation inequalities. The same exact slope formula yields \(P_c>2+c/2\), \(v_s<1\). The relaxed cone therefore holds regardless of \(J_c\). This order works for the entire family of reference widths, avoiding a circular choice of \(\kappa_0\) from a width that itself depends on \(\kappa_0\).

Starting at \(X=100\), multiply the axial prescription by a smooth \(\beta\) from one to zero on a short logarithmic interval. The integrated coefficients change by arbitrarily small errors when the interval is short, while
\[
P_c=p_{1,r}+\beta p_{2,r}^2/p_{1,r}+o(1)>2,
\quad v_s<1,
\]
using \(p_{1,r}>3\) there. Then hold \(D_XU=0\) and interpolate \(a\) convexly from \(\kappa_0p_{1,r}\) to 0.8, after arranging \(\kappa_0\sup p_{1,r}<0.8\). During this second interval \(v_s=a\le0.8\), \(P_c=p_1>2\) by continuity and its strictly positive initial margin. Both intervals can fit strictly between 100 and 110.

On the final interval, \(a=0.8\), \(l=0.6\), \(D_XU=0\). The unchanged leading gradient and the already small field changes give \(S_q>1\): its constant source contribution is \(0.6(-W_*)>1.68\), while the large \(\Lambda L\chi\) term absorbs the previously identified weighted errors. At a hypothetical first downward crossing \(p_1=2\),
\[
D_Xp_1=XS_q/L-0.6p_1>X-1.2>0,
\]
a contradiction. Hence \(p_1>2\) persists through \(X_i=110\).

On a small first collar, all reference coefficients arise from the analytic axis solution, its derivatives and forward integrals. The only cutoffs depend on \(y\), not on \(\eta\). Choose a compact complex neighborhood on which the natural \(\Phi\), \(L\), and the relevant denominators do not vanish. Integration in the real radial variable preserves holomorphy with uniform bounds on that neighborhood. Exponentiating the prescribed logarithm preserves nonvanishing. Radial differentiation changes the cutoff bounds but does not change the holomorphic parameter domain. A smaller fixed \(t_c\in(0,t_1)\), with \(4e^{t_c}<4.1\), therefore gives the common parameter domain for \(F,U,V_0/X,\Pi\) and every fixed radial derivative throughout \([0,X_0e^{t_c}]\). Subsequent changes have later support, so their forward primitives leave this rectangle exactly unchanged.

\subsection{4. Transition length fixed before amplitude}

Let \(\ell_i=\log(CE(X_i,\eta))\), \(G_i=U(X_i,\eta)\), and \(L_0=\log(X_i/X_0)\). Write \(M_k,N_k\) for the reference bounds on \(p_{1,r},n_{s,r}\), which are independent of subsequent \(C\). Before the final interpolation the logarithmic \(\phi\) slope is \(-\kappa p_{1,r}/2\); during it the slope is a convex interpolation with \(-0.4\); after it the slope is exactly \(-0.4\). Hence
\[
\|\log\phi(X_i)\|_{C^k_\eta}
\le\|\log\phi_*+\log\Phi(4,\cdot)\|_{C^k_\eta}
+\tfrac12(M_k+0.8)L_0.
\]
Adding \(\frac12\log(2X_i)\) proves \(\|\ell_i\|_{C^k_\eta}\le B_k\), where \(B_k\) is independent of subsequent \(C\) and smaller transition widths. The axial change on activation is at most \(X_0e^{\bar t}N_kt_1/2\); the remaining small-shear part contributes at most \(X_iN_k\kappa_0/2\). The natural-axis difference is \(O(\Lambda^{-1})\) with a constant independent of \(\Lambda\). Therefore
\[
\|G_i-U_*\|_{C^k_\eta}
\le C_k^{\rm nat}/\Lambda+C_k^{\rm join}(\Lambda)(t_1+\kappa_0+\omega_{\rm fin}),
\]
where \(\omega_{\rm fin}\) is the sum of the two final logarithmic widths. No inverse cutoff width occurs in these integrals.

With \(f=(1+\eta^2)^{-1}\), choose a finite positive \(T_{\rm sh}\) satisfying exactly the source lower bound
\[
T_{\rm sh}\ge20\|\sigma'\|_\infty(B_0+\|\log f\|_\infty).
\]
For \(y_i=\log(X/X_i)\), prescribe
\[
\log E=-\log C+y_i/10+(1-\sigma(y_i/T_{\rm sh}))\ell_i
+\sigma(y_i/T_{\rm sh})\log f,\qquad U=G_i.
\]
At the left endpoint this matches \(E(X_i)\) and its radial jet, because both original and new profiles have \(D_X\log E=0.1\) there and the step is flat. At the right endpoint it matches \(C^{-1}f(X/X_i)^{1/10}\) with its entire radial jet. Direct differentiation gives
\[
l=0.6+T_{\rm sh}^{-1}\sigma'(y_i/T_{\rm sh})(\log f-\ell_i)
\in[0.55,0.65],\quad a\in[0.7,0.9],\quad b_s=0.
\]
The old parameter gradient satisfies
\(-H_c\ell_i'\ge\Lambda L\chi-K\chi-K_{\sigma_*}\sqrt\chi-K/\Lambda-o(1)\).
The new one is
\[
-H_c(\log f)'=\frac{2\eta H_c}{1+\eta^2}
\ge-Kj_0^2-K/\Lambda-o(1).
\]
For the latter bound, the exact identity
\(\eta H_*=(D+4d)\eta^2+dj_0\eta\)
and completion of the square bound its negative part by \(Kj_0^2\); the \(U-U_*\) error is already controlled. These gradients are combined with nonnegative coefficients summing to one. The old large positive term absorbs its correspondingly weighted errors, and the new negative part is small. Since \(-W\) is close to \(-W_*>2.8\) and \(l\ge0.55\), the constant contribution exceeds 1.54 before the small \(h,j_0,\Lambda^{-1}\) errors. Thus \(S_q>1\). The same crossing argument at \(p_1=2\), now \(l\le0.65\), gives \(D_Xp_1>X-1.3>0\). Since \(b_s=0\), this proves \(P_c=p_1>2\), \(v_s=a<1\), hence the relaxed cone, along the full transition and the following ideal-angular interval.

\subsection{5. The five matching equations, including all scaling factors}

Set
\[
X_R=110(CP_*)^{10},\quad x=X/X_R,\quad
X_{\rm sep}=110e^{T_{\rm sh}},\quad
x_{\rm sep}=\frac{e^{T_{\rm sh}}}{(CP_*)^{10}}.
\]
Here \(X_{\rm sep}\) is fixed before \(C\). The ideal pair is \(U_0=4\eta\), \(E_0=P_*fx^{1/10}\). The exact normalized moments are
\[
\widehat M_0=4\eta x,\quad\widehat I_0=\frac{5\sqrt2}{8}P_*fx^{8/5},
\quad\widehat J_0=4\eta\widehat I_0,
\]
\[
\widehat S_0=16\eta^2x-\frac5{12}P_*^2f^2x^{6/5},
\quad C_{p,0}=\frac52P_*^2f^2x^{1/5}.
\]
These follow by direct integration, retaining the factors
\((M,I,J,S,C_p)=(X_R\widehat M,X_R^{3/2}\widehat I,
X_R^{3/2}\widehat J,X_R\widehat S,C_p)\).

On \([0,X_{\rm sep}]\), fixed parameter derivatives of \(U\) and \(CE\) are bounded independently of later \(C\); near zero \(CE=O(\sqrt X)\). Thus
\[
\|\widehat M\|_{C^k}+\|\widehat S\|_{C^k}\le K_kx_{\rm sep},
\quad\|\widehat I\|_{C^k}+\|\widehat J\|_{C^k}
\le K_kC^{-1}x_{\rm sep}^{3/2},
\quad\|C_p\|_{C^k}\le K_kC^{-2}
\]
at \(x=x_{\rm sep}\). For pressure the proof uses the physical integral \(C^{-2}\int_0^{X_{\rm sep}}(CE)^2/(2X)\,dX\), finite because \(CE=O(\sqrt X)\). All constants are fixed before increasing \(C\). The ideal moments also tend to zero with the displayed exact powers. Their pressure increment is \(O(C^{-2})\), because \(x_{\rm sep}^{1/5}=e^{T_{\rm sh}/5}(CP_*)^{-2}\).

Beyond \(X_{\rm sep}\), the angular profile is exactly ideal, since
\(C^{-1}f(X/X_i)^{1/10}=P_*fx^{1/10}\). Choose \(C\) so \(x_{\rm sep}<e^{-8}\). Restore \(G_i\) to \(4\eta\) by a fixed smooth step on \(-8<\log x<-7\), retaining \(E_0\). The size of this change, and its fixed logarithmic radial and parameter derivatives, is bounded by fixed constants times \(\|G_i-4\eta\|\). The moment errors accumulated to the correction interval are bounded by
\[
K_k\|G_i-4\eta\|_{C^{k+1}_\eta}+o_C(1).
\]
Every relevant radial weight is integrable at zero: their actual exponents are \(0,3/5,1/5\) and \(-4/5\) for the pressure integrand, all greater than \(-1\). No divergent contribution is discarded.

On \(-6<\log x<-5\), choose two fixed ordered nonnegative bumps for \(U\), three for \(E\). At the ideal pair make the invertible row operations \(J\mapsto J-4\eta I\), \(S\mapsto S-8\eta M\), with \(4\eta\) held at its base value. The linearized axial block has weights \(1,fx^{3/5}\); the azimuthal block has weights \(x^{1/2},fx^{1/10},fx^{-9/10}\), retaining fixed nonzero factors \(P_*,\sqrt2\) in the matrix. The exponents in each block are distinct. The Chebyshev determinant proof and quadratic contraction in \texttt{core\_exact.md} therefore solve all five moment equations exactly, with small smooth coefficients. The row operations do not impose that the corrected \(U\) remain constant; the full nonlinear differences still include \(\delta U\delta E\), \((\delta U)^2\), and \((\delta E)^2\).

To verify one tolerance can precede \(C\), put \(R=[e^{-8},e^{-5}]\), \(Q_{\min}=\min Q_{s,0}>0\), \(e_* =\min E_0>0\), and \(w_*=1+\max|N_{s,0}/(E_0Q_{s,0})|\) on \(R\times[-1,1]\). The ideal source is positive directly: for this pair
\[
S_{q,0}=\frac95-h+\frac{16}{5}h\eta^2
+\frac{2\eta^2(D+4d)}{1+\eta^2}>0,
\]
as \(W=-3+8h\eta^2\), \(l=3/5\), \(H_c=(D+4d)\eta\). Its regular weighted primitive gives \(Q_{s,0}>0\). The compact minima thus exist without importing the positivity of a separate abbreviated formula. The exact moment map has no \(X_R\) in \(Q_s,N_s\), so a fixed small tolerance in the fields and normalized moments through one extra parameter derivative ensures
\[
Q_s\ge Q_{\min}/2,\quad E\ge e_*/2,\quad |a-0.8|\le0.1,
\quad|N_s/(EQ_s)|\le w_*,\quad |b_s|\le0.1/(1+w_*).
\]
Then, retaining the source numbers,
\[
v_s\le0.9+0.01/0.7<1,
\quad G:=Q_s-b_sN_s/(aE)\ge(1-1/7)Q_s\ge3Q_{\min}/7,
\quad P_c=X_RxG/L.
\]
Choosing \(X_R\) large now gives \(P_c>2\) without shrinking the tolerance. First choose \(j_0\), then \(\Lambda\), so \(j_0+C_k^{\rm nat}/\Lambda\) fits the tolerance. Next fix the \(B_k\) and \(T_{\rm sh}\), then enlarge \(C\) for the vanishing moment contributions. Finally decrease \(\kappa_0,t_1,\omega_{\rm fin}\) for the remaining \(G_i\) error. These choices prove the exact order
\[
(M_d,T_d,P_*,\lambda,h)\to\text{moment tolerance},j_0
\to\delta_*,\sigma_*,\Lambda\to(B_k),T_{\rm sh}
\to C,X_R\to\kappa_0,t_1,\text{final widths}.
\]
Only finitely many parameter derivative orders need to be small. Smooth implicit differentiation supplies all remaining fixed orders after the choices. At \(X_h=X_Re^{-5}\), all five moments and both fields agree with the ideal pair, with the same axis pressure \(\Pi_0\), so the exact exterior-preservation map from \texttt{core\_exact.md} applies.

\subsection{6. The finite-frequency assembly in Section 4}

The outer construction in Appendix A supplies the separate profile, its exact total moments, pressure datum and reserved intervals. The loop construction in Appendix C supplies the periodic shear family. These two existence constructions are dependencies allocated to the other audit lanes. The operations performed with their outputs in Section 4 can be checked exactly here.

The join at \(X_h\) is smooth because both pairs coincide with the same ideal pair on a neighborhood, not only at one radius. Their five matching moments and common axis pressure imply exact agreement of every integrated exterior field. In particular all total moments and the outer heat field persist. The perturbation interval \(I=[X_-,X_+]\) lies after the reserved analytic collar and before the first repair interval; its endpoints already have the admissible cone.

For the loop input \(t_\ell,\rho_\ell\) with mean \(t_s\), variance \(\rho_\ell/a\), and \(v_\ell=v_s+\rho_\ell>2\), Section 4 changes the angular parameter by
\[
\frac{d\varphi}{d\theta'}=\frac{a(1+t_\ell^2)}{2\pi v_\ell},
\quad(a_L,-b_L)=\frac{v_\ell(1,t_\ell)}{1+t_\ell^2}.
\]
The derivative is positive and its integral over \([0,2\pi]\) is
\(a(1+t_s^2+\rho_\ell/a)/v_\ell=1\). Thus it defines an orientation-preserving smooth circle coordinate of period one. Its shear mean is
\[
\int_0^1(a_L,-b_L)\,d\varphi
=\frac a{2\pi}\int_0^{2\pi}(1,t_\ell)\,d\theta'
=(a,-b_s).
\]
The compact positive derivative gives a smooth inverse including all endpoint parameter derivatives. This establishes exactly the reparameterization asserted in Lemma 4.11; the existence of \(t_\ell,\rho_\ell\) itself remains the separately audited Appendix C construction.

Take zero-mean periodic primitives \(\mathscr A,\mathscr B\) with
\[
\partial_\varphi\mathscr A=-(a_L-a_0)/2,
\quad\partial_\varphi\mathscr B=E_0(b_L-b_0)/2.
\]
They exist because the two means vanish, and they vanish near both endpoints where the loop is constant. Set
\[
E_N=E_0\exp(\mathscr A(X,\eta,N\log X)/N),
\quad U_N=U_0+\mathscr B(X,\eta,N\log X)/N.
\]
Holding phase fixed in the following slow derivatives, direct differentiation gives exactly
\[
a_N=a_L-2D_X\mathscr A/N,
\quad b_N=e^{-\mathscr A/N}
\left(b_L+2D_X\mathscr B/(NE_0)\right).
\]
Since the phase has no \(\eta\)-dependence and all intervals are fixed, field differences and fixed parameter derivatives are \(O(N^{-1})\), while the actual shear differs from the desired loop by \(O(N^{-1})\). The five exact integrands of the changes, writing \(u_N=U_N-U_0\), \(e_N=E_N-E_0\), are
\[
u_N,\quad\sqrt{2X}e_N,\quad H_0u_N+U_0\sqrt{2X}e_N+\sqrt{2X}u_Ne_N,
\]
\[
2U_0u_N+u_N^2-E_0e_N-e_N^2/2,
\quad E_0e_N/X+e_N^2/(2X).
\]
Integration on the fixed interval preserves \(O(N^{-1})\) bounds. The moment map, using one extra parameter derivative, gives \(p_{s,N}-p_{s,0}=O(N^{-1})\) without assuming small radial derivatives of the field changes.

The four defining gaps are the components of the smooth map
\[
\Psi(a,b,p)=\left(a,v-2,c-v,2(c-v)^2-(v-2)j^2\right),
\quad v=a+b^2/a,\ c=p_1-bp_2/a,\ j=p_2+bp_1/a.
\]
The loop and unchanged admissible part have positive compact minima \(\mu_L,\mu_R\). Fix a compact neighborhood with \(a>0\), and let \(C_\Psi\) be the supremum norm of its derivative. Choosing \(N\) so the preceding errors are below \(\min(\mu_L,\mu_R)/(2C_\Psi)\) leaves each relevant gap at least half its original minimum. This is a finite uniform choice after every radial scale and loop has been fixed.

On a later repair interval \((Y_0,Y_1)\subset I_1\), the unchanged pair is \(U_0=0\), \(E_0=K(\eta)X^{-1/2-\lambda}\), \(K=c_{\rm patch}(1+\eta^2)^{-1}>0\). Two \(U\) bumps and three \(E\) bumps have the exact quadratic moment changes, in order \((M,J,I,S,C_p)\),
\[
\int u_c,\quad\int(H_0u_c+\sqrt{2X}u_ce_c),\quad\int\sqrt{2X}e_c,
\]
\[
\int(u_c^2-E_0e_c-e_c^2/2),\quad\int(E_0e_c/X+e_c^2/(2X)).
\]
The linear blocks have weights \((1,X^{-\lambda})\) and \((X^{1/2},X^{-1/2-\lambda},X^{-3/2-\lambda})\) with their fixed positive \(K,\sqrt2\) factors. Distinctness follows from the actual \(\lambda>0\). The determinant and quadratic-solve proofs in \texttt{core\_exact.md} apply. With \(\|d_N\|_{C^k}\le D_k/N\), the source's explicit requirement
\[
N\ge8\max_{k=0,2}\nu_k^2q_kD_k
\]
gives a smooth correction with \(\|c\|_{C^2}\le2\nu_2D_2/N\). Fixed bump shapes also make its first radial derivative \(O(N^{-1})\). Increasing \(N\) so \(N\ge4C_\Psi C_{\rm rep}/\mu_R\) leaves the repair gaps at least \(\mu_R/4\). At \(Y_1\) every moment is restored exactly. Below the first perturbation and above \(Y_1\), the full field and stress are therefore exactly the original joined ones. No later modification reaches the analytic collar or the reserved intervals \(I_3,I_4\).

\subsection{7. Endpoint, weight and total-moment conclusions}

The exact inner factor already proved gives \(T_0=e^{-t_1^2/y_a^2}g_aB_a\), where \(y_a=\log(X/X_a)\), \(g_a(0)>0\), and \(B_a(0)=Fs_{\rm shear}\ne0\). The imported outer endpoint computation has the exact factors
\[
T_{0,\theta}=e^{-4/y_b^2}y_b^{-3}b_\theta,
\quad T_{0,z}=e^{-4/y_b^2}y_b^3b_z,
\quad y_b=\log(X_b/X),\quad b_\theta(0)>0.
\]
Consequently the unit direction is \(B_a/|B_a|\) near the inner edge and
\((b_\theta,y_b^6b_z)/(b_\theta^2+y_b^{12}b_z^2)^{1/2}\) near the outer edge. Both extend smoothly. The endpoint directions are respectively \((1,t_s)/(1+t_s^2)^{1/2}\) and \((1,0)\), with \(t_s=0\) at the latter edge.

Let \(n=T_0/|T_0|\), \(A_n=n_\theta+t_sn_z\), \(B_n=n_z-t_sn_\theta\). In the interior,
\[
A_n=(P_c-v_s)/|p_s-s_{\rm shear}|>0,
\quad G_n:=2-(v_s-2)B_n^2/A_n^2>0.
\]
At either edge \(B_n=0\), while \(A_n>0\), so \(G_n=2\). Compactness of the closed annulus gives positive minima for \(A_n,G_n,F,a,v_s-2\). The choice
\(\kappa=\min\{1,\min A_n,\min G_n\}>0\)
proves exactly the two directional bounds in source (4.26), with \(\kappa<2\).

The exact weight
\[
\zeta(X)=\exp(-t_1^2/y_a^2-4/y_b^2)\quad(X_a<X<X_b),\qquad \zeta=0\text{ otherwise}
\]
is smooth and flat at both edges. On either fixed collar its ratio to the local exponential has positive upper and lower bounds. The nonzero vector coefficients and finite powers introduced by differentiation therefore prove
\[
|T_0|\ge c\zeta,\qquad
|\partial^\alpha T_0|\le C_\alpha\zeta\,
\min\{1,y_a,y_b\}^{-m_\alpha}.
\]
On the compact region between the collars this follows from positivity and smoothness directly. Thus all six conclusions concerning stress direction and weight follow from the actual endpoint factors, not from a claim that every flat function has a smooth normalized direction.

Every completed join restores all five cumulative integrals at a finite radius. Therefore total \(M,J,S\) remain zero, \(C_p(\infty)=-\Pi_0\), and the renormalized angular moment remains unchanged because the perturbation of \(H\) is compactly supported and has total integral zero. The pressure is hence exactly
\[
\Pi(X)=\Pi_0+C_p(X)=C_p(X)-C_p(\infty)
=-\int_X^\infty E(x)^2/(2x)\,dx.
\]
Above the retained \(X_v\), the outer data have \(U=M=J=0\); the formula
\(V_0=(2\eta XU-2D\eta M-dM_\eta)/L\)
then gives \(V_0=0\) exactly. The heat formula was proved directly in \texttt{core\_exact.md}. The compensation used only \(I_2\), the last repair only \(I_1\); therefore \(I_{\rm pos}=I_3\) and \(I_{\rm mean}=I_4\) retain their exact specified profiles. These are dependency-preserving assembly statements; they do not by themselves validate any later residual summation or wave realization.


% --- included proof body: derivation.tex ---


\noindent This is a bounded mathematical audit of Section 5 and Appendix A of
\emph{Finite Time Blowup for Navier--Stokes}, cited as [NS]. The source is the
165-page \href{https://cdn.openai.com/pdf/32d9f210-8b73-45e0-91bc-82a30aef8a9a/navier-stokes.pdf}{source PDF}.
Its SHA-256, written on two lines without an intervening character, is
\begin{center}\ttfamily
8c8a94ad9ac824c8b605b9827cadf7be\\
aca48bd10b380de3cfc872a2c37afa81
\end{center}
Page numbers below are both printed and PDF page numbers. Every proof in pp.~45--62
and pp.~126--144 was read. Definitions and proofs in (4.1)--(4.23) were read as
dependencies. The existence and gluing of the leading axis profile in Appendix B,
the cone realization in Appendix C, and the wave and final localization arguments
are outside this lane. Their conclusions are identified at the exact points where
they enter; this document does not assert a global Navier--Stokes theorem.

The derivations below are independent calculations, not a reproduction of the
source prose. Symbols, physical coefficients, signs, and original parameter choices
are retained. We write $\hh$ for the heat factor called $H$ in (A.32) solely to
distinguish it from the angular-momentum profile $H=\sqrt{2X}E$; the exact map
between them is $E_{\rm heat}=c_\infty X^{-A}\hh(2(1-\eta^2)/X)$.

\section{The physical coordinate map and the exact coefficient equations}
Set
\[
\begin{gathered}
 A=\tfrac12+h,\qquad D=\tfrac12-h,\qquad 0<h<\tfrac12,\qquad \tau=1-t,\\
 z=q^D\eta,\qquad \tau=q(1-\eta^2),\qquad s=\tfrac12r^2,\qquad X=s/q,
\end{gathered}
\]
and $d=1-\eta^2$, $L=1-2h\eta^2$. For $\tau>0$, the equation for $q$ is
$q-z^2q^{2h}=\tau$ on $q>|z|^{1/D}$. Its derivative is
$1-2hz^2q^{2h-1}=L\ge1-2h>0$, it starts at zero and tends to infinity.
It consequently defines exactly one $q$, with $|\eta|<1$.
Direct differentiation, with the other physical variables fixed, gives
\[
\begin{array}{lll}
q_t=-L^{-1},&\eta_t=D\eta/(qL),&X_t=X/(qL),\\
q_z=2\eta q^{1-D}/L,&\eta_z=d/(q^DL),&X_z=-2\eta X/(q^DL).
\end{array}
\]
Thus, for every real $b$ and every differentiable profile $f$,
\begin{align*}
 \partial_t(q^bf)&=q^{b-1}T_bf,&
 T_bf&=L^{-1}(-bf+D\eta f_\eta+Xf_X),\\
 \partial_z(q^bf)&=q^{b-D}Z_bf,&
 Z_bf&=L^{-1}(2b\eta f+df_\eta-2\eta Xf_X).
\end{align*}
In particular two axial derivatives mean $Z_{b-D}Z_b$, in that order.
There is no division by $d$, including at $\eta=\pm1$. These are the
operators of [NS, (4.2), pp.~24--25].

Retain the fixed constant $C>1$. Put $\lambda_n=2nh$, $E_n=\sqrt{2X}\phi_n/C$,
and
\[
u_{\theta,n}=rq^{-A-1/2+\lambda_n}\phi_n/C=q^{-A+\lambda_n}E_n,
\quad u_{z,n}=q^{-A+\lambda_n}U_n,\quad
ru_{r,n}=q^{\lambda_n}V_n,\quad p_n=q^{-2A+\lambda_n}\Pi_n.
\]
The zero-index profiles are the leading profiles. All negative-index profiles
are zero. Since $\partial_s(ru_r)+\partial_zu_z=0$ and $A+D=1$,
\[
 (V_n)_X=-Z_{-A+\lambda_n}U_n.
\]
For $F_n=\int_0^XU_n(x,\eta)\,dx=X\AX U_n$ and the regular axis value
$V_n(0)=0$, integration of $x(U_n)_x$ gives
\begin{equation}\label{eq:V}
 V_n=\frac{2\eta XU_n-2\eta(D+\lambda_n)F_n-d(F_n)_\eta}{L}.
\end{equation}
This calculation fixes both coefficients of the moment terms; replacing
$D+\lambda_n$ by $D$ would change incompressibility.

The angular vector Laplacian and the axial scalar Laplacian have different
axis coefficients. Directly using $\partial_rX=r/q$ gives
\[
(\partial_{rr}+r^{-1}\partial_r-r^{-2})(r\phi(X))
 =\frac{2r}{q}(X\phi_{XX}+2\phi_X),
\]
\[
(\partial_{rr}+r^{-1}\partial_r)U(X)
 =\frac2q(XU_{XX}+U_X),\qquad
(\partial_{rr}+r^{-1}\partial_r-r^{-2})\frac{V(X)}r
 =\frac{2X}{qr}V_{XX}.
\]
Angular transport is $u_r(\partial_r+r^{-1})u_\theta$; after removing
$rq^{-A-3/2+\lambda_n}/C$ its coefficient is
$V_i((\phi_j)_X+\phi_j/X)$. Axial radial transport gives $V_i(U_j)_X$.
Multiplication by $u_z$ adds the correct $Z$ operator because $A+D=1$.
Finally $1-2D=2h$: axial viscosity of index $n-1$ occurs at index $n$.
These identities give the following exact coefficient equations, with
$b=-A-1/2$, $c=-A$:
\begin{align}
2(X\phi_{n,XX}+2\phi_{n,X})
 &=T_{b+\lambda_n}\phi_n+
 \sum_{i+j=n}\{V_i(\phi_{j,X}+\phi_j/X)+U_iZ_{b+\lambda_j}\phi_j\}\notag\\
 &\quad-Z_{b+\lambda_{n-1}-D}Z_{b+\lambda_{n-1}}\phi_{n-1},\label{eq:phi}\\
2(XU_{n,XX}+U_{n,X})
 &=T_{c+\lambda_n}U_n+
 \sum_{i+j=n}\{V_iU_{j,X}+U_iZ_{c+\lambda_j}U_j\}\notag\\
 &\quad+Z_{-2A+\lambda_n}\Pi_n
 -Z_{c+\lambda_{n-1}-D}Z_{c+\lambda_{n-1}}U_{n-1},\label{eq:U}\\
\Pi_{n,X}&=C^{-2}\sum_{i+j=n}\phi_i\phi_j-\Omega_{n-1}/(2X),\label{eq:Pi}\\
\Omega_k&=T_{\lambda_k}V_k+
 \sum_{i+j=k}\{V_i(V_{j,X}-V_j/(2X))+U_iZ_{\lambda_j}V_j\}\notag\\
 &\quad-2XV_{k,XX}-Z_{\lambda_{k-1}-D}Z_{\lambda_{k-1}}V_{k-1}.
\label{eq:Omega}
\end{align}
For the last two lines multiply the physical radial equation by $r$.
The pressure term is $2X\Pi_{n,X}$ and the centrifugal term is
$-\sum E_iE_j$. All other radial terms carry $q^{-1+\lambda_k}$,
whereas these carry $q^{-2A+\lambda_n}$; equality requires $k=n-1$.
This proves the sign $-\Omega_{n-1}/(2X)$ rather than choosing it by convention.
The terms in \eqref{eq:Omega} all contain a factor $X$ when $V_j=Xv_j$.
For example $V_i(V_{j,X}-V_j/(2X))=Xv_i(v_j/2+Xv_{j,X})$.
The $T$ and $Z$ operators preserve divisibility by $X$, as does $XV_{XX}$.
Therefore the quotient in \eqref{eq:Pi} is a smooth axis source.
Equations \eqref{eq:V}--\eqref{eq:Omega} verify [NS, (5.2)--(5.6), pp.~46--47].

\section{The analytic radial recursion on an interval independent of the order}
The precise input here is the leading axis solution and its first collar from
[NS, Theorem 4.6(i), Proposition B.2]: there is a fixed interval
$0\le\xi\le a$, $X=\xi^2$, where the leading coefficients are smooth radially
and holomorphic in a complex neighborhood of $[-1,1]$. The collar is also the
location from which one chooses the retained radius $X_-$. This input is owned
by the profiles audit; no existence of that leading solution is inferred here.

At positive order the sums in \eqref{eq:phi}--\eqref{eq:Pi} are linear in the
current profiles. Indeed their only current-index pairs are $(0,n)$ and $(n,0)$.
For a completely local formulation define
\[
 K_n=\AX U_n-U_n,\qquad
 W_n=(\phi_n,U_n,K_n,\Pi_n,\partial_\xi\phi_n,\partial_\xi U_n)^T.
\]
Then $K_{n,\xi}+2K_n/\xi=-U_{n,\xi}$. Formula \eqref{eq:V} becomes
\[
 \frac{V_n}{X}=\frac{2\eta(A-\lambda_n)U_n
 -2\eta(D+\lambda_n)K_n-d(U_{n,\eta}+K_{n,\eta})}{L}.
\]
Also
$2(X\partial_{XX}+2\partial_X)=\tfrac12(\partial_{\xi\xi}+3\xi^{-1}\partial_\xi)$
and
$2(X\partial_{XX}+\partial_X)=\tfrac12(\partial_{\xi\xi}+\xi^{-1}\partial_\xi)$.
The pressure row is
$\Pi_{n,\xi}=4\xi C^{-2}\phi_0\phi_n+f_{n,4}$, with known smooth $f_{n,4}$.
Substituting this row into the term $X\Pi_{n,X}$ in the axial equation gives
\begin{equation}\label{eq:volterra}
 W_{n,\xi}+\xi^{-1}\diag(0,0,2,0,3,1)W_n
   =A_0W_n+A_1\partial_\eta W_n+f_n,
 \qquad W_n(0)=0.
\end{equation}
There are no singular entries left in $A_0$ or $A_1$: $V_0/\xi=\xi v_0$,
$V_n\phi_0/X=v_n\phi_0$, and $\Omega_{n-1}/X$ were explicitly resolved above.
The first four equations and \eqref{eq:phi}--\eqref{eq:U} specify every entry
of $A_0$ exactly. More usefully for convergence, its derivative matrix has precisely
the following potentially nonzero entries, with $H_c=D\eta+dU_0$:
\[
\begin{aligned}
(A_1)_{51}&=2H_c/L,&
(A_1)_{52}=(A_1)_{53}&=-2d(\phi_0+X\phi_{0,X})/L,\\
(A_1)_{62}&=2(H_c-dXU_{0,X})/L,&
(A_1)_{63}&=-2dXU_{0,X}/L,&
(A_1)_{64}&=2d/L.
\end{aligned}
\]
For example the last entry follows from the $d\Pi_{n,\eta}/L$ term and the
factor two introduced when converting the axial radial operator. These entries
verify the sparse matrix in [NS, p.~48], including its factors of two.

Here is a direct convergence proof that allows the order-dependent norms to be
arbitrarily large. For $c_i=(0,0,2,0,3,1)_i$ let
\[
(Gg)_i(\xi)=\int_0^\xi(s/\xi)^{c_i}g_i(s)\,ds,
\qquad K=G(A_0+A_1\partial_\eta).
\]
The kernel is diagonal, has norm at most one, and is independent of $\eta$.
Each $\partial_\eta^jA_1$ has range in coordinates $5,6$ and annihilates
coordinates $5,6$. Thus $A_1(\xi)D_0(\partial_\eta^jA_1)(s)=0$ for every
intervening diagonal kernel $D_0$. Expanding a word of $k$ factors $K$,
any adjacent choice of two derivative factors is zero, including the term where
the left derivative lands on the right matrix. Consequently a surviving word
contains at most $p_k=\lceil k/2\rceil$ derivatives in $\eta$.

Choose $S_\rho=\{\eta:\dist(\eta,[-1,1])<\rho\}$ with no zero of $L$,
on which coefficients and source are bounded uniformly in $\xi$. For a smaller
$S_{\rho'}$, distribute the gap $\Delta=\rho-\rho'$ equally among the at most
$p_k$ differentiations. Repeated Cauchy estimates bound their cost by
$(\max(1,p_k/\Delta))^{p_k}$. The $k+1$ ordered radial integrals have volume
$a^{k+1}/(k+1)!$. The at most $2^k$ word choices can be absorbed into a constant
$C_n^{k+1}$. Therefore
\[
 \|K^kGf_n\|_{S_{\rho'}}
 \le \frac{C_n^{k+1}a^{k+1}}{(k+1)!}
       (\max(1,p_k/\Delta))^{p_k}.
\]
The logarithm divided by $k$ is $-\tfrac12\log k+O(1)$, so the series
$W_n=\sum_{k\ge0}K^kGf_n$ converges for this entire fixed $a$.
For a homogeneous bounded analytic solution, $W=K^kW$ gives the same estimate
with $k$ ordered integrals; it tends to zero. This proves uniqueness in the
stated analytic class without a bound uniform in $n$.

To obtain axis regularity, use the nonsingular identities
\[
(Gg)_i=\xi\int_0^1t^{c_i}g_i(t\xi)\,dt,
\qquad \partial_\xi(Gg)_i=g_i(\xi)-c_i\int_0^1t^{c_i}g_i(t\xi)\,dt.
\]
Start with continuous $W$ holomorphic in $\eta$; Cauchy estimates give continuous
$\eta$ derivatives on a smaller neighborhood, so the right sides give a radial
derivative. Repeating this argument gives every fixed mixed derivative on some
smaller neighborhood. The equations are invariant under reflection of $\xi$
with the first four components even and the last two odd. Uniqueness of the
reflected integral equation imposes these parities. An even $C^\infty$ function
of $\xi$ is $C^\infty$ as a function of $X=\xi^2\ge0$: Taylor's formula to
order $2m+2$ contains only even powers and gives the needed remainders after
each of the first $m$ $X$ derivatives. This works for every $m$, with the same
parameter regularity at that finite order. It proves the local conclusion of
[NS, Lemma 5.1], with the leading-profile existence kept as its exact input.

\section{Power moment maps and exact corrections}
\subsection{An invertible moment map and its quadratic inverse}
For distinct real exponents $\alpha_1,\ldots,\alpha_m$, a nonzero sum
$\sum c_ix^{\alpha_i}$ has at most $m-1$ positive zeros. To prove this by
induction, order the exponents and divide by the smallest power. If the result
had $m$ zeros, its derivative would have $m-1$ zeros by Rolle's theorem. That
derivative is either zero, in which case the original result is constant, or
is a nonzero combination of $m-1$ distinct powers, which has at most $m-2$
zeros by induction. Both alternatives exclude $m$ zeros.

For ordered disjoint intervals $I_1<\cdots<I_m$ and nonnegative nonzero
$\beta_j\in C_c^\infty(I_j)$, the determinant of $(x_j^{\alpha_i})$ is
nonzero at every $x_1<\cdots<x_m$, and its continuous sign is fixed on that
connected set. Multilinearity and Fubini give
\[
\det\left(\int x^{\alpha_i}\beta_j(x)\,dx\right)
=\int_{I_1\times\cdots\times I_m}\det(x_j^{\alpha_i})
 \prod_j\beta_j(x_j)\,dx_1\cdots dx_m\ne0.
\]
Smooth families on compact parameter sets have bounded inverse derivatives,
by the determinant formula and $\partial B^{-1}=-B^{-1}(\partial B)B^{-1}$.
Under $x=e^y$, this is also the exponential-weight statement in [NS, Lemma A.1].
For two identical translated logarithmic bumps at centers separated by
$\Delta>0$, after dividing row $i$ by $e^{s_ic_1}$ the determinant is
$b(s_1)b(s_2)(e^{s_2\Delta}-e^{s_1\Delta})$, where
$b(s)=\int e^{st}\beta(t)\,dt>0$. On a fixed compact slope interval,
the mean value theorem gives an inverse bound $O(|s_1-s_2|^{-1})$.
In particular the near-coincident weights in this construction cost exactly
an allowed $O(\lambda^{-1})$ upper bound, not a constant independent of $\lambda$.

For completeness the finite nonlinear solve can be made fully quantitative.
Let $F_\eta(c)=B(\eta)c+Q_\eta(c,c)$, with $Q$ bilinear. In a chosen
$C^k([-1,1])$ norm let multiplication by $B^{-1}$ have norm $\beta_k$ and
the bilinear map have norm $\kappa_k$. If
$8\beta_k^2\kappa_k\|d\|_{C^k}\le1$, the map
$\Phi(c)=B^{-1}(d-Q(c,c))$ preserves the ball of radius
$r=2\beta_k\|d\|_{C^k}$, since
\[
\|\Phi(c)\|\le\beta_k\|d\|+\beta_k\kappa_kr^2
\le r/2+r/4<r.
\]
Its Lipschitz constant is at most $2\beta_k\kappa_kr\le1/2$.
Iteration from zero therefore converges to the unique solution in this ball.
At $k=0$, $\|B^{-1}D_cQ(c,c)\|\le1/2$, so the Jacobian is invertible.
Pointwise implicit differentiation proves smoothness. Solving in any fixed
$C^k$ ball gives the same branch because the iterates are identical.
This proves [NS, Lemma A.2] and also specifies the branch selected in every use.

The five original leading moments are
\[
M=\int_0^XU\,dx,\quad I=\int_0^X\sqrt{2x}E\,dx,\quad
J=\int_0^XU\sqrt{2x}E\,dx,
\]
\[
S=\int_0^X(U^2-E^2/2)\,dx,\qquad C_p=\int_0^X E^2/(2x)\,dx,
\quad\Pi=\Pi_{\rm ax}+C_p.
\]
The substitution $X=\rho x$ is an exact map: the factors of these five
moments are $(\rho,\rho^{3/2},\rho^{3/2},\rho,1)$, with
$H=\rho^{1/2}\sqrt{2x}E$. No factor is discarded. At
$U=u_c(\eta)$ and $E=e_*f_*(\eta)x^\alpha$, take the differential of
the rows $J-u_cI$ and $S-2u_cM$, where $u_c$ in the row operation stays at
its original value. These differentials are
$\int\sqrt{2x}E\,\delta U\,dx$ and $-\int E\delta E\,dx$.
Together with $M,I,C_p$, they give the two blocks of weights
\[
(1,x^{\alpha+1/2}),\qquad (x^{1/2},x^\alpha,x^{\alpha-1}),
\]
with the factors $e_*f_*$ and $\sqrt2$ retained in their row multipliers.
The exponents are distinct precisely when
$\alpha\notin\{-1/2,1/2,3/2\}$.
The functionals are degree at most two, so their unlisted remainder is
exactly a quadratic bilinear map. This derives [NS, Corollary A.3].

\subsection{The five positive-order moments, with all row signs}
Let $R=\sqrt{2X}$, so $dX=R\,dR$. Choose fixed radii
$X_-<X_{\rm keep}<X_{\rm cut}<a^2<\inf I_{\rm pos}$, with $X_-$ in the
leading inner collar. Take a smooth radial function $\kappa$, equal to one
through $X_{\rm keep}$ and zero from $X_{\rm cut}$ onward. It may be
constructed from the step $\sigma$ defined in Section 6 by affine rescaling.
The preliminary extension is
$\widetilde\phi_n=\kappa\phi_n^{\rm in}$,
$\widetilde U_n=\kappa U_n^{\rm in}$, extended by zero after $a^2$, and
$\widetilde E_n=R\widetilde\phi_n/C$.
Add two $U$ bumps and three $E$ bumps in the reserved interval $I_{\rm pos}$:
$U_n=\widetilde U_n+\sum_{j=1}^2\alpha_jb^U_j(R)$,
$E_n=\widetilde E_n+\sum_{j=1}^3\beta_jb^E_j(R)$,
and $\phi_n=CE_n/R$ by its smooth axis extension. Choose the bumps
nonnegative, of unit integral in $dR$, with disjoint ordered supports.
Here the exact leading values are
$U_0=0$, $E_0=e_*f(\eta)R^{-1-2\lambda}$, $e_*>0$, $f>0$.
The moments to prescribe are
\begin{align*}
m_1&=\int_0^\infty RU_n\,dR,&m_2&=\int_0^\infty R^2E_n\,dR,&
m_3&=\int_0^\infty\Pi_{n,R}\,dR,\\
m_4&=\int_0^\infty R^2\sum_{i+j=n}U_iE_j\,dR,&
m_5&=\int_0^\infty\left(R\sum_{i+j=n}U_iU_j-\tfrac12R^2\Pi_{n,R}\right)dR.
\end{align*}
Reconstruct $F_n,V_n$ by integration and \eqref{eq:V}, and reconstruct pressure
by \eqref{eq:Pi} with zero axis value. The already finalized lower orders are
fixed. Forward integration and the common zero axis values make the new
profiles agree exactly with the inner solution through $X_-$.
Since $\kappa$ is independent of $\eta$ and all bumps lie past $a^2$,
the entire interval $[0,a^2]$ retains parameter analyticity, even where
the radial cutoff has changed the coefficient. Thus the next recursion uses
finalized lower orders on the same analytic rectangle. Then
\[
\Pi_{n,R}=\frac{2E_0E_n+\sum_{i=1}^{n-1}E_iE_{n-i}-\Omega_{n-1}}{R},
\]
\[
m_5=\int_0^\infty\left(\sum U_iU_j-\tfrac12\sum E_iE_j+\tfrac12\Omega_{n-1}\right)dX.
\]
Hence the exact changes under $\delta U_n,\delta E_n$ on this patch are
\[
\begin{array}{lll}
\delta m_1=\int R\delta U_n\,dR,&
\delta m_4=e_*f\int R^{1-2\lambda}\delta U_n\,dR,\\
\delta m_2=\int R^2\delta E_n\,dR,&
\delta m_3=2e_*f\int R^{-2-2\lambda}\delta E_n\,dR,&
\delta m_5=-e_*f\int R^{-2\lambda}\delta E_n\,dR.
\end{array}
\]
There are no current-order quadratic terms because $n>0$ and $i+j=n$.
For the unit-integral ordered bumps, define $B_U$ with row powers
$(1,1-2\lambda)$ and $B_E$ with powers $(2,-2-2\lambda,-2\lambda)$.
All powers in each block are distinct for $\lambda>0$. If $m^0$ are the
moments before these five additions, the exact solution is
\[
\alpha=-B_U^{-1}\binom{m_1^0}{m_4^0/(e_*f)},\qquad
\beta=-B_E^{-1}\begin{pmatrix}m_2^0\\m_3^0/(2e_*f)\\-m_5^0/(e_*f)\end{pmatrix}.
\]
Thus the correction works for arbitrary finite discrepancies; no smallness or
nonlinear inversion is needed at positive order. It is smooth in $\eta$ because
$f$ is bounded away from zero. This independently verifies (5.14)--(5.16).

Beyond all these velocity bumps, $F_n=m_1=0$, hence $V_n=0$ by
\eqref{eq:V}. Choose $X_+$ also beyond the leading axial pulse, so $V_0=U_0=0$
there. Every term of every $\Omega_k$ then vanishes, including its shifted
viscosity term. Since a product $\phi_i\phi_j$ of positive total order has a
positive-order factor, $\Pi_{n,X}=0$ there and $\Pi_n=m_3=0$.
This proves the common support of all five reconstructed profiles. It is essential
to include the leading axial pulse in this choice: $\Omega_0$ need not vanish
immediately after the positive-order correction patch.

On the separate mean patch, the positive-order velocities already vanish:
$E_n=U_n=0$ beyond their bump supports, and $F_n=0$ there by $m_1=0$.
Therefore $V_n=0$ as well. This assertion concerns the velocities, and does not
assert that positive-order pressure has vanished at this earlier patch.

\section{Weighted residual integrals and compactly supported stresses}
For any smooth axisymmetric divergence-free field, expand
$\partial_r(ru_r)+r\partial_zu_z=0$ in the two tangential momentum equations.
One obtains the exact conservative identities
\begin{align*}
r^2\res_\theta={}&\partial_t(r^2u_\theta)+\partial_r(r^2u_ru_\theta)
+\partial_z(r^2u_zu_\theta)-\partial_r(r^2u_{\theta,r}-ru_\theta)
-\partial_{zz}(r^2u_\theta),\\
r\res_z={}&\partial_t(ru_z)+\partial_r(ru_ru_z)
+\partial_z(r(u_z^2+p))-\partial_r(ru_{z,r})-\partial_{zz}(ru_z).
\end{align*}
For example the two surplus terms from the angular flux derivatives are
$2ru_ru_\theta+r^2u_\theta(\partial_ru_r+\partial_zu_z)
=ru_ru_\theta$, precisely the cylindrical-basis term.

At index $n>0$, the moments before physical differentiation are
\begin{align*}
\int r^2u_{\theta,n}\,dr&=q^{3/2-A+\lambda_n}m_2,\qquad
\int r^2\sum u_{z,i}u_{\theta,j}\,dr=q^{3/2-2A+\lambda_n}m_4,\\
\int ru_{z,n}\,dr&=q^{1-A+\lambda_n}m_1,\qquad
\int r\left(\sum u_{z,i}u_{z,j}+p_n\right)dr
=q^{1-2A+\lambda_n}m_5.
\end{align*}
For the last equality,
$\int R\Pi_n\,dR=-\tfrac12\int R^2\Pi_{n,R}\,dR$, with zero boundary
terms because pressure is smooth at zero and compactly supported.
All four moments vanish as identities in $\eta$. Their $t,z$ derivatives
therefore vanish before any coefficient powers are removed. Axis parity and
compact positive-order supports eliminate the radial boundary terms.

The angular axial-viscosity term at $n=1$ needs a separate exact integral,
proved below from the heat construction:
\begin{equation}\label{eq:submoment}
 \int_0^\infty r^2(u_{\theta,0}-c_\infty(r^2/2)^{-A})\,dr=0.
\end{equation}
The pure power is independent of $z$. The second $z$ derivative of this
convergent integral is legitimate since the exterior heat field is independent
of $z$; on compact physical parameter neighborhoods all remaining $z$ derivatives
have bounded radial support. For $n\ge2$ the lower angular moment is $m_{n-1,2}=0$.
The lower axial-viscosity moment is zero also for $n=1$ because the leading
$M(\infty)$ vanishes. Consequently
\[
\int_0^\infty R^2r_{\theta,n}\,dR=0,
\qquad\int_0^\infty Rr_{z,n}\,dR=0.
\]
Here $\res_{j,n}=q^{-A-1+\lambda_n}r_{j,n}$; the shifted viscous terms
have this same physical power, as checked in Section 1.

Define
\[
T_{n,\theta}=-R^{-2}\int_0^R\rho^2r_{\theta,n}\,d\rho
=R^{-2}\int_R^\infty\rho^2r_{\theta,n}\,d\rho,
\]
\[
T_{n,z}=-R^{-1}\int_0^R\rho r_{z,n}\,d\rho
=R^{-1}\int_R^\infty\rho r_{z,n}\,d\rho.
\]
Then $(\partial_R+2/R)T_{n,\theta}=-r_{\theta,n}$ and
$(\partial_R+1/R)T_{n,z}=-r_{z,n}$, including their signs.
The inner equations force zero stress through $X_-$; for $n\ge2$ all exterior
sources vanish past $X_+$, giving common interior stress support. For $n=1$
only angular axial viscosity of the leading field can remain beyond $X_+$.
The exterior is $z$ independent past $X_b$, so that stress still vanishes there.
Its quantitative edge bound is proved after the flat-integral calculation below.

\section{Finite residuals, exact divergence after summation, and uniformity}
For a smooth Cartesian profile $g(x_\perp/\sqrt q,\eta)$, direct repeated use
of the coordinate derivatives gives
\[
|\partial_{x_\perp}^{\beta}\partial_z^k\partial_t^\ell(q^bg)|
\le C_m(1+|b|)^m\|g\|_{C^m}
q^{b-|\beta|/2-Dk-\ell},\qquad m=|\beta|+k+\ell.
\]
The factors produced by differentiating $b$ enter the constant, whereas the
power loss depends only on derivative orders. Cartesian axis regularity follows
from $E_n=\sqrt{2X}\phi_n/C$ and $V_n=Xv_n$:
\[
(u_{1,n},u_{2,n})=
\tfrac12q^{-1+\lambda_n}v_n(x_1,x_2)
+q^{-A-1/2+\lambda_n}(\phi_n/C)(-x_2,x_1).
\]
This equality is componentwise, with the same scalar profiles evaluated at
$(X,\eta)$. It eliminates apparent $r^{-1}$ singularities at the axis.
On stress support, $X\ge X_a>0$, each Cartesian derivative of $r^{-1}$ costs
a fixed further negative power of $q$.

For a truncation through $N$, equations in Sections 1--4 cancel every retained
coefficient of
\[
F_{\rm slow}(u,p,T)=\res(u,p)
+(\partial_r+2/r)T_\theta e_\theta+(\partial_r+1/r)T_ze_z.
\]
The remaining indices are at least $N+1$. Tangential terms have base power
$q^{-A-1}$ and radial terms base power $q^{-2A-1/2}$; the latter is the worse
of the two because $A-1/2=h>0$. Thus one valid explicit loss in (5.25) is
\[
 K_m=2A+\tfrac12+m,
\qquad |F_{\rm slow}(U^{[N]})|_m
 \le C_{N,m}q^{2h(N+1)-K_m},\qquad0<q\le1,
\]
on each compact profile rectangle. This is independent of $N$ as required.

The exact streamfunction and potential for a positive-order velocity are
\[
S_n=q^{1-A+\lambda_n}F_n,\qquad
\mathcal A_n=(S_n/r)e_\theta
=\tfrac12q^{-A+\lambda_n}(F_n/X)(-x_2,x_1,0).
\]
They satisfy $u_{z,n}=\partial_sS_n$ and $ru_{r,n}=-\partial_zS_n$.
For the latter identity use $\partial_z(q^{1-A+\lambda_n}F_n)$ and
$1-A=D$: the resulting bracket is exactly the negative of \eqref{eq:V}.
Because $F_n/X=\int_0^1U_n(tX,\eta)\,dt$, this potential is smooth at the axis.

We give the summation argument with the quantifiers needed for uniform residual
bounds. Suppose potentials and remaining tuple entries satisfy
$|Z_j|_m\le C_{j,m}q^{g_j-\ell_m}(1+|\log q|)^{P_{j,m}}$, with
$g_j>0$ nondecreasing to infinity and $\ell_m$ independent of $j$.
In the present coefficient construction one may take $g_j=2jh$ and
$\ell_m=2A+m$. For a cutoff $\chi(aq)$, equal to one on $aq\le1/2$ and
zero on $aq\ge1$, every differentiated term has the form
$a^k\chi^{(k)}(aq)\prod_{i=1}^k\partial_z^{\alpha_i}\partial_t^{\beta_i}q$.
The bound $|\partial_z^\alpha\partial_t^\beta q|
\le C_{\alpha,\beta}q^{1-D\alpha-\beta}$ supplies $q^k$ to pair with $a^k$.
On support of a nonzero differentiated cutoff, $1/2\le aq\le1$.
Thus its derivatives cost only $q^{-D\sum\alpha_i-\sum\beta_i}$, uniformly
in $a$.

Apply these cutoffs to the potentials, not just to their curls:
\[
\curl(\chi(a_jq)\mathcal A_j)
=\chi(a_jq)\curl\mathcal A_j+a_j\chi'(a_jq)\nabla q\times\mathcal A_j.
\]
At each $j$, choose $a_j\ge2a_{j-1}$ sufficiently large to ensure
\[
\widehat C_{j,m}(1+|\log q|)^{\widehat P_{j,m}}q^{g_j/2}
\le2^{-j},\qquad0<q\le a_j^{-1},\quad0\le m\le j.
\]
Only finitely many inequalities are imposed at this step and each holds near
zero because $g_j>0$. Each cutoff increment then has bound
$2^{-j}q^{g_j/2-\ell'_m}$, with $\ell'_m$ independent of $j$ and accounting
for the extra curl derivative. For $q<(2a_J)^{-1}$ the first $J$ cutoffs
equal one, so
\begin{equation}\label{eq:tail}
 |U-U^{[J]}|_m\le2^{-J}q^{g_{J+1}/2-\ell'_m},\qquad J\ge\max(1,m).
\end{equation}
The sum is locally finite because $q$ is bounded below on compact subsets.
Every curl is divergence free, and an axisymmetric angular field
$b(r,z,t)e_\theta$ remains divergence free when multiplied by $\chi(a_jq)$.
Consequently the resulting velocity is exactly divergence free throughout every
cutoff transition.

For a fixed differential polynomial $F$ of order $s$ and degree $d_F$,
write the difference of each monomial as the sum obtained by changing one
factor at a time. Leibniz's rule and polynomial bounds on its fixed coefficients
give, for $e=U-U^{[J]}$,
\[
|F(U)-F(U^{[J]})|_m
\le C_mq^{-H_m}|e|_{m+s}
(1+|U^{[J]}|_{m+s}+|e|_{m+s})^{d_F-1}.
\]
The loss $H_m$ is independent of $J$. The finite partial sum has a bound
$|U^{[J]}|_k\le C_{J,k}q^{-K_k}(1+|\log q|)^{P_{J,k}}$ with $K_k$
independent of $J$. Fix $m$ and the desired power $N$ first. Choose $J$ so
large that the exponent in \eqref{eq:tail}, after the losses in this displayed
comparison, exceeds $N$, and the finite residual exponent exceeds $N+1$.
Keep this $J$ fixed and decrease the neighborhood of $q=0$ so its fixed
constants and logarithms are bounded by one additional $q^{-1}$ factor.
Then the comparison is $O(q^N)$ and so is the fixed finite residual. Any flat
remainder at that fixed stage remains $O(q^N)$; there is no exchange of limits
between an unknown infinite family of remainders. This proves [NS, Lemma 5.4].

For the actual coefficient potentials, the extra radial velocity is exactly
\begin{equation}\label{eq:cutrad}
ru_{r,n}^{\rm cut}=q^{2nh}\left[\chi(c_nq)V_n
-\frac{2\eta}{L}(c_nq)\chi'(c_nq)F_n\right],\qquad
u_{z,n}^{\rm cut}=q^{-A+2nh}\chi(c_nq)U_n.
\end{equation}
This follows by differentiating $\chi(c_nq)S_n$ in $z$; its minus sign is
fixed by $ru_r=-S_z$. It shows explicitly that the cutoff correction has the
same radial support and vanishes on the reserved mean patch.

The diagonal choice can also impose, through order $n$, the finite collection
of $q\partial_q,\partial_X,\partial_\eta$ seminorms of every normalized
coefficient, of the bracket in \eqref{eq:cutrad}, and for $n\ge2$ of $T_n/\zeta$.
It can therefore make these terms at most $2^{-n}q^{nh}$.
At a fixed derivative order retain the finitely many early terms; they all
have their original power at least $q^{2h}$. Sum the tail starting at an index
at least two and above the derivative order to obtain $O(q^{2h})$.
It follows that the three differences in the first line of (5.42) are
$O(q^{2h})$. In the common $q^A$ velocity scale the leading radial component
is $O(q^h)$ and its positive-order difference is $O(q^{3h})$.
For stress the only noncompact interior quotient is the single $n=1$ edge
term established below; this yields precisely
$|\D^I(\widehat T-T_0)|\le C_Iq^{2h}\zeta\delta^{-N_I}$.

For a full expansion after order $N$, take a tail index
$J\ge\max(2N+2,m+1)$. The intermediate finite block retains the original
orders at least $q^{2h(N+1)}$, whereas
\[
\sum_{n\ge J}2^{-n}q^{nh}\le2^{1-J}q^{Jh}
\le2^{1-J}q^{2h(N+1)}.
\]
Physical derivatives introduce only the already fixed losses. Thus no
asymptotic coefficient is changed by using half-powers for the tail estimates.
The construction is simultaneous on every compact profile range: all positive
velocity and pressure corrections are supported below the single fixed $X_+$,
and all stresses below $X_b$. Outside them the exact leading heat solution has
zero augmented residual. A cutoff sequence selected on one rectangle containing
these supports therefore works for every larger rectangle as well.

On any compact set with $q>0$, only finitely many positive-order profiles occur.
Their parameter neighborhoods can be intersected. The leading heat factor has
the one-sided smoothness proved below, and $L\ge1-2h$ transfers every fixed
physical derivative continuously to $\eta=\pm1$. The resulting residual is
flat as $q\downarrow0$, uniformly for all $|\eta|\le1$ on compact profile
ranges. This is the uniformity asserted in [NS, Proposition 5.5]; it does not
assert that the singular velocity itself extends smoothly at $q=0$.

\section{The complete outer schedule and its exact moment closure}
\subsection{Definitions of all stages and the adjustable coefficients}
Use the fixed step
\[
\sigma(y)=\frac{e^{-1/y^2}}{e^{-1/y^2}+e^{-1/(1-y)^2}}\ (0<y<1),
\qquad \sigma=0\ (y\le0),\quad\sigma=1\ (y\ge1).
\]
The schedule chooses, in order, $M_d$ large, $T_d=e^{M_d}+10$,
$P_*>e^{T_d}$, $\lambda>0$ small depending on these choices, and
$h>0$ sufficiently small relative to both $\lambda$ and $e^{-T_d}$.
Only after these choices is $X_R$ made large. All estimates denoted $C_{\rm pre}$
below depend on the first three choices and the fixed steps, not on the later
$X_R$ or amplitude factor $C$. Smallness requirements use finitely many
derivatives, and parameters can be chosen to meet all of them.

Write $x=X/X_R$ and, initially, $y=\log x$. The reference fields for $y\le0$
are $U=4\eta$, $E=P_*f e^{y/10}$, $f=(1+\eta^2)^{-1}$,
$J_0=\log(1+\eta^2)$. In each next stage $y$ is reset to zero at that
stage's left endpoint, and $l=X\partial_X\log(\sqrt{2X}E)$.
The fields are fully specified by the following data, with values at endpoints
matched by integration of $\partial_y\log E=l-1/2$:
\begin{enumerate}
\item On a unit interval, $U=4\eta$, $l=(3/5)(1-\sigma(y))$.
On the following interval $0\le y\le T_d$, set $l=0$,
$U=k(y)\eta$, $k=4[1-\sigma(\log(1+y)/M_d)]$.
Thus $k$ vanishes for its last eleven units, and
$|k'|\le e_d/(1+y)$ with $e_d=4\|\sigma'\|_\infty/M_d$.
The swirl on this interval is $E=P_1fe^{-y/2}$, where $P_1\asymp P_*$.
\item On a unit interval set $U=0$, $l=-\lambda\sigma(y)$, and then
hold $l=-\lambda$ for $T_w=60\log(1/\lambda)$.
Reserve $(T_w-25,T_w-20)$, $(T_w-20,T_w-15)$,
$(T_w-14,T_w-9)$, and $(T_w-8,T_w-3)$ for the four later corrections.
\item At $X=X_p$ start a pulse with $0\le y\le13/\lambda$,
$E=e_bf e^{-(1/2+\lambda)y}$, $U=ER_b$, $\xi_b=\lambda y$.
Define $\varphi_b(\xi)=\int_0^\xi\sigma(v/.02)\,dv$ for $\xi\ge0$,
extended by zero to negative arguments, and
$R_0(\xi)=\varphi_b(\xi)[1-\sigma(\xi-10)]$.
Set $R_b=\mathrm{Amp}(\eta)R_0(\lambda y)+c_1\beta_1(y)+c_2\beta_2(y)$,
where the two identical width $.3$ bumps have centers $13/\lambda-3$ and
$13/\lambda-1$. Choose $c_1,c_2$ to make $M=J=0$ at the pulse end.
\item From the pulse endpoint $E=e_{\rm end}f$, set $U=0$ permanently and
\[
\log E=\log e_{\rm end}-(1/2+\lambda)y
-\vartheta_f(y)J_0-(1-\vartheta_f(y))\log2,
\quad\vartheta_f=1-\sigma(y/T_f).
\]
A fixed sufficiently large $T_f$ ensures $-\lambda-.1\le l\le-\lambda$.
Next retain the now $\eta$ independent exponential for $30\log(1/\lambda)$.
Two relative swirl bumps of width $.3$, three and one units before its end,
set $I=XH/(1-\lambda)$ there and make their total integral
$\int(E^2-E_{\rm unedited}^2)\,dy$ zero.
\item Starting at $X_{\rm rel}$, vary $l$ from $-\lambda$ to $-1$ in a
unit interval, hold $l=-1$ for $4\log(1/h)$, and vary from $-1$ to $-h$
in a unit interval, all using $\sigma$.
The terminal multiplier is
$f_o(y)=1-\rho_o\psi_o(y)$ on $0\le y\le3$,
$\psi_o=1-\sigma((y-1)/2)$, $\rho_o=c_oh>0$ with $c_o$ fixed so that
$0\le f_o'/f_o<h/4$.
Define
\[
Q_p=\int_0^3\exp\left(\int_0^v(1+l(s))\,ds\right)
\frac{f_o'(v)}{f_o(v)}\,dv,\qquad l=-h+f_o'/f_o
\]
on this terminal interval. Hold $l=-h$ before it until $Q_s$ reaches $Q_p$.
Then apply this terminal interval, and retain $l=-h$ thereafter.
On and after the terminal interval the field is exactly $E_{\rm pow}f_o$,
$E_{\rm pow}=c_\infty X^{-A}$, with $c_\infty$ independent of $\eta$.
\end{enumerate}
All stage endpoints divided by $X_R$ are constants independent of $\eta$.
The following calculations prove that the apparently adjustable stages have
the stated finite solutions, and determine the pulse amplitude.

\subsection{Pulse corrections, angular correction, and amplitude root}
During the long interval $T_w$, the quantities $XE$ and $XHE$ grow at rates
$1/2-\lambda$ and $1/2-2\lambda$, respectively, whereas $XE^2$ decays
at rate $2\lambda$. The preceding $M,J$ are constant once $U=0$.
Since $e^{-T_w/2}=\lambda^{30}$ and
$\lambda T_w=60\lambda\log(1/\lambda)\to0$, it follows, with one
$\eta$ derivative as well, that
\[
e_b\le C_{\rm pre}\lambda^{30},\quad
\left\|\frac{M(X_p)}{X_pe_bf}\right\|_{C^1}
+\left\|\frac{J(X_p)}{X_pH(X_p)e_bf}\right\|_{C^1}
\le C_{\rm pre}\lambda^{29},
\]
\[
\left\|\frac{S(X_p)}{X_pe_b^2f^2}\right\|_{C^1}
\le C_{\rm pre}(1+\log(1/\lambda)).
\]
For the last inequality, express $S$ as an integral of $X(U^2-E^2/2)$
in logarithmic coordinates. The finite preceding stages contribute a fixed
constant in this scale and the long $U=0$ stage contributes its length times
a bounded factor $e^{-2\lambda y}$. These statements prove (A.14) with all
scale factors retained.

The exact two equations for the pulse correction are
\[
\begin{pmatrix}
\int e^{s_1y}\beta_1\,dy&\int e^{s_1y}\beta_2\,dy\\
\int e^{s_2y}\beta_1\,dy&\int e^{s_2y}\beta_2\,dy
\end{pmatrix}\binom{c_1}{c_2}
=-\binom{M(X_p)/(X_pe_bf)+\mathrm{Amp}\int e^{s_1y}R_0(\lambda y)\,dy}
{J(X_p)/(X_pH(X_p)e_bf)+\mathrm{Amp}\int e^{s_2y}R_0(\lambda y)\,dy},
\]
where $s_1=1/2-\lambda$, $s_2=1/2-2\lambda$. Divide each row by its
exponential value at the first bump center. The main pulse is zero after
$11/\lambda$, leaving a gap at least $2/\lambda-3$; both slopes exceed $.4$
for small $\lambda$. Its normalized moments are at most a polynomial in
$\lambda^{-1}$ times $(1+|\mathrm{Amp}|)e^{-.4(2/\lambda-3)}$.
The prior moments have an even larger exponential gap. The matrix inverse costs
$O(\lambda^{-1})$. Reducing the positive exponential constant absorbs every
such fixed polynomial. Hence the exact affine solution
$c=c^{(0)}+\mathrm{Amp}\,c^{(1)}$ satisfies
\[
\|c^{(0)}\|_{C^1}+\|c^{(1)}\|_{C^1}\le C_{\rm pre}e^{-c/\lambda}.
\]
The same argument applies to any fixed derivative, with additional powers of
$\lambda^{-1}$ for $\xi_b$ derivatives. This proves exact $M=J=0$ and (A.15).

For angular momentum put $r_I=I/(XH)$. Since $I_X=H$,
$r_I'+(1+l)r_I=1$. During the uniform interval following interpolation,
$r_I-(1-\lambda)^{-1}$ decays at rate $1-\lambda$.
At the first swirl bump center its discrepancy is
$O_{C^1}(C_{\rm pre}\lambda^{28})$; the exponent $28$ is a valid relaxed
bound for $e^{-(1-\lambda)(30\log(1/\lambda)-3)}$.
For the relative edit $E\mapsto E(1+c_1\beta_1+c_2\beta_2)$,
the linear angular row has logarithmic weight $e^{(1-\lambda)y}$ and
the linear pressure row has weight $2e^{(-1-2\lambda)y}$.
Their slope separation tends to $2$, so their normalized inverse is bounded.
The pressure row alone has its exact quadratic remainder. The contraction
proved above yields coefficients $O_{C^1}(C_{\rm pre}\lambda^{28})$,
with both exact targets met. All fixed radial derivatives of these bumps
are of the same size, so they preserve $E>0$ and $l\le-\lambda/2$.

After the correction, $I=XH/(1-\lambda)$ and $U=M=J=0$ at the start
of the exterior transition. The exact radial identities of (4.16) give
\[
Q_s=-1+(1-h)I/(XH)=(\lambda-h)/(1-\lambda).
\]
Subsequently $Q_s'+(1+l)Q_s=-l-h$. During the first transition the
nonnegative source and bounded integrating factor retain a positive $c\lambda$
lower bound. The $l=-1$ hold gives $Q_s'=1-h$, so its end value is comparable
to $\log(1/h)$. The next unit transition changes this by bounded factors.
On the other hand $Q_p\asymp\rho_o\asymp h$: $f_o$ is bounded above and
below, the interval has fixed length, and $\int f_o'=\rho_o$.
The intervening $l=-h$ solution is $Q_{\rm in}e^{-(1-h)y}$; the exact required
length is $(1-h)^{-1}\log(Q_{\rm in}/Q_p)>0$ for sufficiently small $h$.
On the terminal interval the resulting solution is
\[
Q_s(y)=\int_y^3\exp\left(\int_y^v(1+l(s))\,ds\right)
\frac{f_o'(v)}{f_o(v)}\,dv.
\]
Differentiating this expression proves the equation and its zero right endpoint.
Thus $Q_s=0$ afterward and $I=XH_{\rm pow}/(1-h)$. Since
$\int_0^XH_{\rm pow}\,dx=XH_{\rm pow}/(1-h)$ for $h<1$, this is exactly
$\int_0^\infty(H-H_{\rm pow})\,dX=0$. The subtraction is integrable at
zero, and is identically zero at large $X$ at this pre-heat stage.

For the other moment, note the exact identities on the $l=-1$ hold:
\[
\frac{(XE^2)_{\rm end}}{(XE^2)_{\rm start}}=h^8,
\qquad \frac{E_{\rm end}}{E_{\rm start}}=h^6.
\]
They follow by integrating $(\log(XE^2))'=2l=-2$ and
$(\log E)'=-3/2$ over $4\log(1/h)$.
The first transition and hold have total $\int XE^2dy\le CX_{\rm rel}e_{\rm rel}^2$.
After the hold $l\le-3h/4$, so the remaining integral is at most
$CX_{\rm rel}e_{\rm rel}^2h^8\int_0^\infty e^{-3hy/2}\,dy
\le CX_{\rm rel}e_{\rm rel}^2h^7$. This proves the $h$ uniform bound (A.18).
Interpolation and its next hold cost at most
$C(1+\log(1/\lambda))X_{\rm end}e_{\rm end}^2$; their first parameter
derivatives satisfy the same bounds because $f^{\vartheta}$ and its derivatives
are bounded and the exterior is parameter independent.

There is also an exact algebraic specification of the amplitude root. Substitute
the affine $c$ above and write $R_b=R^{(0)}+\mathrm{Amp}\,R^{(1)}$.
In the full $S(\infty)$, every term outside $\int_{\rm pulse}E^2R_b^2\,dX$
is independent of $\mathrm{Amp}$. Consequently this moment is an explicitly
defined quadratic polynomial $a(\eta)\mathrm{Amp}^2+b(\eta)\mathrm{Amp}+c(\eta)$,
where, without dropping any end bumps,
\[
\begin{aligned}
a&=\int_{\rm pulse}E^2(R^{(1)})^2dX>0,\qquad
b=2\int_{\rm pulse}E^2R^{(0)}R^{(1)}dX,\\
c&=S_{\rm no\ axial\ pulse}(\infty)+\int_{\rm pulse}E^2(R^{(0)})^2dX.
\end{aligned}
\]
Changing variables $\xi_b=\lambda y$ in this polynomial gives
\[
\frac{\lambda S(\infty)}{X_pe_b^2f^2}
=\mathrm{Amp}^2K_b-\frac{1-e^{-26}}4+\mathcal E(\mathrm{Amp},\eta),
\quad K_b=\int_0^{13}e^{-2\xi}R_0(\xi)^2\,d\xi,
\]
\[
\|\mathcal E\|_{C^1([.9,1.2]\times[-1,1])}
\le C_{\rm pre}\lambda(1+\log(1/\lambda)).
\]
Indeed the prior $S$ and the post-pulse negative $E^2/2$ integral have precisely
the bounds just proved after multiplication by $\lambda$, whereas the affine
end bumps are exponentially small; all cross terms are retained in
$\mathcal E$. On $.02\le\xi\le9$, $R_0\ge\xi-.02$, and globally
$0\le R_0\le\xi$. Therefore
\[
\frac{e^{-.04}}4\left[1-e^{-17.96}(1+17.96+2(8.98)^2)\right]
\le K_b\le\frac14,
\]
whose lower bound exceeds $.20$. At $\mathrm{Amp}=.9$ the principal value
is below $-.047$; at $1.2$ it exceeds $.038$; its derivative is at least $.36$.
For sufficiently small $\lambda$ the exact polynomial has a unique root on
this interval. It is the positive branch
$(-b+\sqrt{b^2-4ac})/(2a)$, with its nonzero derivative already proved by
the bracket estimate. Implicit differentiation gives
$|\partial_\eta\mathrm{Amp}|\le C_{\rm pre}\lambda(1+\log(1/\lambda))$.
This constructs the smooth amplitude and proves both required moments, including
all terms in (A.19)--(A.20).

\section{Pressure datum and the cone margins on every outer stage}
\subsection{Exact pressure and its parameter analyticity}
Omit only the swirl bumps whose total $\int E^2dy$ was just proved to be zero.
The remaining scheduled field has the form $E=c(y)f(\eta)^{\vartheta(y)}$,
$0\le\vartheta\le1$, with $c,\vartheta$ and all interval lengths independent
of $\eta$ and $X_R$. On one simply connected complex neighborhood of $[-1,1]$
avoiding $\pm i$, choose the logarithm of $f$. The holomorphic functions
$\exp(\vartheta\log f)$ are uniformly bounded for $0\le\vartheta\le1$.
The squared integrand is bounded by $Ce^{y/5}$ as $y\to-\infty$ and
$Ce^{-(1+2h)y}$ as $y\to\infty$. Uniform convergence of the integral on
compact complex subsets therefore defines a holomorphic function
\[
\Pi_{\rm ax}(\eta)=-\tfrac12\int_{-\infty}^\infty E_{\rm sched}(y,\eta)^2dy.
\]
It is even, independent of $X_R$, and its derivative is an integral of
$\vartheta J_0'c^2f^{2\vartheta}$. This has the sign of $\eta$ and is
strict for $\eta\ne0$ because the reference interval has $\vartheta=1$.
The reference interval contributes exactly
$-\tfrac12P_*^2f^2\int_{-\infty}^0e^{y/5}dy=-\tfrac52P_*^2f^2$.
Restoring the pressure-neutral bumps leaves the total unchanged. Hence the
actual pressure is exactly
\[
\Pi(X,\eta)=-\tfrac12\int_{\log(X/X_R)}^\infty E(v,\eta)^2dv.
\]
This proves all assertions of [NS, Lemma A.5]. It also proves the precise
preservation map: an edit of zero total $\int\delta(E^2)/(2X)dX$ keeps
forward pressure unchanged before and after its support; it need not keep
pressure unchanged inside that support.

\subsection{The exact quantities used for the cone test}
Retain
\[
W=1-2D\eta M/X-dM_\eta/X,\quad H_c=D\eta+dU,
\quad a=2-2l,\quad b_s=2XU_X/E,
\]
\[
Q_s=-W+\frac{(1-h)I-D\eta I_\eta-dJ_\eta+2(h-D)\eta J}{XH},
\]
\[
N_s=-WU+\frac{D(M-\eta M_\eta)+4h\eta S-dS_\eta}{X}
+4A\eta\Pi-d\Pi_\eta.
\]
These formulas follow by integrating the source equations
\[
Q_s'+(1+l)Q_s=-Wl-h(1-2\eta U)-H_c(\log E)_\eta,
\]
\[
N_s'+N_s=-WU'-A(1-2\eta U)U-H_cU_\eta
-d\Pi_\eta+4A\eta\Pi+2\eta X\Pi_X.
\]
For example $(XW)_X=1-2D\eta U-dU_\eta$; integration by parts of
$-WXH_X$ produces $-XWH$ and the displayed $I,J$ terms. In the pressure
terms use $\int_0^X\Pi,dx=X\Pi-\int_0^XE^2/2,dx$.
These operations produce the exact $S$ and $M$ terms above, and establish the
map from the five original moments to stress coefficients in [NS, (4.16)].

Set $w=N_s/(EQ_s)$ where $Q_s>0$, and $p_{s,1}=XQ_s/L$,
$p_{s,2}=wp_{s,1}$. To explain precisely the sufficient test used below, put
\[
c=1-b_sw/a,\quad j=w+b_s/a,\quad v_s=a+b_s^2/a,
\quad P_c=p_{s,1}c,\quad J_c=p_{s,1}j.
\]
Direct expansion gives
\[
2c^2-(v_s-2)j^2
=\left(1+\frac{b_s^2}{a^2}\right)
\left[2-2b_sw-\frac{b_s^2}{a}-(a-2)w^2\right].
\]
Thus on a compact set with $a>0$, $a-b_sw>0$, and the bracket positive,
there are constants $c_*>0$, $G_*>0$, and $B<\infty$ bounding these
expressions away from zero and bounding $v_s,cv_s$. For
$p_{s,1}\ge\max((B+2)/c_*,8B/G_*)$,
$P_c>\max(2,v_s)$ and
\[
\frac{2(P_c-v_s)^2-(v_s-2)J_c^2}{p_{s,1}^2}
\ge G_*-4B/p_{s,1}\ge G_*/2.
\]
For $v_s>2$ these are exactly the admissible cone inequalities; for $v_s\le2$
they imply the relaxed cone. Equivalence follows by solving the quadratic
$2(P_c-v)^2-(v-2)J_c^2=0$, whose smaller root is the function in (4.21),
between $2$ and $P_c$. This verifies the finite threshold used in Lemma 4.5.

\subsection{Reference and axial-decrease stages}
On the ideal reference interval $U=4\eta$, $W=1-4L$, $l=3/5$, so
\[
Q_s=\frac{(3/5)(4L-1)-h(1-8\eta^2)+(D+4d)\eta J_0'}{8/5}>c>0
\]
for the chosen small $h$. Here $a=4/5$, $b_s=0$. Its $N_s$ geometric
source is $O(|\eta|)$ and its pressure source is $O(P_*^2|\eta|)$;
these estimates follow by substituting the constant $U$ and the bounded
backward pressure integral above. The vanishing at zero follows from even
pressure and odd $U$, and one parameter derivative removes the $|\eta|$
factor. The $N_s'+N_s$ convolution has finite bounds on each prescribed
reference interval $x\ge x_->0$.

During the first unit transition, $b_s=0$, $.8\le a\le2$, and
the angular source is at least $-h$ since $W<0$, $l\ge0$, and
$H_cJ_0'\ge0$. The initial positive lower bound and variation of constants
retain $Q_s>c$ for small $h$. With $b_s=0$ and $a\le2$, the second sufficient
cone bracket is automatically at least $2$.

During axial decrease, $l=0$, $a=2$, $b_s=2k'\eta/E$. The angular
source is exactly
$-h(1-2k\eta^2)+(D+dk)\eta J_0'$.
Since $\eta J_0'=2\eta^2/(1+\eta^2)\asymp\eta^2$, convolution with
$e^{-y}$ and $h\ll e^{-T_d}$ gives
$Q_s\ge c(\eta^2+e^{-y})$ for $0\le y\le T_d$.
The geometric axial source is $O(|\eta|)$. The pressure source is
$O(E^2|\eta|)$ because subsequent scheduled swirl decays at least
$e^{-\Delta y/2}$ and $|\partial_\eta\log E|\le |J_0'|$; the later
pressure-neutral angular bumps may be omitted from the backward integral.
Solving the $N_s$ equation yields
\[
|N_s|/E^2\le C|\eta|(e^{T_d}/P_*^2+1+y)
\le C|\eta|(1+y).
\]
Consequently
$|b_sw|\le C|k'|\eta^2(1+y)/(\eta^2+e^{-y})\le Ce_d$ and
$b_s^2\le Ce_d^2e^{T_d}/P_*^2$.
Choose $M_d$ so $Ce_d$ is small, then $P_*>e^{T_d}$ sufficiently large.
Both required strict margins follow, uniformly in $\eta$ and this whole stage.

\subsection{Intermediate power law and the complete pulse}
Write $M/X=\bar k\eta$ during axial decrease. Then $\bar k'=k-\bar k$.
Because $0\le k\le4$ and its final eleven units are zero,
$L\bar k<1/2$ at the next transition. When $U=0$ this gives $W\ge1/2$.
Write $c_\eta=D\eta J_0'\asymp\eta^2$. The angular source is
$(-l)W-h+c_\eta$. During the unit transition $0\ge l\ge-\lambda$,
the prior positive lower bound persists; $w$ has a finite bound depending
only on earlier choices. Thus $\sqrt\lambda|w|$ is small. The cone is
admissible as soon as $l<0$, and relaxed at its initial endpoint.

On the constant-slope interval, $\bar k(y)=\bar k(0)e^{-y}$ and
\[
Q_s'+(1-\lambda)Q_s=\lambda-h+c_\eta-\lambda L\bar k(0)e^{-y}.
\]
The right side is at least $\lambda/4+c_\eta$ for small $h/\lambda$.
Its explicit convolution therefore gives
\[
Q_s\ge c_{\rm pre}(\eta^2+\lambda+e^{-(1-\lambda)y}),\qquad
Q_s-\frac{\lambda-h+c_\eta}{1-\lambda}
=O_{C^1}(C_{\rm pre}(1+y)e^{-(1-\lambda)y}).
\]
There is no geometric $N_s$ source when $U=0$; the pressure source is
$O(E^2|\eta|)$. Its convolution gives
$|N_s/E|\le C_{\rm pre}|\eta|(1+y)e^{-(1/2-\lambda)y}$.
Using $|\eta|/(\eta^2+e^{-(1-\lambda)y})
\le\tfrac12e^{(1-\lambda)y/2}$ proves
\[
\sqrt\lambda|w|\le C_{\rm pre}\sqrt\lambda(1+y)e^{\lambda y/2}=o(1)
\quad(0\le y\le60\log(1/\lambda)).
\]
Since $a=2+2\lambda$, $b_s=0$, this supplies the strict cone bracket.

During the pulse set $m=(M/X)/E$ and $\beta=1/2-\lambda$.
The exact identity $m'+\beta m=R_b$ follows by differentiating $M/X$ and $E$.
With the pulse extended flatly to the left, write its convolution as
$\int_0^\infty e^{-\beta v}R_b(\lambda(y-v))\,dv$ plus
$m(0)e^{-\beta y}$. Taylor's formula under this integral gives
\[
m=R_b/\beta-\lambda\partial_{\xi_b}R_b/\beta^2
+O(\lambda^2)+m(0)e^{-\beta y},
\]
because the integral of $v^2e^{-\beta v}$ is finite uniformly for
$\beta>.4$. The end bumps have exponentially small fixed $\xi_b$ derivatives.
The same estimate holds after one $\eta$ derivative.

Since $E\le C_{\rm pre}\lambda^{30}$, all terms from
$W=1-2D\eta Em-d(Em)_\eta$ and from $dU$ in the angular source have
size $O(C_{\rm pre}E)$, including one parameter derivative. The initial
constant-slope estimate and exponential convolution give
\[
Q_s=\frac{\lambda-h+c_\eta}{1-\lambda}+O(C_{\rm pre}\lambda^{28}).
\]
The backward pressure and its first derivative are $O(E^2)$.
For $S$, integration over $0\le\lambda y\le13$ bounds
$S/(XE^2)$ by $C_{\rm pre}e^{26}/\lambda$ in $C^1$.
The numerator accumulates at most a fixed integral in $\xi_b$ divided by
$\lambda$ and the denominator loses at most $e^{-26}$; this explains the
constant $e^{26}/\lambda$, not $e^{26/\lambda}$.
Substituting $M/X=Em$ in the exact moment formula for $N_s$ gives
\[
N_s/E=-R_b+(D+c_\eta)m-D\eta m_\eta+O(C_{\rm pre}\lambda^{27}).
\]
For example $D(M-\eta M_\eta)/(XE)
=D(m-\eta m_\eta+\eta J_0'm)$, giving exactly $(D+c_\eta)m$.
The $S$ and pressure remainders are at most $CEe^{26}/\lambda$ in this
quotient, which is smaller than the stated bound for sufficiently small
$\lambda$. The difference $W-1$ similarly costs $O(E)$.

The amplitude derivative and the initial-moment bounds show
$|m_\eta|\le C_{\rm pre}\lambda(1+\log(1/\lambda))$.
Since $|\eta|/(\lambda+c_\eta)\le C/\sqrt\lambda$,
the quotient of $\eta m_\eta$ by $Q_s$ is $o(1)$.
The preceding formulas therefore imply, uniformly throughout the pulse,
\[
w=2R_b-C_d\partial_{\xi_b}R_b+o(1),\qquad
C_d=\frac{\lambda(D+c_\eta)(1-\lambda)}
{\beta^2(\lambda-h+c_\eta)}\le2+o(1).
\]
The coefficient of $R_b$ before this last replacement is exactly
$(1-\lambda)/\beta$, since $D+c_\eta-\beta=\lambda-h+c_\eta$.
The quotient defining $C_d$ is maximal at $c_\eta=0$, because its derivative
in $c_\eta$ has sign $\lambda-h-D<0$; $h/\lambda\to0$ gives the bound.
Also $b_s=-(1+2\lambda)R_b+2\lambda\partial_{\xi_b}R_b=-R_b+O(\lambda)$.
Apart from exponentially small end bumps, $R_b\ge0$ and
$\partial_{\xi_b}R_b\le1.2$, since the declining cutoff has a nonpositive
contribution. The exact upper bounds for the two leading quadratics are
\[
\sup_{R\ge0}(-2R^2+2.41R)=2.41^2/8<.74,
\]
\[
\sup_{R\ge0}(-3.5R^2+4.82R)=4.82^2/14<1.68.
\]
The uniformly $o(1)$ remainders can be made smaller than the strict numerical
gaps. Hence $a-b_sw>1$ and
$2b_sw+b_s^2/a+(a-2)w^2<2$ on the complete pulse, including its descending
and correction parts. Here $v_s\ge a=2+2\lambda>2$.

\subsection{Interpolation and exterior transitions}
After the pulse $U=M=J=0$, so $W=1$.
The angular source during interpolation is $-l-h+\vartheta_fc_\eta$,
giving $Q_s\ge c\lambda$. The tiny angular correction preserves this bound.
Since the complete amplitude solve gives $S(\infty)=0$,
\[
S(X)=\tfrac12\int_X^\infty E(X')^2dX'.
\]
Before exterior release, $l\le-\lambda/2$ and the later release has the
$h$ uniform integral bound proved above. Therefore
$(|S|+|S_\eta|)/X\le CE^2/\lambda$ and
$|\Pi|+|\Pi_\eta|\le CE^2$. The exact moment formula yields
$|w|\le CE/\lambda^2\ll1$, because pulse end reduces $E$ exponentially
in $1/\lambda$. Here $b_s=0$ and $2<a\le2+2\lambda+.2+o(1)$,
so both cone margins remain strict.

On the parameter-independent release and tail, the same backward integrals
give the exact expression
\[
N_s=2\eta\int_y^\infty(he^{v-y}-A)E(v)^2dv.
\]
Since $l\le-3h/4$,
$E(v)^2\le E(y)^2e^{-(1+3h/2)(v-y)}$.
Thus the $he^{v-y}$ integral is at most $\tfrac23E(y)^2$ and the
other at most $E(y)^2/(1+3h/2)$. In particular $|N_s|\le CE^2$
with a constant independent of small $h$.
Before the steep hold, $Q_s\ge c\lambda$ and the amplitude is already
small relative to $\lambda$. During the hold it remains above this bound.
After the hold, $E\le Ce_{\rm rel}h^6$, while $Q_s\ge ch$ up to terminal
coordinate $y=1/2$. The latter follows from its explicit positive terminal
integral and its intervening monotone decay to $Q_p\asymp h$.
Consequently $|w|\le CE/Q_s\ll1$, even if $h\ll e_{\rm rel}$.
Here $b_s=0$, $2<a\le4$, so the sufficient cone test has uniform strict slack.

All these intervals are compact once the sequential parameters are fixed.
Under the exact moment scaling $X=X_Rx$, every $X_R$ factor cancels in
$Q_s,N_s$, while $p_{s,1}=X_RxQ_s/L$ grows linearly with $X_R$.
Thus one finite increase of $X_R$ meets the threshold proved above on the
whole designated interval. This completes the local reconstruction of
[NS, Proposition A.4], including each part of its cone verification.

\section{The heat factor, its exact equation, and the compensated replacement}
\subsection{A positive solution of the radial swirl heat equation}
Define for $Z\ge0$,
\[
\hh(Z)=\frac1{\Gamma(1+h)}\int_0^\infty e^{-v}v^h(1+Zv)^{-h}\,dv.
\]
Dominated differentiation gives for every nonnegative integer $m$
\[
\hh^{(m)}(Z)=\frac{(-1)^m(h)_m}{\Gamma(1+h)}
\int_0^\infty e^{-v}v^{h+m}(1+Zv)^{-h-m}\,dv,
\quad \hh^{(m)}(0)=(-1)^m(h)_m(1+h)_m.
\]
For fixed $m$ the integrand after deleting the last factor is integrable
uniformly for $Z\ge0$. Hence $\hh$ is positive, smooth up to zero from the
right, and has bounded fixed derivatives on every bounded nonnegative interval
(in fact on the entire half-line for absolute derivatives).
For $0<h\le h_0<1/2$ and fixed $m\ge1$, the last factorial product is at
most $C_mh$. Taylor's formula therefore proves
\[
\hh(Z)=1-h(1+h)Z+O_h(Z^2),\qquad
|\hh(Z)-1|+|Z\hh'(Z)|\le C hZ
\]
on a fixed bounded nonnegative interval.

The integral solves the required ODE by an exact polynomial calculation.
Multiply
$Z^2\hh''+(1+2(1+h)Z)\hh'+h(1+h)\hh$
by $\Gamma(1+h)/h$, and use common integrand
$e^{-v}v^h(1+Zv)^{-h-2}$. Its remaining polynomial is
\[
(1+h)Z^2v^2-(1+2(1+h)Z)v(1+Zv)+(1+h)(1+Zv)^2
=(1+h)-v-Zv^2.
\]
This is precisely the polynomial from
$\partial_v[e^{-v}v^{1+h}(1+Zv)^{-1-h}]$ in the same common factor.
The boundary values at zero and infinity are both zero. Thus
\begin{equation}\label{eq:heatode}
Z^2\hh''+(1+2(1+h)Z)\hh'+h(1+h)\hh=0.
\end{equation}
Set $K(r,t)=c_\infty s^{-A}\hh(2\tau/s)$, $s=r^2/2$, $Z=2\tau/s$.
Direct differentiation gives
\[
K_t=-2c_\infty s^{-A-1}\hh',
\]
\[
(\partial_{rr}+r^{-1}\partial_r-r^{-2})K
=2c_\infty s^{-A-1}
\{Z^2\hh''+(2A+1)Z\hh'+(A^2-1/4)\hh\}.
\]
Since $2A+1=2(1+h)$ and $A^2-1/4=h(1+h)$,
\eqref{eq:heatode} proves the heat equation, with its swirl term $-r^{-2}K$.

Normalize the positive integrand defining $\hh$ to a probability density.
Then $-Z\hh'/\hh$ is its average of $hZv/(1+Zv)$, so
$0\le-Z\hh'/\hh<h$. It follows that
$rK_r/(2K)=-A-Z\hh'/\hh<-1/2$, in particular $K_r<0$.
At $Z=0$ this is still strict because $h>0$.
Finally $Z=2(1-\eta^2)/X$ gives only factors $-4\eta/X$ and $-4/X$
when differentiated in $\eta$, never a denominator $1-\eta^2$.
Consequently all fixed mixed $X,\eta$ derivatives have one-sided limits at
$\eta=\pm1$, and for each fixed logarithmic radial and parameter derivative,
\[
\frac{E_{\rm heat}}{E_{\rm pow}}
=\hh(2d/X)=1-\frac{2h(1+h)d}{X}+O_{h,j,m}(X^{-2}).
\]
This proves [NS, Lemma A.6] on its stated nonnegative domain. No analytic
continuation through negative $Z$ is used or required.

\subsection{The three heat discrepancies and their exact cancellation}
Let $y=\log(X/X_{\rm tail})$, $X_K=e^{.2}X_{\rm tail}$,
$X_b=e^3X_{\rm tail}$, and choose $\chi_K=0$ for $y\le.2$, $\chi_K=1$
for $y\ge.5$. Replace the reference swirl $E_{\rm cl}$ by
\[
E_{\rm cl}[1+\chi_K(y)(\hh(2d/X)-1)].
\]
For $x=X/X_K\ge1$ and its reference amplitude $e_K$, the preceding
derivative bounds give exactly
\[
|(X\partial_X)^j\partial_\eta^m(E-E_{\rm cl})|
\le C_{j,m}e_KX_K^{-1}x^{-A-1}.
\]
The profile is comparable to $e_Kx^{-A}$ with constants depending on the
fixed schedule. Thus the pressure, $S$, and subtracted angular discrepancies
obey, after $m$ parameter derivatives,
\[
|\Delta C_p|_m\le C_m\frac{e_K^2}{X_K}\int_1^\infty x^{-2A-2}dx,
\quad |\Delta S|_m\le C_me_K^2\int_1^\infty x^{-2A-1}dx,
\]
\[
\left|\Delta\int(H-H_{\rm pow})dX\right|_m
\le C_me_KX_K^{1/2}\int_1^\infty x^{-A-1/2}dx
=C_m e_KX_K^{1/2}/h.
\]
The powers arise respectively from $\delta(E^2)dX/(2X)$,
$-\delta(E^2)dX/2$, and $\sqrt{2X}\delta E,dX$.
The final $1/h$ is retained: $h$ has already been fixed before $X_R$ grows.

At the earlier second reserved patch write $x_*=X/X_*$,
$E=e_*f x_*^{-1/2-\lambda}$, $U=0$, and add
$\delta E=e_*\sum_{i=1}^3 c_i(\eta)\beta_i(x_*)$.
Divide the three rows by $e_*^2$, $X_*e_*^2$, and $X_*^{3/2}e_*$.
The exact linear weights are then
\[
f x_*^{-3/2-\lambda},\qquad -f x_*^{-1/2-\lambda},\qquad
\sqrt2 x_*^{1/2}.
\]
The first two rows have exact quadratic terms and the third is linear.
The moment determinant is nonzero by Section 3; its inverse is uniform in
$\eta$. Since $X_*/X_K$ and $e_*/e_K$ are fixed, all three normalized
discrepancies are $O_m(X_K^{-1})$. The contraction gives the exact correction
with $\|c\|_{C^m}\le C_mX_K^{-1}$ for any prescribed finite $m$, after
choosing $X_R$ large. More generally its higher fixed derivative bounds follow
by implicit differentiation with the same small forcing: the Jacobian stays
invertible and every differentiated equation has only lower derivative products
and an $O_m(X_K^{-1})$ inhomogeneity. This justifies the stated $O_m$ estimates
without selecting infinitely many smallness thresholds at once.

Both edits have $U=0$, so $M,J$ are unchanged exactly. All three other total
moments are restored by this solve. The total pressure integral is therefore
unchanged, so the analytic axis pressure and every profile before the patch
remain unchanged. On the compact logarithmic interval up to $y=.5$,
the exact moment formulas transfer the $O(X_K^{-1})$ changes of profiles and
moments to $Q_s,N_s$, costing one more $\eta$ derivative.
Indeed substituting $I=X_R^{3/2}\widehat I$, $J=X_R^{3/2}\widehat J$,
$M=X_R\widehat M$, $S=X_R\widehat S$, $H=X_R^{1/2}\widehat H$
cancels every $X_R$ in those formulas. Positive lower bounds for $E,Q_s$
and the already strict cone margins persist by continuity, uniformly on this
fixed compact interval. This proves [NS, Proposition A.7] and its preservation
claims; the remaining endpoint is treated next rather than being inferred
from a compact-interior estimate.

\section{Exterior stress, exact flat factors, and the order-one edge estimate}
\subsection{Convergent moment subtraction and backward stress formulas}
The leading regular axis and exact matching of the five moments are inputs
from the profiles lane at this stage, precisely as in [NS, Lemma A.8].
Once supplied, $M(\infty)=J(\infty)=S(\infty)=0$ and the compensated
subtracted angular moment gives
\[
I(X)=\frac{XH_{\rm pow}(X)}{1-h}-\int_X^\infty(H-H_{\rm pow})dx,
\qquad S(X)=\tfrac12\int_X^\infty E^2dx.
\]
The exact pressure is the backward integral already established. These
identities, substituted into the moment formulas, determine the same forward
stress everywhere in the tail.

To prove the backward physical stress representation rather than assume its
boundary constants vanish, put $K_{\rm pow}=c_\infty(r^2/2)^{-A}$.
For $r\to\infty$, uniformly on bounded $\tau\ge0$ intervals, the heat
expansion gives
\[
K-K_{\rm pow}=O_h(\tau r^{-3-2h}),\quad K_t=O_h(r^{-3-2h}),
\quad r^2K_r-rK=O_h(r^{-2h}).
\]
Thus the weighted difference and its time derivative are dominated at infinity
by $Cr^{-1-2h}$, whose integral from $R_*$ is $R_*^{-2h}/(2h)$.
At zero the subtracted pure power has weight $r^{1-2h}$, integrable for
$h<1$, while the regular field is integrable by axis smoothness.
The substitution $r=\sqrt qR$, $dX=R\,dR$ proves exactly
\[
\int_0^\infty r^2(u_\theta-K_{\rm pow})dr
=q^{3/2-A}\int_0^\infty(H-H_{\rm pow})dX=0.
\]
Time differentiation is justified by the displayed integrable majorant on
compact physical parameter neighborhoods, proving \eqref{eq:submoment} and
the zero integrated time term.
The angular transport flux has zero axial integral because it is
$q^{3/2-2A}J(\infty)=0$; its radial boundary vanishes by regularity at
zero and $U=M=V_0=0$ in the exterior. The radial viscous boundary is zero
by the $r^{-2h}$ estimate above.

For the axial component $\int ru_zdr=q^{1-A}M(\infty)=0$.
Fubini with the absolutely convergent positive pressure integrand yields
\[
\int_0^\infty rp(r)dr
=-\int_0^\infty\frac{u_\theta(r')^2}{r'}
\left(\int_0^{r'}r\,dr\right)dr'
=-\tfrac12\int_0^\infty r'u_\theta(r')^2dr'.
\]
The pressure has decay $O(r^{-2-4h})$, hence $r^2p\to0$, and its contribution
to the integrated axial flux is exactly $q^{1-2A}S(\infty)=0$.
All axial radial transport and viscosity boundaries vanish as well.
Therefore both leading residuals have zero weighted totals, so the stresses
defined with a minus sign by integration from the axis are exactly
\begin{equation}\label{eq:backward}
T_\theta(r)=r^{-2}\int_r^\infty r'^2\res_\theta^{(0)}(r')dr',\qquad
T_z(r)=r^{-1}\int_r^\infty r'\res_z^{(0)}(r')dr'.
\end{equation}
For $X\ge X_b$ the field is $(u_r,u_\theta,u_z)=(0,K,0)$ and pressure
is its $z$ independent backward radial integral. The heat equation makes both
leading residuals zero there, proving zero exterior stress.

\subsection{An exact flat integral map}
For $c>0$, integer $j\ge0$, and a smooth $b(u,\eta)$ on a closed compact
parameter rectangle, substitute $u=\delta(1+\delta^2v)^{-1/2}$.
Then $du=-\tfrac12\delta^3(1+\delta^2v)^{-3/2}dv$ and
$c/u^2=c/\delta^2+cv$. Hence
\begin{equation}\label{eq:flat}
\int_0^\delta e^{-c/u^2}u^{-j}b(u,\eta)du
=\tfrac12e^{-c/\delta^2}\delta^{3-j}B(\delta,\eta),
\end{equation}
where the coefficient is explicitly
\[
B(\delta,\eta)=\int_0^\infty e^{-cv}(1+\delta^2v)^{(j-3)/2}
b(\delta/\sqrt{1+\delta^2v},\eta)dv.
\]
Every fixed derivative of the integrand is bounded by $e^{-cv}$ times a
polynomial in $v$, uniformly in the compact parameters and $0\le\delta\le\delta_0$.
This follows directly from differentiating $\delta(1+\delta^2v)^{-1/2}$;
its denominators are at least one and all its numerators are polynomials in
$\delta,v$. Dominated differentiation proves smoothness of $B$ with all
fixed derivatives bounded, and $B(0,\eta)=b(0,\eta)/c$.
This proves the exact map in [NS, Lemma A.9], including the factor $1/2$.

\subsection{Positive angular stress and the limiting direction}
For terminal $y\ge.5$, the heat cutoff is one, so $u_\theta=Kf$, $f=f_o(y)$,
and $u_r=u_z=0$. At fixed $z,t$, $y_r=2/r$, while
$y_t=(qL)^{-1}$, $y_z=-2\eta/(q^DL)$. The leading angular residual,
which omits axial viscosity, is exactly
\[
\res_\theta^{(0)}=Kf'/(qL)-K(f_{rr}+r^{-1}f_r)-2K_rf_r.
\]
Differentiate the physical pressure $p=-\int_r^\infty K(r')^2f(r')^2dr'/r'$.
Since $K_z=0$ and $dr'/r'=dy'/2$, this gives the sign and factor
\[
p_z(r)=\frac{2\eta}{q^DL}\int_y^3K(y')^2f(y')f'(y')dy'.
\]
Inserting the angular residual into \eqref{eq:backward} and integrating its
viscous terms by parts gives
\begin{equation}\label{eq:positive}
T_\theta=Kf_r+\frac1{r^2}\int_r^\infty(r'K-r'^2K_{r'})f_{r'}dr'
+\frac1{qLr^2}\int_r^\infty r'^2Kf'\,dr'.
\end{equation}
To check the integration sign, the term $-\int r'^2Kf_{r'r'}dr'$ contributes
\[
r^2Kf_r+\int(2r'K+r'^2K_{r'})f_{r'}dr'.
\]
The two remaining
viscous terms leave exactly $r'K-r'^2K_{r'}$ in this integral.
Every term in \eqref{eq:positive} is nonnegative since $K>0$, $K_r<0$,
$f'\ge0$. The remaining interval has logarithmic length at most $2.5$;
radii and heat factors there have uniformly bounded ratios. Since
$\int_y^3f'\,dv=\rho_o\psi_o(y)$, its last term implies
\[
T_\theta\ge c\frac{rK}{qL}\rho_o\psi_o(y)>0\quad(y<3).
\]
The pressure formula and its next radial integral similarly give
\[
|p_z|\le Cq^{-D}K^2\rho_o\psi_o,
\quad |T_z|\le Crq^{-D}K^2\rho_o\psi_o,
\quad |T_z/T_\theta|\le Cq^{1-D}K=C E_{\rm pow}\hh(2d/X).
\]
The final equality uses $1-D=A$. The $h^6$ amplitude reduction makes this
ratio small. Furthermore, since $\partial_yZ=-Z$,
\[
a=2+2h+2Z\hh'/\hh-2f'/f.
\]
For sufficiently large $X_R$, $-Z\hh'/\hh<h/4$ uniformly in $\eta$;
also $f'/f<h/4$ by construction. Hence $2+h<a\le2+2h$ and $b_s=0$.
The exact stress identities are
$P_c-a=T_{0,\theta}/F>0$, $J_c=T_{0,z}/F$ with $F=E/\sqrt{2X}>0$.
Thus the cone is exactly $(a-2)(T_z/T_\theta)^2<2$, with uniform strict
slack throughout the terminal interval. This also covers the plateau $.5\le y\le1$:
the positive backward time integral persists there even though $f'(y)=0$.

Near $\delta=3-y=0$, the specified step has the exact factorization
\[
\psi_o=e^{-4/\delta^2}g(\delta),\qquad
g(\delta)=\frac1{e^{-(1-\delta/2)^{-2}}+e^{-4/\delta^2}},\qquad g(0)=e,
\]
\[
f'=\rho_oe^{-4/\delta^2}\delta^{-3}[8g(\delta)+\delta^3g'(\delta)].
\]
The normalized boundary term in \eqref{eq:positive} is
\[
q^{A+1/2}Kf_r=\frac{2E_{\rm pow}(X)\hh(2d/X)}{\sqrt{2X}}f'.
\]
The two remaining angular integrals contain $f'$, so \eqref{eq:flat} with
$c=4,j=3$ makes each $e^{-4/\delta^2}$ times a smooth coefficient.
Relative to the boundary factor $e^{-4/\delta^2}\delta^{-3}$ they vanish
to at least order $\delta^3$. Therefore
\[
T_{0,\theta}=e^{-4/\delta^2}\delta^{-3}b_\theta(\delta,\eta),\qquad
b_\theta(0,\eta)=\frac{16\rho_oE_{\rm pow}(X_b)}{\sqrt{2X_b}}
\hh(2d/X_b)g(0)>0.
\]
The heat factor has a positive minimum on $[0,2/X_b]$, so $b_\theta$
has a positive lower bound on a sufficiently short closed collar uniformly
for $|\eta|\le1$.
The pressure integral contains one $f'$ factor; the same substitution gives
$p_z=q^{-D-2A}e^{-4/\delta^2}$ times a smooth coefficient.
Its further radial integration for $T_z$ uses \eqref{eq:flat} with $j=0$,
giving
\[
T_{0,z}=e^{-4/\delta^2}\delta^3b_z(\delta,\eta).
\]
The powers of $q$ cancel exactly:
$q^{A+1/2}q^{1/2}q^{-D-2A}=q^{1-D-A}=1$.
The remaining coefficients depend only on $X,\eta,\hh(2d/X),L^{-1}$,
with their already proved closed-parameter smoothness. Consequently
\[
T_{0,z}/T_{0,\theta}=\delta^6b_z/b_\theta\longrightarrow0
\]
smoothly, and the stress direction extends to $(1,0)$ at the outer edge.
Every fixed profile derivative costs only a finite inverse power of $\delta$,
and $|T_0|\ge ce^{-4/\delta^2}\delta^{-3}$ there.
This proves (A.48)--(A.51) including the precise leading coefficient and
the independence from the physical scale $q$.

\subsection{The remaining positive-order edge term}
At order $n=1$, on the final collar the only exterior source is
$-\partial_{zz}(Kf_o)$ in the angular equation. With
$a_z=-2\eta/L$ and
$b_z^{\rm coord}=(-2D\eta a_z+d\partial_\eta a_z)/L$,
direct differentiation of $y_z=q^{-D}a_z$ gives
\[
\partial_{zz}(Kf_o)=Kq^{-2D}
\{a_z^2\partial_{yy}f_o+b_z^{\rm coord}\partial_yf_o\}.
\]
All parameter derivatives of the coefficients are bounded because $L>0$.
The two profile terms have factors
$e^{-4/\delta^2}\delta^{-6}$ and $e^{-4/\delta^2}\delta^{-3}$ times
smooth functions. The backward stress integral gains three powers by
\eqref{eq:flat}, producing
$|T_1|\le Ce^{-4/\delta^2}\delta^{-3}$ and, for every fixed derivative,
$|\partial^IT_1|\le C_Ie^{-4/\delta^2}\delta^{-N_I}$.
The physical exponent after radial integration is
$q^{-A-2D+1/2}=q^{-A-1/2+2h}$, exactly the order-one stress scale.
Near this outer edge the other factor in $\zeta$ is a positive smooth bounded
factor, so this is the weighted estimate required in (5.13) and (5.43).
Near the inner edge $T_1$ is identically zero; for $n\ge2$ the common compact
interior support makes every fixed derivative of $T_n/\zeta$ bounded.

\section{Audit conclusion and exact dependency boundary}
The calculations above verify the coordinate operators; every coefficient and
sign in the order recursion; the sparse analytic Volterra mechanism; the
five positive-order moment equations and common supports; conservative residual
cancellation including the subtracted leading angular moment; divergence under
the cutoff curl; the order-independent derivative losses and flat Borel
residual; all outer schedule moment corrections and amplitude selection; the
pressure datum; each outer cone estimate; the heat ODE and endpoint regularity;
the compensation map; and the two outer stress factors.

No mathematical contradiction was found in this assigned subsystem after
these derivations. Two extraction ambiguities required visual resolution:
the pulse estimate is $C_{\rm pre}e^{26}/\lambda$, and the endpoint angular
coefficient has denominator $\sqrt{2X_b}$. Neither is a correction to the source
PDF. The proof of simultaneous compact-range Borel bounds is made explicit
above using the common correction supports; it fills in the domain choice
without changing the statement or the construction.

This result has an exact dependency boundary. The common analytic axis/collar
data used in the positive-order recursion and the final regular-axis matching
used in the backward stress identity are the outputs of [NS, Appendix B and
Theorem 4.6]. Their independent proof is outside this document, as is
the oscillatory realization of the admissible stress and the final exact
Navier--Stokes construction. No acceptance of those parts, or of a claimed
global blowup theorem, follows from this bounded audit.


% --- included proof body: oscillations_body.tex ---
\section{Auxiliary geometry and coefficient algebra: independent derivation}\label{auxiliary-geometry-and-coefficient-algebra-independent-derivation}

Source: supplied \emph{Finite Time Blowup for Navier-Stokes}, Section 6, pp.~62-72. Source identity and audit boundary are recorded in \texttt{LOGBOOK.md} and \texttt{READ\_SCOPE.md}. The formulas below use the source's original cylindrical coordinate orientation, chart scales, integer covering, and support domains. The geometric and algebraic assertions are proved here directly.

\subsection{Coordinate maps and exact derivative intertwining}\label{coordinate-maps-and-exact-derivative-intertwining}

Let \(A=1/2+h\), \(D=1/2-h\), \(0<h<1/100\), \(Q=2^{-\ell}\), \(\varepsilon=Q^h\), \(S_*=\ell^2\), and \[
R=r/Q^{1/2},\quad Z=z/Q^D,\quad T=(1-t)/Q.
\] The profile radius is \(R_{\rm profile}=r/q^{1/2}\), so the exact transition is \(R=\sqrt{q/Q}\,R_{\rm profile}\). No equality of these different radius coordinates is used. The original similarity variables obey \(z=q^D\eta\), \(1-t=q(1-\eta^2)\), with \(L=1-2h\eta^2>0\). On a dyadic shell \(q/Q\in[1/2,2]\), \(q/Q\) and \(\eta\) are smooth functions of \(Z,T\), including one-sided derivatives at \(T=0\). Indeed \(T=s-Z^2s^{2h}\), \(s=q/Q\), has derivative \(1-2hZ^2s^{2h-1}=L\ge1-2h\). Repeated implicit differentiation bounds each fixed chart derivative on the compact shell. Thus the ordinary chart and profile derivatives used below are related by finite sums with uniformly bounded smooth coefficients.

Set \[
J_g=\begin{pmatrix}3&1\\1&5\end{pmatrix},\quad b_g=\sqrt2-1,
\quad v_r=(1,-b_g),\quad v_t=(b_g,1),
\] \[
\Lambda_g=4-\sqrt2,\quad T_g=4+\sqrt2,\quad
\rho_g=\frac{\log\Lambda_g}{\log T_g},\quad \kappa_s=10^{-5},
\quad d_r=2((1+h)\rho_g-h\kappa_s).
\] Direct multiplication gives \(J_gv_r=\Lambda_gv_r\), \(J_gv_t=T_gv_t\), \(v_r\cdot v_t=0\), and \(\det J_g=14\). Also \(\rho_g>0\) and \(\rho_g>\kappa_s\), hence \(d_r>0\). For example \(\Lambda_g>2\), \(T_g<6\) imply \(\rho_g>\log2/\log6>1/3>\kappa_s\).

On the extended domain with independent \(Y\in\mathbb R^2/\mathbb Z^2\), define \[
\iota(r,\theta,z,t)=(r,\theta,z,t,v_rr^{d_r}+v_tt\bmod\mathbb Z^2).
\] For every smooth extended scalar \(F\), the ordinary chain rule proves the exact identities \[
\partial_t(\iota^*F)=\iota^*((\partial_t+v_t\cdot\partial_Y)F),\qquad
\partial_r(\iota^*F)=\iota^*((\partial_r+d_rr^{d_r-1}v_r\cdot\partial_Y)F).
\] The \(z,\theta\) identities have no extra term. These are morphisms of differential algebras: pullback preserves sums and products and intertwines the four stated derivations. Their commutators vanish because the only variable coefficient depends on \(r\), the extra directions are constant and independent of \(Y,t,z,\theta\), and the radial derivation is the only derivation differentiating that coefficient. The fields depending on \(Y\) vanish in a neighborhood of the axis at each \(q>0\), so no differentiability of \(r^{d_r}\) at \(r=0\) is invoked.

For \(i=\lfloor\log_{T_g}(Q^{-1-h}/S_*)\rfloor\ge0\), let \(\pi_i(Y)=J_g^iY\bmod\mathbb Z^2\). This is a surjective group homomorphism of tori with finite kernel of cardinality \(14^i\). On a pullback \(f\circ\pi_i\), the two directional operators are exactly \(T_g^iN_if\) and \(\Lambda_g^iL_if\), where \(N_i=v_t\cdot\partial_{Y_i}\), \(L_i=v_r\cdot\partial_{Y_i}\). Thus \[
\mathsf t_*=Q^{1+h}\mathsf t=-\varepsilon\partial_T+c_iN_i,
\quad D_r=\sqrt Q\,\mathsf r=\partial_R+M_id_rR^{d_r-1}L_i,
\] \[
c_i=T_g^iQ^{1+h},\quad M_i=\Lambda_g^iQ^{d_r/2},\quad
D_z=\sqrt Q\partial_z=\varepsilon\partial_Z,
\quad D_\theta=R^{-1}\partial_\theta.
\] The floor inequality is strict on its left: \[
T_g^{-1}Q^{-1-h}/S_*<T_g^i\le Q^{-1-h}/S_*.
\] It follows that \(T_g^{-1}S_*^{-1}<c_i\le S_*^{-1}\), and, using \(\Lambda_g^i=(T_g^i)^{\rho_g}\), \[
T_g^{-\rho_g}\varepsilon^{-\kappa_s}S_*^{-\rho_g}
<M_i\le\varepsilon^{-\kappa_s}S_*^{-\rho_g}.
\] Every radial derivative of \(d_rR^{d_r-1}\) is bounded on the fixed positive radial interval. This proves the claimed additional radial cost \(\varepsilon^{-\kappa_s}\), with its retained power of \(S_*\).

For physical velocity \(Q^{-A}u_*\), pressure \(Q^{-2A}p_*\), and residual multiplier \(Q^{2A+1/2}\), the time, spatial first-derivative, and Laplacian coefficients are respectively \[
Q^{2A+1/2-A}=Q^{1+h},\qquad
Q^{1/2},\qquad
Q^{2A+1/2-A-1}=Q^h=\varepsilon.
\] These identities retain the source's viscosity-one chart convention; scaling to arbitrary viscosity is outside this lane.

\subsection{Fourier inversion and descent}\label{fourier-inversion-and-descent}

For \(n=(n_1,n_2)\ne0\), \[
v_r\cdot n=(n_1+n_2)-\sqrt2n_2.
\] Its product with the algebraic conjugate is the nonzero integer \((n_1+n_2)^2-2n_2^2\). It cannot be zero unless \(n_2=0=n_1\), by irrationality of \(\sqrt2\). The conjugate has absolute value at most \(C|n|\), so \(|v_r\cdot n|\ge1/(C|n|)\). For \(v_t\cdot n=(n_2-n_1)+\sqrt2n_1\), the same argument uses \((n_2-n_1)^2-2n_1^2\). Therefore both directional inverses on zero-mean smooth torus functions have Fourier multipliers bounded by \(C(1+|n|)\), retaining the Fourier derivative factor \(2\pi i\): the inverse of \(v\cdot\partial_Y\) multiplies the \(n\)-coefficient by \((2\pi i\,v\cdot n)^{-1}\). Rapid Fourier decay is preserved. In Sobolev norms this gives \(\|L^{-1}f\|_{H^s}\le C_s\|f\|_{H^{s+1}}\), with the zero coefficient defined to be zero. Taking a further fixed Sobolev embedding loss gives all fixed classical derivatives. No infinite derivative loss occurs.

If \(f\) descends through \(\pi_i\), its Fourier support lies in \((J_g^i)^T\mathbb Z^2\). Translation multiplies each coefficient by a character, averaging deletes nonzero coefficients, and a directional Fourier inverse multiplies them by scalar values. Each operation preserves this sublattice. Radial integration at fixed \(z,t\) preserves deck invariance pointwise in the integral. These are the precise maps proving the descent claims.

\subsection{Separation of all interacting labels}\label{separation-of-all-interacting-labels}

Use the dyadic squared partition \(\sum_\ell\chi_\ell(q)^2=1\), the product grid squared partitions of mesh \(S_*^{-3}\), and the duplicated signs in the source's labels \(\gamma=(\ell,a,\sigma)\). On the active annulus select every grid support meeting the active shell and a representative in that intersection. Then \(\eta_\gamma=\chi_\ell\chi_{\ell,a}\) gives \(\sum_{\beta=(\ell,a)}\eta_\beta^2=1\), counting each box once and the two signs separately only when assigning rectangles.

Overlapping dyadic supports imply \(|\ell-\ell'|\le2\). Join labels whenever the fixed-factor enlarged physical grid boxes intersect and the band difference is at most four, and join the two signs of each box. A chart scale ratio for bands differing by four is bounded above and below; \(\ell^2/(\ell')^2\) has the same property for \(\ell,\ell'\ge\ell_0>4\). A grid box thus meets a uniformly bounded number of boxes in each of finitely many bands. The countable graph has maximum degree \(d<\infty\). Enumerating its vertices and assigning the first unused color in \(\{1,\ldots,d+1\}\) gives a proper coloring.

Writing \(f(\ell)=((1+h)\ell\log2-2\log\ell)/\log T_g\), the mean-value bound and the floor operation give \[
|i(\ell)-i(\ell')|\le1+\frac{4(1+h)\log2+8/\ell_0}{\log T_g}.
\] Fix an integer upper bound \(\Delta_{\max}\). There are only finitely many ordered color constraints \[
c_\mu\ne J_g^\Delta c_\nu\pmod {\mathbb Z^2},\qquad0\le\Delta\le\Delta_{\max},
\] excluding only the tautological constraint \((\Delta,\mu,\nu)=(0,\nu,\nu)\). For different colors, the forbidden equality is a proper closed subset with empty interior of the finite product of tori. For a self-pair with \(\Delta>0\), \(J_g^\Delta-I\) is nonsingular because its eigenvalues are \(\Lambda_g^\Delta-1,T_g^\Delta-1>0\); its kernel on the torus is finite, so this forbidden subset also has empty interior. A finite union of these closed sets has empty interior. The complement is open and nonempty, and rational tuples are dense, giving the rational centers specified in the source. There is consequently a strictly positive minimum distance \(d_c\) of every forbidden difference from the lattice.

Let \(C_J=\max_{0\le\Delta\le\Delta_{\max}}\|J_g^\Delta\|\) and \(C_v=\|v_r\|+\|v_t\|\). Choose \(r_0>0\) sufficiently small that all enlarged rectangles are injective and \(2r_0 C_v(C_J+1)<d_c\). If two absolute preimages of enlarged rectangles at levels \(i,i+\Delta\) intersect, their common point produces \[
c_\mu-J_g^\Delta c_\nu=J_g^\Delta e_\nu-e_\mu\pmod{\mathbb Z^2},
\quad |e_\nu|,|e_\mu|\le2r_0C_v.
\] The right side has distance from zero less than \(d_c\), a contradiction for joined labels. All labels whose slow supports meet were joined. Therefore their enlarged auxiliary preimages are disjoint. For smoothly extended fields supported in these products, the support of every derivative is contained in the same closed support, so all cross-label derivative products vanish before and after \(\iota^*\). This proves the required nonlinear support assertion, including boundaries.

On a band rectangle, write \(Y_i-c_\gamma-k=\xi_gv_r+\eta_gv_t\) and \(v=(\eta_g+r_0)/c_i\), \(L_s=2r_0/c_i\). Orthogonality gives \(L_i\eta_g=N_i\xi_g=0\), \(L_i\xi_g=N_i\eta_g=1\). Hence \(D_rv=D_zv=0\), \(\mathsf t_*v=1\). This is the exact reason spatial differentiation of the pulse does not incur a time-direction torus frequency.

\subsection{Common torus, lifts, integration paths, and label counts}\label{common-torus-lifts-integration-paths-and-label-counts}

At a fixed \((z_0,t_0)\), put \(q_{\mathrm{loc}}=q(z_0,t_0)>0\).
By continuity choose a neighborhood \(U\) with
\(q(U)\subset(q_{\mathrm{loc}}/2,2q_{\mathrm{loc}})\).
An active dyadic band obeys \(q/2\le Q_\ell\le2q\).
Every band meeting \(U\) therefore has
\(q_{\mathrm{loc}}/4<Q_\ell<4q_{\mathrm{loc}}\), so its dyadic
indices lie in a finite set \(\mathcal L(U)\) with span at most four.
The previously proved four-band bound on \(i(\ell)-i(\ell')\)
therefore applies to this entire set. Take \(i_0=\min_{\ell\in\mathcal L(U)}i(\ell)\), \(H=J_g^{i_0}Y\). Then \(Y_i=J_g^\Delta H\), \(0\le\Delta\le\Delta_{\max}\). The preimage of an injective rectangle consists exactly of the \(14^\Delta\) images \[
J_g^{-\Delta}(c_\gamma+k+\xi_gv_r+\eta_gv_t)\pmod {\mathbb Z^2},
\quad[k]\in\mathbb Z^2/J_g^\Delta\mathbb Z^2.
\] Different cosets give disjoint copies: an overlap would imply that their \(k\)-difference is in \(J_g^\Delta\mathbb Z^2\), after using injectivity in the band rectangle. This also proves there are no omitted copies. Derivative factors of \(J_g^\Delta\) are uniformly bounded. Surjectivity of \(\pi_{i_0}\) transfers the previously proved support separation to \(H\).

For a torus character \(e^{2\pi i n\cdot Y_i}\), pullback has frequency \((J_g^\Delta)^Tn\), zero exactly when \(n=0\). Integrating a Fourier series therefore proves exact preservation of normalized Haar measure, not an approximation depending on frequency. This proves agreement of the absolute, common, and band averages. On overlaps the representatives pull back to the identical extended field; their compatibility is equality after the respective covering maps. Those maps are independent of \(R,Z,T\), so slow derivatives preserve compatibility. The normalized fast derivatives are the same absolute operators \(Q^{1+h}N_{\rm abs}\) and \(Q^{d_r/2}L_{\rm abs}\); they also preserve it.

A mesh of spacing \(S_*^{-3}\) gives \(O(1)\) overlap at a point, \(O(S_*^9)\) boxes in a compact three-dimensional chart set, and \(O(S_*^3)\) boxes along a bounded radial line. The two signs and bounded covering multiplicities contribute fixed factors. Radial integration fixes \(q(z,t)\); it can collect radial boxes but no additional bands. The collected sum therefore introduces only a polynomial in \(S_*\).

Define \(\lambda_t(w)=v_t\cdot w/(1+b_g^2)\). In a lifted common rectangle, the exact path holding slow and transverse coordinates fixed is \[
H_w=H+T_g^{-\Delta}c_i(w-v(H))v_t.
\] Because \(\lambda_tJ_g^\Delta=T_g^\Delta\lambda_t\), \[
D_Hv=T_g^\Delta\lambda_t/c_i,\quad
D_HH_w=I-v_t\lambda_t,
\] and all higher derivatives of these affine expressions vanish. The first is \(O(S_*)\), the second uniformly bounded. No phase-averaging operation is hidden in this path: values at distinct preimages remain distinct. Deck translations commute with the path and merely reindex the rectangle copy. A uniquely prescribed initial-value solution therefore descends to the common torus.

\subsection{Weighted coefficient classes and complete product rules}\label{weighted-coefficient-classes-and-complete-product-rules}

On \(X_a<X<X_b\), let \[
\zeta=\exp\left(-\frac{a_a}{\log^2(X/X_a)}-\frac{a_b}{\log^2(X_b/X)}\right),
\quad\delta=\min\{1,\log(X/X_a),\log(X_b/X)\},
\] with fixed \(a_a,a_b>0\), and extend \(\zeta\) by zero. For any \(b>0,N\ge0\), \(\zeta^b\delta^{-N}\to0\) faster than every power at either edge: put \(x=1/\delta\to\infty\) and use \(x^Ne^{-cx^2}\to0\), proved for example by comparing \(N\log x-cx^2\) with \(-cx^2/2\) for large \(x\). Repeated derivatives of the explicit exponential have the same exponential times a finite sum of inverse edge-distance powers. Thus all stated derivative bounds imply smooth zero extension at radial boundaries.

The mean class \(\mathcal M_\alpha\) consists of angular-independent compatible coefficients with derivatives bounded by \(C_I\varepsilon^\alpha S_*^{b_I}\zeta\delta^{-d_I}\), smooth zero extension outside the active shell, and no rectangle or envelope condition. The wave class \(\mathcal W_\alpha\) uses the same derivative variables \((R,Z,T,H_1,H_2)\), retains a label and nonzero harmonic, has the prescribed closed slow/transverse/shell supports, and bounds \(C_I\varepsilon^\alpha S_*^{b_I}\sqrt\zeta\delta^{-d_I}P(v)\). Actual cutoff coefficients also have compact pulse support and smooth extension there. Equal label-harmonic contributions must be added before testing the bounds. \(P(v)\) is a pointwise weight; these definitions already require all coefficient derivatives to obey it, so no differentiation of a mere inequality for \(P\) is being assumed.

Leibniz's formula has finitely many terms at each order. The retained products of weights obey \[
\zeta^2\le\zeta,\quad\zeta^{3/2}P\le\sqrt\zeta P,
\quad\zeta P^2\le\zeta,\quad\zeta P^2\le\sqrt\zeta P
\] because \(0\le\zeta,P\le1\). These prove \(\mathcal M_\alpha\mathcal M_\beta\subset\mathcal M_{\alpha+\beta}\), \(\mathcal M_\alpha\mathcal W_\beta\subset\mathcal W_{\alpha+\beta}\), and the two same-label wave-product rules. A zero-harmonic product extends to a mean after temporal cutoff. Before cutoff it is only a mean-bounded local coefficient. Multiplication by a mean cannot enlarge the wave support; products of different actual labels are zero by separation. Finite sums retain the same exponent and suitable maxima of the polynomial degrees.

For a nonzero angular frequency \(kp_\gamma\in\mathbb Z\setminus\{0\}\), angular integration of \(e^{ik(m+m')\Phi_\gamma}\) is zero unless \(m+m'=0\), in which case it is one. This proves exactly which product coefficients survive. It does not identify angular averaging with auxiliary averaging. Products add harmonic integers and differentiation preserves them; a finite sequence of such operations leaves a finite band-independent set at each finite stage.

Ordinary coefficient derivatives shift the derivative index in the defining estimate, preserving \(\alpha\). Multiplication by \(M_id_rR^{d_r-1}\), with all its derivatives bounded above, proves \(D_r:\mathcal C_\alpha\to\mathcal C_{\alpha-\kappa_s}\), \(\mathcal C=\mathcal M,\mathcal W\). The other exact maps are \(D_z=\varepsilon\partial_Z:\mathcal C_\alpha\to\mathcal C_{\alpha+1}\), \(-\varepsilon\partial_T:\mathcal C_\alpha\to\mathcal C_{\alpha+1}\), and \(c_{i_0}N_{i_0}:\mathcal C_\alpha\to\mathcal C_\alpha\). Here \(c_{i_0}=T_g^{-\Delta}c_i\) and \(M_{i_0}=\Lambda_g^{-\Delta}M_i\), so a change to a common torus changes constants, not exponents. All these claims concern coefficients with the exponential factored out. Differentiating the exponential itself produces its exact \(ikm\nabla_*\Phi\) or \(ikm\mathsf t_*\Phi\) factor and is accounted for separately in the pulse analysis.

Finally, for \(f_*=Q^{a_f}f_{\rm phys}\), substitution \(r=\sqrt Q R\) gives the exact moment map \[
\int R^e\langle f_*\rangle_Y\,dR
=Q^{a_f-(e+1)/2}\int r^e\langle f_{\rm phys}\rangle_Y\,dr.
\] The measure and power weight are both retained. Normalized Haar compatibility proves that the average may be computed in any of the three torus representations.

\section{Phase, constrained evolution, and every fixed derivative}\label{phase-constrained-evolution-and-every-fixed-derivative}

This file independently derives the claims of source pp.~73-81, Lemma 7.1, Proposition 7.2, Corollary 7.3, and Lemma 7.4. The exact auxiliary maps and operators are proved in the preceding file. The imported background assertion is Proposition 5.5: on the fixed enlarged annulus its chart radial component \(b\) is \(O(\varepsilon)\), its tangential components \(V,G\) differ from their leading profiles by \(O(\varepsilon^2)\), and these statements hold at every fixed slow derivative. The construction of that background belongs to another audit lane. Here the consequences for the pulse system are established explicitly.

\subsection{Cone coordinates and the phase gradient}\label{cone-coordinates-and-the-phase-gradient}

Write \(F=V/R\), \(g=(RF_R,G_R)\) in component order \((\theta,z)\). At the fixed representative of a slow box, the leading shear is \[
g_0=F_0(-a,b_s),\quad N=g_0/|g_0|,\quad K=(-N_z,N_\theta).
\] Let \(t_s=-b_s/a\) and \(v_s=a(1+t_s^2)\). Since \(a,F_0>0\), \[
N=-\frac{(1,t_s)}{\sqrt{1+t_s^2}},\quad
K=\frac{(t_s,-1)}{\sqrt{1+t_s^2}},\quad
|g_0|=aF_0\sqrt{1+t_s^2}.
\] Retain the source's \[
\lambda_0^2=-2F_0N_\theta(2F_0N_\theta+|g_0|),
\qquad c_0=\lambda_0/(2F_0N_\theta).
\] Substitution gives \[
\lambda_0^2
=2F_0^2\left(a-\frac2{1+t_s^2}\right)
=2aF_0^2(1-2/v_s)>0,\qquad
c_0^2=(v_s-2)/2,\quad c_0<0.
\] The positive lower bounds on \(F_0,a,v_s-2\) and compact upper bounds give positive lower bounds on \(\lambda_0,|c_0|\), and finite upper bounds.

At the same, unfrozen profile point, \(T_0=F(p_s-(a,-b_s))\), \(P_c=p_{s,1}+t_sp_{s,2}\), \(J_c=p_{s,2}-t_sp_{s,1}\). Consequently \[
T_0\cdot N=-\frac{F(P_c-v_s)}{\sqrt{1+t_s^2}},
\qquad
T_0\cdot K=-\frac{FJ_c}{\sqrt{1+t_s^2}}.
\] Multiplication of \(T_0\) by the positive chart factor \((Q/q)^{A+1/2}\) changes neither sign nor quotient. Thus the two stress cone inequalities are exactly \(T_{0,*}\cdot N<0\) and \[
\left|c_0\frac{T_{0,*}\cdot K}{T_{0,*}\cdot N}\right|<1.
\] The weighted direction at a vanishing edge is used, not a division of zero values. Its smooth extension and strict compact margin give a supremum \(\mu<1\). Choosing \(u_*>0\) with \(u_*/\sqrt{1+u_*^2}>\mu\) is possible because that function increases to one. All quantities now frozen at representatives are kept fixed under differentiation. Uniform continuity in the compact chart domain makes the same strict cone margin valid on every enlarged box of diameter \(CS_*^{-3}\), after one lower bound on \(\ell\).

Retain \(k=\lceil\varepsilon^{-1/2}\rceil\), so for \(0<\varepsilon\le1\), \[
1\le\varepsilon k^2\le(1+\sqrt\varepsilon)^2\le4,\qquad
k^{-1}\le\sqrt\varepsilon.
\] Define \[
B_s^2=\frac{\lambda_0}{\varepsilon k^2(1+u_*^2)^{3/2}},
\quad s(v)=\sigma(u_*/2+u_*v/L_s),\quad x_0=\sigma B_su_*/2.
\] The \(B_s\) have fixed positive lower and upper bounds. Set \[
(\widetilde p/R_0,p_z)
=B_s\left(K-\frac{\sigma u_*g_0}{L_s|g_0|^2}\right).
\] Choose \(kp\) as the specified nearest nonzero integer to \(k\widetilde p\). The error \(|p-\widetilde p|\le k^{-1}\) suffices, including when the nearest unrestricted integer would have been zero. The integer is constant with the label. Put \[
\Phi=p\theta+p_zZ/\varepsilon+x_0R-vH_\Phi,\qquad H_\Phi=pF+p_zG.
\] Because \(D_rv=D_zv=0\) and the base is auxiliary-independent, \[
n_\Phi=\nabla_*\Phi
=\bigl(x_0-v(H_\Phi)_R,\ p/R,\ p_z-\varepsilon v(H_\Phi)_Z\bigr).
\] At the representative, before rounding, \((\widetilde p/R_0,p_z)\cdot g_0=-\sigma B_su_*/L_s\). Variation across the slow box changes this by \(O(S_*^{-3})\); the background change contributes \(O(\varepsilon^2)\); rounding contributes \(O(k^{-1})\). Hence \[
|(H_\Phi)_R+\sigma B_su_*/L_s|
\le C(S_*^{-3}+\varepsilon^2+k^{-1}).
\] Multiplication by \(v\le L_s=O(S_*)\) gives the radial comparison with \(B_ss(v)\). The tangential comparison has contributions \(O(S_*^{-1})\) from the deliberate tilt, \(O(S_*^{-3}+\varepsilon^2+k^{-1})\) from freezing and rounding, and \(O(\varepsilon S_*)\) from the axial phase term. A single sufficiently large \(\ell_*\) satisfies \[
S_*^2(\varepsilon+\varepsilon^2+k^{-1})\le1
\] for every \(\ell\ge\ell_*\), since \(S_*=\ell^2\) and \(\varepsilon=2^{-h\ell}\). Therefore \[
|n_\Phi-B_s(s(v),K)|\le C/S_*.
\] Differentiation in \(v\) gives \(n_\Phi'=(-(H_\Phi)_R,0,-\varepsilon(H_\Phi)_Z)\), also \(O(S_*^{-1})\). Each fixed number of slow derivatives of these coefficients is bounded by a polynomial in \(S_*\): the only growing factors in these formulas are \(v\), \(\varepsilon^{-1}\) in the phase itself, and the fixed chart derivatives. The \(\varepsilon^{-1}\) term of \(\Phi\) disappears in \(n_\Phi\); it is not discarded from the full exponential when physical derivatives are later taken.

The exact transport defect follows term by term: \[
\mathsf t_*\Phi=-H_\Phi+\varepsilon v(H_\Phi)_T,\qquad
F\partial_\theta\Phi+GD_z\Phi=H_\Phi-\varepsilon vG(H_\Phi)_Z.
\] Thus \[
E_{\rm ik}:=(\mathsf t_*+bD_r+F\partial_\theta+GD_z)\Phi
=\varepsilon v((H_\Phi)_T-G(H_\Phi)_Z)+bn_{\Phi,r}.
\] The \(b\) term is \(O(\varepsilon)\) at every fixed slow derivative. Auxiliary coefficient derivatives of \(v\) are \(O(S_*)\), while derivatives of the fixed affine path of order two or greater vanish. Leibniz's rule therefore proves \(D^IE_{\rm ik}=O_I(\varepsilon S_*^{b_I})\) for every fixed \(I\), on the same domain. It imposes no smallness condition on the higher derivatives.

\subsection{Exact pressure elimination and plane isomorphism}\label{exact-pressure-elimination-and-plane-isomorphism}

Set \[
\mathcal K=
\begin{pmatrix}0&-2F&0\\2F+RF_R&0&0\\G_R&0&0\end{pmatrix},
\quad d=\varepsilon k^2|n_\Phi|^2,\quad
\mathcal A_\Phi=-\mathcal K+
\frac{n_\Phi(n_\Phi^T\mathcal K-(n_\Phi')^T)}{|n_\Phi|^2}.
\] For the original equation \[
t_m'+\mathcal Kt_m+m^2dt_m+ikm n_\Phi\pi_m=-f_m,
\quad n_\Phi\cdot t_m=0,
\] differentiate its constraint to get \(n_\Phi\cdot t_m'=-n_\Phi'\cdot t_m\). Taking its normal component then forces \[
\pi_m=\frac{i}{km}\,
\frac{n_\Phi\cdot\mathcal Kt_m-n_\Phi'\cdot t_m+n_\Phi\cdot f_m}{|n_\Phi|^2}.
\] The factor has sign \(+i/(km)\): multiplying it by \(ikm n_\Phi\) gives the negative of the required normal vector. Substituting gives the projected equation \[
t_m'=\mathcal A_\Phi t_m-m^2dt_m-\operatorname{proj}_{n_\Phi^\perp}f_m.
\] Conversely this projected equation and the displayed pressure imply the original equation provided the constraint holds. For any solution of the projected equation, \[
(n_\Phi\cdot t_m)'=-m^2d(n_\Phi\cdot t_m).
\] An initially zero constraint remains zero, proving equivalence, existence, and uniqueness once the finite-dimensional initial-value equation is solved.

Since \(n_{\Phi,\tan}\) stays close to \(B_sK\), its norm has a uniform positive lower bound. The nonzero \(p/R\) also prevents it from being literally zero. Define \[
K_a=n_{\Phi,\tan}/|n_{\Phi,\tan}|,\quad
N_a=((K_a)_z,-(K_a)_\theta),\quad
s_a=n_{\Phi,r}/|n_{\Phi,\tan}|.
\] Then \(N_a\perp K_a\), \(|N_a|=|K_a|=1\), and \[
U=[e_r-s_aK_a,N_a],\quad
M=\begin{pmatrix}1&1\\c_0\sqrt{1+s^2}&-c_0\sqrt{1+s^2}\end{pmatrix},
\quad B=UM
\] maps \(\mathbb C^2\) isomorphically onto the actual complex plane \(n_\Phi^\perp\). The exact left inverse is \[
B^\ell=M^{-1}\begin{pmatrix}e_r^T\\N_a^T\end{pmatrix}
\] because the last displayed row matrix times \(U\) is \(I_2\). The determinant of \(M\) is \(-2c_0\sqrt{1+s^2}\), bounded away from zero. All factors and their inverses are uniformly bounded in value and polynomially bounded in each fixed coefficient derivative.

For the reference normal \(n_{\rm ref}=B_s(s,K)\), use \(e=e_r-sK\) and \(N\) as an orthogonal plane basis, of lengths \(\sqrt{1+s^2}\) and \(1\). The actual source matrix \(\mathcal K_0\) gives \[
\mathcal K_0e=(2F_0sK_\theta,\ 2F_0e_\theta+g_0),\quad
\mathcal K_0N=(-2F_0N_\theta,0).
\] Consequently \[
e\cdot(-\mathcal K_0e)=sK\cdot g_0=0,\quad
e\cdot(-\mathcal K_0N)=2F_0N_\theta,
\] \[
N\cdot(-\mathcal K_0e)=-(2F_0N_\theta+|g_0|),\quad
N\cdot(-\mathcal K_0N)=0.
\] Orthogonal projection onto \(n_{\rm ref}^\perp\) leaves these scalar products unchanged. Dividing the first coordinate by \(1+s^2\) produces \[
\mathcal D=
\begin{pmatrix}
0&2F_0N_\theta/(1+s^2)\\
-(2F_0N_\theta+|g_0|)&0
\end{pmatrix}.
\] Its two eigenvectors are exactly the columns of \(M\), with eigenvalues \(\lambda=\lambda_0/\sqrt{1+s^2}\) and \(-\lambda\), since \(c_0=\lambda_0/(2F_0N_\theta)\) and \(\lambda_0^2=-2F_0N_\theta(2F_0N_\theta+|g_0|)\). This verifies the orientation and signs of the growing coordinate.

Replacing the background and normal by the actual ones changes the algebraic projected matrix by \(O(S_*^{-1})\). This follows from the explicit rational expression for the projection, the uniform positive denominators, the value bounds already proved, and the mean-value formula on their common compact range. Also \(K_a-K,s_a-s=O(S_*^{-1})\), while their \(v\)-derivatives and \(s'\) are \(O(S_*^{-1})\). Hence \(B'=O(S_*^{-1})\) and the normal-motion term in \(\mathcal A_\Phi\) is \(O(S_*^{-1})\). Therefore \[
B^\ell(\mathcal A_\Phi B-B')=
\operatorname{diag}(\lambda,-\lambda)+E,\qquad |E|\le C/S_*.
\] Every fixed derivative of \(E\) is polynomially bounded; no derivative-smallness claim is needed. The damping comparison is \[
d_{\rm ref}=\varepsilon k^2B_s^2(1+s^2),\qquad
|d-d_{\rm ref}|\le C/S_*,
\] because \(\varepsilon k^2\) is bounded and \(|n_\Phi|^2-B_s^2(1+s^2)\) is \(O(S_*^{-1})\).

\subsection{Envelope and uniform forward inverse}\label{envelope-and-uniform-forward-inverse}

Let \[
a_{\rm net}(y)=\frac{\lambda_0}{\sqrt{1+y^2}}-
\frac{\lambda_0(1+y^2)}{(1+u_*^2)^{3/2}},
\qquad y=|s(v)|=u_*/2+u_*v/L_s.
\] It is zero at \(v=L_s/2\), and direct differentiation gives \[
\frac{d}{dv}a_{\rm net}(|s(v)|)
=-\frac{u_*\lambda_0|s(v)|}{L_s}
\left((1+s(v)^2)^{-3/2}+2(1+u_*^2)^{-3/2}\right).
\] Since \(u_*/2\le|s|\le3u_*/2\), its value lies between \(-C/L_s\) and \(-c/L_s\), for fixed positive \(c,C\). Define \(P(v)=\exp(\int_{L_s/2}^v a_{\rm net}(|s(w)|)\,dw)\). Integrating the derivative inequality on either side of the midpoint, keeping the direction of each integral, gives \[
-C_1(v-L_s/2)^2/L_s\le\log P(v)\le-c_1(v-L_s/2)^2/L_s.
\] This proves the two Gaussian bounds and \(0<P\le1\) with the source's integral normalization.

Writing \(t_m=Bz_m\) gives \[
z_m'=A_mz_m+g_m,\quad
A_m=\operatorname{diag}(\lambda,-\lambda)+E-m^2dI_2,
\quad g_m=-B^\ell\operatorname{proj}_{n_\Phi^\perp}f_m.
\] For the complex Euclidean norm, taking the real part of \(z^*A_1z\) bounds the upper right derivative of \(|z|\) by \((\lambda-d_{\rm ref}+C/S_*)|z|\). Integration gives \[
\|V_1(v,w)\|\le
\exp(C(v-w)/S_*)\,P(v)/P(w)\le C'P(v)/P(w).
\] The factor \(C'\) is fixed because \(0\le v-w\le L_s\asymp S_*\). Since the difference of \(A_m\) from \(A_1\) is the scalar matrix \(-(m^2-1)dI_2\), differentiation directly verifies \[
V_m(v,w)=\exp\left(-(m^2-1)\int_w^v d(a)\,da\right)V_1(v,w).
\] For every nonzero integer \(m\), \(m^2-1\ge0\) and \(d>0\), so the same value bound is valid for all such harmonics. This is the exact mechanism preventing a new small-domain requirement when the finite harmonic set increases.

If the source obeys \[
|D^If_m|\le C_I\varepsilon^\alpha S_*^{b_I}
\sqrt\zeta\,\delta^{-a_I}P(v),
\] the zero-data Duhamel integral gives \[
z_m(v)=\int_0^vV_m(v,w)g_m(w)\,dw,\quad
|z_m(v)|\le C\varepsilon^\alpha S_*^{b+1}
\sqrt\zeta\,\delta^{-a}P(v).
\] The source's \(P(w)\) cancels the propagator's denominator exactly. The shell weights are constant along the integration path because \(X\) is constant there.

For completeness, parameter differentiation can be performed without differentiating a possibly exponentially large propagator in isolation. Hold \(v\) fixed and use slow variables and the initial transverse coordinate as parameters. They commute with \(\partial_v\); the zero initial data and all their parameter derivatives vanish. For each multi-index \(I\), \[
(\partial^Iz_m)'=A_m\partial^Iz_m+\partial^Ig_m+
\sum_{0<J\le I}\binom IJ(\partial^JA_m)\partial^{I-J}z_m.
\] At fixed finite harmonic bound \(M\), each \(\partial^JA_m\) is bounded by \(C_{J,M}S_*^{c_J}\). Induction on \(|I|\), using the already proved propagator bound, shows that every source term on the right has the same \(\varepsilon^\alpha\sqrt\zeta P\) weight with a larger inverse-distance exponent and polynomial degree. A further time integral contributes only \(L_s\). One may choose the degree at step \(I\) to exceed \[
1+\max\{b_I,\max_{0<J\le I}(c_J+b'_{I-J})\}.
\] Pulse derivatives follow from the differential equation, and mixed pulse/parameter derivatives follow by repeated differentiation of that equation. Returning to the current \(H\) uses \(D_Hv=O(S_*)\), \(D_HH_w=I-v_t\lambda_t\), with higher affine derivatives zero, and therefore introduces only additional polynomial powers. Multiplying by \(B\) yields the stated bound for \(t_m\). The explicit pressure has the extra factor \((km)^{-1}=O(\varepsilon^{1/2})\); its other factors have the proved derivative bounds. Thus the exponent for pressure is \(\alpha+1/2\).

The support argument also follows from this equation. If a fixed slow/transverse point is outside the source's closed support, the entire source path is zero; uniqueness gives a zero solution. At a boundary all source jets vanish by the prescribed smooth zero extension, so the differentiated equations inductively give zero solution jets. Across a shell edge the retained \(\sqrt\zeta\delta^{-a_I}\) bounds give the same conclusion. Deck translations preserve the path and the prescribed data, giving common-torus descent by uniqueness. Smooth extension to the rectangle enlargement follows from smooth finite-dimensional ODE dependence on its coefficients and data; no temporal zero extension is asserted before cutoff. On compact sets with \(q>0\), one-sided background/source derivatives at \(\tau=0\) pass through the same parameter equations and pressure formula.

Only the lower bounds for denominators and finitely many value comparisons determine \(q_*\). The induction above changes constants and polynomial degrees for each fixed \(M,I\); it does not impose further smallness inequalities. This proves the quantifier statement of Corollary 7.3.

\subsection{The real homogeneous pulse and energy transfer}\label{the-real-homogeneous-pulse-and-energy-transfer}

For \(m=1,f=0\), specify \(z_+(0)=P(0)>0,z_-(0)=0\). While \(z_+>0\), the ratio \(r=z_-/z_+\) obeys exactly \[
r'=E_{21}+(-2\lambda+E_{22}-E_{11})r-E_{12}r^2.
\] Let \(c_\lambda=\inf\lambda>0\). On \(r=\pm K_b/S_*\), the outward derivative is at most \[
\frac{-2c_\lambda K_b+C(1+K_b/S_*+K_b^2/S_*^2)}{S_*}.
\] Choose fixed \(K_b>C/c_\lambda\), then enlarge the fixed band threshold to make the numerator negative. The closed interval \([-K_b/S_*,K_b/S_*]\) is invariant until any putative first zero of \(z_+\). On this interval, \[
\left(\log(z_+/P)\right)'=-(d-d_{\rm ref})+E_{11}+E_{12}r=O(S_*^{-1}).
\] Since \(z_+(0)=P(0)\) and \(L_s/S_*\) is bounded, integration proves \(cP\le z_+\le CP\), ruling out that putative zero. For all \(v\), \[
x=z_++z_-=z_+(1+r),\quad
y=c_0\sqrt{1+s^2}(z_+-z_-),
\] so \[
cP\le x\le CP,\qquad
y/x=c_0\sqrt{1+s^2}+O(S_*^{-1}).
\] The same differentiated equation used for the inverse proves every fixed derivative bound, with the nonzero initial value retained: it is fixed with the label, so all slow/transverse derivatives of that datum vanish. The value solution contributes \(\|V_1(v,0)\|P(0)\le CP(v)\). Pulse derivatives are obtained from the equation. Pressure again has the extra \(\varepsilon^{1/2}\).

For the real homogeneous amplitude \(t\), \(n_\Phi\cdot t=0\) makes the normal component of \(\mathcal A_\Phi t\) perpendicular to \(t\). Direct multiplication, keeping both cylindrical \(2F\) terms, gives \[
t\cdot\mathcal Kt
=-2Ft_rt_\theta+(2F+RF_R)t_rt_\theta+G_Rt_rt_z
=g\cdot(t_r(t_\theta,t_z)).
\] Taking \(t\cdot t'\) therefore proves the exact identity \[
\frac12\frac{d}{dv}|t|^2
=-g\cdot(t_r(t_\theta,t_z))-\varepsilon k^2|n_\Phi|^2|t|^2.
\] At the representative the first term is \(-|g_0|T_N\) for the flux \(T=T_NN+T_KK\). It is positive for the negative \(N\)-component of the realized stress. The carrier physical wavelength is \(\sqrt Q/k\asymp Q^{1/2+h/2}\), while the leading physical wave amplitude has power \(Q^{-A}\sqrt\varepsilon=Q^{-1/2-h/2}\), up to polynomial factors in \(S_*\). Their product has power \(Q^0\); this agrees with the retained leading damping \(\varepsilon k^2\asymp1\).

\section{Covariance, signed stress increments, exact curls, and cutoff residuals}\label{covariance-signed-stress-increments-exact-curls-and-cutoff-residuals}

This file gives direct proofs for source pp.~81-87, Propositions 7.5-7.6, Lemma 7.7, and Corollary 7.8. The background inputs and pulse estimates are recorded and derived in the preceding files. All averages below are on the extended domain at fixed slow coordinates, before physical evaluation.

\subsection{Two columns and their positive coefficients}\label{two-columns-and-their-positive-coefficients}

For real extended vector fields define \[
\mathcal C(w)=\langle\langle w_rw_{\tan}\rangle_\theta\rangle_Y,\qquad
\mathcal B(w,v)=
\langle\langle w_rv_{\tan}+v_rw_{\tan}\rangle_\theta\rangle_Y,
\quad w_{\tan}=(w_\theta,w_z).
\] Expanding the product proves \(\mathcal C(w+v)=\mathcal C(w)+\mathcal B(w,v)+\mathcal C(v)\) and \(D\mathcal C(w)[v]=\mathcal B(w,v)\). These maps preserve the original component order \((r\theta,rz)\).

For a slow box \(\beta=(\ell,a)\), its two signs have disjoint auxiliary rectangles and the same slow cutoff. Write \[
b_\sigma=\chi_g(\xi_g)\psi(v)t_\sigma^h\cos(k\Phi_\sigma),\qquad
H=[\mathcal C(b_+)\mid\mathcal C(b_-)].
\] Each \(b_\sigma\) is extended by zero outside its rectangle, using its smooth transverse and temporal cutoffs. The integer \(kp\ne0\) ensures that \[
\frac1{2\pi}\int_0^{2\pi}\cos^2(k\Phi_\sigma)\,d\theta=\frac12
\] exactly, whatever the non-angular part of the phase. Haar pullback permits averaging on the band torus. The coordinate map \((\xi_g,\eta_g)\mapsto c_\gamma+\xi_gv_r+\eta_gv_t\) has the retained Jacobian \(|\det(v_r,v_t)|=1+b_g^2\), and \(d\eta_g=c_i\,dv\). The homogeneous amplitude is independent of \(\xi_g\), since its ODE coefficients depend on slow coordinates and \(v\), and its initial value is fixed with the label. With radial component \(x\), this gives \[
H_\sigma=D_g c_i\int_0^{L_s}\psi^2x\,t_{\sigma,\tan}^h\,dv,\qquad
h_\sigma=D_g c_i\int_0^{L_s}\psi^2x^2\,dv,
\] \[
D_g=\frac{|\det(v_r,v_t)|}{2}\int_{\mathbb R}\chi_g(\xi)^2\,d\xi>0.
\] Thus the angular factor \(1/2\), the rectangular Jacobian, the time measure \(c_i\,dv\), and the transverse cutoff integral are all present.

The proved bounds \(cP\le x\le CP\) and \(P(v)\asymp\exp(-c(v-L_s/2)^2/L_s)\), with separate upper and lower constants, imply \[
c\sqrt{L_s}\le\int_0^{L_s}\psi^2x^2\,dv\le C\sqrt{L_s}.
\] For the lower bound, integrate over \(|v-L_s/2|\le\sqrt{L_s}\), which lies in the region \(\psi=1\) for all sufficiently large \(L_s\); \(P\) there has a fixed positive lower bound. The upper bound follows by extending the Gaussian integral to \(\mathbb R\). The same substitution \(w=(v-L_s/2)/\sqrt{L_s}\) proves \[
\int_0^{L_s}\frac{|v-L_s/2|}{L_s}\psi^2x^2\,dv\le C.
\] It follows that \(h_\sigma\asymp c_i\sqrt{L_s}\asymp S_*^{-1/2}\), and the probability measure proportional to \(\psi^2x^2\,dv\) has expectation of \(|v-L_s/2|/L_s\) at most \(CS_*^{-1/2}\).

Since \(t_{\sigma,\tan}^h/x=-s_aK_a+(y/x)N_a\), the previous file proves its uniform difference from \(-s(v)K+c_0\sqrt{1+s(v)^2}N\) is \(O(S_*^{-1})\). The latter function is Lipschitz in \(s\) on the fixed compact range, and \(|s(v)-\sigma u_*|=u_*|v-L_s/2|/L_s\). Taking the weighted average therefore gives \[
H_\sigma=h_\sigma(-A_cN-\sigma u_*K+e_\sigma),\quad
A_c=-c_0\sqrt{1+u_*^2}>0,\quad |e_\sigma|\le CS_*^{-1/2}.
\] Every fixed slow derivative of \(H_\sigma\) is polynomially bounded: differentiate the finite integral using the already proved pulse derivatives, fixed \(L_s\), and fixed label cutoffs. None of these parameter derivatives differentiates a discrete label or moves the integration endpoints.

In coordinates \(N,K\), omit \(e_\sigma\) temporarily and solve the exact two-by-two system: \[
h_+y_+=\frac12(-T_N/A_c-T_K/u_*),\qquad
h_-y_-=\frac12(-T_N/A_c+T_K/u_*).
\] The cone choice made in the previous file gives a fixed \(\eta_c>0\) such that \[
\frac{|T_K|}{u_*}\le(1-\eta_c)\frac{-T_N}{A_c},
\qquad T=T_{0,*}.
\] Since \(N,K\) are orthonormal and \(T_N<0\), this also gives \(-T_N\ge c|T|\). Hence the two displayed reference coefficients are at least \(c|T|\) and at most \(C|T|\).

The reference column matrix in these coordinates is \[
H_{\rm ref}=
\begin{pmatrix}-A_ch_+&-A_ch_-\\-u_*h_+&u_*h_-\end{pmatrix},
\quad
\det H_{\rm ref}=-2A_cu_*h_+h_-.
\] Writing \(H=(H_{\rm ref}\operatorname{diag}(h_+^{-1},h_-^{-1})+
\mathcal E)\operatorname{diag}(h_+,h_-)\), the first factor differs from a fixed invertible matrix by \(O(S_*^{-1/2})\). For \(\|H_0^{-1}\mathcal E\|<1/2\), the exact inverse of
\(I+H_0^{-1}\mathcal E\) is the convergent series
\(\sum_{n=0}^{\infty}(-H_0^{-1}\mathcal E)^n\).
Indeed multiplication of its finite partial sum through \(n=N\) by
\(I+H_0^{-1}\mathcal E\), in either order, gives
\(I-(-H_0^{-1}\mathcal E)^{N+1}\), and the last power tends to zero.
The sum has norm at most \(2\), and its difference from \(I\) has norm
at most \(2\|H_0^{-1}\mathcal E\|\), by the geometric-series bounds.
Thus the exact first-factor inverse is
\((I+H_0^{-1}\mathcal E)^{-1}H_0^{-1}\); the original diagonal
\(\operatorname{diag}(h_+,h_-)^{-1}\) remains on its left. The coefficient perturbation is then bounded by \(CS_*^{-1/2}|T|\) before dividing by \(h_\sigma\). For a further fixed lower band threshold, it is smaller than half the positive reference coefficient margin. Therefore \[
|\det H|\ge c/S_*,\quad \|H^{-1}\|\le C\sqrt{S_*},\quad
c\sqrt{S_*}|T_{0,*}|\le y_\sigma\le C\sqrt{S_*}|T_{0,*}|,
\quad y=H^{-1}T_{0,*}.
\] This proves positivity uniformly up to limiting stress directions at the shell edges. At an edge the actual target is zero; the positive inequalities for \(y\) are asserted in the open shell.

The background profile theorem gives \(|T_{0,*}|\ge c\zeta\) and \(|D^IT_{0,*}|\le C_I\zeta\delta^{-a_I}\), using the exact chart/profile transition from the first file. Differentiating \(HH^{-1}=I\) expresses each fixed derivative of \(H^{-1}\) as a finite sum of products of \(H^{-1}\) and derivatives of \(H\), hence gives polynomial \(S_*\) bounds. Leibniz's rule proves \[
y_\sigma\ge c\sqrt{S_*}\zeta,\qquad
|D^Iy_\sigma|\le C_IS_*^{b_I}\zeta\delta^{-a_I}.
\] For \(a_\sigma=\sqrt{y_\sigma}\), every derivative of positive order is a finite sum \[
c\,y_\sigma^{1/2-n}\prod_{\nu=1}^nD^{I_\nu}y_\sigma,
\qquad |I_\nu|\ge1,\quad \sum I_\nu=I.
\] The numerator contributes \(\zeta^n\), while the inverse power contributes \(\zeta^{1/2-n}\); the product has exactly the remaining \(\sqrt\zeta\). The \(S_*\)-powers and inverse edge-distance powers are finite at every fixed order. The value bound uses the upper bound for \(y_\sigma\). Thus \[
|D^Ia_\sigma|\le C_IS_*^{b'_I}\sqrt\zeta\,\delta^{-a'_I},
\] which proves smooth zero extension of every coefficient at both shell edges.

Set \(W_0=\sqrt\varepsilon\sum_{\sigma=\pm}a_\sigma b_\sigma\). The two disjoint supports remove the cross products, and the scalar amplitudes are independent of both averaging variables. Consequently \[
\mathcal C(W_0)=\varepsilon H(a_+^2,a_-^2)^T
=\varepsilon T_{0,*}
\] exactly. The pulse and square-root estimates give the local \(\mathcal W_{1/2}\) bounds. Multiplying by \(\eta_\beta\) gives the required slow support.

The physical covariance of the assembled field is obtained with the original factors: \[
Q^{-2A}\varepsilon T_{0,*}
=Q^{-2A+h}(Q/q)^{A+1/2}T_0
=Q^{-A+h+1/2}q^{-A-1/2}T_0
=q^{-A-1/2}T_0.
\] The final equality uses \(A=1/2+h\). Each physical slow point has bounded overlap, and cross-label products vanish. Summing the squared partition \(\sum_\beta\eta_\beta^2=1\) therefore supplies precisely the leading physical stress, with the sign suitable to cancel the negative stress divergence in Proposition 5.5.

\subsection{Signed differential inverse and assembly across bands}\label{signed-differential-inverse-and-assembly-across-bands}

Let \(\Sigma(R,Z,T)\) be a real auxiliary-independent two-vector in \(\mathcal M_\alpha\). Keep the previously chosen positive amplitudes fixed. Define \[
d_\Sigma=H^{-1}(\Sigma/\varepsilon),\qquad
\delta a_\sigma=(d_\Sigma)_\sigma/(2a_\sigma),\qquad
\mathcal L\Sigma=\sqrt\varepsilon\sum_{\sigma=\pm}\delta a_\sigma b_\sigma.
\] Here the division is taken only on the open shell before smooth zero extension. Every derivative of \(y_\sigma^{-1/2}\) is a finite sum \[
c\,y_\sigma^{-n-1/2}\prod_{\nu=1}^nD^{I_\nu}y_\sigma,
\] so it has the weight \(\zeta^{-1/2}\) with allowed polynomial and edge-distance factors. Since \(d_\Sigma\) has weight \(\varepsilon^{\alpha-1}\zeta\), the product defining \(\delta a_\sigma\) has \[
|D^I\delta a_\sigma|
\le C_I\varepsilon^{\alpha-1}S_*^{b_I}\sqrt\zeta\,\delta^{-a_I}.
\] This proves smooth zero extension and the local \(\mathcal W_{\alpha-1/2}\) bound for \(\mathcal L\Sigma\), with the envelope supplied by \(b_\sigma\).

Disjointness and the retained bilinear factor two give \[
\mathcal B(W_0,\mathcal L\Sigma)
=\varepsilon H(2a_+\delta a_+,2a_-\delta a_-)^T
=\Sigma.
\] This is an exact linear right inverse of \(D\mathcal C(W_0)\), without positivity or smallness assumptions on \(\Sigma\). It is not an exact nonlinear inverse: the retained identity is \[
\mathcal C(W_0+\mathcal L\Sigma)
=\varepsilon T_{0,*}+\Sigma+\mathcal C(\mathcal L\Sigma).
\] The last term lies in \(\mathcal M_{2\alpha-1}\), as both wave coefficients have exponent \(\alpha-1/2\) and the product has the mean weight \(\zeta\).

For a real physical auxiliary-independent stress \(\sigma(r,z,t)\)
supported in the active shell, whose representatives
\(Q_\beta^{2A}\sigma\) belong to the stated local \(\mathcal M_\alpha\)
domain, use the source's assembled maps. This domain includes every fixed
derivative bound with the weight \(\zeta\delta^{-a_I}\), compatibility on
overlapping charts, and smooth zero extension at the shell edges. The
previously chosen positive amplitudes remain fixed. The following identity
holds throughout the active partition domain and extends by zero outside it
because the input stress vanishes there: \[
W_{0,\rm phys}^{\rm as}
=\sum_\beta Q_\beta^{-A}\eta_\beta W_{0,\beta},\qquad
\mathcal L_{\rm phys}^{\rm as}\sigma
=\sum_\beta Q_\beta^{-A}\eta_\beta
\mathcal L_\beta(Q_\beta^{2A}\sigma).
\] The different \(Q_\beta\) are retained. The contribution of box \(\beta\) to the bilinear covariance is \[
Q_\beta^{-2A}\eta_\beta^2Q_\beta^{2A}\sigma
=\eta_\beta^2\sigma.
\] There are no cross-box terms, so summation gives \(\mathcal B(W_{0,\rm phys}^{\rm as},\mathcal L_{\rm phys}^{\rm as}\sigma)=\sigma\). In a reference chart of scale \(Q\), the representatives are \[
W_0^{\rm as}=Q^AW_{0,\rm phys}^{\rm as},\quad
\mathcal L^{\rm as}\Sigma=Q^A\mathcal L_{\rm phys}^{\rm as}(Q^{-2A}\Sigma).
\] Thus the exact chart identity is \(\mathcal B(W_0^{\rm as},\mathcal L^{\rm as}\Sigma)=\Sigma\). Overlapping \(Q_\beta/Q\) lie between fixed positive bounds, so chart changes preserve all claimed exponents; the bounded covering changes preserve derivatives and averages. Equal label-harmonic contributions are collected before class membership is tested. Auxiliary independence of \(\Sigma\) is essential to taking the amplitudes outside the integral and is explicitly a hypothesis here.

\subsection{Curl identity with all cylindrical terms}\label{curl-identity-with-all-cylindrical-terms}

The normalized cylindrical curl is \[
\operatorname{curl}_*A=
\left(R^{-1}\partial_\theta A_z-D_zA_\theta,\,
D_zA_r-D_rA_z,\,
(D_r+R^{-1})A_\theta-R^{-1}\partial_\theta A_r\right).
\] For a smoothly supported \(t_m\in\mathcal W_\alpha\) with \(n_\Phi\cdot t_m=0\), put \[
C_m=\frac{i\,n_\Phi\times t_m}{km|n_\Phi|^2},
\quad A_m=C_me^{ikm\Phi}.
\] Using the right-handed \((e_r,e_\theta,e_z)\) orientation, the exponential differentiation contributes \[
ikm\,n_\Phi\times C_m
=-\frac{n_\Phi\times(n_\Phi\times t_m)}{|n_\Phi|^2}
=t_m-\frac{n_\Phi(n_\Phi\cdot t_m)}{|n_\Phi|^2}=t_m.
\] The exact remainder is \[
r_m=\left(-D_z(C_m)_\theta,\,
D_z(C_m)_r-D_r(C_m)_z,\,
(D_r+R^{-1})(C_m)_\theta\right).
\] The multiplier \(n_\Phi/|n_\Phi|^2\) has all polynomial coefficient bounds. Hence \(C_m\in\mathcal W_{\alpha+1/2}\), and \(r_m\in\mathcal W_{\alpha+1/2-\kappa_s}\). The less costly \(D_z\) and frame terms satisfy the same weaker common bound. Supports are preserved by local differentiation and smooth extension.

For any smooth normalized vector potential, retain \[
\operatorname{div}_*a=(D_r+R^{-1})a_r+R^{-1}\partial_\theta a_\theta+D_za_z.
\] Because \(D_r(R^{-1})=-R^{-2}\), one has \((D_r+R^{-1})(R^{-1}f)=R^{-1}D_rf\). Substitution of the three curl entries gives the six terms \[
R^{-1}D_r\partial_\theta A_z
-(D_r+R^{-1})D_zA_\theta
+R^{-1}D_z\partial_\theta A_r
-R^{-1}\partial_\theta D_rA_z
+D_z(D_r+R^{-1})A_\theta
-R^{-1}D_z\partial_\theta A_r.
\] They cancel in pairs by the commuting coordinate derivations and \(D_z(R^{-1})=0\). Thus divergence of the exact curl vanishes. The pullback intertwining proved in the first file transfers this identity exactly to physical space.

For the full amplitude \(a_m=t_m+r_m\), its coefficient is angular-independent. Expanding the divergence of \(a_me^{ikm\Phi}\) proves \[
ikm\,n_\Phi\cdot a_m
=-\big((D_r+R^{-1})(a_m)_r+D_z(a_m)_z\big).
\] The scalar \(n_\Phi\cdot a_m=n_\Phi\cdot r_m\) belongs to \(\mathcal W_{\alpha+1/2-\kappa_s}\), since \(n_\Phi\cdot t_m=0\) and the full phase vector has polynomial coefficient bounds. The displayed left side includes the additional factor \(ikm\), with \(k=O(\varepsilon^{-1/2})\) and fixed nonzero harmonic \(m\); its class is therefore \(\mathcal W_{\alpha-\kappa_s}\). This is also the class supplied by the amplitude derivative on the right. In particular, the complete amplitude is not asserted to be transverse: its retained longitudinal component is precisely what the divergence identity controls. Physical vector potentials are \(Q^{1/2-A}A_*\), because physical curl contributes \(Q^{-1/2}\); physical pressures are \(Q^{-2A}p_*\). These factors give exactly the prescribed physical velocity and normalized pressure.

Real waves use conjugate harmonic pairs. For real \(t\cos(k\Phi)\), the velocity coefficients are \(t/2\) for \(m=\pm1\); the factor \(i/(km)\) changes by complex conjugation when \(m\) changes sign. Thus the potentials and pressures also form real physical fields. No missing factor is introduced by replacing the cosine with exponentials.

\subsection{Exact cutoff residual and terminal flatness at each fixed stage}\label{exact-cutoff-residual-and-terminal-flatness-at-each-fixed-stage}

Let the zero-data inverse solve the principal equation for an actual supported source \(f_m\), and set \(\widehat t_m=\psi t_m,\widehat\pi_m=\psi\pi_m\). The principal equation yields exactly \[
\widehat t_m'+\mathcal K\widehat t_m+m^2d\widehat t_m
+ikmn_\Phi\widehat\pi_m+f_m
=(1-\psi)f_m+\psi't_m.
\] The source on the left remains \(f_m\), not \(\psi f_m\). This is why the uncancelled \((1-\psi)f_m\) term is required. The support of either right-hand term lies in \(|v-L_s/2|\ge L_s/5\), since \(\psi=1\) on the complementary central region. There the proved Gaussian estimate is \(P(v)\le e^{-cL_s/25}\le e^{-c'S_*}\). For every fixed coefficient derivative, Leibniz's formula gives the same exponential factor times allowed polynomial powers: derivatives of \(\psi\) are supported in the same exterior region, and all derivatives of \(f_m,t_m\) already obey the envelope bound. The transverse/slow/shell weights stay unchanged.

When the corresponding exponential and physical factors are differentiated any fixed number of times, only a fixed power \(Q^{-M}\) and a polynomial \(S_*^C\) are added. This can be verified directly: on band pullbacks, \(T_g^i\le Q^{-1-h}/S_*\), \(\Lambda_g^i\le Q^{-(1+h)\rho_g}S_*^{-\rho_g}\), \(k\le2Q^{-h/2}\); chart derivatives contribute fixed powers \(Q^{-1},Q^{-D},Q^{-1/2}\); and the phase itself has only the displayed \(p_zZ/\varepsilon\), \(vH_\Phi\), \(p\theta\), \(x_0R\) terms, with \(v=O(S_*)\) and all relevant differentiated coefficients bounded polynomially. The cylindrical frame has \(r^{-1}=Q^{-1/2}R^{-1}\), with \(R\) bounded away from zero. A fixed product/chain-rule expansion contains finitely many such factors. The exponents \(M,C\) may depend on the derivative order, stage, and finite harmonic set, but not on the band.

For every fixed \(M,C,N\), \(q\asymp Q=2^{-\ell}\) gives \[
q^{-N}Q^{-M}S_*^Ce^{-c'S_*}
\le C_N\exp(-c'\ell^2+(M+N)\ell\log2+2C\log\ell)\longrightarrow0.
\] For a direct bound, choose \(\ell\) so that the last two positive terms are at most \(c'\ell^2/2\); the remaining Gaussian is bounded and tends to zero. This proves every physical derivative of this fixed-stage cutoff residual is \(O(q^N)\) for every \(N\). Homogeneous pulse cutoff tails obey the same argument, with source zero and the retained \(\psi't^h\) term.

At each physical point only a bounded number of slow labels occur, so local summation retains the same bound. On their annular supports, \(r^2/(2q)\ge X_a>0\) implies \(q\le r^2/(2X_a)\). Thus these derivative bounds give arbitrary vanishing order at the singular space-time point \(r=z=0,t=1\). Every one-sided derivative jet there has value zero; the derivative identities persist at that boundary point by induction and the fundamental theorem of calculus along coordinate segments lying in the time half-space. Outside the annular supports the fields vanish smoothly. At \(t=1,z\ne0\), however, \(q=|z|^{1/D}>0\), so this estimate does not assert zero values on the whole terminal plane. One-sided smooth jets there are supplied by the parameter-dependent ODE proof. Any continuation across that plane, and the infinite-stage diagonal sum, are the later sections' responsibilities and are not inferred merely from the fixed-stage tail estimate.

\subsection{Actual covariance increment with every remainder retained}\label{actual-covariance-increment-with-every-remainder-retained}

Let \(U=W_0^{\rm as}+E\) be the actual real divergence-free wave field, with \(U\in\mathcal W_{1/2}\), \(E\in\mathcal W_\rho\). Given \(\Sigma\in\mathcal M_\alpha\) independent of the averaging variables, applying the exact curl construction to \(\mathcal L^{\rm as}\Sigma\) gives \[
V=\mathcal L^{\rm as}\Sigma+R,\qquad R\in\mathcal W_{\alpha-\kappa_s}.
\] The latter exponent is \((\alpha-1/2)+(1/2-\kappa_s)\). Then \[
\begin{aligned}
\mathcal C(U+V)-\mathcal C(U)
&=\mathcal B(U,V)+\mathcal C(V)\\
&=\Sigma+\mathcal B(E,\mathcal L^{\rm as}\Sigma)
+\mathcal B(U,R)+\mathcal C(V).
\end{aligned}
\] The exact bilinear identity at the fixed \(W_0^{\rm as}\) is the only cancellation used. Products yield the respective mean exponents \[
\rho+\alpha-\tfrac12,\qquad
\alpha+\tfrac12-\kappa_s,\qquad
2\alpha-1.
\] For the last exponent, \(\kappa_s<1/2\) and \(0<\varepsilon\le1\) imply \(R\in\mathcal W_{\alpha-1/2}\); hence \(V\) is in that class. The radial divergence is the original componentwise operator \(D_r+2/R\) or \(D_r+1/R\), respectively. Its derivative term loses at most \(\kappa_s\), and its bounded frame terms cause no further loss. Thus all three remainder exponents after radial divergence are their displayed values minus at most \(\kappa_s\). These are precise local estimates. Their compatibility with the residual-improvement exponents selected in Section 9 must be checked where those exponents are selected.

\section{Admissible shear loop, radial modulation, and moment restoration}\label{admissible-shear-loop-radial-modulation-and-moment-restoration}

This file independently proves the operations in Appendix C, pp.~157-165. It retains the original logarithmic coordinate \(D_X=X\partial_X\), profile \(E>0,U\), pressure integration constant \(\Pi_0(\eta)\), component signs, interval ordering, and all five moments. The existence and numerical parameter margins of the input profile from Appendices A and B are audited by the profiles lane. This file verifies what Appendix C does to that input, including the exact restoration map; it does not claim that its imported input construction has been independently proved here.

\subsection{Scalar cone and its four gaps}\label{scalar-cone-and-its-four-gaps}

At \(a>0\), set \[
t_s=-b_s/a,\quad v_s=a(1+t_s^2),\quad
P_c=p_{s,1}+t_sp_{s,2},\quad
J_c=p_{s,2}-t_sp_{s,1}.
\] The relaxed condition is \(P_c>2\), \(v_s<U(P_c,J_c)\), with \[
U(P,J)=P+\frac{J^2}{4}
-|J|\sqrt{\frac{P-2}{2}+\frac{J^2}{16}}.
\] The admissible condition also requires \(v_s>2\). The quadratic expression \[
G(v)=2(P-v)^2-(v-2)J^2
\] has roots \(U(P,J)\) and the same expression with a plus before the square root. For \(P>2\), \(U>2\): writing \(w=P-2>0\), the inequality \[
\left(w+\frac{J^2}{4}\right)^2
-J^2\left(\frac w2+\frac{J^2}{16}\right)=w^2>0
\] proves \(w+J^2/4>|J|\sqrt{w/2+J^2/16}\). Also \(U\le P\), because the square-root term is at least \(J^2/4\), with equality only for \(J=0\). When \(J=0\), \(U=P>2\). Therefore for \(v>2\), admissibility is exactly \[
P-v>0,\qquad 2(P-v)^2-(v-2)J^2>0.
\] The four continuous smooth gaps on \(a>0\) are \(a\), \(v_s-2\), \(P_c-v_s\), and \(2(P_c-v_s)^2-(v_s-2)J_c^2\). Positivity on a compact parameter set gives one positive lower margin for all four. The root \(U\) need only be continuous; no differentiation of its \(|J|\) term at zero is used.

\subsection{A smooth loop whose average is exactly the original shear}\label{a-smooth-loop-whose-average-is-exactly-the-original-shear}

Let \(I\times[-1,1]\) be the fixed compact domain from Appendix C, with relaxed data and admissibility on radial boundary collars. The positive minimum of \(P_c(t_s)-2\) permits a constant \[
0<d_0<\tfrac12\min(P_c(t_s)-2).
\] Here \(P_c(t)=p_{s,1}+p_{s,2}t\) and \(J_c(t)=p_{s,2}-p_{s,1}t\) for a variable slope \(t\). With normalized \(\theta'\)-average on \([0,2\pi]\), define \[
M_e(z)=\frac1{2\pi}\int_0^{2\pi}e^{z\sin\theta'}\,d\theta',
\quad
t(\theta';\mu)=t_s+
\frac{d_0}{p}\left(\frac{e^{\mu p\sin\theta'}}{M_e(\mu p)}-1\right),
\quad p=p_{s,2}.
\] For \(p=0\) the quotient is defined by continuity. Its numerator is smooth in \(p,\mu,\theta'\), zero at \(p=0\), and the identity \[
\frac{G(p)}p=\int_0^1G'(up)\,du
\] proves joint smoothness, including all parameters and derivatives. Since \(M_e(0)=1,M_e'(0)=0\), the value there is \(t=t_s+d_0\mu\sin\theta'\).

The definitions give \[
\langle t\rangle_{\theta'}=t_s,\qquad
P_c(t)=P_c(t_s)+d_0\left(\frac{e^{\mu p\sin\theta'}}{M_e(\mu p)}-1\right)
\ge P_c(t_s)-d_0>2.
\] The variance is \[
V(\mu,p)=\langle(t-t_s)^2\rangle
=\frac{d_0^2}{p^2}
\left(\frac{M_e(2\mu p)}{M_e(\mu p)^2}-1\right),
\] with value \(d_0^2\mu^2/2\) at \(p=0\). To show strict increase, let \(g(z)=\log M_e(z)\). Differentiating under the finite integral gives \[
g''(z)=
\langle(\sin\theta'-\langle\sin\theta'\rangle_z)^2\rangle_z>0,
\] where the tilted probability density is \(e^{z\sin\theta'}/(2\pi M_e(z))\). It is strictly positive on the circle and \(\sin\theta'\) is not constant, so the variance is strictly positive. For \(p\ne0,\mu>0\), \[
\partial_\mu(g(2\mu p)-2g(\mu p))
=2p(g'(2\mu p)-g'(\mu p))>0.
\] For \(p>0\) both factors on the right are positive; for \(p<0\) both are negative. Thus \(V\) strictly increases in \(\mu>0\), including its separately proved formula at \(p=0\).

Here is the required growth bound, without an unproved asymptotic invocation. For \(z\ge1\), write \(\theta'=\pi/2+x\) near the maximum. On a fixed small interval \(|x|\le a<\pi\), there are \(c_1,c_2>0\) with \(1-c_2x^2\le\cos x\le1-c_1x^2\). A lower bound integrates over \(|x|\le z^{-1/2}\min(a,1)\), giving \(ce^zz^{-1/2}\). An upper bound on the same interval uses the Gaussian integral and gives \(Ce^zz^{-1/2}\). On its complement \(\sin\theta'\le1-c_a\), so the contribution is at most \(e^{(1-c_a)z}\), also bounded by \(Ce^zz^{-1/2}\). For \(z<0\), shift \(\theta'\) by \(\pi\) to obtain \(M_e(z)=M_e(-z)\). Hence \[
c\,e^{|z|}/\sqrt{|z|}
\le M_e(z)\le C\,e^{|z|}/\sqrt{|z|}\quad (|z|\ge1).
\] Taking the ratio yields positive constants bounding \(M_e(2z)/M_e(z)^2\) above and below by multiples of \(\sqrt{|z|}\). Therefore \(V(\mu,p)\to\infty\) for each fixed \(p\ne0\); for \(p=0\) its explicit quadratic formula proves the same.

The compact range of \(p\) admits a uniform finite \(\mu_{\max}\) such that \(V(\mu_{\max},p)>3/\min a\). To prove uniformity, for each \(p_0\) select a finite \(\mu_{p_0}\) giving the strict bound there; continuity gives a neighborhood in \(p\); a finite subcover gives finitely many such \(\mu\), and their maximum works by monotonicity. This step is valid at \(p=0\), so no divergent factor \(p^{-2}\) is being used near zero.

All slopes \(t(\theta';\mu)\) for \(0\le\mu\le\mu_{\max}\) form a compact smooth family with \(P_c(t)>2\). Continuity of \(U\) and the strict inequality \(U(P,J)>2\) established above give a positive minimum of \(U-2\). Choose \(0<\delta_L<1\) smaller than that minimum, and smaller than the positive minimum of \(v_s-2\) on chosen smaller radial boundary collars.

Choose a smooth \(0\le\zeta_L\le1\) as a function of \(v_s\), equal to one for \(v_s\le2+\delta_L/8\) and zero for \(v_s\ge2+\delta_L/4\). Put \[
v_*=2+\delta_L/2,\qquad
\rho=\zeta_L(v_s)^2(v_*-v_s),\qquad
v=v_s+\rho,
\] and define \(\rho=0\) where \(\zeta_L=0\). On its possible nonzero support, \(v_*-v_s\ge\delta_L/4\), so \[
\sqrt\rho=\zeta_L(v_s)\sqrt{v_*-v_s}
\] extends smoothly by zero. Since \(v_s>0\), \(0\le\rho<v_*<3\).

There is a unique \(0\le\mu<\mu_{\max}\) solving \(V(\mu,p)=\rho/a\), by the preceding strict monotonicity and endpoint bounds. Its smoothness at \(\rho=0\) needs a distinct verification. Expanding the finite integral in its convergent power series gives \(M_e(z)=1+z^2/4+O(z^4)\), and therefore \[
V(\mu,p)=\tfrac12d_0^2\mu^2+O(\mu^4p^2)
\] uniformly for bounded \(p\). Thus the signed square root is \(\mu(d_0/\sqrt2+O(\mu^2p^2))\), a smooth odd function of \(\mu\) near zero with positive derivative \(d_0/\sqrt2\). Applying the implicit-function theorem to this function and \(\sqrt{\rho/a}\), already proved smooth, gives smooth \(\mu\) there. For \(\rho>0\), strict positivity of \(\partial_\mu V\) gives the same conclusion. On the overlap the two local inverses agree by uniqueness.

The three cutoff regimes show \(v>2\) everywhere. Where \(\zeta_L\ne0\), \(v\le v_*<2+\delta_L<U(P_c(t),J_c(t))\). Where \(\zeta_L=0\), \(\mu=0,t=t_s,v=v_s>2\), and the original relaxed condition gives \(v<U\). Thus each proposed state lies in the admissible cone.

Define a lifted phase coordinate by \[
\phi(0)=0,\qquad
\frac{d\phi}{d\theta'}=\frac{a(1+t^2)}{2\pi v},
\quad
(a_L,-b_L)=\frac{v(1,t)}{1+t^2}.
\] Since \(\langle t\rangle=t_s\) and \(\langle(t-t_s)^2\rangle=V\), \[
a\langle1+t^2\rangle=a(1+t_s^2)+aV=v_s+\rho=v.
\] Therefore \(\phi(\theta'+2\pi)=\phi(\theta')+1\). Its derivative has a uniform positive lower bound on the compact parameter set; it defines an orientation-preserving circle diffeomorphism with a smooth parameter-dependent inverse. Changing variables proves both original mean components exactly: \[
\int_0^1a_L\,d\phi
=\frac a{2\pi}\int_0^{2\pi}1\,d\theta'=a,
\quad
\int_0^1(-b_L)\,d\phi
=\frac a{2\pi}\int_0^{2\pi}t\,d\theta'=at_s=-b_s.
\] Also \(-b_L/a_L=t\), \(a_L(1+t^2)=v\), so the verified scalar cone inequalities are the desired vector cone gaps. At boundary collars \(\zeta_L=0\), and the reparametrized shear is exactly the original constant-in-\(\phi\) shear. Compactness gives one \(\kappa_L>0\) for all four cone gaps. This completes the complete loop construction, including the mean, boundary, and smoothness claims.

\subsection{Exact insertion into the unchanged original radial coordinate}\label{exact-insertion-into-the-unchanged-original-radial-coordinate}

Let \(\mathcal A,\mathcal B\) be the zero-mean periodic antiderivatives \[
\partial_\phi\mathcal A=-\tfrac12(a_L-a),\qquad
\partial_\phi\mathcal B=\tfrac12E(b_L-b_s).
\] They exist because the right sides have zero \(\phi\)-mean. One explicit construction for any zero-mean \(g\) is \[
\mathcal I g(\phi)=\int_0^\phi g(w)\,dw-
\int_0^1\int_0^s g(w)\,dw\,ds.
\] It is periodic, zero mean, smooth in all parameters, and has derivative \(g\). It vanishes identically when \(g\equiv0\). Therefore these antiderivatives vanish on the boundary collars and extend smoothly by zero outside \(I\).

For an integer \(N>0\), keep the source's exact definitions \[
E_N=E\exp(\mathcal A(X,\eta,N\log X)/N),\qquad
U_N=U+\mathcal B(X,\eta,N\log X)/N.
\] The modified \(E_N\) is positive without approximation. The chain rule in the logarithmic coordinate is \(X\partial_X=D_X+N\partial_\phi\) after substitution. Hence \[
\begin{aligned}
a_N
&=1-2X\partial_X\log E_N\\
&=a-2D_X\mathcal A/N-2\partial_\phi\mathcal A
=a_L-2D_X\mathcal A/N,
\end{aligned}
\] \[
\begin{aligned}
b_N
&=\frac{2X\partial_XU_N}{E_N}\\
&=e^{-\mathcal A/N}
\left(\frac{2D_XU+2\partial_\phi\mathcal B}{E}
+\frac{2D_X\mathcal B}{NE}\right)\\
&=e^{-\mathcal A/N}
\left(b_L+\frac{2D_X\mathcal B}{NE}\right).
\end{aligned}
\] The axial exponential denominator is retained; replacing it by one would destroy this exact shear formula.

On the fixed compact interval from \(X_-\) through the first reserved patch, \(X,E,H=\sqrt{2X}E,L\) have positive lower bounds. The phase \(N\log X\) is independent of \(\eta\), so every fixed number \(m\) of \(\eta\)-derivatives of \(E_N-E,U_N-U,a_N-a_L,b_N-b_L\) is bounded by \(C_m/N\). For the exponential difference, use \[
e^{\mathcal A/N}-1
=\frac{\mathcal A}N\int_0^1e^{s\mathcal A/N}\,ds
\] and then differentiate under its fixed finite integral; all factors of \(\mathcal A\) and its parameter derivatives are uniformly bounded on the periodic compact domain. For radial derivatives, \(r\ge1\) differentiations produce at most \(N^r\) from the phase and retain the initial \(N^{-1}\), giving \(C_{r,m}N^{r-1}\). Thus no small-radial-derivative assertion is used for the modulation.

\subsection{All five moments and the dependence of the integrated stress}\label{all-five-moments-and-the-dependence-of-the-integrated-stress}

Retain exactly \[
M=\int_0^XU\,dx,\quad I=\int_0^X\sqrt{2x}E\,dx,\quad
J=\int_0^XU\sqrt{2x}E\,dx,
\] \[
S=\int_0^X(U^2-E^2/2)\,dx,\qquad
C_p=\int_0^X E^2/(2x)\,dx,\qquad \Pi=\Pi_0+C_p.
\] On the support of the differences, all scalar weights are bounded, since it is a compact interval with positive left endpoint. For instance \[
U_N^2-U^2=(U_N-U)(U_N+U),\qquad
E_N^2-E^2=(E_N-E)(E_N+E).
\] These identities, the product rule, and integration of the uniform \(C_m/N\) profile bounds prove every fixed \(\eta\)-derivative of the full prefix moment difference is \(O_m(N^{-1})\), uniformly in the endpoint \(X\). The pressure datum is the same, so \(\Pi_N-\Pi=C_{p,N}-C_p\) has exactly that bound.

The formulas whose continuity matters are \[
W=1-\frac{2D\eta M+(1-\eta^2)M_\eta}{X},
\] \[
Q_s=-W+
\frac{(1-h)I-D\eta I_\eta-(1-\eta^2)J_\eta+2(h-D)\eta J}{XH},
\] \[
N_s=-WU+
\frac{D(M-\eta M_\eta)+4h\eta S-(1-\eta^2)S_\eta}{X}
+4A\eta\Pi-(1-\eta^2)\Pi_\eta,
\] \[
p_s=(XQ_s/L,\ XN_s/(LE)).
\] All denominators are bounded away from zero on the fixed interval. Differences of reciprocals are exactly \[
E_N^{-1}-E^{-1}=-(E_N-E)/(E_NE),\qquad
H_N^{-1}-H^{-1}=-(H_N-H)/(H_NH).
\] Applying the product rule to the displayed formulas proves the \(m\)-th parameter derivative bound for \(p_{s,N}-p_s\) from the profile and moment bounds through order \(m+1\), with no radial derivative of the differences. This proves the precise assertion needed for the cone stability.

The compact loop has positive four-gap margin. Its values and the original \(p_s\) range in a compact subset of \(a>0\); take a fixed small neighborhood still in \(a>0\), on which the four-gap map has a finite Lipschitz constant. The just-proved \(O(N^{-1})\) errors in \(a,b,p_s\) preserve each gap for sufficiently large finite \(N\). Between \(I\) and the first correction patch, \(E_N=E,U_N=U\), so the shears themselves equal the originals; the \(p_s\) error remains \(O(N^{-1})\). The original strict admissible margins on that compact intervening interval are preserved by the same argument.

\subsection{Nonlinear moment correction with an exact endpoint map}\label{nonlinear-moment-correction-with-an-exact-endpoint-map}

On the first reserved patch, the original profile is \(U=0,E=K(\eta)X^{-1/2-\lambda}\), where fixed \(\lambda>0\) and positive \(K\) have been selected before \(N\). The earlier modulation is zero on this patch. Add two smooth compactly supported \(U\)-bumps and three \(E\)-bumps with five coefficient functions \(c(\eta)\). Each block's bumps are nonnegative, nonzero, and supported on ordered disjoint intervals.

Differentiate the five original moment functionals at zero correction. In row order \((M,J;I,S,C_p)\), the two \(U\)-columns have weights \(1,H=\sqrt{2X}E\), and the three \(E\)-columns have weights \(\sqrt{2X},-E,E/X\). The off-diagonal blocks are zero because \(U=0\) before the correction. Up to the displayed nonzero factors, the exponent lists are \[
(0,-\lambda),\qquad
(1/2,-1/2-\lambda,-3/2-\lambda).
\] All powers within each block are distinct for the fixed \(\lambda>0\). To establish invertibility directly, a nonzero linear combination of \(m\) distinct powers has at most \(m-1\) distinct positive zeros: divide by the least power, use Rolle's theorem on \(m\) putative zeros, and apply induction to the derivative, which has \(m-1\) distinct powers and nonzero coefficients wherever the original nonconstant coefficients were nonzero. Thus the evaluation determinant \(\det[x_j^{\alpha_i}]\) never vanishes for ordered \(0<x_1<\cdots<x_m\). Its sign is constant on that connected set. Multilinearity and Fubini give \[
\det\left[\int x^{\alpha_i}\beta_j(x)\,dx\right]
=\int\det[x_j^{\alpha_i}]\prod_j\beta_j(x_j)\,dx_1\cdots dx_m\ne0.
\] The integrand has a constant nonzero sign on a set of positive measure. This proves each block's rank and hence the full moment Jacobian's rank. Smooth dependence on \(\eta\), compactness of \([-1,1]\), and the derivative identity for the inverse give uniform finite bounds at every fixed parameter order. No uniform estimate as \(\lambda\to0\) is claimed or needed.

The exact moment increment map is \(B(\eta)c+Q_\eta(c,c)\), because \(M,I\) are linear, \(J\) is bilinear in \(U,E\), and \(S,C_p\) are quadratic. Let \(d_N(\eta)\) be the negative moment discrepancy accumulated before the patch; it obeys \(\|d_N\|_{C^m}\le C_m/N\). On the \(C^0\) ball of radius \(r=2\beta_0\|d_N\|_{C^0}\), where \(\beta_0=\|B^{-1}\|_{C^0\to C^0}\), the map \[
\mathcal T(c)=B^{-1}(d_N-Q(c,c))
\] maps the ball to itself and has Lipschitz constant at most \(2\beta_0\kappa_0r\le1/2\), provided \(8\beta_0^2\kappa_0\|d_N\|_{C^0}\le1\). Thus it has a unique fixed point selected by iteration from zero. The derivative \(B+D_cQ(c,c)\) is invertible by the same estimate, giving smooth parameter dependence including one-sided endpoint derivatives. In each preselected finite \(C^m\) norm the identical contraction proof works after increasing \(N\), with the relevant multiplication and bilinear norms, and yields \(C_m/N\) for the same branch.

The assertion that every other fixed derivative is \(O_m(N^{-1})\) can also be justified without demanding simultaneous smallness of infinitely many norms. Differentiate the pointwise equation \(m\) times. The highest coefficient derivative is multiplied by the uniformly invertible \(B+D_cQ(c,c)\). All other terms contain either a derivative of \(d_N\), \(O_m(N^{-1})\), or a lower derivative of \(c\); induction gives \(O_m(N^{-1})\) for them, with the fixed data derivatives included in the constants. The same inverse bound therefore proves the estimate for the highest derivative. Only the finite orders needed for positivity and the cone gaps determine the actual choice of \(N\).

The bump shapes and widths are fixed, so their radial derivatives multiply these small coefficients by fixed constants. The correction is therefore small in the radial derivatives entering \(a,b_s\), in the first parameter derivatives entering \(p_s\), and in all prefix moments. The compact input cone margin on the correction patch persists. Positivity of \(E\) there also persists, after one further finite lower bound on \(N\).

At the right endpoint \(X_{\rm rep}\), the five corrected moments equal the original moments exactly as functions of \(\eta\), while the corrected profiles equal the original profiles for every \(X\ge X_{\rm rep}\). Their derivatives as functions of \(\eta\) also agree. Each difference moment has zero derivative in \(X\) there, by the original moment integrands, so it stays zero at every later radius. The common axis pressure datum gives identical \(\Pi\). The displayed formulas for \(Q_s,N_s,p_s\) and the original incompressibility formula \[
V_0=\frac X L\left(2\eta U-2D\eta M/X-(1-\eta^2)\partial_\eta(M/X)\right)
\] then prove exact equality of these fields at every later radius. This is the exact morphism from the five-moment endpoint data to exterior pressure, velocity, and stress; it proves preservation of the terminal compensation, both unused reserved correction patches, and the entire exterior. The inner collar is unchanged because both modifications start later and pressure was integrated from the same axis datum.

\subsection{Edge direction and the common weight}\label{edge-direction-and-the-common-weight}

Appendix C leaves both original collars unchanged. The input edge facts it uses are: at the inner edge, with \(y_a=\log(X/X_a)\), the stress is \[
T_0=e^{-t_1^2/y_a^2}g_a(y_a)B_0(X,\eta),\qquad
g_a(0)>0,\quad B_0(X_a,\eta)=F p_{s,r}\ne0;
\] its limiting direction is a positive multiple of the limiting shear \((a,-b_s)=a(1,t_s)\). At the outer edge, with \(y_b=\log(X_b/X)\), the component factorization has common exponential \(e^{-4/y_b^2}\), angular prefactor \(y_b^{-3}b_\theta\) with \(b_\theta(0,\eta)>0\), and the direction \[
n=\frac{(b_\theta,y_b^6b_z)}
{\sqrt{b_\theta^2+y_b^{12}b_z^2}}.
\] The original collar analysis supplies smooth factors and the fixed-order derivative bounds. Their derivation is assigned to the profiles lane; their exact use here is as follows.

The positive factors cancel in \(n=T_0/|T_0|\), giving a smooth direction through both endpoints. At the inner edge \(n_z-t_sn_\theta=0\) and \(n_\theta+t_sn_z>0\). At the outer edge \(n=(1,0)\), \(t_s=0\), giving the same two conclusions. In the interior, \[
n_\theta+t_sn_z=\frac{P_c-v_s}{|p_s-(a,-b_s)|}>0,\qquad
n_z-t_sn_\theta=\frac{J_c}{|p_s-(a,-b_s)|}.
\] The admissible cone makes \[
2-(v_s-2)\frac{(n_z-t_sn_\theta)^2}{(n_\theta+t_sn_z)^2}>0.
\] Both the denominator and this gap extend continuously to the closed annulus, where the gap has value two at both radial edges. Their positive compact minima permit one \(0<\kappa<2\) with \[
n_\theta+t_sn_z\ge\kappa,\quad
(v_s-2)(n_z-t_sn_\theta)^2
\le(2-\kappa)(n_\theta+t_sn_z)^2.
\] The positive lower bounds on \(F,a,v_s-2\) in the interior and the original endpoint bounds likewise extend uniformly by compactness.

The original stress identity \(T_0=F(p_s-(a,-b_s))\) proves nonvanishing throughout the open shell: if it vanished, \(P_c=v_s\), contrary to admissibility. The outside stress remains zero by the exact endpoint restoration and unchanged inner region. Define the original proposed common edge weight \[
\zeta(X)=\exp(-t_1^2/y_a^2-4/y_b^2),\qquad
\delta=\min(1,y_a,y_b),
\] extended by zero outside the open shell. On a sufficiently short inner collar, the second exponential is smooth and bounded above and below by positive constants, while \(g_a,B_0\) have the required nonzero value bounds and fixed derivative bounds. The inner factorization therefore proves \(|T_0|\ge c\zeta\) and \(|\partial^IT_0|\le C_I\zeta\delta^{-m_I}\). On a short outer collar the first exponential is similarly a smooth positive factor; \(y_b^{-3}b_\theta\) has a lower positive bound and its derivatives have only finite inverse powers of \(y_b\). The supplied outer factor derivative bounds give the same conclusion. On the remaining compact interior both \(\zeta\) and \(|T_0|\) have positive minima and every fixed derivative is bounded, so the same pair of inequalities follows after increasing constants. These three regions cover the annulus. The elementary exponential argument in the first file proves smooth zero extension of all stress derivatives at both boundaries.

All remaining completion claims in C.3 are exact preservation statements: the axis Cartesian formulas remain unchanged on the analytic rectangle; the pressure at infinity and exterior moment values remain unchanged by the proved five-moment map; the exterior heat function is unchanged pointwise; and on \(X\ge X_v=X_{\rm end}\) the original \(U=M=0\) gives \(V_0=0\) directly by its formula. The third and fourth reserved intervals remain the exact original \(U=0,E=c_{\rm patch}(1+\eta^2)^{-1}X^{-1/2-\lambda}\) intervals, with their original order. Fixing this one finite \(N\) before the physical scale \(q\), all dyadic bands, and all later correction stages proves the claimed independence of the resulting derivative constants from those later parameters.


% --- included proof body: mean_corrections_body.tex ---
% Standalone inclusion body. Requires amsmath and amssymb only.
% No custom macros; source/source-code dependency ledger is external.
% Generated from the independently authored AUDIT.md by prepare_tex.py.
\section{Independent audit of compactly supported mean corrections}

Source: \emph{Finite Time Blowup for Navier-Stokes}, Section 8, printed pp.87-100, released at the source PDF; its URL is recorded in the accompanying provenance ledger. Verified SHA-256: \texttt{8c8a94ad\allowbreak 9ac824c8\allowbreak b605b982\allowbreak 7cadf7be\allowbreak aca48bd1\allowbreak 0b380de3\allowbreak cfc872a2\allowbreak c37afa81}. All statements and proofs in this section were read, and every assigned page was rendered and visually inspected.

This audit finds no contradiction in the internal Section 8 calculations. The proofs below establish the local operator facts and algebra explicitly. The dependency ledger identifies the exact earlier data used; the existence of the actual background and the full iteration are audited elsewhere. In particular, the checks here do not establish the Navier-Stokes endpoint.

\subsection{Original coordinates, operators, and spaces}

Keep the paper's \(A=1/2+h\), \(D=1/2-h\), fixed \(h>0\), \(Q=2^{-\ell}\), \(\epsilon=Q^h\), \(S_*=\ell^2\), and

\[
R=Q^{-1/2}r,\quad Z=Q^{-D}z,\quad T=Q^{-1}(1-t).
\]

The positive function \(q=q(z,t)\) is independent of radius. The active shell is \(X_a<X<X_b\), where \(X=QR^2/(2q)\). In a fixed band \(q/Q\asymp1\); hence the radii are bounded above and bounded away from zero. All radial integrals below run from \(0\) to \(\infty\), except where their bounds are written otherwise. Auxiliary Haar measure has total mass one on \(\mathbb T^2=\mathbb R^2/\mathbb Z^2\). Angular averaging means componentwise averaging in the cylindrical frame.

The exact constants and vectors are

\[
J_g=\begin{pmatrix}3&1\\1&5\end{pmatrix},\quad
\Lambda_g=4-\sqrt2,\quad T_g=4+\sqrt2,\quad
b_g=\sqrt2-1,
\]
\[
v_r=(1,-b_g),\quad v_t=(b_g,1),\quad
\rho_g=\frac{\log\Lambda_g}{\log T_g},\quad
\kappa_s=10^{-5},\quad
d_r=2((1+h)\rho_g-h\kappa_s)>0.
\]

The phase map is \(Y=v_r r^{d_r}+v_t t\pmod {\mathbb Z^2}\). For an independent auxiliary variable \(Y\), differentiation after evaluation is

\[
\mathsf r=\partial_r+d_r r^{d_r-1}v_r\cdot\partial_Y,
\qquad
\mathsf t=\partial_t+v_t\cdot\partial_Y.
\]

Both commute with \(\partial_z\), and with each other: the coefficient of the auxiliary radial derivative depends only on \(r\), while the coefficient of the auxiliary temporal derivative is constant. Direct multiplication gives \(J_gv_r=\Lambda_gv_r\) and \(J_gv_t=T_gv_t\).

For a fixed common covering level \(i_0\), let \(y=J_g^{i_0}Y\), \(L=v_r\cdot\partial_y\), \(N=v_t\cdot\partial_y\), \(M=\Lambda_g^{i_0}Q^{d_r/2}\), and \(c=T_g^{i_0}Q^{1+h}\). Then

\[
D_r=\partial_R+Md_rR^{d_r-1}L,\quad
D_z=\epsilon\partial_Z,\quad
t_*=-\epsilon\partial_T+cN,
\quad \Delta_0=D_r^2+R^{-1}D_r+D_z^2.
\]

The integer \(i_0\), \(Q\), \(\epsilon\), \(M\), and \(c\) remain fixed under coefficient differentiation. At the band level,

\[
i(\ell)=\left\lfloor\log_{T_g}(Q^{-1-h}/S_*)\right\rfloor
\]

implies \(T_g^{-1}Q^{-1-h}/S_*<T_g^{i(\ell)}\le Q^{-1-h}/S_*\). Raising to \(\rho_g\) and using the displayed definition of \(d_r\) proves

\[
c\asymp S_*^{-1},\qquad
M\asymp\epsilon^{-\kappa_s}S_*^{-\rho_g},
\qquad M^{-1}\lesssim\epsilon^{\kappa_s}S_*^{\rho_g}.
\]

Bounded changes of covering level multiply these expressions by fixed constants.

The class \(\mathcal M_\alpha\) consists of angularly independent, compatible torus coefficients supported in the shell with smooth zero extension and, for every fixed multi-index \(I\) in \((R,Z,T,y_1,y_2)\), bounds

\[
|\partial^If|\le C_{j,I}\epsilon^\alpha S_*^{B_{j,I}}\zeta\delta^{-C_{j,I}'},
\quad
\zeta=\exp(-a_a/d_a^2-a_b/d_b^2),
\]
\[
d_a=\log(X/X_a),\quad d_b=\log(X_b/X),\quad
\delta=\min(1,d_a,d_b),\qquad a_a,a_b>0.
\]

The powers may depend on the fixed derivative order and correction stage, but not on the band or point. The class \(\mathcal S_\alpha\) has the same slow derivative bounds for functions of \((Z,T)\), without the radial weight. Combining powers into a common larger power is legitimate since \(S_*\ge1\), \(\delta\le1\). This does not change the source object or its exponent.

Leibniz' rule, \(0<\zeta\le1\), and the bounded shell imply \(\mathcal M_a\mathcal M_b\subset\mathcal M_{a+b}\). Slow coefficient derivatives preserve the exponent; \(D_r\) loses at most \(\kappa_s\), \(D_z\) gains one, the slow part of \(t_*\) gains one, and \(cN\) preserves it. A slow field has no \(cN\) contribution. Each assertion follows term by term from the displayed operators. Multiplying an interior-supported coefficient by a bounded smooth base coefficient preserves its weighted class even when that base coefficient has no global shell weight.

\subsection{Exact radial range, inverse, obstruction, and projection}

Fix \((Z,T)\), a covering level, and \(e\in\{0,1,2\}\). Write

\[
\mathcal D_e=D_r+e/R.
\]

Let \(\mathcal C\) be the smooth real functions of \(R>0,y\in\mathbb T^2\) supported in the closed active shell and smoothly zero across its edges. Define the conjugation

\[
(\Gamma H)(R,w)=H(R,w+MR^{d_r}v_r).
\]

Its inverse replaces the plus sign by a minus sign. The chain rule gives the exact intertwining identity

\[
\partial_R\Gamma(R^e\phi)=\Gamma(R^e\mathcal D_e\phi).
\tag{MC1}
\]

Consequently the full obstruction is the function on the entire torus

\[
(K_ef)(w)=\int R^ef(R,w+MR^{d_r}v_r)\,dR.
\tag{MC2}
\]

A compactly supported solution of \(\mathcal D_e\phi=f\) exists exactly when \(K_ef=0\) as a torus function. Necessity follows by integrating (MC1), since \(R^e\phi\) vanishes at both radial endpoints. For sufficiency define

\[
\phi(R,y)=R^{-e}\int_0^R(R')^e
f(R',y+M((R')^{d_r}-R^{d_r})v_r)\,dR'.
\tag{MC3}
\]

It vanishes below the shell. Above the shell it is \(R^{-e}(K_ef)(y-MR^{d_r}v_r)=0\). All derivatives extend by zero, and (MC1) proves the equation. If two compactly supported solutions exist, their difference has \(\partial_R\Gamma(R^e\phi)=0\); the constant is zero below the shell, proving uniqueness.

The obstruction map is onto \(C^\infty(\mathbb T^2)\): given \(F(w)\), choose any smooth interior bump \(h(R)\) with integral one and take \(f(R,y)=R^{-e}h(R)F(y-MR^{d_r}v_r)\). Substitution in (MC2) gives \(K_ef=F\). Thus the strongest exact relationship, at each fixed slow point, is the exact sequence

\[
0\longrightarrow\mathcal C
\xrightarrow{\mathcal D_e}\mathcal C
\xrightarrow{K_e}C^\infty(\mathbb T^2)
\longrightarrow0.
\tag{MC4}
\]

This identifies all obstructions, not only their Haar averages.

Choose the paper's slow cutoff \(\chi_m(X)\), equal to zero near the inner shell edge and one near the outer edge, with derivative supported on a fixed interior interval. For \(g=R^ef\), define

\[
Ig(R,y)=\int_0^R g(R',y+M((R')^{d_r}-R^{d_r})v_r)\,dR',
\]
\[
Jg(R,y)=\int_0^\infty g(R',y+M((R')^{d_r}-R^{d_r})v_r)\,dR',
\]
\[
T_ef=R^{-e}(Ig-\chi_mJg),\qquad
A_ef=R^{-e}(\partial_R\chi_m)Jg.
\tag{MC5}
\]

The conjugation yields \(D_rIg=g\), \(D_rJg=0\). Since \(D_r\chi_m=\partial_R\chi_m\), the product rule yields

\[
\mathcal D_eT_e=1-A_e.
\tag{MC6}
\]

Moreover \(K_eA_ef=K_ef\), because after conjugation the factor \(K_ef(w)\) is independent of \(R\) and \(\int\partial_R\chi_m\,dR=1\). Thus \(A_e^2=A_e\), \(T_e\mathcal D_e=1\) on \(\mathcal C\), and \(T_e\) is the exact inverse of \(\mathcal D_e\) on \(\ker K_e\). The paper's inverse is an explicitly split inverse with a specified obstruction projector.

Haar averaging (MC2) gives

\[
\langle K_ef\rangle=\int R^e\langle f\rangle\,dR.
\tag{MC7}
\]

The zero weighted Haar moment is therefore necessary but in general insufficient for exact compact inversion. As an explicit smooth counterexample to that stronger, unclaimed assertion, take a nonzero integer \(k\), the bump \(h\) above, and

\[
f_M(R,y)=e^{-M^2}R^{-e}h(R)
\cos(2\pi k\cdot(y-MR^{d_r}v_r)).
\]

Its Haar mean vanishes at every radius, while \(K_ef_M=e^{-M^2}\cos(2\pi k\cdot w)\ne0\). The family has all fixed coefficient derivative bounds: each derivative introduces at most a fixed power of \(M\), absorbed by \(e^{-M^2}\). With the original scale relation for \(M\), these bounds are flat in every power of \(\epsilon\). The counterexample does not contradict Lemma 8.2; it proves why its explicitly retained flat remainder cannot be deleted.

\subsection{Primitive estimates and the full flatness argument}

Set \(U=R^{d_r}\) and

\[
F(U,Z,T,y)=\frac{g(R,Z,T,y)}{d_rR^{d_r-1}}.
\]

Use its smooth zero extension. The change of variables has uniformly bounded derivatives on the shell. Put

\[
A_-(U,y)=\int_{-\infty}^0 F(U+u,y+Muv_r)\,du,
\qquad
A_+(U,y)=\int_0^\infty F(U+u,y+Muv_r)\,du.
\]

Then \(Ig=A_-\), \(Jg=A_-+A_+\), and \(Ig-\chi_mJg=(1-\chi_m)A_--\chi_mA_+\). In these integrals the shift \(Muv_r\) does not depend on \(U,Z,T\). Hence every fixed coefficient derivative differentiates only \(F\) and the cutoff; it introduces no factor \(M\). All integrations have bounded effective length. Since the source is smoothly zero extended, differentiating moving support creates no boundary term.

For completeness, the precise edge inequality is as follows. For any fixed real \(B\ge0\),

\[
\frac{d}{ds}\log(e^{-a/s^2}s^{-B})=2a/s^3-B/s>0
\]

whenever \(0<s<\sqrt{2a/B}\), with no upper restriction needed when \(B=0\). Near the inner edge, only \(A_-\) occurs and \(0<d_a(R')\le d_a(R)\) on its effective segment. The factor associated with \(d_b\) stays between positive constants there. This inequality bounds the complete differentiated input weight by \(C\zeta(R)\delta(R)^{-B'}\). The outer edge uses only \(A_+\) and the same argument with \(d_b\). On the fixed interior region both radial edge distances are bounded below, so the output weight is bounded below and the unweighted bound suffices. Multiplication by \(R^{\pm e}\) and the \(R\leftrightarrow U\) chain rule cost fixed constants. This proves \(T_e:\mathcal M_\alpha\to\mathcal M_\alpha\), including every fixed derivative, shell support, and projection of slow support.

The directional small divisors are controlled without a numerical approximation. For \(k=(k_1,k_2)\ne0\),

\[
v_r\cdot k=(k_1+k_2)-\sqrt2 k_2,
\]

and its product with \((k_1+k_2)+\sqrt2 k_2\) is the integer \((k_1+k_2)^2-2k_2^2\ne0\). Nonzero follows from the irrationality of \(\sqrt2\); if \(k_2=0\), it follows from \(k_1\ne0\). The conjugate has absolute value at most \(C|k|\). Therefore \(|v_r\cdot k|^{-1}\le C(1+|k|)\). The same calculation for \(v_t\cdot k=(k_2-k_1)+\sqrt2 k_1\) gives the identical bound.

For \(p\ge1\), the inverse \(L^{-p}\) on zero-mean torus functions is the multiplier \((2\pi i v_r\cdot k)^{-p}\) at \(k\ne0\), zero at \(k=0\). For a \(C^{m+p+3}\) function, integrate its Fourier coefficient \(m+p+3\) times in a coordinate with \(|k_j|\ge |k|/\sqrt2\). A derivative of order at most \(m\), followed by the multiplier bound, gives a summand bounded by \(C\|H\|_{C^{m+p+3}}(1+|k|)^{-3}\). The series over \(\mathbb Z^2\) converges, proving

\[
\|L^{-p}H\|_{C^m}\le C_{m,p}\|H\|_{C^{m+p+3}}.
\tag{MC8}
\]

It also proves commutation with slow derivatives, smoothness, and reality. There is no assertion of constants uniform in \(p\).

Under \(\int R^e\langle f\rangle\,dR=0\), the integral of the zero Fourier mode of \(F\) is zero. For \(F^\circ=F-\langle F\rangle\), differentiation in \(u\) gives

\[
\frac d{du}L^{-1}F^\circ(U+u,y+Muv_r)
=\partial_UL^{-1}F^\circ(U+u,y+Muv_r)+MF^\circ(U+u,y+Muv_r).
\]

Its full-line integral is zero by smooth compact support. Iterating this identity gives the exact equality

\[
Jg=(-M^{-1})^p\int_{\mathbb R}
\partial_U^pL^{-p}F^\circ(U+u,y+Muv_r)\,du.
\tag{MC9}
\]

Each repetition has the same minus sign. All omitted zero-mode integrals, including their slow derivatives, are exactly zero. Fixed-shift differentiation and (MC8) bound any further derivative of (MC9). For an output derivative \(I\), radial input derivatives through \(I_R+p\) and torus derivatives through \(|I_y|+p+3\), with the requested slow derivatives, suffice. The coefficient \(\partial_R\chi_m\) has fixed interior support. Thus

\[
|\partial^IA_ef|\le C_{j,I,p}
\epsilon^{\alpha+p\kappa_s}S_*^{B_{j,I,p}}
\zeta\delta^{-B'_{j,I,p}},\qquad
\langle A_ef\rangle=0.
\tag{MC10}
\]

If \(f\) is independent of \(y\), its full characteristic integral is its ordinary weighted radial integral. Under the moment condition the remainder is identically zero, not merely small.

To pass from (MC10) to physical flatness, fix the physical output normalization and derivative order before choosing \(p\). Its graph derivatives require only finitely many coefficient derivatives and cost some finite power \(q^{-L}\) times a power of \(S_*\); \(L\) depends on the output order, not on \(p\). The cutoff remainder is interior supported, so no boundary-weight loss remains there. For prescribed \(N\ge0\), choose \(p\) with \(h(\alpha+p\kappa_s)>N+L\). Since \(q\asymp Q\) and \(S_*=\ell^2\), the positive excess in the exponent dominates every fixed power of \(\ell\), proving an \(O(q^N)\) bound. This proves every fixed physical derivative is flat. Finite regularity alone would not imply this all-orders conclusion.

\subsection{Mean momentum equations derived in the cylindrical frame}

Let \(e_r=(\cos\theta,\sin\theta,0)\), \(e_\theta=(-\sin\theta,\cos\theta,0)\), \(e_z=(0,0,1)\). Then \(\partial_\theta e_r=e_\theta\), \(\partial_\theta e_\theta=-e_r\). For a divergence-free vector \(U=(U_r,U_\theta,U_z)\),

\[
(D_r+R^{-1})U_r+R^{-1}\partial_\theta U_\theta+D_zU_z=0.
\]

Expand transport using this identity. For example,

\[
(D_r+R^{-1})U_r^2+R^{-1}\partial_\theta(U_\theta U_r)+D_z(U_zU_r)
=U_rD_rU_r+R^{-1}U_\theta\partial_\theta U_r+U_zD_zU_r,
\]

because the extra factor multiplying \(U_r\) is the displayed divergence. The frame derivative then adds \(-U_\theta^2/R\). In the angular component it adds \(U_rU_\theta/R\), so its radial flux has weight two. In the axial component the frame contributes nothing. Therefore the transport components are

\[
\begin{aligned}
(U\cdot\nabla U)_r&=(D_r+R^{-1})U_r^2+R^{-1}\partial_\theta(U_\theta U_r)+D_z(U_zU_r)-U_\theta^2/R,\\
(U\cdot\nabla U)_\theta&=(D_r+2R^{-1})(U_rU_\theta)+R^{-1}\partial_\theta U_\theta^2+D_z(U_zU_\theta),\\
(U\cdot\nabla U)_z&=(D_r+R^{-1})(U_rU_z)+R^{-1}\partial_\theta(U_\theta U_z)+D_zU_z^2.
\end{aligned}
\tag{MC11}
\]

The cylindrical vector Laplacian is

\[
\begin{aligned}
(\Delta U)_r&=(\Delta_0+R^{-2}\partial_\theta^2-R^{-2})U_r-2R^{-2}\partial_\theta U_\theta,\\
(\Delta U)_\theta&=(\Delta_0+R^{-2}\partial_\theta^2-R^{-2})U_\theta+2R^{-2}\partial_\theta U_r,\\
(\Delta U)_z&=(\Delta_0+R^{-2}\partial_\theta^2)U_z.
\end{aligned}
\]

These follow by applying \(\partial_\theta^2\) to each component times its frame vector, which produces both the \(\mp2\partial_\theta\) terms and \(-U_{r,\theta}\). Their angular averages retain exactly the radial/angular \(-R^{-2}\) terms.

Now set \(U=(b+\beta,V+v,G+\gamma)+w\), \(\langle w\rangle_\theta=0\), and \(W^{ab}=\langle w_aw_b\rangle_\theta\). The identity

\[
\langle U_aU_b\rangle_\theta=(U_a-w_a)(U_b-w_b)+W^{ab}
\]

keeps the full actual wave covariance. No curl term can be removed from \(W\). Subtract the pure base residual, whose stress contribution is \(-(D_r+2/R)\Sigma_\theta\) in the angular row and \(-(D_r+1/R)\Sigma_z\) in the axial row. Retain its separate flat part. The remaining residual is \((D_rp_m-g_r,E_\theta,E_z)\), with

\[
\begin{aligned}
E_\theta={}&t_*v+(D_r+2/R)(bv+\beta V+\beta v+W^{r\theta})\\
&+D_z(Gv+V\gamma+\gamma v+W^{z\theta})
-\epsilon(\Delta_0-R^{-2})v-(D_r+2/R)\Sigma_\theta,\\
E_z={}&t_*\gamma+(D_r+1/R)(b\gamma+\beta G+\beta\gamma+W^{rz})\\
&+D_z(2G\gamma+\gamma^2+W^{zz}+p_m)
-\epsilon\Delta_0\gamma-(D_r+1/R)\Sigma_z,\\
g_r={}&-t_*\beta-(D_r+1/R)(2b\beta+\beta^2+W^{rr})\\
&-D_z(b\gamma+G\beta+\beta\gamma+W^{zr})
+(2Vv+v^2+W^{\theta\theta})/R
+\epsilon(\Delta_0-R^{-2})\beta.
\end{aligned}
\tag{MC12}
\]

Every sign follows from (MC11) and subtracting the radial pressure derivative to define \(g_r\). These are precisely the checked mathematical identities in source \(8.3\).

\subsection{Pressure, solenoidal axial realization, and dimensions}

Choose the reserved interior bump \(\widehat\rho\in C_c^\infty(I_m)\), \(I_m=\{\sqrt{2X}:X\in I_{\rm mean}\}\), with integral one. In physical variables \(\rho_{\rm phys}=q^{-1/2}\widehat\rho(r/\sqrt q)\), and in a chart \(\rho=Q^{1/2}\rho_{\rm phys}\). Changing variable \(r=\sqrt q\,x\) proves both physical and chart radial integrals of the bump equal one.

Define
\[
P=\int\langle g_r\rangle\,dR,\qquad p_m=T_0(g_r-\rho P).
\]

The input has zero radial Haar integral. Equation (MC6) gives exactly

\[
D_rp_m-g_r=-\rho P-A_0(g_r-\rho P),
\qquad
\partial_R\langle p_m\rangle=\langle g_r\rangle-\rho P.
\tag{MC13}
\]

The map \(g_r\mapsto p_m\) is linear because the bump is fixed. Radial integration maps \(\mathcal M_\alpha\to\mathcal S_\alpha\): for each fixed derivative, \(\zeta\delta^{-B}\) is bounded and the effective interval has bounded length. The interior bump then maps \(P\) back into \(\mathcal M_\alpha\), and Section 3 proves the pressure estimate and its difference estimate. Its cutoff remainder has zero Haar mean and is flat by (MC10).

For a desired axial increment \(\gamma_d\) with \(\int R\langle\gamma_d\rangle\,dR=0\), define \(\Psi_*=T_1\gamma_d\) and

\[
\Delta\beta=-\epsilon\partial_Z\Psi_*,\qquad
\Delta\gamma=(D_r+1/R)\Psi_*=\gamma_d-A_1\gamma_d.
\tag{MC14}
\]

Its divergence is

\[
(D_r+1/R)\Delta\beta+D_z\Delta\gamma
=-\epsilon(D_r+1/R)\partial_Z\Psi_*+
\epsilon\partial_Z(D_r+1/R)\Psi_*=0.
\]

This cancellation uses commutation of the differential operators on the already defined potential. It does not commute \(\partial_Z\) through the integral operator, whose moving cutoff may depend on \(Z\). The averaged axial increment equals \(\langle\gamma_d\rangle\) at each radius, and its weighted integral remains zero. The class bounds are \(\Psi_*,\Delta\gamma\in\mathcal M_\alpha\), \(\Delta\beta\in\mathcal M_{\alpha+1}\). If \(\gamma_d\) is torus independent, the axial realization is pointwise exact.

All physical scales can be checked from one change of variable. If \(f_*=Q^af_{\rm phys}\), then

\[
\int R^ef_*\,dR=Q^{a-(e+1)/2}\int r^ef_{\rm phys}\,dr,
\qquad
T_ef_*=Q^{a-1/2}r^{-e}I_{c,\rm phys}(r^ef_{\rm phys}).
\tag{MC15}
\]

Use \(u_*=Q^Au_{\rm phys}\), \(g_*=Q^{2A+1/2}g_{\rm phys}\), \(p_*=Q^{2A}p_{\rm phys}\), \(\Psi_*=Q^{A-1/2}\Psi_{\rm phys}\). In particular, normalizing \(-\partial_z\Psi_{\rm phys}\) yields \(-Q^{1/2-D}\partial_Z\Psi_*=-\epsilon\partial_Z\Psi_*\), preserving the minus sign. The pressure primitive scale is \(Q^{2A}\). The resulting moment exponents are

\begin{description}
\item[] Quantity: Chart value relative to physical value.
\item[] \(M_\theta=\int R^2\langle v\rangle\,dR\): \(Q^{A-3/2}\).
\item[] \(M_z=\int R\langle\gamma\rangle\,dR\): \(Q^{A-1}\).
\item[] \(P=\int\langle g_r\rangle\,dR\): \(Q^{2A}\).
\item[] \(c_\rho=\frac12\int R^2\rho\,dR\): \(Q^{-1}\).
\item[] \(J_\theta\) defined below: \(Q^{2A-3/2}\).
\item[] \(J_z\) defined below: \(Q^{2A-1}\).
\end{description}

The second term in \(J_z\), involving \(R^2g_r\), has exponent \(2A+1/2-3/2=2A-1\), exactly matching its velocity-product term. Thus the fifth moment row has a consistent physical dimension.

\subsection{Integrated residuals and both covariance targets}

Let bars denote Haar averages. Translation invariance on the torus proves \(\overline{D_rf}=\partial_R\bar f\), \(\overline{t_*f}=-\epsilon\partial_T\bar f\), and the corresponding second-derivative statements. For any smoothly supported scalar \(h\),

\[
\int R^e(\partial_R+e/R)h\,dR=\int\partial_R(R^eh)\,dR=0.
\]

For the angular radial viscosity, retaining all three terms,

\[
\int R^2h_{RR}\,dR=2\int h\,dR,\quad
\int Rh_R\,dR=-\int h\,dR,\quad
\int(-h)\,dR=-\int h\,dR.
\]

Their sum is zero. For the axial viscosity, \(\int Rh_{RR}\,dR=-\int h_R\,dR\), which cancels the remaining \(\int h_R\,dR\). The time and axial-viscosity contributions to the two residual moments are respectively

\[
-\epsilon\partial_TM_\theta-\epsilon^3\partial_Z^2M_\theta,
\qquad
-\epsilon\partial_TM_z-\epsilon^3\partial_Z^2M_z.
\]

They vanish when the two moments are zero at every slow point. The stress divergences vanish by their own compact support. Finally, compact support and (MC13) give

\[
\int R\bar p_m\,dR=-\frac12\int R^2\partial_R\bar p_m\,dR
=-\frac12\int R^2\bar g_r\,dR+c_\rho P.
\tag{MC16}
\]

Set

\[
J_\theta=\int R^2\overline{Gv+V\gamma+\gamma v+W^{z\theta}}\,dR,
\]
\[
J_z=\int R\overline{2G\gamma+\gamma^2+W^{zz}}\,dR
-\frac12\int R^2\bar g_r\,dR,
\qquad c_\rho=\frac12\int R^2\rho\,dR.
\]

Every surviving term of (MC12) is now accounted for, and the exact identities are

\[
\int R^2\bar E_\theta\,dR=\epsilon\partial_ZJ_\theta,
\qquad
\int R\bar E_z\,dR=\epsilon\partial_Z(J_z+c_\rho P).
\tag{MC17}
\]

The product \(c_\rho P\) must stay inside the derivative; \(c_\rho\) changes with the moving scale \(q\). These are the complete compatibility moments of the two torus-independent weighted radial equations.

Use the fixed interior bumps \(\sigma_\theta,\sigma_z\) with \(\int R^2\sigma_\theta\,dR=\int R\sigma_z\,dR=1\). Their physical versions have the powers \(q^{-3/2}\widehat\sigma_\theta(r/\sqrt q)\) and \(q^{-1}\widehat\sigma_z(r/\sqrt q)\). Define
\[
H_\theta=-T_2(\bar E_\theta-\sigma_\theta\epsilon\partial_ZJ_\theta),
\qquad
H_z=-T_1(\bar E_z-\sigma_z\epsilon\partial_Z(J_z+c_\rho P)).
\]

Their sources have zero weighted radial integrals by (MC17), and are torus independent, so their \(A_e\) remainders vanish exactly. Therefore

\[
(D_r+2/R)H_\theta=-\bar E_\theta+\sigma_\theta\epsilon\partial_ZJ_\theta,
\quad
(D_r+1/R)H_z=-\bar E_z+\sigma_z\epsilon\partial_Z(J_z+c_\rho P).
\tag{MC18}
\]

The minus signs are required because adding \(W^{r\theta}\) or \(W^{rz}\) contributes a positive divergence in (MC12). Support at a fixed \((Z,T)\) cannot be created when the whole radial source vanishes there. These are exact radial stress inverses with the two explicitly retained bump residuals.

\subsection{Fast-time inverse and chart agreement}

The same Diophantine proof used for \(L\) gives the unique smooth zero-Haar-mean solution of \(N\phi=F\), for \(\bar F=0\):

\[
N^{-1}F=\sum_{k\in\mathbb Z^2\setminus\{0\}}
\frac{\widehat F(k)}{2\pi i\,v_t\cdot k}e^{2\pi ik\cdot y},
\qquad
\|N^{-1}F\|_{C^m}\le C_m\|F\|_{C^{m+4}}.
\tag{MC19}
\]

All nonzero modes have nonzero divisors. A kernel element therefore has only a constant mode, and the zero-mean condition removes it. The four-derivative count is the \(p=1\) instance of (MC8). This operator is local in the independent variables \(R,Z,T\), hence preserves their closed support and smooth zero extension. It need not preserve support within a proper subset of the torus, which is allowed for \(\mathcal M_\alpha\).

Define \(P_iF(Y)=F(J_g^iY)\). A mode \(k\) pulls back to \((J_g^i)^Tk\). This lattice map is injective because \(\det J_g=14\ne0\), and hence it preserves the zero Fourier coefficient and Haar average. The exact chain rule and the eigenvector identity prove

\[
N_{\rm abs}P_i=T_g^iP_iN_i,
\qquad N_{\rm abs}^{-1}P_i=T_g^{-i}P_iN_i^{-1}
\tag{MC20}
\]

on zero-mean functions. The inverse identity also follows from uniqueness, so it does not require the covering to be one-to-one.

Define the physical tangential update as \(-N_{\rm abs}^{-1}(E_{\theta,\rm phys}^\circ,E_{z,\rm phys}^\circ)\). Its chart representative has multiplier

\[
Q^{A-(2A+1/2)}T_g^{-i_0}=Q^{-1-h}T_g^{-i_0}=c^{-1}.
\]

Thus \(\Delta v=-c^{-1}N^{-1}E_\theta^\circ\), \(\gamma_d=-c^{-1}N^{-1}E_z^\circ\). They belong to \(\mathcal M_\alpha\) when \(E_\theta,E_z\) do, since \(c^{-1}\lesssim S_*\). Both have zero Haar mean at every radius. Realize \(\gamma_d\) by (MC14), and put \(a=-A_1\gamma_d\). Then

\[
cN\Delta v=-E_\theta^\circ,
\qquad cN\Delta\gamma=-E_z^\circ+cNa.
\tag{MC21}
\]

Translations commute with \(N\), and the cutoff is torus independent, so there is no additional commutator. The remainder \(cNa\) is flat by (MC10), with one more requested torus derivative. The slow-time contribution \(-\epsilon\partial_T\) is retained and has one extra power of \(\epsilon\). The actual velocity is divergence-free and preserves both moments because its angular and axial Haar means vanish pointwise.

Radial chart agreement is equally explicit. Pulling a physical integrand \(F(J_g^iY)\) along the absolute radial shift gives

\[
F(J_g^iY+\Lambda_g^i((r')^{d_r}-r^{d_r})v_r),
\]

which is exactly the chart shift after \(r=\sqrt Q R\). Together with (MC15), this proves the radial operators in different charts define one physical field. Their integer indices are held fixed on each common neighborhood. No derivative of a floor function is taken.

\subsection{Explicit five-dimensional moment correction}

Preserve the exact reserved-patch base in the paper's variable \(x=r/\sqrt q\in I_m\):

\[
G_q=0,\qquad V_q=a(\eta)x^{-1-2\lambda},\qquad
a(\eta)=2^{1/2+\lambda}c_{\rm patch}(1+\eta^2)^{-1},
\qquad |a(\eta)|\ge a_0>0,\quad\lambda>0.
\tag{MC22}
\]

This is the specific input base whose construction is assigned to the base audit. The following gives the complete inverse on that specified patch, including a formula beyond just a nonzero determinant.

Choose an interior \(x_0>0\), a small fixed \(d>0\), and a nonnegative nonzero smooth bump \(\eta_0\) sufficiently narrow that its three copies

\[
\eta_j(x)=e^{-jd}\eta_0(e^{-jd}x),\qquad j=0,1,2,
\]

have disjoint supports in \(I_m\). Such a choice exists: first choose \(d\) so \(x_0,x_0e^d,x_0e^{2d}\) remain in the open interval; their finite positive separations then permit a sufficiently small common relative support width. For every real \(p\), define \(\mu_p=\int x^p\eta_0(x)\,dx>0\). The exact change of variable gives

\[
\int x^p\eta_j(x)\,dx=\mu_pe^{jdp}.
\tag{MC23}
\]

The required five rows are

\[
\int R^2\Delta v=0,\quad\int R\gamma_d=0,\quad
\int(2V/R)\Delta v=-P,
\]
\[
\int R^2(G\Delta v+V\gamma_d)=-J_\theta,
\qquad
\int(2RG\gamma_d-RV\Delta v)=-J_z,
\tag{MC24}
\]

with \(dR\) understood in each integral. After using physical targets scaled by \(q\) with the exact powers in Section 5, the angular powers are \(p_0=2,p_1=-2-2\lambda,p_2=-2\lambda\). Set \(t_i=e^{dp_i}\). All three are distinct for \(\lambda>0\).

Let \(C(s)=u_0+u_1s+u_2s^2\). The angular equations are exactly

\[
C(t_0)=0,\qquad
C(t_1)=-\frac{P_q}{2a\mu_{p_1}}=:b_1,
\qquad
C(t_2)=\frac{J_{z,q}}{a\mu_{p_2}}=:b_2.
\]

Therefore the unique polynomial and all three coefficients are given explicitly by

\[
C(s)=b_1\frac{(s-t_0)(s-t_2)}{(t_1-t_0)(t_1-t_2)}
+b_2\frac{(s-t_0)(s-t_1)}{(t_2-t_0)(t_2-t_1)}.
\tag{MC25}
\]

Evaluation at the three nodes proves all three equations exactly. Uniqueness follows since a nonzero polynomial of degree at most two cannot have three distinct roots: division by each linear factor would force degree at least three. Equivalently the displayed source matrix has determinant

\[
-2a^2\mu_2\mu_{-2-2\lambda}\mu_{-2\lambda}
(t_1-t_0)(t_2-t_0)(t_2-t_1)\ne0.
\]

For the axial block set \(t_b=e^d\), \(t_a=e^{d(1-2\lambda)}\). They are distinct. Its two coefficients are

\[
s_1=-\frac{J_{\theta,q}}{a\mu_{1-2\lambda}(t_a-t_b)},
\qquad s_0=-t_bs_1.
\tag{MC26}
\]

Then \(\mu_1(s_0+t_bs_1)=0\), and

\[
a\mu_{1-2\lambda}(s_0+t_as_1)=-J_{\theta,q}.
\]

The determinant is \(a\mu_1\mu_{1-2\lambda}(t_a-t_b)\ne0\). Define the physical corrections, with no change to their powers,

\[
\Delta v_{\rm phys}=q^{-A}\sum_{j=0}^2u_j\eta_j(r/\sqrt q),
\qquad
\gamma_{d,\rm phys}=q^{-A}\sum_{j=0}^1s_j\eta_j(r/\sqrt q).
\tag{MC27}
\]

Equations (MC23)-(MC26) prove all five rows (MC24), including the negative \(RV\Delta v\) in the last row. The profile-to-chart changes multiply every row and its target by the identical factor from (MC15), so the equations also hold in the fixed \(Q\) chart. This is an exact map from the three target functions into the five-dimensional space spanned by the chosen profiles; uniqueness concerns this chosen space, not every possible smooth correction.

All denominators except \(a\) are fixed nonzero constants; \(a^{-1}\) has bounded slow derivatives for the prescribed smooth base and nonzero lower bound. Equations (MC25)-(MC27), the smooth bounded powers of \(q/Q\), and Leibniz' rule prove \(\mathcal S_\alpha\to\mathcal M_\alpha\) at every fixed derivative order. On the interior support the weight is bounded below. Constants may deteriorate as \(\lambda\downarrow0\); the construction needs no estimate uniform in that limit. At \(\lambda=0\) the two axial powers coincide, so there is a real degeneracy outside the stated parameter domain.

The axial correction has zero ordinary weighted integral and is torus independent, so (MC14) realizes it exactly: \(\Delta\gamma=\gamma_d\). Its potential vanishes below the first bump and above the last bump but may be nonzero between them. Its entire support remains in the reserved interval. The induced radial correction is in \(\mathcal M_{\alpha+1}\). This explains precisely why local base bounds must hold throughout that interval, rather than only on the union of bump supports.

\subsection{Every nonlinear change and its exponent}

Keep \(w\), its complete covariance, the base, and the stress fixed. Write \(B=\Delta\beta\), \(u=\Delta v\), \(d=\Delta\gamma=\gamma_d\) only within this calculation; these are the actual fields just constructed, not altered objects. The product differences are

\[
2b(\beta+B)+(\beta+B)^2-(2b\beta+\beta^2)
=2bB+2\beta B+B^2,
\]
\[
b(\gamma+d)+G(\beta+B)+(\beta+B)(\gamma+d)
-(b\gamma+G\beta+\beta\gamma)
=bd+GB+\beta d+\gamma B+Bd,
\]
\[
2V(v+u)+(v+u)^2-(2Vv+v^2)=2Vu+2vu+u^2.
\]

Substituting these three identities into (MC12) proves the exact radial-source difference

\[
g_{r,\rm new}-g_r=(2V/R)u+R_g,
\]
\[
\begin{aligned}
R_g={}&-t_*B-(D_r+1/R)(2bB+2\beta B+B^2)\\
&-D_z(bd+GB+\beta d+\gamma B+Bd)
+(2vu+u^2)/R+\epsilon(\Delta_0-R^{-2})B.
\end{aligned}
\tag{MC28}
\]

No covariance difference is present because \(w\) is unchanged. The finite-dimensional rows cancel only the displayed terms linear in the fixed base. Integrating the radial difference and the expanded flux products gives

\[
P_{\rm new}=\int\langle R_g\rangle\,dR,
\]
\[
(J_\theta)_{\rm new}=\int R^2\langle\gamma u+vd+du\rangle\,dR,
\]
\[
(J_z)_{\rm new}=\int R\langle2\gamma d+d^2\rangle\,dR
-\frac12\int R^2\langle R_g\rangle\,dR.
\tag{MC29}
\]

To verify the last cancellation, the linear part of the pressure-source moment is

\[
-\frac12\int R^2(2V/R)u\,dR=-\int RVu\,dR,
\]

which combines with \(\int2RGd\,dR\) to form the last row of (MC24). This proves the signs in (MC29) without assuming that new defects vanish.

Take the original hypotheses \(v,\gamma\in\mathcal M_{0.9}\), \(\beta\in\mathcal M_{1.9}\), targets in \(\mathcal S_\alpha\), \(\alpha\ge0.9\). Then \(u,d\in\mathcal M_\alpha\), \(B\in\mathcal M_{\alpha+1}\). All increments are torus independent. On their potential support, retain the stated base bounds \(b=O(\epsilon S_*^B)\), \(V,G=O(S_*^B)\) for every fixed coefficient derivative. The complete term list is:

\begin{description}
\item[] Term in \(R_g\): Guaranteed exponent.
\item[] \(-t_*B=+\epsilon\partial_TB\): \(\alpha+2\).
\item[] \(-(D_r+1/R)(2bB)\): \(\alpha+2-\kappa_s\).
\item[] \(-(D_r+1/R)(2\beta B)\): \(\alpha+2.9-\kappa_s\).
\item[] \(-(D_r+1/R)B^2\): \(2\alpha+2-\kappa_s\).
\item[] \(-D_z(bd)\): \(\alpha+2\).
\item[] \(-D_z(GB)\): \(\alpha+2\).
\item[] \(-D_z(\beta d)\): \(\alpha+2.9\).
\item[] \(-D_z(\gamma B)\): \(\alpha+2.9\).
\item[] \(-D_z(Bd)\): \(2\alpha+2\).
\item[] \(2vu/R\): \(\alpha+0.9\).
\item[] \(u^2/R\): \(2\alpha\).
\item[] \(\epsilon D_r^2B\): \(\alpha+2-2\kappa_s\).
\item[] \(\epsilon R^{-1}D_rB\): \(\alpha+2-\kappa_s\).
\item[] \(\epsilon D_z^2B\): \(\alpha+4\).
\item[] \(-\epsilon R^{-2}B\): \(\alpha+2\).
\end{description}

Every exponent is at least \(\alpha+0.9-2\kappa_s\): the only potentially tight quadratic one is \(2\alpha\), and \(2\alpha-(\alpha+0.9-2\kappa_s)=\alpha-0.9+2\kappa_s\ge2\kappa_s>0\). All remaining comparisons follow by subtraction and \(\alpha\ge0.9\). Thus

\[
R_g\in\mathcal M_{\alpha+0.9-2\kappa_s},\qquad
(P_{\rm new},(J_\theta)_{\rm new},(J_z)_{\rm new})
\in\mathcal S_{\alpha+0.9-2\kappa_s}.
\tag{MC30}
\]

The two product-only defects in (MC29) actually have bounds \(\alpha+0.9\) and \(2\alpha\) before the \(R_g\) moment is added. The table records every transport and viscous term, including those compressed into the source proof's remaining-transport statement.

\subsection{Recalculation, endpoint scope, and audit disposition}

The exact construction order is forced by the definitions: calculate the full \(W^{ab}\) and \(g_r\) from the current actual velocity; calculate \(P\) and \(p_m\); calculate \(E_\theta,E_z\) with that pressure; calculate \(J_\theta,J_z\) from the same \(g_r\) and velocity. Repeat this sequence after every update. In particular, the temporal correction's \(a=-A_1\gamma_d\) must occur in all new products. It is insufficient that \(\bar a=0\): for example \(\overline{a^2}\) can be positive. The exact difference formulas (MC28)-(MC29) apply to the slow five-dimensional update, whose axial remainder is identically zero.

All local operations preserve one-sided slow derivatives on closed positive-\(q\) regions because the radial integrals have fixed-shift formulas, Fourier inverse estimates use finitely many derivatives for each requested derivative, and matrix coefficients are smooth functions of the slow data. This statement carries the actual source regularity through those operators; it does not construct missing source regularity. On \(q\ge c_0>0\), only finitely many bands occur, so no accumulation of covering changes is involved. Near \(q=0\), physical flatness of the radial cutoff remainders was proved in Section 3.

The verified results of this bounded audit are the full radial exact sequence and its cutoff splitting; all derivative losses for the primitive and time inverse; the exact solenoidal velocity realization; the complete integrated mean residual identities; an explicit inverse of the five moment rows; and the exact nonlinear recomputation and stated exponent gain. No Section 8 repair is required by these calculations. This conclusion is limited to these maps and estimates. The actual base input, wave construction, iteration, and passage to an exact solution remain separate portions of the coordinated source audit, listed precisely in \texttt{\detokenize{DEPENDENCIES.md}}.


% --- included proof body: stage9_body.tex ---
% Independent reconstruction of Section 9 (printed pp.100--116).
% Inclusion body: requires amsmath, amssymb, and standard theorem environments.
% This file records the finite-stage proof architecture and does not claim a
% theorem-complete Navier--Stokes construction.
\section{Residual improvement and the local field (independent reconstruction)}
\label{sec:stage9}

\paragraph{Scope and provenance.}
The formulas and page assignments in this section were transcribed from the
released manuscript, printed pp.~100--116, and then rederived below.  Data
imported from Sections~5--8 are hypotheses here: the fixed background and
primary profiles, the pulse inverse, curl realization, oscillatory covariance,
and the temporal and five-equation mean maps. Those earlier constructions
are proved in the accompanying bodies at the locations in the dependency
ledger. The present text proves their composition, the displayed residual
identities, exponent arithmetic, and single-stage summation argument.
It does not independently reprove existence of all the imported data.

Put $\varepsilon=Q^h$, $Q=2^{-\ell}$, $S_*=\ell^2$, and retain the source
normalization $R_*=Q^{2A+1/2}R$.  Throughout
\[
  \kappa_s=10^{-5},\qquad \sigma_j=\frac{1}{5}+\frac{j}{10},\qquad
  B_j=\frac{1}{2}+\sigma_j,\qquad C_j=B_j+\frac{1}{2}=1+\sigma_j .
\]
The classes $W^\alpha,M^\alpha,S^\alpha$ are exactly those fixed in
Section~6.4; polynomial factors in $S_*$ and the radial edge factors
$\sqrt\zeta\,\delta^{-M}$ or $\zeta\,\delta^{-M}$ are retained in every
estimate.  The angular mean $\langle\cdot\rangle_\theta$ and auxiliary mean
$\langle\cdot\rangle_Y$ are different operations.
Since $0<\varepsilon\le1$, larger exponents give smaller classes:
$W^a\subset W^b$ and $M^a\subset M^b$ for $a\ge b$. Thus the quantities
$W^{0.68}$ and $M^{0.9}$ are class bounds; they are not powers of large
parameters called $W$ or $M$.

\subsection{Curl identities and the linear residual}
\label{subsec:stage9-linear}

For a fixed nonzero harmonic $m\in\mathbb Z$, set $k_m=km$, where
$k=\lceil\varepsilon^{-1/2}\rceil$, and write $a\,e^{ik_m\Phi}$.
Retain the full phase covector $n_\Phi=\nabla_*\Phi$, without unit
normalization. In the fixed label the phase and covector are
\[
\Phi=p\theta+p_zZ/\varepsilon+x_0R-vH_\Phi,\qquad
H_\Phi=pF+p_zG,\qquad F=V/R,
\]
\[
n_\Phi=(x_0-v\partial_RH_\Phi,\ p/R,\
p_z-\varepsilon v\partial_ZH_\Phi).
\]
Here $v$ in the phase is the pulse variable; later $v$ in the mean tuple
denotes the azimuthal mean correction, as in the source. The displayed
covector follows from $D_rv=D_zv=0$ and $D_z=\varepsilon\partial_Z$
on auxiliary-independent coefficients. Its fixed positive lower length
bound is supplied by the pulse geometry.
The imported pulse inverse supplies a transverse amplitude $t_m$ and pressure
$\pi_m$ solving
\begin{equation}
 \begin{aligned}
 L_m(t_m,\pi_m):=Q^{1+h}\!N_{\rm abs}t_m+K t_m
 &+\varepsilon k_m^2|n_\Phi|^2t_m\\
 &+ik_m n_\Phi\pi_m=-f_m,\qquad n_\Phi\cdot t_m=0 .
 \end{aligned}                         \tag{9.1}
\end{equation}
Its normalized potential coefficient and curl are
\[
 C_m=\frac{i n_\Phi\times t_m}{k_m|n_\Phi|^2},\qquad
\operatorname{curl}_*(C_me^{ik_m\Phi})=(t_m+r_m)e^{ik_m\Phi},
\]
Indeed, in the retained right-handed $(e_r,e_\theta,e_z)$ orientation,
\[
ik_m n_\Phi\times C_m
=-\frac{n_\Phi\times(n_\Phi\times t_m)}{|n_\Phi|^2}
=t_m-\frac{n_\Phi(n_\Phi\cdot t_m)}{|n_\Phi|^2}=t_m.
\]
The remaining terms are exactly
\[
r_m=(-D_z(C_m)_\theta,\
D_z(C_m)_r-D_r(C_m)_z,\
(D_r+R^{-1})(C_m)_\theta).
\]
The physical potential is $Q^{1/2-A}C_me^{ik_m\Phi}$.
where the imported curl lemma gives $r_m\in W^{\alpha+1/2-\kappa_s}$ if
$t_m\in W^\alpha$.  Thus the complete divergence-free amplitude is
$a_m=t_m+r_m$; the transport factor in quadratic products is always this
complete amplitude.

For an arbitrary amplitude $a$ and pressure coefficient $\pi$, direct
expansion in cylindrical coordinates gives the remainder after (9.1):
\begin{align}
 L^{\rm rem}_m(a,\pi)={}&(-\varepsilon\partial_T+bD_r+GD_z)a
 +ik_m E_{ik}a\\
 &+a_r(\partial_Rb)e_r+a_zD_z(b,V,G)+\frac{b}{R}a_\theta e_\theta
 +(D_r\pi,0,D_z\pi)\\
 &-\varepsilon\big\{(D_r^2+R^{-1}D_r+D_z^2)a-R^{-2}(a_r,a_\theta,0)\\
 &\qquad +2ik_m(n_{\Phi,r}D_r+n_{\Phi,z}D_z)a
 +ik_m(D_rn_{\Phi,r}+R^{-1}n_{\Phi,r}+D_zn_{\Phi,z})a\\
 &\qquad +2ik_m(n_{\Phi,\theta}/R)(-a_\theta,a_r,0)\big\}.       \tag{9.2}
\end{align}
Here $E_{ik}=(t_*+bD_r+(V/R)\partial_\theta+GD_z)\Phi$ is the imported
phase-transport defect.  The displayed formula follows by differentiating the
cylindrical basis: transport of $v$ contributes
$(-u_\theta v_\theta,u_\theta v_r,0)/R$, while the scalar and vector Laplacians
give respectively the mixed phase terms and the last angular connection.

\paragraph{Exponent audit.}
If $f_m\in W^\alpha$, the seven groups in (9.2) gain, in their written order,
\[
1-\kappa_s,\quad \frac{1}{2},\quad 1,\quad \frac{1}{2}-\kappa_s,
\quad1-2\kappa_s,\quad \frac{1}{2}-\kappa_s,\quad \frac{1}{2}.
\]
The principal velocity operator applied to the curl correction preserves the
same exponent (the fast derivative is already in (9.1), and
$\varepsilon k_m^2=O(1)$).  Hence the normalized linear residual belongs to
$W^{\alpha+1/2-3\kappa_s}$.  Multiplication by a pulse cutoff is done before
the curl.  The uncancelled $(1-\psi)f_m$ and $\psi'_t$ terms satisfy, for all
fixed Cartesian orders $m,N$,
\[
 |E_{\rm cut}|_m\le C_{m,N}Q^{-M}S_*^P e^{-cS_*}\le C_{m,N,N'}q^{N'} .
\]
They are retained in the additive flat remainder and are never reintroduced as
forward-pulse sources.

\subsection{Nonlinear products and finite-stage source decomposition}
\label{subsec:stage9-products}

If $w\in W^\alpha$ is a wave and a mean field has tangential class
$M^\mu$ and radial class $M^{\mu+1}$, then every nonzero harmonic of
$(w\!\cdot\!\nabla_*)v+(v\!\cdot\!\nabla_*)w$ lies in
$W^{\alpha+\mu-1/2}$.  For two curl-generated waves of classes
$W^\alpha,W^{\alpha'}$, the nonzero harmonics lie in
$W^{\alpha+\alpha'-\kappa_s}$ and the zero harmonic in
$M^{\alpha+\alpha'-\kappa_s}$.  Indeed, phase differentiation costs
$k=O(\varepsilon^{-1/2})$. The complete amplitude is not asserted
transverse. Expanding its exact zero divergence gives
\[
ik_m(n_\Phi\cdot a_m)
=-\big((D_r+R^{-1})(a_m)_r+D_z(a_m)_z\big),
\]
\[
n_\Phi\cdot a_m=n_\Phi\cdot r_m\in
W^{\alpha+1/2-\kappa_s}.
\]
Thus phase differentiation of an $m'$ harmonic by this wave contributes
$ik_{m'}(a_m\cdot n_\Phi)a_{m'}$, which belongs to
$W^{\alpha+\alpha'-\kappa_s}$. Equivalently the divergence identity
replaces $ik_{m'}n_\Phi\cdot a_m$ by $m'/m$ times the displayed amplitude
divergence; $m'/m$ is a fixed constant at each finite stage. Amplitude
derivatives cost at most $\kappa_s$ and cylindrical connections cost none.
Since $P_v^2\le P_v$ and supports of distinct labels
are disjoint, no other loss occurs.  Leibniz' rule proves the same statements
at every fixed coefficient derivative order.

For a finite stage $j$, write the physical residual as
\begin{equation}
 R(u^{[j]},p^{[j]})=G^{[j]}+F^{[j]},\qquad
 G_*^{[j]}=Q^{2A+1/2}G^{[j]},\quad F_*^{[j]}=Q^{2A+1/2}F^{[j]} .       \tag{9.3}
\end{equation}
On the common auxiliary torus,
\begin{equation}
 G_*^{[j]}-\langle G_*^{[j]}\rangle_\theta
 =\sum_{\gamma\in\Gamma}\sum_{m\in H_j}f^{[j]}_{\gamma,m}e^{ik_{\gamma m}\Phi_\gamma},
 \quad H_j\Subset\mathbb Z\setminus\{0\},                         \tag{9.4}
\end{equation}
with $\operatorname{supp}(f^{[j]}_{\gamma,m}|_{E_\gamma})\subset\Omega_\gamma$,
smooth zero extension to $E_\gamma^+$, and $f^{[j]}_{\gamma,m}\in W^\alpha$.
The angular mean belongs to $M^\mu$, with the exact exponents supplied by the
finite preceding construction.  Radial integration is at fixed $(z,t)$, so
$q(z,t)$ and the local band set do not change; pressure reconstruction and the
mean vector potential preserve support and classes.  For the stress correction,
subtract an interior profile with equal weighted moment.  If $s$ is logarithmic
distance to an edge and $\zeta\asymp e^{-a/s^2}$, then
\[
\int_0^\delta s^{-M}e^{-a/s^2}ds\le C_{M,a}\delta^{3-M}e^{-a/\delta^2},\qquad
\left|\partial_s^k e^{-a/s^2}\right|\le C_{k,a}s^{-3k}e^{-a/s^2}.
\]
For the first inequality substitute $u=a/s^2$, obtaining
$\frac12a^{(1-M)/2}\int_{a/\delta^2}^\infty
u^{(M-3)/2}e^{-u}du$. With $u=u_0+y$ and $u_0\ge1$, a nonpositive
power is bounded by $u_0^{(M-3)/2}$, while a positive power is bounded
by this factor times $(1+y)^{(M-3)/2}$. Integration in $y$ gives the
claimed bound. For the second, induction expresses the derivative as
$e^{-a/s^2}$ times a polynomial in $s^{-1}$ of degree at most $3k$.
This also proves that every fixed inverse-distance power times either
edge factor tends to zero with all derivatives.
Thus the correction remains in $M^\alpha$ with only polynomial $S_*$ losses.
The amplitude division is exact:
\[
 d\Sigma=H^{-1}(\Sigma/\varepsilon),\qquad
 \delta a_\sigma=\frac{(d\Sigma)_\sigma}{2a_\sigma},\qquad
 |D_I\delta a_\sigma|\lesssim \varepsilon^{\alpha-1}S_*^{P_I}\sqrt\zeta\,\delta^{-M_I},
\]
and $\sqrt\varepsilon\,\delta a_\sigma b_\sigma\in W^{\alpha-1/2}$ with the
retained factor $P_v$. This proves the pulse inverse hypotheses at every finite
stage without shrinking the common threshold.

The flat term satisfies, for every fixed $j,m,N$,
\begin{equation}
 |F^{[j]}|_m\le C_{j,m,N}q^N .                                      \tag{9.5}
\end{equation}
For an increment $(w,\pi)$ the exact residual identity is
\[
 R(u+w,p+\pi)=R(u,p)+\partial_tw-\Delta w+\nabla\pi
 +(u\cdot\nabla)w+(w\cdot\nabla)u+(w\cdot\nabla)w .                \tag{9.5a}
\]
This proves inductively that all products use complete cutoff curls and that
new cutoff tails are added to $F$. The exact mean identities, with
$g_r,E_\theta,E_z$ defined by the imported formulas (8.3), are
\begin{equation}
 \begin{aligned}
 W_{ab}=\langle w_aw_b\rangle_\theta,\quad
 P&=\int\langle g_r\rangle_YdR,\quad p_m=\mathcal T^0(g_r-\rho P),\\
 \langle G_*+F_*-E_{B,*}\rangle_\theta
 &=(D_rp_m-g_r,E_\theta,E_z).
 \end{aligned}                                      \tag{9.6}
\end{equation}

\subsection{Finite correction state and initialization}
\label{subsec:stage9-state}

Write the normalized finite state as
\begin{equation}
 u_*=(b+\beta,V+v,G+\gamma)+w,\qquad
 p_*=p_{B,*}+p_m+p_w,\qquad \langle w\rangle_\theta=\langle p_w\rangle_\theta=0 . \tag{9.7}
\end{equation}
At stage $j$, the nonzero wave residual, tangential mean residual, and defects
obey
\begin{equation}
 \begin{aligned}
 G_{\rm wave}^{[j]}\in W^{B_j},\quad E_\theta,E_z\in M^{C_j},\quad
 &(P,J_\theta,J_z)\in S^{C_j},\\
 B_j=\tfrac12+\sigma_j,\quad &C_j=1+\sigma_j .
 \end{aligned}                         \tag{9.8}
\end{equation}
The cumulative fields satisfy the deliberately retained budgets
\begin{equation}
 w\in W^{1/2},\quad w-w_0^{\rm tan}\in W^{0.68},\quad
 v,\gamma,p_m\in M^{0.9},\quad \beta\in M^{1.9},                    \tag{9.9}
\end{equation}
and the moments remain exact:
\begin{equation}
 \int_0^\infty R^2\langle v\rangle_YdR=0,\qquad
 \int_0^\infty R\langle\gamma\rangle_YdR=0 .                     \tag{9.10}
\end{equation}

The primary transverse amplitudes have class $W^{1/2}$; their linear residual
has $1-3\kappa_s$, and wave nonlinearities $1-\kappa_s$.  Before mean updates,
the exact balances and defects are $M^{1-\kappa_s}$ and $S^{1-\kappa_s}$.
The leading auxiliary covariance cancels exactly because the squared slow
partition functions sum to one.  Writing $\Sigma_a^{(0)}$ for this tensor and
$\Sigma_a$ for the full tensor,
\begin{align}
\langle E_\theta\rangle_Y&=(\partial_R+2/R)(\langle W_{r\theta}\rangle_Y-\Sigma_\theta^{(0)})
 -(\partial_R+2/R)(\Sigma_\theta-\Sigma_\theta^{(0)})+\varepsilon\partial_Z\langle W_{z\theta}\rangle_Y,\\
\langle E_z\rangle_Y&=(\partial_R+1/R)(\langle W_{rz}\rangle_Y-\Sigma_z^{(0)})
 -(\partial_R+1/R)(\Sigma_z-\Sigma_z^{(0)})+\varepsilon\partial_Z\langle W_{zz}+p_m\rangle_Y . 
\end{align}
Cross terms contain a curl correction and are $M^{3/2-\kappa_s}$; axial and
higher covariance terms are at least $M^{2-\kappa_s}$.  Hence the averaged
balances are $M^{1.49}$.  The temporal inverse at order
$H_0=1-\kappa_s$ gives tangential increment $M^{H_0}$ and radial increment
$M^{H_0+1}$. After pressure reconstruction the new tangential mean
residual terms have order at least $H_0+1-2\kappa_s$, while their
unchanged auxiliary averages retain order $1.49$. Wave--mean
interaction has wave order at least $H_0$; defects remain $S^{H_0}$.
The five-equation map
cancels the three linear defect contributions, leaving defects above
$1.8-\kappa_s$.  Since $1-\kappa_s>0.68$ and $H_0>0.9$, this is stage $0$:
$B_0=0.7$, $C_0=1.2$, with (9.9)--(9.10).

\subsection{One correction cycle}
\label{subsec:stage9-cycle}

Assume a state at stage $j$ and abbreviate $B=B_j$, $C=C_j$.  Recompute $g_r$
from complete fields and set $P=\int\langle g_r\rangle_YdR$,
$p_m=\mathcal T^0(g_r-\rho P)$ after each substep.

\paragraph{(i) Supported nonzero harmonics.}
For each coefficient in (9.4), solve the pulse equation with zero data on the
entrance path and set
\[
 \Delta w_m=\operatorname{curl}_*\!\left(\frac{in_\Phi\times(\psi t_m)}{k_m|n_\Phi|^2}e^{ik_m\Phi}\right),
 \qquad \Delta p_{w,m}=\psi\pi_me^{ik_m\Phi}.
\]
The increment is $W^B$, its pressure is $W^{B+1/2}$, and the four resulting
nonzero-harmonic errors have exponents
\[
 B+\tfrac12-3\kappa_s,\quad B+\tfrac12-\kappa_s,\quad 2B-\kappa_s,\quad B+0.4 .
\]
The last $B+0.4$ estimate is a nonzero-harmonic (wave) estimate. Separately, the mean covariance changes have exponent at least $B+1/2=C$; a radial divergence costs $\kappa_s$. Thus $E_\theta,E_z,\Delta g_r,\Delta p_m\in M^{C-\kappa_s}$ and $(P,J_\theta,J_z)\in S^{C-\kappa_s}$. The exact integrated identities improve
the weighted moments by one:
\[
 \int R^2\langle E_\theta\rangle_YdR,\quad
\int R\langle E_z\rangle_YdR\in S^{C+1-\kappa_s}.                \tag{9.12}
\]
Precisely, compact radial support and the two moment constraints give
\[
\int R^2\langle E_\theta\rangle_YdR
=\varepsilon\partial_ZJ_\theta,\qquad
\int R\langle E_z\rangle_YdR
=\varepsilon\partial_Z(J_z+c_\rho P),\qquad
c_\rho=\tfrac12\int R^2\rho\,dR .
\]
Multiplying a radial flux $(\partial_R+e/R)h$ by $R^e$ turns its
integral into a zero boundary term. For the angular viscosity the
integrals of $R^2h_{RR}$, $Rh_R$, and $-h$ are respectively
$2\int h$, $-\int h$, and $-\int h$; axial radial viscosity cancels
in the same way. Time and axial derivatives of the two zero moments
vanish. Finally
$\int R\langle p_m\rangle_YdR
=-\frac12\int R^2\langle g_r\rangle_YdR+c_\rho P$.
These calculations leave exactly the two displayed axial flux
derivatives and account for the $c_\rho P$ term without removing it
from its derivative. The $\varepsilon\partial_Z$ supplies the claimed
gain of one.

\paragraph{(ii) Auxiliary-averaged stress correction.}
In physical variables let $F_2=Q^{-2A-1/2}\langle E_\theta\rangle_Y$,
$F_1=Q^{-2A-1/2}\langle E_z\rangle_Y$.  Choose fixed profiles $\widehat b_e$
with $\int x^e\widehat b_e(x)dx=1$, set
$b_e(r)=q^{-(e+1)/2}\widehat b_e(r/\sqrt q)$, and define
\[
 M_e=\int_0^\infty r^eF_e(r)dr,\qquad
 \sigma_e(r)=-r^{-e}\int_0^r(r')^e(F_e(r')-b_e(r')M_e)dr'.
\]
The equal-moment subtraction implies $\sigma_e$ vanishes outside the annulus and
$(\partial_r+e/r)\sigma_e=-F_e+b_eM_e$.  Put
$\Sigma=Q^{2A}(\sigma_2,\sigma_1)$.  The edge estimates above give
$\Sigma\in M^{C-\kappa_s}$.  Applying the physical signed amplitude map divides
by the fixed positive $2\sqrt{y_\sigma}$ and realizes the prescribed covariance
with $w_0^{\rm tan}$; no square root of a perturbed covariance is taken.

Write the signed increment as $s=t_s+r_s$.  Its exact covariance change is
\[
B(a,b)=\langle\langle a\otimes b+b\otimes a\rangle_\theta\rangle_Y,
\qquad
\langle\Delta W\rangle_Y
=\langle\langle (w+s)\otimes(w+s)-w\otimes w\rangle_\theta\rangle_Y.
\]
\begin{equation}
 \begin{aligned}
 \langle\Delta W\rangle_Y={}&B(w_0^{\rm tan},t_s)+B(w-w_0^{\rm tan},t_s)\\
 &+B(w,r_s)+\langle\langle s\otimes s\rangle_\theta\rangle_Y .
 \end{aligned}\tag{9.13}
\end{equation}
This is the distributive law with $w=w_0^{\rm tan}+(w-w_0^{\rm tan})$
and $s=t_s+r_s$; in particular the complete self-interaction is retained.
The four new wave errors from this signed increment are, respectively,
\[
B+\tfrac12-4\kappa_s,\quad B+0.4-\kappa_s,\quad
B+\tfrac12-2\kappa_s,\quad 2B-3\kappa_s .
\]
The first term has radial-tangential component $\Sigma$ and cancels the target
average.  The remaining terms before divergence have exponents
$C+0.18-\kappa_s$, $C+1/2-2\kappa_s$, and $C+\sigma_j-2\kappa_s$; radial divergence costs
$\kappa_s$.  Together with (9.12), pressure reconstruction yields
\begin{equation}
 \langle E_\theta\rangle_Y,\langle E_z\rangle_Y\in M^{C+0.17},\qquad
 E_\theta,E_z\in M^H,\quad (P,J_\theta,J_z)\in S^H,\quad H=C-2\kappa_s . \tag{9.14}
\end{equation}

\paragraph{(iii) Temporal removal of auxiliary oscillation.}
The imported radial mean inverse is the split operator for the full torus
obstruction $K_e f(w)=\int R^e f(R,w+M R^{d_r}v_r)\,dR$:
$D_e\mathcal T_e=1-\mathcal A_e$, $K_e\mathcal A_e=K_e$, and
$\mathcal A_e^2=\mathcal A_e$.  The vanishing Haar moment alone is not used as
an inversion hypothesis; the projector contribution is retained as the flat
remainder. The temporal inverse $N^{-1}$ acts on the zero auxiliary
average subspace. The radial inverse is used with its projector term;
zero auxiliary average does not assert $K_e f=0$.
Set $E_a^\circ=E_a-\langle E_a\rangle_Y$ and use the zero-average inverse:
\[
 \Delta v=-c_{i0}^{-1}N_{i0}^{-1}E_\theta^\circ,\quad
 \gamma_d=-c_{i0}^{-1}N_{i0}^{-1}E_z^\circ,\quad
 \Psi_*=\mathcal T^1\gamma_d,\quad
 \Delta\beta=-\varepsilon\partial_Z\Psi_*,\quad
 \Delta\gamma=(D_r+R^{-1})\Psi_* .
\]
The actual $\Delta\gamma$ includes the flat reconstruction remainder
$F_{\rm ax}=c_{i0}N_{i0}(\Delta\gamma-\gamma_d)$.  Fast-time identities are
$c_{i0}N_{i0}\Delta v=-E_\theta^\circ$ and
$c_{i0}N_{i0}\Delta\gamma=-E_z^\circ+F_{\rm ax}$.  Subtracting old from new
balances gives exactly
\begin{align}
 E_{\theta,\rm new}-\langle E_\theta\rangle_Y={}&-\varepsilon\partial_T\Delta v
 -\varepsilon(\Delta_0-R^{-2})\Delta v
 +(D_r+2/R)((b+\beta)\Delta v+(V+v)\Delta\beta+\Delta\beta\Delta v)\\
 &+D_z((G+\gamma)\Delta v+(V+v)\Delta\gamma+\Delta\gamma\Delta v),\\
 E_{z,\rm new}-\langle E_z\rangle_Y={}&-\varepsilon\partial_T\Delta\gamma-\varepsilon\Delta_0\Delta\gamma+F_{\rm ax}
 +(D_r+R^{-1})((b+\beta)\Delta\gamma+(G+\gamma)\Delta\beta+\Delta\beta\Delta\gamma)\\
 &+D_z(2(G+\gamma)\Delta\gamma+(\Delta\gamma)^2+\Delta p_m).
\end{align}
Slow time, radial flux, axial flux, and viscosity have gains respectively
$H+1,H+1-\kappa_s,H+1,H+1-2\kappa_s$.  Thus, after placing $F_{\rm ax}$ in
the flat remainder, $E_\theta,E_z\in M^{\min(C+0.17,H+1-2\kappa_s)}$ and
the wave change is $W^H$.

\paragraph{(iv) Compatibility defects.}
Apply the five-equation map to $(P,J_\theta,J_z)$ at order $H$.  The azimuthal
and prescribed axial increments are $M^H$, their induced radial increment is
$M^{H+1}$, and the first two rows preserve (9.10).  The term
$2V\Delta v/R$ and its two moment contributions are exactly the linear terms
canceled by the last three rows.  Every remaining term has exponent at least
$H+0.9-2\kappa_s$, hence
\[
 (P,J_\theta,J_z)\in S^{H+0.9-2\kappa_s}=S^{C+0.9-4\kappa_s}.
\]
Here is the complete calculation behind the defect estimate. Write
$u=\Delta v$, $d=\Delta\gamma=\gamma_d$, and $B_r=\Delta\beta$ in this
paragraph only. The slow moment correction has zero axial projector
remainder, because its prescribed axial mean satisfies the exact weighted
moment constraint. Since the wave is unchanged, subtraction gives
\[
g_{r,\mathrm{new}}-g_r=(2V/R)u+R_g,
\]
\begin{align*}
R_g={}&-t_*B_r-(D_r+R^{-1})(2bB_r+2\beta B_r+B_r^2)\\
&-D_z(bd+GB_r+\beta d+\gamma B_r+B_rd)\\
&+R^{-1}(2vu+u^2)+\varepsilon(\Delta_0-R^{-2})B_r .
\end{align*}
The last three rows of the fixed linear system imply exactly
\[
P_{\rm new}=\int\langle R_g\rangle_YdR,\qquad
J_{\theta,\rm new}=\int R^2\langle\gamma u+vd+du\rangle_YdR,
\]
\[
J_{z,\rm new}=\int R\langle2\gamma d+d^2\rangle_YdR
-\tfrac12\int R^2\langle R_g\rangle_YdR .
\]
For $u,d\in M^H$ and $B_r\in M^{H+1}$ the termwise exponents are
\[
\begin{array}{c|c}
\text{term}&\text{exponent}\\ \hline
-t_*B_r&H+2\\
(D_r+R^{-1})bB_r&H+2-\kappa_s\\
(D_r+R^{-1})\beta B_r&H+2.9-\kappa_s\\
(D_r+R^{-1})B_r^2&2H+2-\kappa_s\\
D_zbd,\ D_zGB_r&H+2\\
D_z\beta d,\ D_z\gamma B_r&H+2.9\\
D_zB_rd&2H+2\\
2vu/R,\ u^2/R&H+0.9,\ 2H\\
\varepsilon D_r^2B_r&H+2-2\kappa_s\\
\varepsilon R^{-1}D_rB_r&H+2-\kappa_s\\
\varepsilon D_z^2B_r,\ -\varepsilon R^{-2}B_r&H+4,\ H+2
\end{array}
\]
Every entry is at least $H+0.9-2\kappa_s$, since
$H=C-2\kappa_s\ge1.19998>0.9$. The two product-only defect expressions
have exponents at least $\min(H+0.9,2H)$ before the $R_g$ moment.
This establishes the displayed defect order with every retained term.
The closure inequalities are
\[
\min\{\tfrac12-3\kappa_s,\tfrac12-\kappa_s,B-\kappa_s,0.4\}\ge0.4,\quad
\min\{\tfrac12-4\kappa_s,0.4-\kappa_s,\tfrac12-2\kappa_s,B-3\kappa_s\}\ge0.4-\kappa_s,
\]
\[
 H-B=\tfrac12-2\kappa_s>0.1,\qquad
 \min\{0.17,1-4\kappa_s\}=0.17>0.1,\qquad 0.9-4\kappa_s>0.1 .
\]
The numerical margins are exact decimal rationals:
$0.18-2\kappa_s=0.17998>0.17$,
$\sigma_j-3\kappa_s\ge0.19997>0.17$,
$0.4-\kappa_s=0.39999$,
$H-B=0.49998$,
$1-4\kappa_s=0.99996$, and $0.9-4\kappa_s=0.89996$.
All exceed the needed $0.1$ gain in the relevant comparison.
Since $B\ge0.7$, the next state has $B_{j+1}=B_j+0.1$ and
$C_{j+1}=C_j+0.1$. The smallest signed increment is
$B-\kappa_s\ge0.69999>0.68$; mean and radial increments retain $0.9$ and
$1.9$. This proves Proposition~9.6 by induction, with pressure recomputed
after every substep and all retained flat terms in $F$.

\subsection{Common domain and Cartesian derivative budgets}
\label{subsec:stage9-domain}

Choose one $q_{\rm big}\le q_*$ on which the fixed background, positive primary
amplitudes, pulse propagators, and mean correction profiles are defined. The
propagator for every later harmonic is the fixed fundamental propagator times a
forward damping factor of modulus at most one, so no stage-dependent threshold
is needed. Signed updates divide by the original $2\sqrt{y_\sigma}$; the
five-equation map, pressure map, temporal inverse, and modified radial integrals
are fixed linear maps. Thus all finite states exist on $0<q<q_{\rm big}$.

For a fixed stage, let $D_{v,a}(m)$ count source derivatives needed at vertex $v$
of the finite directed acyclic graph of sums, products, derivatives, and inverses:
\[
 D_{v,a}(m)=\max_{w\to v}D_{w,a}(m+d_{v,w}),\qquad
 D_{b,a}(m)=m\ (b=a),\quad D_{b,a}(m)=0\ (b\ne a).
\]
The graph is finite at each stage. Extra derivatives of $N^{-1}D_rf$ do not
repeat the explicit radial loss; they only increase $D_{v,a}(m)$ and polynomial
$S_*$ powers.

On perturbation supports $r\asymp\sqrt Q$, $q\asymp Q$. In a fixed chart the
physical derivative losses are
\[
 \partial_r:\tfrac12+h\kappa_s,\qquad D_z:\tfrac12-h,\qquad
 \partial_t:1+h,\qquad \text{frame derivative}:\tfrac12 .              \tag{9.19}
\]
The phase obeys
$|\nabla_x^a\partial_t^b\Phi_\gamma|\lesssim
Q^{-a/2-b(1+h)}S_*^{P_{a,b}}$, while $k\le2Q^{-h/2}$ gives, with
$s_x=1/2+h/2$ and $s_t=1+3h/2$,
\[
 |\nabla_x^a\partial_t^b e^{ikm\Phi}|\lesssim
 Q^{-as_x-bs_t}S_*^{P_{a,b,m}} .                                  \tag{9.19a}
\]
These dominate (9.19), including chart transitions.  For every derivative order
$m$, a sufficient common loss is
\[
 \ell_m=2A+(m+1)(1+3h/2).
\]
Therefore, with $g_j=hj/10$, each retained stage increment obeys
\begin{equation}
 \begin{aligned}
 |Z_j|_m+|\Delta u_j|_m&\le C_{j,m}q^{g_j-\ell_m}(1+|\log q|)^{P_{j,m}},\\
 Z_j=(A_j,B_j,p_j),\quad &\Delta u_j=\operatorname{curl}A_j+B_j .
 \end{aligned}\tag{9.17}
\end{equation}
The finite sums satisfy
\[
 |(u^{[j]},p^{[j]})|_m\le C_{j,m}q^{-K_m}(1+|\log q|)^{P_{j,m}},
\]
and their residuals
\begin{equation}
 |R(u^{[j]},p^{[j]})|_m\le C_{j,m}q^{h\sigma_j-K_m}(1+|\log q|)^{P_{j,m}}+E_{j,m},
 \qquad |E_{j,m}|\le C_{j,m,N}q^N .                                 \tag{9.18}
\end{equation}
The fixed physical rescalings $Q^{-A},Q^{-2A},Q^{1/2-A}$ are bounded by
$Q^{-2A}$; one extra spatial derivative accounts for recovering velocity from a
potential.  On overlaps, prescribed torus coordinate changes and smooth zero
extensions prove the same estimates across support and band boundaries.

\subsection{Realization by a diagonal single-stage sum}
\label{subsec:stage9-realization}

Let $U_0=(u^{[0]},p^{[0]})$ and, for $j\ge1$, let $A_j^w$ be the physical
wave-potential increment, $\Psi_j$ its azimuthal mean potential, and $b_j$ its
direct azimuthal velocity.  Set
\[
 A_j=A_j^w+\Psi_je_\theta,\qquad B_j=b_je_\theta,\qquad
 p_j=p^{[j]}-p^{[j-1]},\qquad Z_j=(A_j,B_j,p_j),
\]
so $\Delta u_j=\operatorname{curl}A_j+B_j$.  In a fixed chart the pressure
increment is the exact telescoping identity
\begin{equation}
 Q^{2A}p_j=p_w^{[j]}-p_w^{[j-1]}+
 \mathcal T^0\!\left(g_r^{[j]}-g_r^{[j-1]}-\rho(P^{[j]}-P^{[j-1]})\right). \tag{9.16}
\end{equation}
Choose a fixed smooth $\chi$ with $\chi(s)=1$ for $s\le1/2$,
$\chi(s)=0$ for $s\ge1$, and $0\le\chi\le1$. We now construct, rather
than assume, a sequence $a_j\uparrow\infty$ which gives the claimed
residual bounds. It will define
\begin{equation}
 \begin{aligned}
 A&=A_0+\sum_{j\ge1}\chi(a_jq)A_j,\\
 Be_\theta&=B_0e_\theta+\sum_{j\ge1}\chi(a_jq)B_j,\\
 p_{\rm loc}&=p_0+\sum_{j\ge1}\chi(a_jq)p_j,\\
 u_{\rm loc}&=\operatorname{curl}A+Be_\theta .
 \end{aligned}                \tag{9.21}
\end{equation}
For a multi-index $\beta$ in physical space and time, the similarity map
gives $|\partial^\beta q|\le C_\beta q^{1-|\beta|}$ on $0<q<q_{\rm big}<1$;
the sharper axial loss is $D<1$. The multivariate chain rule expresses
each derivative of a cutoff as a sum of
\[
a_j^k\chi^{(k)}(a_jq)
\prod_{\nu=1}^k\partial^{\beta_\nu}q,\qquad
\sum_\nu\beta_\nu=\beta,\quad |\beta_\nu|\ge1 .
\]
For $k\ge1$, its support has $1/2\le a_jq\le1$, so pairing
$a_j^k$ with the $q^k$ from the product proves
$|\partial^\beta\chi(a_jq)|\le C_\beta q^{-|\beta|}$, independently
of $a_j$. The undifferentiated cutoff is bounded by one.

Define the actual cutoff increment
\[
V_j=\operatorname{curl}(\chi(a_jq)A_j)+\chi(a_jq)B_j,\qquad
\Pi_j=\chi(a_jq)p_j .
\]
The exact product rule retains
\[
V_j=\chi(a_jq)\Delta u_j
+a_j\chi'(a_jq)\nabla q\times A_j .
\]
Thus (9.17) and Leibniz's rule give constants
$\widehat C_{j,m},\widehat P_{j,m}$ independent of $a_j$ such that
\[
|(V_j,\Pi_j)|_m\le
\widehat C_{j,m}q^{g_j-L_m}(1+|\log q|)^{\widehat P_{j,m}},
\qquad L_m=\ell_{m+1}+m+1 ,
\]
on their support $q\le a_j^{-1}$. This deliberately generous
stage-independent loss includes every derivative of the cutoff and
the additional curl.

Recursively choose $a_j\ge\max(2a_{j-1},2/q_{\rm big},j)$ sufficiently
large that the finitely many inequalities
\[
\widehat C_{j,m}(1+|\log q|)^{\widehat P_{j,m}}q^{g_j/2}
\le2^{-j}\qquad(0<q\le a_j^{-1},\ 0\le m\le j)
\]
hold. Such a choice exists because $g_j=hj/10>0$ and
$q^\delta(1+|\log q|)^P\to0$ for every $\delta>0$ and finite $P$.
Consequently
$|(V_j,\Pi_j)|_m\le2^{-j}q^{g_j/2-L_m}$ for $m\le j$.
On a compact subset of the local domain, $q$ has a positive lower
bound, so only finitely many cutoff terms survive. This proves
smoothness of (9.21).

Each representative has smooth zero extension and vanishes near the axis for
fixed $t<1$. Each $B_j=b_je_\theta$ is an axisymmetric angular vector,
so in physical cylindrical coordinates
$\operatorname{div}(\chi(a_jq)B_j)=r^{-1}\partial_\theta
(\chi(a_jq)b_j)=0$. Divergence of each curl vanishes. Thus
$\operatorname{div}u_{\rm loc}=0$, including all cutoff transitions.

Fix a finite $J\ge\max(1,m)$ and restrict to
$q<(2a_J)^{-1}$. Then all the first $J$ cutoffs equal one, and for
$e_J=(u_{\rm loc},p_{\rm loc})-(u^{[J]},p^{[J]})$,
\[
|e_J|_m\le\sum_{j>J}2^{-j}q^{g_j/2-L_m}
\le2^{-J}q^{g_{J+1}/2-L_m}.                                     \tag{9.23}
\]
This is a comparison with one fixed stage. It does not estimate an
infinite sum of stage residuals.

Write $e_J=(v_J,\pi_J)$. The residual difference is exactly
\[
R(u^{[J]}+v_J,p^{[J]}+\pi_J)-R(u^{[J]},p^{[J]})
=\partial_tv_J-\Delta v_J+\nabla\pi_J
+(u^{[J]}\cdot\nabla)v_J+(v_J\cdot\nabla)u^{[J]}
+(v_J\cdot\nabla)v_J .
\]
For the Cartesian jet norm, Leibniz's rule bounds its order $m$ by
$C_m|e_J|_{m+2}(1+|u^{[J]}|_{m+2}+|e_J|_{m+2})$.
Given $m,N$, choose one $J\ge m+2$ so large that
\[
g_{J+1}/2-L_{m+2}>N+K_{m+2}+2,\qquad
h\sigma_J-K_m>N+1 .
\]
This is possible since both $g_j$ and $\sigma_j$ tend to infinity.
Keep this $J$ fixed. The finite-stage logarithm is
$O(q^{-1})$, with a constant depending on $J,m$, so (9.23) and the
finite-stage growth bound imply that the residual difference is
$O(q^N)$. Equation (9.18) gives the same estimate for the fixed
residual, including $E_{J,m}$ by its fixed-stage flatness. Adding
these two estimates proves for all fixed $m,N$
\[
 |R(u_{\rm loc},p_{\rm loc})|_m=O(q^N)\qquad(q\downarrow0).          \tag{9.20}
\]
No bound uniform in $J$ on the flat remainder constants was assumed or
used. On compact sets $c\le q\le c'$ only finitely many stages and bands are
active. The imported endpoint bounds apply to each of these finite
coefficients with the discrete label choices held fixed. Since
$L=1-2h\eta^2\ge1-2h>0$, coordinate conversion preserves bounds through
every fixed order. Bounding one additional time derivative and
integrating it gives uniform one-sided limits of each derivative at
$t=1$. Integrating spatial derivatives on compact coordinate segments
shows that the limiting jets are compatible with spatial
differentiation; integrating in time gives compatibility of successive
time derivatives. Thus the finite sum has the asserted endpoint
regularity away from the singular point.
Outside a fixed $X_{\rm ext}$ all annular corrections vanish; the remaining field
is the exact exterior heat profile with centrifugal pressure. Its precise
formula is
\[
K(r,\tau)=r^{-1-2h}H_{\rm ext}(\tau/r^2),\qquad
H_{\rm ext}(s)=2^A c_\infty H(4s),
\]
\[
p_{\rm ext}(r,\tau)=-\int_r^\infty K(\rho,\tau)^2\,\frac{d\rho}{\rho}.
\]
Indeed the imported leading field is
$c_\infty(r^2/2)^{-A}H(4\tau/r^2)$ and $2A=1+2h$, which gives
this exact prefactor. The integral defining pressure converges because
$K=O(r^{-1-2h})$ at infinity. It has derivative
$\partial_rp_{\rm ext}=K^2/r$. Consequently the radial component of
the momentum residual is $-K^2/r+\partial_rp_{\rm ext}=0$; the
axial component is zero, and the angular component is zero by the
imported exact heat equation
$\partial_tK=(\partial_r^2+r^{-1}\partial_r-r^{-2})K$.

The exterior profile integral is
\[
H(Z)=\frac1{\Gamma(1+h)}\int_0^\infty
e^{-v}v^h(1+Zv)^{-h}\,dv,\qquad Z\ge0 .
\]
Its $m$th derivative is the same integral with factor
$(-1)^m(h)_m v^m(1+Zv)^{-m}$. The absolute integrand is bounded
by $e^{-v}v^{h+m}(h)_m/\Gamma(1+h)$, an integrable function
independent of $Z$. Differentiation under the integral and the Gamma
recurrence therefore prove
\[
\sup_{s\ge0}|H_{\rm ext}^{(m)}(s)|
\le2^A c_\infty4^m(h)_m(1+h)_m .
\]
Here $(a)_0=1$ and $(a)_m=\prod_{\nu=0}^{m-1}(a+\nu)$.
Thus the exterior derivative estimate is uniform on the whole
nonnegative half-line, including its endpoint.

Near the axis the
base streamfunction is $r^2$ times a smooth function and the base angular field
is $r$ times a smooth scalar, so (9.21) has a smooth Cartesian representative.
At $(X,\eta)=(X_{\rm in},0)$, all annular corrections vanish and the leading
profile yields, with $q=\tau=1-t$,
\[
 u_{\theta,\rm loc}(\sqrt{2X_{\rm in}\tau},0,0,1-\tau)
 =\tau^{-A}\big(e_0+O(\tau^{2h})\big),\qquad e_0>0 .              \tag{9.22}
\]
These endpoint, exterior, and inner statements are consequences of the imported
profile, base, oscillation and mean results plus the finite-stage proof above; they
are recorded here without a theorem-complete claim.





% --- included proof body: endpoint_derivation.tex ---
\section{The terminal force, the comparison argument, and the two official domains}
\label{sec:released-endpoint}

This chapter reconstructs the final argument of the released manuscript,
Section 10, pp.~117--126. It uses the actual local fields constructed in its
preceding sections. The complete independent review of those constructions
is recorded separately; this chapter is not a declaration that the entire
manuscript has already passed that review. The maps, cutoff terms, limiting
arguments and comparison estimates used here are written out explicitly.

\subsection{The local representatives and the localization map}

The retained parameters are
\[
 0<h<\frac1{100},\qquad A=\frac12+h,\qquad D=\frac12-h,\qquad
 \tau=1-t=q(1-\eta^2),\qquad z=q^D\eta,\qquad X=\frac{r^2}{2q}.
\]
Here $r=(x_1^2+x_2^2)^{1/2}$, $q>0$, $|\eta|<1$ before the
terminal time. The constant $q_*>0$ is the fixed local-domain threshold
from Theorem 3.1 (chosen there uniformly for the finite correction stages).
The local domain is $\Omega_*=\{\tau>0,\ q<q_*\}$.
Write the actual local solution representatives as
\[
 u_{\rm loc}=\operatorname{curl}\mathcal A+B e_\theta,\qquad p_{\rm loc}.
\]
The separate scalar $B$ is independent of $\theta$. This notation
$\mathcal A$ for the vector potential avoids confusing it with the
retained scalar exponent $A$.

The background contribution to the potential is determined by the actual
axial coefficients $U_n$:
\[
 S_n=q^{1-A+2nh}\int_0^XU_n(X',\eta)\,dX',\qquad
 S=S_0+\sum_{n\ge1}\chi(c_nq)S_n,\qquad
 \mathcal A_{\rm base}=\frac{S}{r}e_\theta.
\]
Each full radial integral of $U_n$ is zero. Consequently $S_n=0$
beyond the common axial support, with its cutoff retained in $S$.
Near the axis $S=r^2a(r^2,z,t)$ and hence
\[
 \mathcal A_{\rm base}=a(r^2,z,t)(-x_2,x_1,0).
\]
This is a smooth Cartesian vector field; no division by the coordinate
$r$ remains in this representative. The wave potentials are supported
in the active annulus. A mean axial increment $\gamma_{d,\rm phys}$
has potential coefficient
\[
 \Psi=r^{-1}I_c(r\gamma_{d,\rm phys}),\qquad
 I_cg=Ig-\chi_m Jg,\quad
 Ig(r)=\int_0^rg(s)\,ds,\quad Jg=\int_0^\infty g(s)\,ds.
\]
Beyond the source and transition of $\chi_m$, $Ig=Jg$ and $\chi_m=1$;
therefore this potential vanishes there as well. These facts give the
actual representative $\mathcal A=0$ for $X\ge X_{\rm ext}$.
All annular terms vanish near the axis for each fixed $t<1$.

The coordinates give an elementary useful estimate without an
asymptotic replacement. If $1-\eta^2\ge1/2$, then $q\le2\tau$.
Otherwise $|z|=q^D|\eta|>q^D/\sqrt2$, so
$q<(\sqrt2|z|)^{1/D}$. Thus, for a fixed $C_0>0$,
\[
 q\le C_0\big(\tau+|z|^{1/D}\big).
\]
Choose $0<\tau_0<1/2$ and $z_0>0$ such that
$C_0(\tau_0+z_0^{1/D})<q_*/2$. Choose a smooth axisymmetric
$\chi_x$, smooth in $(r^2,z)$, equal to one for
$r\le r_0/2,\ |z|\le z_0/2$, with compact support in
$\{r<r_0,\ |z|<z_0\}$. Choose $\chi_t$ zero for
$1-t\ge\tau_0$ and one for $0\le1-t\le\tau_0/2$.
Put $c=\chi_x\chi_t$, and define, before time one,
\[
 u=\operatorname{curl}(c\mathcal A)+cB e_\theta,\qquad p=cp_{\rm loc}.
 \tag{E.1}
\]
These expressions are extended by zero off the cutoff support. The
support is contained in the fixed compact set $K=\operatorname{supp}\chi_x$,
and its part in the local domain has the strict margin $q<q_*/2$.
Smoothness of the cutoffs, the Cartesian axis formulas and this margin
prove smoothness of the extended fields.

The exact product formula is
\[
 u=c u_{\rm loc}+\nabla c\times\mathcal A.
 \tag{E.2}
\]
Incompressibility follows directly:
\[
 \operatorname{div}u
 =\operatorname{div}\operatorname{curl}(c\mathcal A)
   +r^{-1}\partial_\theta(cB)=0.
\]
The axis identity extends by Cartesian smoothness. The extra term in
(E.2) is retained throughout; localization has not been asserted to
commute with curl. The time cutoff gives $u=p=0$ on an interval
after time zero.

Define the full residual
\[
 f=\partial_tu+(u\cdot\nabla)u-\Delta u+\nabla p .
 \tag{E.3}
\]
This defines the required force before time one and includes every
derivative and quadratic product of the cutoffs in (E.1).
Taking the divergence of this exact identity gives
\[
 -\Delta p=\sum_{i,j=1}^3\partial_i\partial_j(u_i u_j)
                   -\operatorname{div}f.
\]
There is no divergence-free restriction on the prescribed force in this
identity or in the official forced alternatives.

\subsection{All terminal derivatives, including the exterior pressure}

The local construction supplies bounds on every derivative of
$\mathcal A,B e_\theta,p_{\rm loc}$ on compact sets on which $q$ is
bounded away from zero and $q_*$. These are bounds for the actual
representatives on the closed parameter interval $-1\le\eta\le1$.
For one such set $K_0$, a component $G$, a spatial multi-index $\alpha$
and an integer $j\ge0$, the identity
\[
 \partial_x^\alpha\partial_t^jG(x,t')
 -\partial_x^\alpha\partial_t^jG(x,t)
 =\int_t^{t'}\partial_x^\alpha\partial_s^{j+1}G(x,s)\,ds
\]
shows that this derivative is uniformly Cauchy on $K_0$ as $t,t'\uparrow1$.
Multiplication by the fixed cutoffs and the finite differential
polynomial (E.3) preserves this property.

For positive radius and sufficiently small $q$, the representative is the
exact heat exterior. With $s=r^2/2$, its azimuthal velocity is
\[
 K_{\rm heat}(r,\tau)=c_\infty s^{-A}H(2\tau/s)
       =r^{-1-2h}H_{\rm ext}(\tau/r^2),
 \qquad H_{\rm ext}(Z)=2^Ac_\infty H(4Z).
 \tag{E.4}
\]
The integral defining $H$ in the exterior construction can be
differentiated $j$ times using an integrable majorant proportional to
$e^{-v}v^{h+j}$, uniformly for its nonnegative argument. Thus
\[
 |\partial_\tau^jK_{\rm heat}(r,\tau)|
       \le C_j r^{-1-2h-2j},\qquad
 \left|\partial_\tau^j
       \frac{K_{\rm heat}(\rho,\tau)^2}{\rho}\right|
       \le C_j\rho^{-3-4h-2j}.
 \tag{E.5}
\]
The second bound also follows by writing the complete Leibniz sum
$\sum_{a=0}^j\binom ja
(\partial_\tau^aK_{\rm heat})(\partial_\tau^{j-a}K_{\rm heat})/\rho$;
each summand has exactly the same displayed radial power.

The pressure representative is fixed by radial infinity:
\[
 p_{\rm loc}(r,\tau)=-\int_r^\infty
                 \frac{K_{\rm heat}(\rho,\tau)^2}{\rho}\,d\rho.
 \tag{E.6}
\]
On $r\in[r_-,r_+]$ with $r_->0$, the majorant in (E.5) is integrable
over $[r_-,\infty)$. It justifies all time differentiations in (E.6),
and dominated convergence gives their limits at $\tau=0$ uniformly
on this radial interval. Differentiation of the lower endpoint gives
the radial derivatives; differentiating (E.4) gives the velocity's
radial derivatives. The sign $\partial_t=-\partial_\tau$ is retained.
Cartesian derivatives follow from these radial derivatives with bounded
coefficients on a positive-radius compact set. Since $\mathcal A=0$
here, the localized expressions involve only these fields and the fixed
cutoffs. The heat equation and
$\partial_r p_{\rm loc}=K_{\rm heat}^2/r$ give zero uncut residual.

These two arguments cover the terminal spatial slice away from the
origin. At $z\ne0$, $q\ge |z|^{1/D}>0$. At $z=0$ and $x\ne0$,
the radius is positive and the heat exterior applies in a neighborhood.
There are therefore limits for all mixed derivatives of (E.3), uniform
on each compact spatial set excluding the origin.

Near the origin both cutoffs are one. For $X\le X_{\rm ext}$ the
constructed local residual has every derivative of order $O(q^N)$
for every integer $N\ge0$, and beyond $X_{\rm ext}$ it is exactly zero.
The coordinate bound above therefore gives
\[
 |\partial_x^\alpha\partial_t^jf(x,t)|
 \le C_{\alpha,j,N}\big(\tau+|z|^{1/D}\big)^N
 \tag{E.7}
\]
on a fixed neighborhood of $(0,1)$. For any tolerance choose first a
small ball and then a late time making (E.7) smaller than that tolerance.
On the remaining compact part of $K$, the previous uniform convergence
applies using a finite cover. All derivatives vanish off $K$.
They are consequently uniformly Cauchy on all of $\mathbb R^3$.

Let $F_j(x)=\lim_{t\uparrow1}\partial_t^jf(x,t)$. To identify spatial
derivatives, pass to the uniform limit in
\[
 \partial_x^\alpha\partial_t^jf(x+se_i,t)
 -\partial_x^\alpha\partial_t^jf(x,t)
 =\int_0^s\partial_{x_i}\partial_x^\alpha\partial_t^jf(x+ve_i,t)\,dv.
\]
Induction on $|\alpha|$ shows that all the derivative limits are
$\partial_x^\alpha F_j$. Thus $F_j\in C_c^\infty(\mathbb R^3;\mathbb R^3)$,
$\operatorname{supp}F_j\subset K$, and all its spatial derivatives
vanish at the origin. Passing to the limit in the time integral gives
the exact compatibility identity
\[
 \partial_x^\alpha\partial_t^jf(x,t)
 =\partial_x^\alpha F_j(x)
       -\int_t^1\partial_x^\alpha\partial_s^{j+1}f(x,s)\,ds.
 \tag{E.8}
\]

\subsection{A concrete smooth extension of the prescribed force}

Fix $\chi_0\in C_c^\infty(\mathbb R)$ equal to one near zero and
zero on $[1,\infty)$. For $\sigma=t-1\ge0$ define
\[
 f(x,1+\sigma)=
 \sum_{j=0}^\infty \chi_0(b_j\sigma)\frac{\sigma^j}{j!}F_j(x).
 \tag{E.9}
\]
The positive integers $b_j$ will be specified by finite inequalities.
For a temporal derivative of order $m$, the complete product rule is
\[
 \partial_\sigma^m\left(\chi_0(b\sigma)\frac{\sigma^j}{j!}\right)
 =\sum_{\substack{0\le a\le m\\m-a\le j}}
   \binom ma b^a\chi_0^{(a)}(b\sigma)
                    \frac{\sigma^{j-m+a}}{(j-m+a)!}.
\]
On $\sigma\ge0$ every term is supported in $0\le\sigma\le b^{-1}$.
Its supremum is bounded by its fixed coefficient times $b^{m-j}$.
For example one can use the finite explicit constant
\[
 C_{j,m}=\sum_{\substack{0\le a\le m\\m-a\le j}}
        \binom ma\frac{\|\chi_0^{(a)}\|_\infty}{(j-m+a)!}.
\]
Thus, with the usual maximum norm over derivatives up to order $a_0$,
\[
 \left\|\partial_x^\alpha\partial_\sigma^m
       \left(\chi_0(b\sigma)\frac{\sigma^j}{j!}F_j\right)\right\|_\infty
 \le C_{j,m}b^{m-j}\|F_j\|_{C^{|\alpha|}}.
 \tag{E.10}
\]
Take $b_0=1$. For $j\ge1$ take the least integer $b>b_{j-1}$
for which
\[
 C_{j,m}b^{m-j}\|F_j\|_{C^{a_0}}\le2^{-j}
 \quad\hbox{for every }a_0,m\ge0,\quad
                   a_0+m\le\lfloor j/2\rfloor.
 \tag{E.11}
\]
There are finitely many constraints, and each exponent $m-j$ is
strictly negative. Therefore the set defining this least integer is
nonempty. This specifies the sequence without a uniform bound on the
size of the jets $F_j$.

For fixed $\alpha,m$, all sufficiently late summands in (E.9) satisfy
the summable bound $2^{-j}$. The series of each differentiated order
converges uniformly. Repeated integration of this uniform series
shows that it is the corresponding derivative of (E.9). At $\sigma=0$
the cutoff in each fixed summand is constant on a neighborhood, so its
$m$th derivative contributes only when $j=m$, with value $F_m$.
The mixed derivatives therefore match (E.8) exactly, proving smooth
gluing across $t=1$. The same uniform convergence proves smoothness
at the spatial support boundary and at $t=2$. The original field
vanishes near $t=0$. Hence
\[
 f\in C_c^\infty(\mathbb R^3\times(0,\infty);\mathbb R^3).
 \tag{E.12}
\]
If $K\subset B(0,R_0)$ and
$M_{\alpha,m}=\|\partial_x^\alpha\partial_t^mf\|_\infty$, its
actual support implies, for every integer $k\ge0$,
\[
 |\partial_x^\alpha\partial_t^mf(x,t)|
 \le M_{\alpha,m}(3+R_0)^k(1+|x|+t)^{-k},\qquad t\ge0.
 \tag{E.13}
\]
On the support $1+|x|+t\le3+R_0$; off the support the left side is
zero. This proves all the original whole-space force decay conditions.

\subsection{The energy estimate retains the dissipation}

On every $[0,T]$ with $T<1$, the fields in (E.1) are smooth with fixed
compact support. Integration by parts and incompressibility give
\[
 \frac12\frac{d}{dt}\|u(t)\|_2^2+\|\nabla u(t)\|_2^2
                =\langle f(t),u(t)\rangle.
 \tag{E.14}
\]
The transport integral is the integral of
$\operatorname{div}(u|u|^2/2)$ and the pressure integral is the integral
of $\operatorname{div}(pu)$, so both vanish. Put
$F(t)=\int_0^t\|f(s)\|_2\,ds$. For $\delta>0$ division of (E.14) by
$(\|u\|_2^2+\delta^2)^{1/2}$ yields
\[
 \frac{d}{dt}(\|u\|_2^2+\delta^2)^{1/2}
 \le\|f(t)\|_2.
\]
After integration from zero and $\delta\downarrow0$, $\|u(t)\|_2\le F(t)$.
Returning to the original equality, without discarding its dissipation,
gives
\[
 \|u(t)\|_2^2+2\int_0^t\|\nabla u(s)\|_2^2\,ds
 \le2\int_0^t F'(s)F(s)\,ds=F(t)^2.
 \tag{E.15}
\]
The smooth compact force has $F(1)<\infty$. This proves a uniform
kinetic-energy bound and finite total dissipation on $[0,1)$ by
monotone convergence.

\subsection{Comparison without a pressure condition at infinity}

Fix $T<1$. Let $(v,P)$ be any smooth solution with the same force and
zero datum, with $\sup_{0\le t\le T}\|v(t)\|_2<\infty$. Set
\[
 w=v-u,\quad \pi=P-p,\quad
 g_{ij}=v_iv_j-u_iu_j=w_iw_j+w_iu_j+u_iw_j.
\]
Then $w$ is uniformly in $L^2$ and $g$ uniformly in $L^1$.
The difference equation, with the tensor convention
$(\operatorname{div}g)_i=\sum_j\partial_jg_{ij}$, is
\[
 \partial_tw+(v\cdot\nabla)w+(w\cdot\nabla)u
       =\Delta w-\nabla\pi,\qquad \operatorname{div}w=0.
 \tag{E.16}
\]
Its conservative version has $\operatorname{div}g$ as the transport
term.

Use the Fourier convention with derivative multiplier $i\xi_i$ and
Riesz multiplier $i\xi_i/|\xi|$. Define
\[
 \pi_*=\sum_{i,j}R_iR_jg_{ij}.
 \tag{E.17}
\]
For $s>3/2$, boundedness of the Fourier transform of each $L^1$
component gives
\[
 \|\pi_*\|_{H^{-s}}^2
 \le C_s\left(\sum_{i,j}\|g_{ij}\|_1\right)^2
          \int_{\mathbb R^3}(1+|\xi|^2)^{-s}\,d\xi<\infty.
\]
The value of the multiplier at the origin does not affect this
distribution. Its sign gives
$\Delta\pi_*=-\sum_{i,j}\partial_i\partial_jg_{ij}$, the same
Poisson equation as $\pi$.

The pressure gradient itself is fixed by the equation. For
$a\in C_c^\infty(0,T)$ its time integral satisfies
\[
 \int_0^T a\nabla\pi\,dt
 =\Delta\int_0^T aw\,dt+\int_0^T a'w\,dt
                -\operatorname{div}\int_0^T ag\,dt.
\]
These terms belong respectively to $H^{-2}$, $L^2$ and $H^{-3}$.
The last inclusion follows from $L^1\subset H^{-2}$ by the preceding
Fourier bound and one derivative. Also $\pi_*\in L^\infty_tH^{-2}_x$.
Consequently
\[
 H_a=\int_0^T a(\nabla\pi-\nabla\pi_*)\,dt\in H^{-3},
 \qquad \Delta H_a=0.
\]
The Fourier transform of $H_a$ is a function with finite weighted
$L^2$ norm. The equation $|\xi|^2\widehat H_a=0$ makes that function
zero off the singleton $\{0\}$; a function supported on a set of
Lebesgue measure zero is zero almost everywhere. Thus $H_a=0$.
Testing in space and time yields $\nabla\pi=\nabla\pi_*$ in the
interior. No spatial growth hypothesis on $P$ has been inserted.

Choose $0\le\phi\le1$ smooth and compactly supported, equal to one on
the unit ball. For $R\ge1$ put
\[
 \phi_R(x)=\phi(x/R),\quad \chi_R=\phi_R^8,\quad
 E_R=\int\chi_R|w|^2,\quad
 A_R=\left(\int\chi_R|\nabla w|^2\right)^{1/2},\quad
 B_R=\|\phi_R^4w\|_6.
\]
The ordinary three-dimensional homogeneous Sobolev inequality and
the full product rule give
\[
 B_R\le C\|\nabla(\phi_R^4w)\|_2
     \le C(A_R+R^{-1}\|w\|_2).
 \tag{E.18}
\]
Using the classical $L^p$ boundedness of Riesz transforms for
$1<p<\infty$, write exactly
\[
 \phi_R^4\pi_*=
 \sum_{i,j}R_iR_j(\phi_R^4g_{ij})+
 \sum_{i,j}[\phi_R^4,R_iR_j]g_{ij}.
 \tag{E.19}
\]
The first sum has $L^{3/2}$ norm at most $C_T(B_R+1)$, since
\[
 \|\phi_R^4w_iw_j\|_{3/2}\le B_R\|w\|_2,\quad
 \|\phi_R^4w_iu_j\|_{3/2}
       \le\|w\|_2\|u\|_6,
\]
and the same estimate holds for $u_iw_j$.

The nonlocal kernel of $R_iR_j$ is bounded by $C|x-y|^{-3}$.
Its possible local identity term cancels in the commutator. The
Lipschitz and supremum bounds on $\phi_R^4$ therefore majorize its
kernel by
\[
 K_R(x-y)=C|x-y|^{-3}\min\{|x-y|/R,1\}.
\]
The radial integral is
\[
 \|K_R\|_{4/3}^{4/3}
 \le C\left(R^{-4/3}\int_0^Rr^{-2/3}\,dr
                      +\int_R^\infty r^{-2}\,dr\right)
 \le CR^{-1}.
 \tag{E.20}
\]
Convolution of this kernel with $g\in L^1$ bounds the second sum of
(E.19) in $L^{4/3}$ by $C_TR^{-3/4}$. One may obtain (E.19) first for
compact smooth cutoffs of $g$, then pass to the limit: $g$ converges
in $L^1$, its Riesz transform in $H^{-s}$, and the commutator in
$L^{4/3}$ by (E.20). On every compact set the first term is in
$L^{3/2}$ because $w$ is smooth there. Thus $\pi_*$ is locally
integrable and may be used in the following pressure pairing.

Testing the equality of pressure gradients against the compact field
$a(t)\chi_Rw$ and using $\operatorname{div}w=0$ gives, for almost
every $t$,
\[
 \int\pi w\cdot\nabla\chi_R=\int\pi_*w\cdot\nabla\chi_R.
\]
Since $|\nabla\chi_R|\le CR^{-1}\phi_R^7$, (E.19) pairs with
$CR^{-1}\phi_R^3|w|$. The exact weighted interpolations are
\[
 \|\phi_R^3w\|_3\le\|\phi_R^2w\|_3
                 \le B_R^{1/2}\|w\|_2^{1/2},\qquad
 \|\phi_R^3w\|_4\le B_R^{3/4}\|w\|_2^{1/4}.
\]
For the first use $\phi_R^2|w|=(\phi_R^4|w|)^{1/2}|w|^{1/2}$;
for the second use powers $3/4$ and $1/4$. Hölder then proves
\[
 \left|\int\pi w\cdot\nabla\chi_R\right|
 \le C_TR^{-1}
       \left((B_R+1)B_R^{1/2}+R^{-3/4}B_R^{3/4}\right).
 \tag{E.21}
\]

All integrations in the difference energy identity are compact:
\[
 \begin{split}
 \frac12E_R'+A_R^2={}&-\int\chi_R(w\cdot\nabla)u\cdot w
       +\frac12\int|w|^2\Delta\chi_R\\
 &+\frac12\int|w|^2v\cdot\nabla\chi_R
       +\int\pi w\cdot\nabla\chi_R .
 \end{split}
 \tag{E.22}
\]
For $R$ sufficiently large, $\chi_R=1$ on a neighborhood of the
fixed support of $u$, hence $v=w$ on $\operatorname{supp}\nabla\chi_R$.
The transport term is at most
\[
 CR^{-1}\int\phi_R^6|w|^3
 \le C_TR^{-1}B_R^{3/2},
\]
where $0\le\phi_R\le1$ and the preceding interpolation were used.
The Laplacian term is bounded by $C_TR^{-2}$.
The remaining non-flux term is at most $\|\nabla u\|_\infty E_R$.
Substitute (E.18) into (E.21). Every resulting positive power of
$A_R$ is at most $3/2<2$. For example the scalar Young inequality
gives $R^{-1}A_R^{3/2}\le\delta A_R^2+C_\delta R^{-4}$;
the smaller powers are absorbed by the same method. Choosing the
finite sum of dissipation coefficients below $1/2$ proves
\[
 \frac12E_R'+\frac12A_R^2
       \le\|\nabla u\|_\infty E_R+\frac{C_T}{R}.
\]
Here $\|\nabla u\|_\infty$ is bounded on $[0,T]$, and $E_R(0)=0$.
Integrating the differential inequality gives explicitly
\[
 E_R(t)\le\frac{2C_T}{R}\int_0^t
     \exp\left(2\int_s^t\|\nabla u(a)\|_\infty\,da\right)\,ds
       \le\frac{C'_T}{R}.
\]
Each fixed ball has $\chi_R=1$ for all sufficiently large $R$.
Letting $R\to\infty$ shows its integral of $|w|^2$ is zero.
Smoothness gives $w=0$ everywhere on $[0,T]$.

\subsection{The actual path of velocity growth}

The local inner asymptotic retained from the constructed profiles is
\[
 x_\tau=(\sqrt{2X_{\rm in}\tau},0,0),\qquad
 u_\theta(x_\tau,1-\tau)=
       \tau^{-A}\big(e_0+O(\tau^{2h})\big),\qquad e_0>0.
 \tag{E.23}
\]
Along this path $z=0$, $q=\tau$, $X=X_{\rm in}$; the annular
corrections vanish, and both localization cutoffs are one for small
$\tau$. The cutoff sum of all positive-order background terms is
included in the displayed error. Since $A>0$ and $h>0$, the bracket
is at least $e_0/2$ for sufficiently small $\tau$ and the velocity
diverges. Its spatial points converge to the origin.

Every global smooth competitor satisfying the official uniform
kinetic-energy condition agrees with $u$ on each $[0,T]$, $T<1$,
by the comparison calculation. Such a competitor is bounded on a
compact neighborhood of $(0,1)$ by continuity, contradicting
(E.23). This conclusion uses the actual earlier profile and residual
construction; it does not follow from a singular expression alone.
The prescribed force is nonzero, since $f=0$ would make (E.15) imply
$u=0$.

\subsection{Exact viscosity and parabolic maps}

For the residual convention
$\mathcal R_\nu(u,p,f)=\partial_tu+(u\cdot\nabla)u-\nu\Delta u
 +\nabla p-f$, define the map from viscosity one to viscosity $\nu>0$:
\[
 \mathscr A_\nu(u,p,f)(x,t)=
 \big(\sqrt\nu\,u(x/\sqrt\nu,t),\
       \nu p(x/\sqrt\nu,t),\
       \sqrt\nu\,f(x/\sqrt\nu,t)\big).
 \tag{E.24}
\]
With $y=x/\sqrt\nu$, its time, convection, viscous and pressure
terms are respectively
$\sqrt\nu\,\partial_tu$, $\sqrt\nu\,(u\cdot\nabla_y)u$,
$-\sqrt\nu\,\Delta_yu$, $\sqrt\nu\,\nabla_yp$. Therefore
\[
 \mathcal R_\nu(\mathscr A_\nu(u,p,f))(x,t)
       =\sqrt\nu\,\mathcal R_1(u,p,f)(y,t),\qquad
 \operatorname{div}_x u^\nu=\operatorname{div}_yu.
\]
The map is invertible on its image:
$(v,P,F)\mapsto(\nu^{-1/2}v(\sqrt\nu y,t),
\nu^{-1}P(\sqrt\nu y,t),\nu^{-1/2}F(\sqrt\nu y,t))$.
It preserves time and zero initial data, sends $K$ to $\sqrt\nu K$,
and obeys, for every $\alpha,m$,
\[
 \partial_x^\alpha\partial_t^mf^\nu
       =\nu^{(1-|\alpha|)/2}
          (\partial_y^\alpha\partial_t^mf)(y,t).
\]
The volume Jacobian is $\nu^{3/2}$, so
\[
 \|u^\nu(t)\|_2^2=\nu^{5/2}\|u(t)\|_2^2,\qquad
 \nu\int_0^1\|\nabla u^\nu(t)\|_2^2dt
       =\nu^{5/2}\int_0^1\|\nabla u(t)\|_2^2dt.
 \tag{E.25}
\]
The same inverse map sends any bounded-energy competitor to the
already excluded viscosity-one competitor.

At a fixed viscosity define the parabolic map
\[
 \mathscr S_\lambda(v,P,F)(x,t)=
   \big(\lambda v(\lambda x,\lambda^2t),\
        \lambda^2P(\lambda x,\lambda^2t),\
        \lambda^3F(\lambda x,\lambda^2t)\big),\quad\lambda>0.
 \tag{E.26}
\]
Its momentum residual is $\lambda^3$ times the original residual
at $(\lambda x,\lambda^2t)$, and divergence is $\lambda^2$ times
the original divergence. Its inverse is $\mathscr S_{1/\lambda}$.
The instantaneous energy factor is $\lambda^{-1}$.

The formalization uses a different viscosity adapter,
\[
 \mathscr T_\nu(u,p,f)(x,t)=
       \big(\nu u(x,\nu t),\nu^2p(x,\nu t),\nu^2f(x,\nu t)\big).
 \tag{E.27}
\]
Every momentum term at viscosity $\nu$ has factor $\nu^2$ times
its viscosity-one counterpart. Its singular time is $1/\nu$,
its energy factor is $\nu^2$, and it preserves the unit spatial
period. These two actual adapters are related by the exact identity
on the entire velocity--pressure--force triple
\[
                  \mathscr T_\nu
                       =\mathscr S_{\sqrt\nu}\circ\mathscr A_\nu.
 \tag{E.28}
\]
Indeed the arguments become
$(\sqrt\nu x/\sqrt\nu,(\sqrt\nu)^2t)=(x,\nu t)$,
and the three prefactors become $\nu,\nu^2,\nu^2$.
The energy factors compose as $\nu^{-1/2}\nu^{5/2}=\nu^2$.
For a unit-periodic input, $\mathscr A_\nu$ gives period $\sqrt\nu$,
and $\mathscr S_{\sqrt\nu}$ returns it to period one.
Thus their different intermediate periods and singular times are
accounted for by a proved morphism. This identity concerns the two
adapters; it does not assert without inspection that their input
representatives in the paper and Lean source are literally identical.

\subsection{The unit-periodic construction, with periodic pressure}

Fix $\nu>0$ and use the fields at that viscosity just constructed.
Choose $\lambda>1$ with
$\lambda^{-1}\sqrt\nu K\Subset Q_0=(-1/2,1/2)^3$, and let
$t_0=1-\lambda^{-2}$. For $t\ge t_0$ put
\[
 \widetilde u(x,t)=\lambda u^\nu(\lambda x,\lambda^2(t-t_0)),\quad
 \widetilde p(x,t)=\lambda^2p^\nu(\lambda x,\lambda^2(t-t_0)),\quad
 \widetilde f(x,t)=\lambda^3f^\nu(\lambda x,\lambda^2(t-t_0)).
 \tag{E.29}
\]
Use the first two formulas before time one and the force formula for
all subsequent time; extend all three by zero for $t<t_0$.
The original fields vanish on an initial time interval, so this
extension is smooth. All momentum terms scale by $\lambda^3$ and
the viscosity stays $\nu$. Original time one is reached at time one
because $\lambda^2(1-t_0)=1$.

Define on $\mathbb R^3/\mathbb Z^3$
\[
 U(x,t)=\sum_{k\in\mathbb Z^3}\widetilde u(x+k,t),\quad
 P_{\rm per}(x,t)=\sum_{k\in\mathbb Z^3}\widetilde p(x+k,t),\quad
 F_{\rm per}(x,t)=\sum_{k\in\mathbb Z^3}\widetilde f(x+k,t).
 \tag{E.30}
\]
The supports of distinct translates have a positive gap, since
the unshifted compact support lies strictly inside $Q_0$.
The sums are locally finite and at each point at most one
translate is nonzero. The support of each derivative remains in
the same closed support. Consequently the full convection term is
\[
 (U\cdot\nabla)U
       =\sum_{k\in\mathbb Z^3}
             (\widetilde u(x+k,t)\cdot\nabla)\widetilde u(x+k,t).
\]
All cross terms vanish by disjointness, not by linearity of the
equation. Termwise differentiation now proves
\[
 \partial_tU+(U\cdot\nabla)U-\nu\Delta U+\nabla P_{\rm per}
        =F_{\rm per},\qquad
 \operatorname{div}U=0,\qquad U(\cdot,0)=0.
\]
Integer translation reindexes each sum, proving unit periodicity of
velocity, pressure and force separately. The force has time support
in $[t_0,1+\lambda^{-2}]$, and hence every prescribed periodic
force derivative has arbitrary polynomial time decay.

The precise path of growth is
\[
 \widetilde x_\tau=\lambda^{-1}\sqrt\nu\,x_\tau,\qquad
 \widetilde t_\tau=1-\lambda^{-2}\tau,\qquad
 U_\theta(\widetilde x_\tau,\widetilde t_\tau)
 =\lambda\sqrt\nu\,\tau^{-A}(e_0+O(\tau^{2h})).
\]
For small $\tau$ this lies inside $Q_0$ and only its zeroth
translate contributes. The value diverges as before.
Two smooth periodic solutions with the same data have difference
$w$ satisfying, on each $[0,T]$, $T<1$,
\[
 \frac12\frac{d}{dt}\|w\|_{L^2(\mathbb T^3)}^2
       +\nu\|\nabla w\|_{L^2(\mathbb T^3)}^2
 \le\|\nabla U\|_\infty\|w\|_{L^2(\mathbb T^3)}^2.
\]
Periodic integration removes the transport divergence and pressure
gradient; the remaining quadratic term has the displayed bound.
The zero datum and integration of this inequality force $w=0$.
A global smooth periodic competitor therefore contradicts the same
growth path. Both the original unit periods and the pressure
periodicity required by the official correction have been preserved.


% --- included proof body: profile_assembly_body.tex ---
% Integration body. Requires amsmath, amssymb, amsthm, mathtools, mathrsfs, hyperref.
% All labels and references use pa:; custom macros use pa-prefixed names.
\providecommand{\paHeat}{\mathcal H}
\providecommand{\paResidual}{\mathcal R}
\providecommand{\paReal}{\mathbb R}
\providecommand{\paSupport}{\operatorname{supp}}
\section{The exact constructed objects and the order of the argument}
Retain the original right-handed cylindrical coordinates
\[
 x=(r\cos\theta,r\sin\theta,z),\quad s=r^2/2,\quad
 \tau=1-t=q(1-\eta^2),\quad z=q^D\eta,\quad X=s/q,
\]
where
\[
 0<h<1/100,\quad A=\tfrac12+h,\quad D=\tfrac12-h,
 \quad d=1-\eta^2,\quad L=1-2h\eta^2,\quad -1\le\eta\le1.
\]
The physical chart for $t<1$ has $|\eta|<1$; the profile endpoints have one-sided
derivatives. The coordinate and Cartesian extension proofs are in
Section~\ref{pa:sec:component-core}. No radial dilation below changes these coordinates.
The fields are
\[
 u_\theta=q^{-A}E,\quad u_z=q^{-A}U,\quad ru_r=V_0,\quad
 p=q^{-2A}\Pi,\quad E=\sqrt{2X}F=\sqrt{2X}\phi/C,
 \quad H=\sqrt{2X}E.
\]
Every integral and component order in this document is the original one:
\[
 \begin{aligned}
 M&=\int_0^XU\,dx,&I&=\int_0^XH\,dx,&J&=\int_0^XUH\,dx,\\
 S&=\int_0^X(U^2-E^2/2)\,dx,&
 C_p&=\int_0^XE^2/(2x)\,dx,&\Pi&=\Pi_0+C_p.
 \end{aligned}
\]
The two tangential components are ordered $(r\theta,rz)$. In particular,
\[
\begin{split}
 V_0&=(2\eta XU-2D\eta M-dM_\eta)/L,\\
 W&=1-(2D\eta M+dM_\eta)/X,\\
 Q_s&=-W+\frac{(1-h)I-D\eta I_\eta-dJ_\eta+2(h-D)\eta J}{XH},\\
 N_s&=-WU+\frac{D(M-\eta M_\eta)+4h\eta S-dS_\eta}{X}
            +4A\eta\Pi-d\Pi_\eta,\\
 p_s&=(XQ_s/L,XN_s/(LE)),\quad
 a=1-2X\partial_X\log E,\quad b_s=2XU_X/E,\\
 T_0&=F\{p_s-(a,-b_s)\}.
\end{split}                                                     \tag{PA.1}
\]
At the regular axis the formulas are extended by the proved radial primitive
identities, rather than evaluated as indeterminate quotients.

The construction uses the following directed chain:
\[
\begin{gathered}
\text{outer schedule and moment adjustments}
\longrightarrow \Pi_0\ \text{and the outer family},\\
\Pi_0\longrightarrow\text{analytic axis and inner continuation}
\longrightarrow\text{five-row join},\\
\text{regular join and total moments}
\longrightarrow\text{backward stress and endpoint directions},\\
\text{joined relaxed shear and admissible collars}
\longrightarrow\text{variance loop},\\
\text{variance loop}\longrightarrow\text{finite-frequency field and five-row repair}.
\end{gathered}                                                    \tag{PA.2}
\]
The outer schedule does not use a regular axis field. Its temporary inner branch
$U_o=4\eta$, $E_o=P_*f(X/X_R)^{1/10}$, $f=(1+\eta^2)^{-1}$,
has convergent five prefix moments. The later backward-stress calculation is
applied to the regular joined field actually constructed below. Thus it is a
later use of the schedule, and supplies no premise needed to construct that schedule.

\section{The outer pressure is the actual axis datum}
The outer schedule in Section~\ref{pa:sec:component-outer} first fixes
\[
 M_d>0,\quad T_d=e^{M_d}+10,\quad P_*>e^{T_d},\quad
 \lambda>0,\quad 0<h<\min\{1/100,\lambda,e^{-T_d}\}.                 \tag{PA.3}
\]
Its successive slope changes, finite pulse, affine two-bump axial solve,
pressure-neutral angular solve, amplitude root, and finite tail waiting length
are proved there. In particular no stage length depends on $\eta$ or on $X_R$.
After omitting only the angular edit of zero total pressure increment, write
its swirl exactly as
\[
 E_{\mathrm{sched}}(y,\eta)=c(y)f(\eta)^{\vartheta(y)},\qquad
 y=\log(X/X_R),\quad 0\le\vartheta\le1.
\]
Here $\vartheta=1$ and $c(y)=P_*e^{y/10}$ for $y\le0$; in the
final tail $\vartheta=0$ and $c(y)$ decays with exponent $-A$.
Define, before selecting $X_R$,
\[
 \Pi_0(\eta)=-\frac12\int_{-\infty}^{\infty}
                  c(y)^2f(\eta)^{2\vartheta(y)}\,dy.               \tag{PA.4}
\]
For $z=x+iu$, $|u|<1/4$, the real part of $1+z^2$ is
$1+x^2-u^2>15/16$. The analytic principal logarithm of $1+z^2$ therefore
defines $f(z)^{2\vartheta}=\exp[-2\vartheta\log(1+z^2)]$ throughout
this strip. On each compact subset it is bounded uniformly for
$0\le\vartheta\le1$. The integrable bounds $C e^{y/5}$ at $-\infty$
and $C e^{-(1+2h)y}$ at $+\infty$, and the finite intervening length,
give uniform convergence on compact complex subsets. Integrating on rectangles
and applying Fubini and Cauchy's theorem proves holomorphy of (PA.4); alternatively,
uniform convergence of truncated holomorphic integrals proves the same statement.
The integrand is real positive for real $\eta$ and even. As $f'=-2\eta f^2$,
dominated differentiation gives the stronger exact bounds
\[
 \Pi_0\le-\frac52P_*^2f^2,\qquad
 \eta\Pi_0'=2\eta^2 f\int\vartheta c^2f^{2\vartheta}\,dy
       \ge10P_*^2\eta^2 f^3>0\quad(\eta\ne0).                    \tag{PA.5}
\]
The reference half-line contributes $\int_{-\infty}^0e^{y/5}dy=5$;
all other contributions to the derivative have the same sign.

The scale independence and its survival under the actual modifications are exact.
The substitution $X=X_Re^y$ gives $dX/(2X)=dy/2$, so (PA.4) contains no
$X_R$. The two angular bumps restore their total $\int E^2dy$ to its pre-edit
value. The heat replacement and its earlier three-row correction restore together
the total increments in $C_p,S,I$ (with $I$ understood as its convergent
subtracted total). Their equations are proved in
Section~\ref{pa:sec:component-outer}; $M,J$ are unchanged because $U=0$ on both
supports. Thus the complete prepared outer profile has
\[
 \int_0^\infty E_o^2/(2X)\,dX=-\Pi_0                           \tag{PA.6}
\]
for every sufficiently large $X_R$, with this very same function $\Pi_0$.
The pressure used in the axis equation is (PA.4), not a later approximation to it.

For completeness, positivity at the turning point of the axis equation has an
explicit bound. Choose $0<j_0\le1/20$, put
\[
 U_*=4\eta+j_0,\quad H_*=D\eta+dU_*,\quad
 Z_*=-A(1-2\eta U_*)U_*-H_*U_*'-d\Pi_0'+4A\eta\Pi_0.
\]
The function $H_*/d=D\eta/(1-\eta^2)+4\eta+j_0$ has derivative
$D(1+\eta^2)/(1-\eta^2)^2+4>0$, with endpoint limits of opposite
infinite sign. Its unique zero is $\eta_0=-s$, where
\[
 0<s\le1/80,\quad j_0/5<s<j_0/4,\quad
 d_0=1-s^2\ge6399/6400,\quad U_*(\eta_0)=Ds/d_0.
\]
The first bound follows from $s=j_0/(4+D/d_0)$; it then gives
$D/d_0\le3200/6399<1$, proving the strict lower bound.
At this zero, $-H_*U_*'=0$ and $-d_0\Pi_0'(\eta_0)>0$.
Using $f^2\ge1/4$ and (PA.5), the positive pressure term is at least
$A j_0P_*^2/2$. On the other hand,
\[
 U_*(\eta_0)/j_0\le800/6399,\quad
 2sU_*(\eta_0)\le1/6399,
\]
so its only negative term has absolute value at most
\[
 A j_0\frac{5120000}{40947201}<A j_0/7.
\]
Consequently
\[
 Z_*(\eta_0)>A j_0(P_*^2/2-1/7)>0.                              \tag{PA.7}
\]
This proves the exact hypothesis of the axis construction for the already selected
$P_*>1$, without a further dependence of $P_*$ on $j_0$ or $X_R$.
Choose a neighborhood $V$ of $\eta_0$ with $Z_*>Z_*(\eta_0)/2$ and
$0<\delta_*<Z_*(\eta_0)/4$. On the compact set
$K=\{|Z_*|\le\delta_*\}$, the continuous function $|H_*|$ has a
positive minimum $m$ when $K$ is nonempty. Choose
$0<\sigma_*^2<m^2/99$; for empty $K$ take any sufficiently small positive
$\sigma_*$. Then
\[
 \chi=H_*^2/(H_*^2+\sigma_*^2)>.99\quad\hbox{on }K.              \tag{PA.8}
\]
After a small real and complex enlargement, $L$ and $H_*^2+\sigma_*^2$
remain nonzero and $\Pi_0$ remains holomorphic. These are all the analytic
datum hypotheses used by the coefficient-space contraction in
Section~\ref{pa:sec:component-axis}.
Denote the selected simply connected complex neighborhood, with compact closure
inside that nonvanishing domain, by $\Omega$.

\section{One finite choice of radius, amplitude, and joining widths}\label{pa:sec:assembly-parameters}
The outer schedule estimates use finite constants formed from fixed smooth steps,
finite intervals, and their displayed convolution integrals. The outer proof first
chooses $M_d$ and $P_*$ so the axial-decrease cone gaps are strictly positive,
then chooses $\lambda$ so that
\[
 C_{\rm pre}\sqrt\lambda(1+\log(1/\lambda))
 \quad\hbox{and}\quad C_{\rm pre}\lambda(1+\log(1/\lambda))
\]
are smaller than the fixed strict allowances in its pulse and power-law estimates.
Those allowances are positive: the two pulse quadratic upper bounds are
$2.41^2/8<.74$ and $4.82^2/14<1.68$, and the amplitude-root endpoint
gaps exceed $.047$ and $.038$ before the small error. Choosing $h$ subsequently
small relative to both $\lambda$ and $e^{-T_d}$ gives (PA.3) and the remaining
positive outer gaps. These assertions use the explicit estimates proved in the
outer component, not an assumed cone theorem.

Here are precise inequalities for the later choices. All constants in them are
finite norms, minima, or Lipschitz constants of the displayed maps on specified
compact sets. Fix the finite derivative orders required for the joining inequalities,
including one extra $\eta$ derivative for (PA.1). On
$R=[e^{-8},e^{-5}]\times[-1,1]$, the ideal pair is
\[
 U_{\rm id}=4\eta,\quad E_{\rm id}=P_*fx^{1/10}.
\]
Its angular source is exactly
\[
 S_{q,\rm id}=\frac95-h+\frac{16}{5}h\eta^2
     +\frac{2\eta^2(D+4d)}{1+\eta^2}>0,
\]
and the positive regular weighted primitive gives $Q_{\min}=\min_RQ_{s,\rm id}>0$.
Put $e_*=\min_RE_{\rm id}>0$ and
$w_*=1+\max_R|N_{s,\rm id}/(E_{\rm id}Q_{s,\rm id})|$.
Choose a positive tolerance $\varepsilon_m$ in the normalized field/moment map
and fixed restoration/repair maps sufficiently small to imply
\[
 Q_s\ge Q_{\min}/2,\quad E\ge e_*/2,\quad .7\le a\le.9,
 \quad |N_s/(EQ_s)|\le w_*,\quad |b_s|\le .1/(1+w_*).            \tag{PA.9}
\]
Such a positive tolerance follows directly from the rational formulas (PA.1), the
fixed bump Jacobian inverse and quadratic contraction proved below. The
denominators have positive minima at the ideal pair. In particular the tolerance
contains no $X_R$.

Select $j_0$ and then $\Lambda$ so the finite bounds
$C_kj_0+C_k^{\rm nat}/\Lambda$ are less than one quarter of the required
joining-error allowance for each of the selected orders. In the radial variable
$Y=\Lambda X$, write $\phi=\phi_*\Phi$. Also impose all the
finitely many axis-contraction and endpoint bounds proved in the axis component:
\[
 \Lambda\ge\max\{1,2M,2K\},\quad \Phi\ge c_0>0,\quad
 S_q\ge2.5+.95\Lambda L\chi,
\]
and its two exit alternatives. Here $M,K$ are the actual radius-one-ball norm
and Lipschitz bounds for the rescaled remainder operators, which are uniform in
the later amplitude after its complex bound is imposed.
Choose the amplitude-independent bounds $B_k$ and then
\[
 T_{\rm sh}\ge20\|\sigma'\|_\infty(B_0+\|\log f\|_\infty),
 \qquad X_{\rm sep}=110e^{T_{\rm sh}}.                           \tag{PA.10}
\]
The complete proof of their independence from subsequent $C$ and smaller joining
widths is in Section~\ref{pa:sec:component-continuation}.

Now choose $C$, once, to exceed all the following finite bounds:
\[
\begin{gathered}
 C>1,\qquad C\ge\sup_{\overline\Omega}|\phi_*|,\qquad
 C\ge C_0,\qquad
 C\ge (R_*/110)^{1/10}/P_*,\\
 (CP_*)^{10}>e^{T_{\rm sh}+8},\qquad
 X_R:=110(CP_*)^{10}>14e^8/(3Q_{\min}).                          \tag{PA.11}
\end{gathered}
\]
The threshold $R_*$ includes the prepared outer cone, heat compensation, and
terminal $-Z\paHeat'/\paHeat<h/4$ requirements, all depending only on the fixed outer
data. The threshold $C_0$ includes the uniform exit and reference inequalities
$p_{1,r}+p_{2,r}^2/p_{1,r}>2+c$. Also require, for the finite selected orders,
the vanishing early normalized moment errors and ideal early moments to use less
than one quarter of the remaining matching allowance. This is possible because
\[
 x_{\rm sep}=e^{T_{\rm sh}}/(CP_*)^{10}\longrightarrow0,
\]
\[
 \|\widehat M\|_k+\|\widehat S\|_k\le K_kx_{\rm sep},\quad
 \|\widehat I\|_k+\|\widehat J\|_k\le K_kC^{-1}x_{\rm sep}^{3/2},
 \quad\|C_p\|_k\le K_kC^{-2},                                  \tag{PA.12}
\]
with constants already fixed; their exact integral proof is included below.
The complex supremum in (PA.11) gives uniform $\eta$ derivatives by Cauchy's
estimate; a real supremum alone would not do so.

After fixing this $C$, first choose $\kappa_0>0$ so
\[
 \kappa_0<1/2,\quad \kappa_0V_{\max}<1/4,\quad
 \kappa_0\sup p_{1,r}<.8,
\]
and so the $O(\kappa_0)$ value and moment errors fit their positive allowances.
The number $V_{\max}$ is a common bound on the family $0<t_1\le t_*$,
with $t_*$ fixed before $\kappa_0$. Now choose $0<t_1\le t_*$ and the two final
widths, of sum $\omega_{\rm fin}>0$, sufficiently small that the remaining
$C_k^{\rm join}(\Lambda)(t_1+\kappa_0+\omega_{\rm fin})$ errors fit the
matching allowance and all the activation comparison gaps. The full derivation
has no factor $1/\kappa_0$ in the activation slope. It proves the exact
$e^{-t_1^2/y_a^2}$ factor before any division by it. Thus none of these choices
forces a new earlier choice of $T_{\rm sh}$ or $C$.

The relaxed-cone verification on the matching patch is quantitative. From (PA.9),
\[
 v_s=a+b_s^2/a\le .9+.01/.7<1,\quad
 G=Q_s-b_sN_s/(aE)\ge6Q_s/7\ge3Q_{\min}/7.
\]
Since $L\le1$ and $x\ge e^{-8}$, (PA.11) gives
\[
 P_c=X_RxG/L>2.                                                 \tag{PA.13}
\]
The smaller root $U(P_c,J_c)$ of the cone polynomial is strictly greater than
$2$, as proved in the core and loop components. Thus (PA.13) proves the strict
relaxed cone throughout the actual restoration and correction, with the single
previously fixed tolerance.

\section{Every row of both five-moment repairs and the exterior map}
Write $u=\delta U$, $e=\delta E$. At any fixed reference pair the complete
five increment integrands, in order $(M,I,J,S,C_p)$, are
\[
 \begin{pmatrix}
 u\\ \sqrt{2X}e\\ Hu+U\sqrt{2X}e+\sqrt{2X}ue\\
 2Uu+u^2-Ee-e^2/2\\ Ee/X+e^2/(2X)
 \end{pmatrix}.                                                \tag{PA.14}
\]
This follows by multiplying out the two quadratic and one bilinear original
integrands. In particular the $C_p$ row includes its factor $1/X$ and the
$S$ row its negative angular quadratic term.
Under $X=X_Rx$, the exact moment transformation is
\[
 (M,I,J,S,C_p)
  =(X_R\widehat M,X_R^{3/2}\widehat I,X_R^{3/2}\widehat J,
     X_R\widehat S,\widehat C_p),\quad H=X_R^{1/2}\widehat H.     \tag{PA.15}
\]
Substitution in (PA.1) cancels $X_R$ in $Q_s,N_s$; it retains the factor
$X_R$ in $p_s$.

For the inner join, on $e^{-6}<x<e^{-5}$ the base pair is
$U=4\eta$, $E=P_*fx^{1/10}$. Take two nonnegative ordered disjoint
bumps for $u$ and three for $e$. The invertible row changes
$J\mapsto J-4\eta I$, $S\mapsto S-8\eta M$ use their displayed base
values, while retaining all quadratic terms in (PA.14). The derivative blocks are
\[
 (1,\sqrt2P_*f x^{3/5}),\qquad
 (\sqrt2x^{1/2},-P_*f x^{1/10},P_*f x^{-9/10}).                 \tag{PA.16}
\]
For the later shear repair on a compact subinterval of $I_1$, the base pair is
$U=0$, $E=K(\eta)X^{-1/2-\lambda}$ with $K>0$. In order $(M,J;I,S,C_p)$,
the blocks are exactly
\[
 (1,\sqrt2KX^{-\lambda}),\qquad
 (\sqrt{2X},-KX^{-1/2-\lambda},KX^{-3/2-\lambda}).              \tag{PA.17}
\]
The powers in (PA.16) and (PA.17) are distinct within each block; (PA.17) uses the
already fixed $\lambda>0$. Their matrices are invertible. Indeed a nonzero
combination of $m$ distinct powers has at most $m-1$ positive zeros: divide
by the lowest power and apply Rolle's theorem and induction to its derivative.
The evaluation determinant consequently has constant nonzero sign on the
connected set of ordered positive evaluation points. Multilinearity expresses
the bump-moment determinant as its integral against the nonnegative product of
the bumps. This product is positive on a set of positive measure, so the integral
is nonzero. Compactness in $\eta$ and the adjugate identity bound every fixed
derivative of the inverse.

The exact increment map is $B(\eta)c+Q_\eta(c,c)$. Let $\nu_j$ denote
the $C^j$ multiplication norm of $B^{-1}$ and $q_j$ the bilinear norm of $Q$.
For each of the finitely selected orders $j$, impose
\[
 8\nu_j^2q_j\|d\|_{C^j}\le1.                                  \tag{PA.18}
\]
The map $c\mapsto B^{-1}(d-Q(c,c))$ is a contraction on the radius
$2\nu_j\|d\|_{C^j}$ ball: its Lipschitz constant is at most
$4\nu_j^2q_j\|d\|_{C^j}\le1/2$, and its image norm is at most
three quarters of that radius. Starting from zero in $C^0$ and $C^j$ selects
the same solution. The pointwise Jacobian differs from $B$ by a perturbation
whose $B^{-1}$-norm is at most $1/2$; hence it is invertible. The implicit
function theorem then gives smoothness at every remaining fixed order, with
one-sided endpoint derivatives. These statements prove the complete nonlinear
repair maps, rather than replacing them by their linearizations.

The join takes place at $X_h=X_Re^{-5}$. Its two field definitions agree with
the identical ideal pair on a neighborhood, so every radial jet agrees. The
first repair sets all five moments equal there as functions of $\eta$.
For every $X\ge X_h$, integrating (PA.14) from $X_h$ gives zero, since $u=e=0$
there. Hence every moment and every $\eta$ derivative agrees for all such $X$.
The common $\Pi_0$ gives identical $\Pi$. Formula (PA.1) then gives identical
$V_0,W,Q_s,N_s,p_s$; equality of field jets gives identical $a,b_s,F,T_0$.
This proves the exact map from the five matching values to the entire outer
field and stress. The identical argument applies after the second repair.

In particular the actual regular joined pair has
\[
 M(\infty)=J(\infty)=S(\infty)=0,\qquad
 \int_0^\infty(H-H_{\rm pow})\,dX=0,\quad
 H_{\rm pow}=\sqrt{2X}c_\infty X^{-A},                           \tag{PA.19}
\]
because its difference from the prepared outer pair is supported before $X_h$,
and its $I$ difference at $X_h$ is zero. The difference $H-H_{\rm pow}$ is
integrable at zero since $h<1$, and at infinity since the heat difference is
$O(X^{-1-h})$. Its pressure is exactly
\[
 \Pi=\Pi_0+C_p=C_p-C_p(\infty)
     =-\int_X^\infty E(x,\eta)^2/(2x)\,dx.                    \tag{PA.20}
\]

\section{Applying the outer endpoint calculation to the actual join}
The joined field is Cartesian smooth at the axis by the axis component, agrees
pointwise with the heat-corrected outer profile for $X\ge X_h<X_{\rm tail}$,
and has (PA.19)--(PA.20). We now prove precisely the backward-stress input formerly
listed separately from the axis construction.

Put $K_{\rm pow}=c_\infty(r^2/2)^{-A}$. The exact substitution
$r=\sqrt{2qX}$, with $r\,dr=q\,dX$, gives
\begin{equation}\label{pa:eq:submoment}\tag{PA.SUB}
 \int_0^\infty r^2(u_\theta-K_{\rm pow})\,dr
  =q^{3/2-A}\int_0^\infty(H-H_{\rm pow})\,dX=0.
\end{equation}
The heat expansion gives $u_\theta-K_{\rm pow}=O_h(\tau r^{-3-2h})$,
$\partial_tu_\theta=O_h(r^{-3-2h})$, and
$r^2\partial_ru_\theta-ru_\theta=O_h(r^{-2h})$ in the exterior.
Their weighted tails have integrable majorant $Cr^{-1-2h}$ for fixed $h>0$.
At zero the subtracted weight is $r^{1-2h}$ and the field is regular.
Dominated differentiation therefore gives zero integrated angular time term.
The axial angular-momentum flux has integral $q^{3/2-2A}J(\infty)=0$.
The radial transport boundary vanishes at zero by regularity and at infinity
because $U=M=V_0=0$. Its viscosity boundary vanishes by the displayed heat
estimate. Thus $\int_0^\infty r^2\paResidual_\theta^{(0)}\,dr=0$.

For the axial component $\int ru_zdr=q^{1-A}M(\infty)=0$. The canonical
pressure obeys, by absolute convergence and Fubini,
\[
 \int_0^\infty rp(r)\,dr
 =-\frac12\int_0^\infty ru_\theta(r)^2\,dr.
\]
The integrated axial momentum flux is consequently
$q^{1-2A}S(\infty)=0$. Its transport and viscosity radial boundaries vanish,
and the pressure boundary $r^2p\to0$ follows from $p=O(r^{-2-4h})$.
Hence $\int_0^\infty r\paResidual_z^{(0)}\,dr=0$. Here
$\paResidual_j^{(0)}=\mathcal R(u,p)\cdot e_j+\partial_{zz}u_j$, with viscosity
one, exactly as in the core calculation. The leading residuals exclude only the
axial viscosity term in this identity. Their regular forward primitives are
therefore exactly
\[
 T_\theta=r^{-2}\int_r^\infty r'^2\paResidual_\theta^{(0)}(r')\,dr',
 \quad T_z=r^{-1}\int_r^\infty r'\paResidual_z^{(0)}(r')\,dr',
 \quad T=q^{-A-1/2}T_0.                                        \tag{PA.21}
\]
All original hypotheses of the outer endpoint calculation have now been
proved for this specific joined field.

Section~\ref{pa:sec:component-edge} derives (PA.21) directly once more, evaluates
the heat-tail residual including its pressure term and both integration signs,
and proves
\[
 b_s=0,\quad 2+h<a\le2+2h,\quad T_\theta>0,\quad
 2-(a-2)(T_z/T_\theta)^2\ge\kappa_o>0                           \tag{PA.22}
\]
on $e^{1/2}X_{\rm tail}\le X<X_b=e^3X_{\rm tail}$. Its exact endpoint
factors, with $y_b=\log(X_b/X)$, are
\[
 T_{0,\theta}=e^{-4/y_b^2}y_b^{-3}b_\theta(y_b,\eta),\quad
 T_{0,z}=e^{-4/y_b^2}y_b^3b_z(y_b,\eta),                         \tag{PA.23}
\]
with smooth factors and
\[
 b_\theta(0,\eta)=\frac{16\rho_o E_{\rm pow}(X_b)}{\sqrt{2X_b}}
                   \paHeat(2d/X_b)e>0.
\]
This coefficient keeps the cutoff, heat, radial, and factor $16$ constants.
The stress is zero for $X\ge X_b$ by the exact radial heat equation and (PA.21).
The inner continuation component proves, for $y_a=\log(X/X_a)$,
\[
 T_0=e^{-t_1^2/y_a^2}g_a(y_a)B_a(y_a,\eta),\quad
 g_a(0)=(1-\kappa_0)e>0,\quad
 B_a(0,\eta)=F(X_a,\eta)(a,-b_s)(X_a,\eta)\ne0.                 \tag{PA.24}
\]
Its stress is zero through $X_a=4/\Lambda$, and its exit has $v_s>2+.2$.
Thus the actual joined profile supplies the two admissible boundary collars
needed for the next construction.

\section{A single finite loop insertion and restoration}\label{pa:sec:assembly-loop}
Choose $X_-$ strictly inside the first admissible collar, and
$X_{\rm good}<X_+<\inf I_1$. The rectangle $I=[X_-,X_+]\times[-1,1]$
contains all potentially nonadmissible points and has admissible radial
boundary collars. Shrink $X_{\rm an}$ so $X_a<X_{\rm an}<X_-$.
Choose $[Y_0,Y_1]\subset I_1$ with $X_+<Y_0<Y_1$.
The exact relaxed data $a>0$, $P_c>2$, $v_s<U(P_c,J_c)$ follow from
the preceding inner/outer construction on all of $I$. The required compact
strict inequalities are actual positive continuous minima. The admissible
condition holds throughout $[X_+,Y_1]$.

The component proof in Section~\ref{pa:sec:component-loop} constructs the smooth
variance loop from these very data. It proves its $p_{s,2}=0$ extension, strict
variance monotonicity, finite uniform variance range, smooth inverse at zero
variance, and exact mean after the positive circle reparameterization.
The output satisfies, in phase $\varphi\in\paReal/\mathbb Z$,
\[
 \int_0^1(a_L,-b_L)\,d\varphi=(a,-b_s),
\]
with four strictly positive cone gaps on the compact loop family, and equals
the original shear on the two boundary collars. The positive compact derivative
of the phase map gives a smooth inverse on the same parameter rectangle.

The zero-mean periodic primitives $\mathscr A,\mathscr B$ have
\[
 \partial_\varphi\mathscr A=-(a_L-a)/2,\qquad
 \partial_\varphi\mathscr B=E(b_L-b_s)/2.
\]
They are zero on those collars. Set
\[
 E_N=E\exp[\mathscr A(X,\eta,N\log X)/N],\qquad
 U_N=U+\mathscr B(X,\eta,N\log X)/N.
\]
The exact chain rule, with $D_X$ holding phase fixed, is
\[
 a_N=a_L-2D_X\mathscr A/N,\qquad
 b_N=e^{-\mathscr A/N}\{b_L+2D_X\mathscr B/(NE)\}.               \tag{PA.25}
\]
For each fixed number of $\eta$ derivatives, field differences, shear errors
relative to the loop, and all prefix-moment differences are $O(N^{-1})$ on
the fixed interval $[X_-,Y_1]$. This follows directly from (PA.14) and the identity
$e^v-1=v\int_0^1e^{sv}ds$. The phase has no $\eta$ dependence. Formula (PA.1)
then bounds $p_{s,N}-p_s$ by $C/N$, using one extra $\eta$ derivative and
no radial derivative of the field differences. Actual high radial derivatives
can grow as $N^{r-1}$; $N$ is fixed before they are used.

Here is one simultaneous finite bound for $N$. Let $\mu_L$ be the minimum
of the four cone gaps on the loop and $\mu_R$ their minimum on the original
compact segment $[X_+,Y_1]$. Choose fixed neighborhoods of these states in
$a>0$, and let $L_\Psi$ be a common Lipschitz constant for
\[
 \Psi(a,b,p)=\left(a,v-2,c-v,2(c-v)^2-(v-2)j^2\right),
 \quad v=a+b^2/a,\quad c=p_1-bp_2/a,\quad j=p_2+bp_1/a.
\]
Let $B_{\rm mod}/N$ bound the modulation errors in its three inputs, and
let $B_{\rm rep}/N$ bound the repair's changes of those inputs, including all
its prefix moments. Let $D_j/N$ bound the defect for the repair (PA.17),
$j=0,2$, and let $B_E/N$ bound its change in $E$. These constants are fixed
by the loop, bump shapes, and all the parameters chosen above. Take an integer
$N$ satisfying
\[
\begin{gathered}
 N>1,\qquad
 N>2L_\Psi B_{\rm mod}/\min(\mu_L,\mu_R),\\
 N>4L_\Psi B_{\rm rep}/\mu_R,\qquad
 N>8\max_{j=0,2}\nu_j^2q_jD_j,\\
 N>2B_E/\min_{[Y_0,Y_1]\times[-1,1]}E .
\end{gathered}                                                     \tag{PA.26}
\]
Increase it finitely, when necessary, to put the input changes inside the chosen
neighborhoods. The first bound preserves half the loop and intervening margins.
The $C^2$ contraction (PA.18) gives the exact repair with
$\|c\|_{C^2}\le2\nu_2D_2/N$; its fixed bump shapes give the bound defining
$B_{\rm rep}$. The next bound preserves at least $\mu_R/4$ on the repair.
The last preserves $E>0$. Equation (PA.26) imposes only finitely many inequalities.
Every other fixed parameter derivative is finite by implicit differentiation of
the same coefficient equation, without reselecting $N$.

At $Y_1$ all five moments are restored. The exterior map proved after (PA.18)
therefore proves exact equality of $\Pi,V_0,Q_s,N_s,p_s,a,b_s,T_0$ beyond
$Y_1$. The inner analytic rectangle is unchanged because both perturbations
start later and use the same axis pressure. The heat compensation in $I_2$,
and the unused $I_3,I_4$, consequently remain exactly their original profiles.

\section{Closed leading-profile conclusion}
\begin{quote}\textbf{Leading-profile assembly theorem.}
For the particular $h$ selected in (PA.3) with the finite smallness choices in
Section~\ref{pa:sec:assembly-parameters}, the preceding construction produces the profiles and every property
of [NS, Theorem 4.6]. This chosen $h$ lies in the original range $0<h<1/100$.
The original coordinates and component signs are retained. All leading-profile constants are independent of physical
$q$, later dyadic bands, and later correction stages.
\end{quote}
\begin{proof}
Cartesian smoothness, positivity of $\phi$ and $E$ for $X>0$, and the common
complex parameter neighborhood for all fixed radial derivatives through
$[0,X_{\rm an}]$ follow from the explicit analytic axis and continuation
construction, unchanged by both later supports. The average identity
$M/X=\int_0^1U(vX,\eta)dv$ supplies the smooth extension of $V_0/X$.
Equation (PA.20) proves the requested canonical pressure. Differentiating it gives
$\Pi_X=E^2/(2X)=F^2$, including at $X=0$. The exact tangential divergence
identities in the core component, retaining radial viscosity, hold for every
constructed profile; through $X_a$ its two stress-free equations were solved
by contraction. Equations (PA.21)--(PA.24) and the final restoration prove zero
stress outside $[X_a,X_b]$. On its interior the admissible cone is now strict.
Since $T_0=0$ would imply $P_c=v_s$, it is nonzero there.

The unit direction is smooth on the open annulus. At the inner edge (PA.24) gives
$n=B_a/|B_a|$, with direction parallel to $(a,-b_s)$; at the outer edge (PA.23)
gives
\[
 n=\frac{(b_\theta,y_b^6b_z)}{(b_\theta^2+y_b^{12}b_z^2)^{1/2}},
\]
which is smooth through the edge and equals $(1,0)$ there. Put
$A_n=n_\theta+t_sn_z$ and $B_n=n_z-t_sn_\theta$, where $t_s=-b_s/a$.
In the interior, the exact scalar products in (PA.1) give
\[
 A_n=(P_c-v_s)/|p_s-(a,-b_s)|>0,\quad
 B_n=J_c/|p_s-(a,-b_s)|,
\]
so $G_n=2-(v_s-2)B_n^2/A_n^2>0$. At both edges $B_n=0$, $A_n>0$,
and $G_n=2$. Compactness therefore gives positive minima of $A_n,G_n$.
The choice $\kappa=\min\{1,\min A_n,\min G_n\}$ is in $(0,2)$ and
gives exactly (4.26). Positivity of $F,a,v_s-2$ has positive compact minima,
using the inner $v_s>2+.2$ bound, outer $a>2+h$ bound, and the finite-loop
gaps. Thus the lower bounds extend over the closed annulus.

Define, with zero extension outside the open annulus,
\[
 \zeta(X)=\exp(-t_1^2/y_a^2-4/y_b^2),\quad
 \delta=\min\{1,y_a,y_b\}.
\]
Differentiating either exponential introduces only finitely many inverse
powers of its logarithmic distance. The exponential dominates each such
power, so the extension is smooth with every radial derivative zero at the
edges. On a fixed small inner collar (PA.24) gives
$|T_0|\ge c\zeta$, $|\partial^\alpha T_0|\le C_\alpha\zeta\delta^{-m_\alpha}$.
The same holds on a fixed outer collar by (PA.23); its $y_b^{-3}$ factor is
bounded below there, and all its derivatives cost finite inverse powers.
On the intervening compact interior, $|T_0|$ and $\zeta$ have positive
minima and all fixed derivatives are bounded. Taking the minimum of the
three lower constants gives one $c>0$ independent of $\alpha$, and increasing
the upper constants proves (4.27) on the whole open annulus.

The second repair restores (PA.19)--(PA.20). In the unchanged exterior the field is
\[
 U=V_0=0,\quad E=c_\infty X^{-A}\paHeat(2d/X),\quad
 \paHeat(Z)=\Gamma(1+h)^{-1}\int_0^\infty e^{-v}v^h(1+Zv)^{-h}dv.
\]
The core and outer calculations prove this positive heat factor, all its
one-sided derivatives at zero, and the exact radial heat equation for
$q^{-A}E=c_\infty s^{-A}\paHeat(2\tau/s)$. Above the original $X_v$,
$U=M=J=0$, so (PA.1) gives $V_0=0$ on the entire outer collar as well.
The exact exterior map preserves $I_{\rm pos}=I_3$ and $I_{\rm mean}=I_4$,
with $\sup I_3<\inf I_4<X_v$ and
$U=0$, $E=c_{\rm patch}(1+\eta^2)^{-1}X^{-1/2-\lambda}$ there.
All radii and widths, the circle loop, its bumps, and the single integer $N$
were selected in (PA.3)--(PA.26) before physical $q$ or any later derivative stage.
Every resulting fixed derivative is a finite derivative of these fixed smooth
functions. This proves the final independence assertion and all six parts.
\end{proof}

\medskip
% Appendix numbering is intentionally not reset in the cumulative integration body.
\noindent The remaining sections supply the complete independently written component
calculations used above. They are embedded in this file; compilation and reading require
no external TeX inputs. The proof order within the main construction distinguishes the
outer schedule from its later use in the backward-stress calculation. References to [NS]
below identify equations in the fixed source, not additional unproved inputs.


% Source body: core_exact.md
\section{Exact leading-field maps and radial primitives}\label{pa:sec:component-core}

\subsection{1. Coordinates and their exact Jacobian}

Take right-handed cylindrical coordinates
\[
x=(r\cos\theta,r\sin\theta,z),\quad e_\theta=(-\sin\theta,\cos\theta,0),
\quad \theta\in\mathbb R/(2\pi\mathbb Z),\quad s=r^2/2.
\]
Put
\[
\begin{aligned}
&\tau=1-t>0,\quad A=\tfrac12+h,\quad D=\tfrac12-h,\quad 0<h<\tfrac12,\\
&z=q^D\eta,\quad \tau=q(1-\eta^2),\quad X=s/q,\\
&d=1-\eta^2,\quad L=1-2h\eta^2.
\end{aligned}
\]
For fixed \(\tau>0,z\in\mathbb R\), eliminating \(\eta\) gives
\(g(q)=q-z^2q^{2h}=\tau\). At \(q_*=|z|^{1/D}\), \(g(q_*)=0\); for \(z=0\) take \(q_*=0\). For \(q>q_*\), \(z^2q^{2h-1}<1\), and
\[
g'(q)=1-2hz^2q^{2h-1}\ge1-2h>0.
\]
Moreover \(g(q)\to\infty\). Thus exactly one \(q>q_*\) exists, \(|\eta|<1\), and the implicit function theorem makes \(q,\eta\) smooth for \(t<1\). The values \(\eta=\pm1\) used for profiles are one-sided parameter endpoints; they are not additional points with \(\tau>0\).

Differentiation at fixed physical coordinates gives
\[
q_t=-L^{-1},\quad \eta_t=\frac{D\eta}{qL},\quad X_t=\frac X{qL},
\qquad q_z=\frac{2\eta q^{1-D}}L,
\quad \eta_z=\frac d{q^DL},\quad X_z=-\frac{2\eta X}{q^DL}.
\]
The formula for \(\eta_z\) uses the exact equality \(L-2D\eta^2=d\). Consequently, for every real \(b\) and smooth profile \(f\),
\[
\partial_t(q^bf)=q^{b-1}L^{-1}(-bf+D\eta f_\eta+Xf_X),
\]
\[
\partial_z(q^bf)=q^{b-D}L^{-1}(2b\eta f+df_\eta-2\eta Xf_X).
\]
No constant or power of \(q\) is discarded in these maps.

\subsection{2. Cartesian axis extension and incompressibility}

Let \(C>1\), \(F=\phi/C\), \(E=\sqrt{2X}F\), and prescribe
\[
u_\theta=q^{-A}E,\quad u_z=q^{-A}U,\quad ru_r=V_0=Xv_0,
\quad p=q^{-2A}\Pi.
\]
For \(F,U,v_0,\Pi\) smooth on \([0,R]\times[-1,1]\), the Cartesian components are exactly
\[
u_1=\frac{v_0}{2q}x_1-q^{-A-1/2}Fx_2,
\quad u_2=\frac{v_0}{2q}x_2+q^{-A-1/2}Fx_1,
\quad u_3=q^{-A}U.
\]
Since \(X=(x_1^2+x_2^2)/(2q)\), these formulas prove smoothness across \(r=0\) for each \(t<1\). Smoothness is imposed on \(F\), not on \(E\) as a function of \(X\). This explicitly proves the map from scalar profile regularity to Cartesian velocity regularity.

Incompressibility is
\(\partial_s(ru_r)+\partial_zu_z=0\). Because \(A+D=1\), its exact scalar equation is
\[
(V_0)_X=L^{-1}(2A\eta U-dU_\eta+2\eta XU_X).
\]
Integrate with \(V_0(0,\eta)=0\). Writing \(M(X,\eta)=\int_0^XU(x,\eta)\,dx\), the result is
\[
V_0=\frac{2\eta XU-2D\eta M-dM_\eta}{L}.
\]
The extension of \(M/X\) at zero is smooth because
\(M/X=\int_0^1U(vX,\eta)\,dv\); differentiating this integral proves every mixed derivative assertion. Thus the formula also proves smoothness of \(V_0/X\).

The radial derivative of pressure and centrifugal term have the common factor \(q^{-2A-1/2}\), with respective coefficients \(\sqrt{2X}\Pi_X\) and \(E^2/\sqrt{2X}\). Their exact equality is
\[
\Pi_X=\frac{E^2}{2X}=F^2.
\]
The other radial terms are still present in the full Navier-Stokes residual. Radial time differentiation, both material-transport contributions and radial viscosity have scale \(q^{-3/2}=q^{2h}q^{-2A-1/2}\). Axial viscosity of the radial velocity has scale \(q^{-1/2-2D}=q^{-3/2+2h}=q^{4h}q^{-2A-1/2}\). Axial viscosity has scale \(q^{-A-2D}=q^{2h}q^{-A-1}\) in the tangential equations. These equalities identify every remaining type of term without setting any of them to zero.

\subsection{3. Exact transport and stress integration}

All residual identities here have viscosity one, as defined on source p. 6. Define
\[
H=\sqrt{2X}E=2XF,\quad l=X\partial_X\log H,
\quad W=1-\frac{2D\eta M+dM_\eta}{X},\quad H_c=D\eta+dU.
\]
Applying the preceding derivative maps and \(u_r=V_0/r\) gives, for arbitrary \(b,f\),
\[
(\partial_t+u_r\partial_r+u_z\partial_z)(q^bf)
=\frac{q^{b-1}}L\{W Xf_X+H_cf_\eta-b(1-2\eta U)f\}.
\]
Indeed the coefficient of \(Xf_X\) before substituting \(V_0\) is \(1+LV_0/X-2\eta U\), exactly \(W\). Angular momentum is \(ru_\theta=q^{-h}H\). Hence the angular material derivative equals \(-q^{-h-1}HS_q/L\), where
\[
S_q=-Wl-h(1-2\eta U)-H_c(\log E)_\eta.
\]
For the axial velocity, add \(\partial_zp\). Since \(-2A-D=-A-1\), the result is \(-q^{-A-1}S_n/L\), with
\[
S_n=-WXU_X-A(1-2\eta U)U-H_cU_\eta
-d\Pi_\eta+4A\eta\Pi+2\eta X\Pi_X.
\]

Let \(D_X=X\partial_X\). Solve
\[
D_XQ_s+(1+l)Q_s=S_q,\qquad D_XN_s+N_s=S_n.
\]
Their integrating factors are \(XH\) and \(X\). For \(\phi>0\) smooth at the axis, \(H=2X\phi/C\), and any nonzero integration constant produces an \(X^{-2}\) singularity in \(Q_s\) or an \(X^{-1}\) singularity in \(N_s\). Therefore the unique regular choices are
\[
Q_s=\frac{\int_0^XHS_q\,dx}{XH},\qquad
N_s=\frac1X\int_0^XS_n\,dx.
\]
In particular,
\[
Q_s=\frac1{F(X,\eta)}\int_0^1vF(vX,\eta)S_q(vX,\eta)\,dv,
\quad N_s=\int_0^1S_n(vX,\eta)\,dv,
\]
which prove the scalar smooth extensions and give \(Q_s(0)=S_q(0)/2\), \(N_s(0)=S_n(0)\).

Define
\[
\begin{aligned}
&p_s=\left(\frac{XQ_s}L,\frac{XN_s}{LE}\right),\quad a=1-2D_X\log E=2-2l,\\
&b_s=\frac{2D_XU}{E},\quad s_{\rm shear}=(a,-b_s),\quad T_0=F(p_s-s_{\rm shear}).
\end{aligned}
\]
Here \(s_{\rm shear}\) is only a typographical disambiguation from the unchanged coordinate \(s=r^2/2\); it is exactly the manuscript's vector \(s\). The stress components remain ordered \((r\theta,rz)\).

For any profile \(G\), at fixed \(t,z\),
\[
(\partial_r+k/r)(q^{-A-1/2}G)
=q^{-A-1}\left(\sqrt{2X}\,G_X+\frac k{\sqrt{2X}}G\right).
\]
Substituting \(G=FXQ_s/L=HQ_s/(2L)\), \(k=2\), produces \(q^{-A-1}ES_q/L\). Substituting \(G=FXN_s/(LE)=\sqrt{X/2}N_s/L\), \(k=1\), produces \(q^{-A-1}S_n/L\). The remaining stress terms are exactly
\[
\partial_ru_\theta-u_\theta/r=-q^{-A-1/2}Fa,
\qquad \partial_ru_z=q^{-A-1/2}Fb_s.
\]
The differential identities
\[
(\partial_r+2/r)(\partial_ru_\theta-u_\theta/r)
=\partial_{rr}u_\theta+r^{-1}\partial_ru_\theta-r^{-2}u_\theta,
\]
\[
(\partial_r+1/r)\partial_ru_z=\partial_{rr}u_z+r^{-1}\partial_ru_z
\]
now prove, including both signs,
\[
R^{(0)}_\theta=-(\partial_r+2/r)(q^{-A-1/2}T_{0,\theta}),
\quad R^{(0)}_z=-(\partial_r+1/r)(q^{-A-1/2}T_{0,z}),
\]
where \(R^{(0)}_j=\mathcal R(u,p)\cdot e_j+\partial_{zz}u_j\). Thus only axial viscosity has been removed from these exact identities.

Zero stress is equivalent to
\(Q_s=-2L\phi_X/\phi\) and \(N_s=-2LU_X\). Applying the two radial operators and canceling the two opposite quadratic logarithmic-derivative terms gives exactly
\[
-2L\frac{X\phi_{XX}+2\phi_X}{\phi}=S_q,
\qquad -2L(XU_{XX}+U_X)=S_n.
\]
Conversely, these equations and regularity imply the same \(Q_s,N_s\), by uniqueness of the regular radial primitives, hence zero stress. The first displayed fraction was checked directly against PDF p. 27.

\subsection{4. The five moments and their exact matching map}

Retain all five quantities
\[
\begin{aligned}
&M=\int_0^XU\,dx,\quad I=\int_0^XH\,dx,\quad J=\int_0^XUH\,dx,\\
&S=\int_0^X(U^2-E^2/2)\,dx,\quad C_p=\int_0^X\frac{E^2}{2x}\,dx,\quad \Pi=\Pi_{\rm ax}+C_p.
\end{aligned}
\]
Since \((XW)_X=1-2D\eta U-dU_\eta\), integration by parts in \(-WXH_X\) gives \(-XWH\) and the remaining integrand
\((1-h)H-D\eta H_\eta-d(UH)_\eta+2(h-D)\eta UH\). Hence
\[
Q_s=-W+\frac{(1-h)I-D\eta I_\eta-dJ_\eta+2(h-D)\eta J}{XH}.
\]
For the axial primitive, integrate \(-WXU_X\) in the same way. Its velocity terms become
\(D(U-\eta U_\eta)+4h\eta U^2-d(U^2)_\eta\). Pressure contributes
\[
X(4A\eta\Pi-d\Pi_\eta)-2h\eta\int_0^XE^2\,dx
+\frac d2\partial_\eta\int_0^XE^2\,dx,
\]
because \(\int_0^X\Pi\,dx=X\Pi-\frac12\int_0^XE^2\,dx\). Thus
\[
N_s=-WU+\frac{D(M-\eta M_\eta)+4h\eta S-dS_\eta}{X}
+4A\eta\Pi-d\Pi_\eta.
\]
The independent exact check differentiates both right-hand primitives, retains \(\Pi_X=H^2/(4X^2)\), and obtains precisely \(HS_q\) and \(S_n\).

If two pairs \((U_i,E_i)\) agree for \(X\ge X_h>0\) and their five moment values at \(X_h\) agree as functions of \(\eta\), all five moment differences are identically zero beyond \(X_h\): their radial derivatives are zero and their initial values are zero. All parameter derivatives also agree. With the same \(\Pi_{\rm ax}\), the explicit formulas then give identical \(\Pi,V_0,Q_s,N_s,p_s\). Agreement of profiles on an interval gives agreement of radial derivatives, hence of \(a,b_s,T_0\). This proves the complete morphism from matching data to the exterior field and stress; pressure matching alone would not suffice.

On a compact \(X\)-rectangle bounded away from zero and with \(E_i\ge e_{\min}>0\), the displayed formulas are compositions of addition, multiplication and reciprocal on denominators \(X,L,E,H\) bounded away from zero. For example
\(H_2^{-1}-H_1^{-1}=-(H_2-H_1)/(H_1H_2)\). The product rule through \(k\) parameter derivatives bounds the output difference in \((Q_s,N_s,p_s)\) by the input differences in \((U,E,M,I,J,S,C_p)\) through order \(k+1\). This uses no radial derivative of the difference. The one extra parameter derivative is necessary because \(W\) contains \(M_\eta\) and the primitive formulas contain moment derivatives.

Under the exact coordinate substitution \(X=\rho x\), use
\[
\widehat U(x,\eta)=U(\rho x,\eta),\quad\widehat E(x,\eta)=E(\rho x,\eta),
\quad H=\rho^{1/2}\widehat H,
\]
\[
(M,I,J,S,C_p)=(\rho\widehat M,\rho^{3/2}\widehat I,
\rho^{3/2}\widehat J,\rho\widehat S,\widehat C_p),
\quad\widehat Q_s=Q_s,\quad\widehat N_s=N_s,\quad\widehat p_s=\rho^{-1}p_s.
\]
Direct substitution proves that precisely the same rational formulas hold on the fixed \(x\)-rectangle. This is an invertible change of variables retaining all factors, not a replacement of the original moment data.

\subsection{5. Cone equivalence and finite-dimensional solve}

For \(a>0\), define
\[
t_s=-b_s/a,\quad v_s=a(1+t_s^2),\quad P_c=p_{s,1}+t_sp_{s,2},\quad
J_c=p_{s,2}-t_sp_{s,1}.
\]
The roots in \(v\) of \(2(P_c-v)^2-(v-2)J_c^2\) are
\[
v_\pm=P_c+J_c^2/4\pm |J_c|\sqrt{(P_c-2)/2+J_c^2/16}.
\]
For \(P_c>2\), evaluate the polynomial at \(2\) and \(P_c\): it is positive at \(2\), nonpositive at \(P_c\), and has its vertex at or above \(P_c\). Consequently \(2<v_-\le P_c\). This includes \(J_c=0\), where the two roots are \(P_c\). Therefore for \(v_s>2\), the admissible inequality \(v_s<v_-\) is exactly
\[
P_c>v_s,\qquad (v_s-2)J_c^2<2(P_c-v_s)^2.
\]
In particular \(v_s\le2\), \(P_c>2\) always implies the strict relaxed inequality \(v_s<v_-\).

For \(p_{s,2}=wp_{s,1}\), put \(c=1-b_sw/a\), \(j=w+b_s/a\), \(v=a+b_s^2/a\). Direct expansion gives
\[
(v-2)j^2-2c^2=(1+b_s^2/a^2)((a-2)w^2+2b_sw+b_s^2/a-2).
\]
On the compact set specified in Lemma 4.5, choose actual minima \(c_K>0\), \(\gamma_K=\min(2c^2-(v-2)j^2)>0\), and \(B_K\ge\max\{1,v,cv\}\). For
\[
p_{s,1}\ge P_K=\max\{(B_K+2)/c_K,8B_K/\gamma_K\},
\]
one has \(P_c>\max\{2,v\}\) and
\[
\frac{2(P_c-v)^2-(v-2)J_c^2}{p_{s,1}^2}
\ge\gamma_K-4B_K/p_{s,1}\ge\gamma_K/2.
\]
These are exactly the source constants, with the endpoint \(v=2\) covered by the preceding root calculation.

Stress coordinates satisfy, by direct scalar products,
\[
T_{0,\theta}+t_sT_{0,z}=F(P_c-v_s),\quad
T_{0,z}-t_sT_{0,\theta}=FJ_c.
\]
Thus a positive margin for the two stress inequalities proves nonvanishing and controls direction independently of magnitude.

For completeness, the moment matrix invertibility used later is proved as follows. Distinct real exponents \(\alpha_i\) form a Chebyshev family on \((0,\infty)\). For \(m\) exponents, divide a nonzero combination by its smallest power and differentiate. If the original combination had \(m\) distinct positive zeros, Rolle's theorem would give \(m-1\) zeros of a combination of \(m-1\) distinct powers, contradicting the induction starting with one nonzero power. Therefore the determinant \(\det[x_j^{\alpha_i}]\) never vanishes for ordered \(x_1<\cdots<x_m\), and its sign is constant on that connected set. Multilinearity expresses the bump moment determinant as the integral of this determinant against \(\prod_j\beta_j(x_j)\ge0\). Ordered disjoint nontrivial supports make the integral nonzero. Compact smooth parameter families then have bounded inverse matrices and bounded fixed derivatives by the adjugate formula.

For the exact quadratic moment equation \(Bc+Q(c,c)=d\), multiplication by \(B^{-1}\) has norm \(\nu_k\), and \(Q\) has bilinear norm \(q_k\) on \(C^k\). On the ball of radius \(r_k=2\nu_k\|d\|_{C^k}\), the map \(c\mapsto B^{-1}(d-Q(c,c))\) has image norm at most \(r_k/2+\nu_kq_kr_k^2\) and Lipschitz constant at most \(2\nu_kq_kr_k\). The source condition \(8\nu_k^2q_k\|d\|_{C^k}\le1\) gives invariance and Lipschitz constant at most \(1/2\). Iteration from zero in \(C^0\) and \(C^k\) gives the same fixed point. Its pointwise derivative matrix is \(B+Q(c,\cdot)+Q(\cdot,c)\); the corresponding \(B^{-1}\)-perturbation has norm at most \(1/2\), proving invertibility. The finite-dimensional implicit function theorem therefore gives all higher fixed parameter derivatives without imposing infinitely many smallness conditions.

\subsection{6. Direct heat-exterior verification}

The function named \(\mathcal H\) here is exactly the source's \(H(Z)\) in (4.29), distinguished typographically from angular momentum \(H(X,\eta)\):
\[
\mathcal H(Z)=\frac1{\Gamma(1+h)}\int_0^\infty e^{-v}v^h(1+Zv)^{-h}\,dv,
\qquad Z\ge0.
\]
For each \(k\), its \(k\)-th derivative is the integral with factor \((-1)^k(h)_kv^k(1+Zv)^{-h-k}\). The integrable majorant is \((h)_ke^{-v}v^{h+k}\), proving smooth one-sided derivatives at \(Z=0\) and positivity. Differentiating
\(e^{-v}v^{h+1}(1+Zv)^{-h-1}\) with respect to \(v\) shows
\[
Z^2\mathcal H''+((2+2h)Z+1)\mathcal H'+h(1+h)\mathcal H=0.
\]
In fact the numerator of the ODE integrand after taking the common denominator \((1+Zv)^{h+2}\) is \(h(h+1-v-Zv^2)\), exactly \(h\) times that derivative numerator; both boundary values vanish. This is a full integration-by-parts proof.

For \(K=c_\infty s^{-A}\mathcal H(2\tau/s)\), let \(Z=2\tau/s\). Then
\[
\partial_tK=-2c_\infty s^{-A-1}\mathcal H',\quad
\partial_sK=c_\infty s^{-A-1}(-A\mathcal H-Z\mathcal H'),
\]
\[
(\partial_{rr}+r^{-1}\partial_r-r^{-2})K
=c_\infty s^{-A-1}\{(2A^2-\tfrac12)\mathcal H
+(4A+2)Z\mathcal H'+2Z^2\mathcal H''\}.
\]
Using \(2A^2-1/2=2h(1+h)\) and the preceding ODE proves the exact radial heat equation. Also
\[
A\mathcal H+Z\mathcal H'
=\frac1{\Gamma(1+h)}\int_0^\infty e^{-v}v^h
\frac{A+Zv/2}{(1+Zv)^{h+1}}\,dv>0,
\]
so \(K_r=rK_s<0\) for \(r>0\). Finally
\(q^{-A}c_\infty X^{-A}\mathcal H(2d/X)=c_\infty s^{-A}\mathcal H(2\tau/s)\) exactly; the exterior has no residual hidden by a coordinate substitution.



\section{Power moment maps and exact corrections}
\subsection{An invertible moment map and its quadratic inverse}
For distinct real exponents $\alpha_1,\ldots,\alpha_m$, a nonzero sum
$\sum c_ix^{\alpha_i}$ has at most $m-1$ positive zeros. To prove this by
induction, order the exponents and divide by the smallest power. If the result
had $m$ zeros, its derivative would have $m-1$ zeros by Rolle's theorem. That
derivative is either zero, in which case the original result is constant, or
is a nonzero combination of $m-1$ distinct powers, which has at most $m-2$
zeros by induction. Both alternatives exclude $m$ zeros.

For ordered disjoint intervals $I_1<\cdots<I_m$ and nonnegative nonzero
$\beta_j\in C_c^\infty(I_j)$, the determinant of $(x_j^{\alpha_i})$ is
nonzero at every $x_1<\cdots<x_m$, and its continuous sign is fixed on that
connected set. Multilinearity and Fubini give
\[
\det\left(\int x^{\alpha_i}\beta_j(x)\,dx\right)
=\int_{I_1\times\cdots\times I_m}\det(x_j^{\alpha_i})
 \prod_j\beta_j(x_j)\,dx_1\cdots dx_m\ne0.
\]
Smooth families on compact parameter sets have bounded inverse derivatives,
by the determinant formula and $\partial B^{-1}=-B^{-1}(\partial B)B^{-1}$.
Under $x=e^y$, this is also the exponential-weight statement in [NS, Lemma A.1].
For two identical translated logarithmic bumps at centers separated by
$\Delta>0$, after dividing row $i$ by $e^{s_ic_1}$ the determinant is
$b(s_1)b(s_2)(e^{s_2\Delta}-e^{s_1\Delta})$, where
$b(s)=\int e^{st}\beta(t)\,dt>0$. On a fixed compact slope interval,
the mean value theorem gives an inverse bound $O(|s_1-s_2|^{-1})$.
In particular the near-coincident weights in this construction cost exactly
an allowed $O(\lambda^{-1})$ upper bound, not a constant independent of $\lambda$.

For completeness the finite nonlinear solve can be made fully quantitative.
Let $F_\eta(c)=B(\eta)c+Q_\eta(c,c)$, with $Q$ bilinear. In a chosen
$C^k([-1,1])$ norm let multiplication by $B^{-1}$ have norm $\beta_k$ and
the bilinear map have norm $\kappa_k$. If
$8\beta_k^2\kappa_k\|d\|_{C^k}\le1$, the map
$\Phi(c)=B^{-1}(d-Q(c,c))$ preserves the ball of radius
$r=2\beta_k\|d\|_{C^k}$, since
\[
\|\Phi(c)\|\le\beta_k\|d\|+\beta_k\kappa_kr^2
\le r/2+r/4<r.
\]
Its Lipschitz constant is at most $2\beta_k\kappa_kr\le1/2$.
Iteration from zero therefore converges to the unique solution in this ball.
At $k=0$, $\|B^{-1}D_cQ(c,c)\|\le1/2$, so the Jacobian is invertible.
Pointwise implicit differentiation proves smoothness. Solving in any fixed
$C^k$ ball gives the same branch because the iterates are identical.
This proves [NS, Lemma A.2] and also specifies the branch selected in every use.

The five original leading moments are
\[
M=\int_0^XU\,dx,\quad I=\int_0^X\sqrt{2x}E\,dx,\quad
J=\int_0^XU\sqrt{2x}E\,dx,
\]
\[
S=\int_0^X(U^2-E^2/2)\,dx,\qquad C_p=\int_0^X E^2/(2x)\,dx,
\quad\Pi=\Pi_{\rm ax}+C_p.
\]
The substitution $X=\rho x$ is an exact map: the factors of these five
moments are $(\rho,\rho^{3/2},\rho^{3/2},\rho,1)$, with
$H=\rho^{1/2}\sqrt{2x}E$. No factor is discarded. At
$U=u_c(\eta)$ and $E=e_*f_*(\eta)x^\alpha$, take the differential of
the rows $J-u_cI$ and $S-2u_cM$, where $u_c$ in the row operation stays at
its original value. These differentials are
$\int\sqrt{2x}E\,\delta U\,dx$ and $-\int E\delta E\,dx$.
Together with $M,I,C_p$, they give the two blocks of weights
\[
(1,x^{\alpha+1/2}),\qquad (x^{1/2},x^\alpha,x^{\alpha-1}),
\]
with the factors $e_*f_*$ and $\sqrt2$ retained in their row multipliers.
The exponents are distinct precisely when
$\alpha\notin\{-1/2,1/2,3/2\}$.
The functionals are degree at most two, so their unlisted remainder is
exactly a quadratic bilinear map. This derives [NS, Corollary A.3].


\section{The complete outer schedule and its exact moment closure}\label{pa:sec:component-outer}
\subsection{Definitions of all stages and the adjustable coefficients}
Use the fixed step
\[
\sigma(y)=\frac{e^{-1/y^2}}{e^{-1/y^2}+e^{-1/(1-y)^2}}\ (0<y<1),
\qquad \sigma=0\ (y\le0),\quad\sigma=1\ (y\ge1).
\]
The schedule chooses, in order, $M_d$ large, $T_d=e^{M_d}+10$,
$P_*>e^{T_d}$, $\lambda>0$ small depending on these choices, and
$h>0$ sufficiently small relative to both $\lambda$ and $e^{-T_d}$.
Only after these choices is $X_R$ made large. All estimates denoted $C_{\rm pre}$
below depend on the first three choices and the fixed steps, not on the later
$X_R$ or amplitude factor $C$. Smallness requirements use finitely many
derivatives, and parameters can be chosen to meet all of them.

Write $x=X/X_R$ and, initially, $y=\log x$. The reference fields for $y\le0$
are $U=4\eta$, $E=P_*f e^{y/10}$, $f=(1+\eta^2)^{-1}$,
$J_0=\log(1+\eta^2)$. In each next stage $y$ is reset to zero at that
stage's left endpoint, and $l=X\partial_X\log(\sqrt{2X}E)$.
The fields are fully specified by the following data, with values at endpoints
matched by integration of $\partial_y\log E=l-1/2$:
\begin{enumerate}
\item On a unit interval, $U=4\eta$, $l=(3/5)(1-\sigma(y))$.
On the following interval $0\le y\le T_d$, set $l=0$,
$U=k(y)\eta$, $k=4[1-\sigma(\log(1+y)/M_d)]$.
Thus $k$ vanishes for its last eleven units, and
$|k'|\le e_d/(1+y)$ with $e_d=4\|\sigma'\|_\infty/M_d$.
The swirl on this interval is $E=P_1fe^{-y/2}$, where $P_1\asymp P_*$.
\item On a unit interval set $U=0$, $l=-\lambda\sigma(y)$, and then
hold $l=-\lambda$ for $T_w=60\log(1/\lambda)$.
Reserve $(T_w-25,T_w-20)$, $(T_w-20,T_w-15)$,
$(T_w-14,T_w-9)$, and $(T_w-8,T_w-3)$ for the four later corrections.
\item At $X=X_p$ start a pulse with $0\le y\le13/\lambda$,
$E=e_bf e^{-(1/2+\lambda)y}$, $U=ER_b$, $\xi_b=\lambda y$.
Define $\varphi_b(\xi)=\int_0^\xi\sigma(v/.02)\,dv$ for $\xi\ge0$,
extended by zero to negative arguments, and
$R_0(\xi)=\varphi_b(\xi)[1-\sigma(\xi-10)]$.
Set $R_b=\mathrm{Amp}(\eta)R_0(\lambda y)+c_1\beta_1(y)+c_2\beta_2(y)$,
where the two identical width $.3$ bumps have centers $13/\lambda-3$ and
$13/\lambda-1$. Choose $c_1,c_2$ to make $M=J=0$ at the pulse end.
\item From the pulse endpoint $E=e_{\rm end}f$, set $U=0$ permanently and
\[
\log E=\log e_{\rm end}-(1/2+\lambda)y
-\vartheta_f(y)J_0-(1-\vartheta_f(y))\log2,
\quad\vartheta_f=1-\sigma(y/T_f).
\]
A fixed sufficiently large $T_f$ ensures $-\lambda-.1\le l\le-\lambda$.
Next retain the now $\eta$ independent exponential for $30\log(1/\lambda)$.
Two relative swirl bumps of width $.3$, three and one units before its end,
set $I=XH/(1-\lambda)$ there and make their total integral
$\int(E^2-E_{\rm unedited}^2)\,dy$ zero.
\item Starting at $X_{\rm rel}$, vary $l$ from $-\lambda$ to $-1$ in a
unit interval, hold $l=-1$ for $4\log(1/h)$, and vary from $-1$ to $-h$
in a unit interval, all using $\sigma$.
The terminal multiplier is
$f_o(y)=1-\rho_o\psi_o(y)$ on $0\le y\le3$,
$\psi_o=1-\sigma((y-1)/2)$, $\rho_o=c_oh>0$ with $c_o$ fixed so that
$0\le f_o'/f_o<h/4$.
Define
\[
Q_p=\int_0^3\exp\left(\int_0^v(1+l(s))\,ds\right)
\frac{f_o'(v)}{f_o(v)}\,dv,\qquad l=-h+f_o'/f_o
\]
on this terminal interval. Hold $l=-h$ before it until $Q_s$ reaches $Q_p$.
Then apply this terminal interval, and retain $l=-h$ thereafter.
On and after the terminal interval the field is exactly $E_{\rm pow}f_o$,
$E_{\rm pow}=c_\infty X^{-A}$, with $c_\infty$ independent of $\eta$.
\end{enumerate}
All stage endpoints divided by $X_R$ are constants independent of $\eta$.
The following calculations prove that the apparently adjustable stages have
the stated finite solutions, and determine the pulse amplitude.

\subsection{Pulse corrections, angular correction, and amplitude root}
During the long interval $T_w$, the quantities $XE$ and $XHE$ grow at rates
$1/2-\lambda$ and $1/2-2\lambda$, respectively, whereas $XE^2$ decays
at rate $2\lambda$. The preceding $M,J$ are constant once $U=0$.
Since $e^{-T_w/2}=\lambda^{30}$ and
$\lambda T_w=60\lambda\log(1/\lambda)\to0$, it follows, with one
$\eta$ derivative as well, that
\[
e_b\le C_{\rm pre}\lambda^{30},\quad
\left\|\frac{M(X_p)}{X_pe_bf}\right\|_{C^1}
+\left\|\frac{J(X_p)}{X_pH(X_p)e_bf}\right\|_{C^1}
\le C_{\rm pre}\lambda^{29},
\]
\[
\left\|\frac{S(X_p)}{X_pe_b^2f^2}\right\|_{C^1}
\le C_{\rm pre}(1+\log(1/\lambda)).
\]
For the last inequality, express $S$ as an integral of $X(U^2-E^2/2)$
in logarithmic coordinates. The finite preceding stages contribute a fixed
constant in this scale and the long $U=0$ stage contributes its length times
a bounded factor $e^{-2\lambda y}$. These statements prove (A.14) with all
scale factors retained.

The exact two equations for the pulse correction are
\[
\begin{pmatrix}
\int e^{s_1y}\beta_1\,dy&\int e^{s_1y}\beta_2\,dy\\
\int e^{s_2y}\beta_1\,dy&\int e^{s_2y}\beta_2\,dy
\end{pmatrix}\binom{c_1}{c_2}
=-\binom{M(X_p)/(X_pe_bf)+\mathrm{Amp}\int e^{s_1y}R_0(\lambda y)\,dy}
{J(X_p)/(X_pH(X_p)e_bf)+\mathrm{Amp}\int e^{s_2y}R_0(\lambda y)\,dy},
\]
where $s_1=1/2-\lambda$, $s_2=1/2-2\lambda$. Divide each row by its
exponential value at the first bump center. The main pulse is zero after
$11/\lambda$, leaving a gap at least $2/\lambda-3$; both slopes exceed $.4$
for small $\lambda$. Its normalized moments are at most a polynomial in
$\lambda^{-1}$ times $(1+|\mathrm{Amp}|)e^{-.4(2/\lambda-3)}$.
The prior moments have an even larger exponential gap. The matrix inverse costs
$O(\lambda^{-1})$. Reducing the positive exponential constant absorbs every
such fixed polynomial. Hence the exact affine solution
$c=c^{(0)}+\mathrm{Amp}\,c^{(1)}$ satisfies
\[
\|c^{(0)}\|_{C^1}+\|c^{(1)}\|_{C^1}\le C_{\rm pre}e^{-c/\lambda}.
\]
The same argument applies to any fixed derivative, with additional powers of
$\lambda^{-1}$ for $\xi_b$ derivatives. This proves exact $M=J=0$ and (A.15).

For angular momentum put $r_I=I/(XH)$. Since $I_X=H$,
$r_I'+(1+l)r_I=1$. During the uniform interval following interpolation,
$r_I-(1-\lambda)^{-1}$ decays at rate $1-\lambda$.
At the first swirl bump center its discrepancy is
$O_{C^1}(C_{\rm pre}\lambda^{28})$; the exponent $28$ is a valid relaxed
bound for $e^{-(1-\lambda)(30\log(1/\lambda)-3)}$.
For the relative edit $E\mapsto E(1+c_1\beta_1+c_2\beta_2)$,
the linear angular row has logarithmic weight $e^{(1-\lambda)y}$ and
the linear pressure row has weight $2e^{(-1-2\lambda)y}$.
Their slope separation tends to $2$, so their normalized inverse is bounded.
The pressure row alone has its exact quadratic remainder. The contraction
proved above yields coefficients $O_{C^1}(C_{\rm pre}\lambda^{28})$,
with both exact targets met. All fixed radial derivatives of these bumps
are of the same size, so they preserve $E>0$ and $l\le-\lambda/2$.

After the correction, $I=XH/(1-\lambda)$ and $U=M=J=0$ at the start
of the exterior transition. The exact radial identities of (4.16) give
\[
Q_s=-1+(1-h)I/(XH)=(\lambda-h)/(1-\lambda).
\]
Subsequently $Q_s'+(1+l)Q_s=-l-h$. During the first transition the
nonnegative source and bounded integrating factor retain a positive $c\lambda$
lower bound. The $l=-1$ hold gives $Q_s'=1-h$, so its end value is comparable
to $\log(1/h)$. The next unit transition changes this by bounded factors.
On the other hand $Q_p\asymp\rho_o\asymp h$: $f_o$ is bounded above and
below, the interval has fixed length, and $\int f_o'=\rho_o$.
The intervening $l=-h$ solution is $Q_{\rm in}e^{-(1-h)y}$; the exact required
length is $(1-h)^{-1}\log(Q_{\rm in}/Q_p)>0$ for sufficiently small $h$.
On the terminal interval the resulting solution is
\[
Q_s(y)=\int_y^3\exp\left(\int_y^v(1+l(s))\,ds\right)
\frac{f_o'(v)}{f_o(v)}\,dv.
\]
Differentiating this expression proves the equation and its zero right endpoint.
Thus $Q_s=0$ afterward and $I=XH_{\rm pow}/(1-h)$. Since
$\int_0^XH_{\rm pow}\,dx=XH_{\rm pow}/(1-h)$ for $h<1$, this is exactly
$\int_0^\infty(H-H_{\rm pow})\,dX=0$. The subtraction is integrable at
zero, and is identically zero at large $X$ at this pre-heat stage.

For the other moment, note the exact identities on the $l=-1$ hold:
\[
\frac{(XE^2)_{\rm end}}{(XE^2)_{\rm start}}=h^8,
\qquad \frac{E_{\rm end}}{E_{\rm start}}=h^6.
\]
They follow by integrating $(\log(XE^2))'=2l=-2$ and
$(\log E)'=-3/2$ over $4\log(1/h)$.
The first transition and hold have total $\int XE^2dy\le CX_{\rm rel}e_{\rm rel}^2$.
After the hold $l\le-3h/4$, so the remaining integral is at most
$CX_{\rm rel}e_{\rm rel}^2h^8\int_0^\infty e^{-3hy/2}\,dy
\le CX_{\rm rel}e_{\rm rel}^2h^7$. This proves the $h$ uniform bound (A.18).
Interpolation and its next hold cost at most
$C(1+\log(1/\lambda))X_{\rm end}e_{\rm end}^2$; their first parameter
derivatives satisfy the same bounds because $f^{\vartheta}$ and its derivatives
are bounded and the exterior is parameter independent.

There is also an exact algebraic specification of the amplitude root. Substitute
the affine $c$ above and write $R_b=R^{(0)}+\mathrm{Amp}\,R^{(1)}$.
In the full $S(\infty)$, every term outside $\int_{\rm pulse}E^2R_b^2\,dX$
is independent of $\mathrm{Amp}$. Consequently this moment is an explicitly
defined quadratic polynomial $a(\eta)\mathrm{Amp}^2+b(\eta)\mathrm{Amp}+c(\eta)$,
where, without dropping any end bumps,
\[
\begin{aligned}
a&=\int_{\rm pulse}E^2(R^{(1)})^2dX>0,\qquad
b=2\int_{\rm pulse}E^2R^{(0)}R^{(1)}dX,\\
c&=S_{\rm no\ axial\ pulse}(\infty)+\int_{\rm pulse}E^2(R^{(0)})^2dX.
\end{aligned}
\]
Changing variables $\xi_b=\lambda y$ in this polynomial gives
\[
\frac{\lambda S(\infty)}{X_pe_b^2f^2}
=\mathrm{Amp}^2K_b-\frac{1-e^{-26}}4+\mathcal E(\mathrm{Amp},\eta),
\quad K_b=\int_0^{13}e^{-2\xi}R_0(\xi)^2\,d\xi,
\]
\[
\|\mathcal E\|_{C^1([.9,1.2]\times[-1,1])}
\le C_{\rm pre}\lambda(1+\log(1/\lambda)).
\]
Indeed the prior $S$ and the post-pulse negative $E^2/2$ integral have precisely
the bounds just proved after multiplication by $\lambda$, whereas the affine
end bumps are exponentially small; all cross terms are retained in
$\mathcal E$. On $.02\le\xi\le9$, $R_0\ge\xi-.02$, and globally
$0\le R_0\le\xi$. Therefore
\[
\frac{e^{-.04}}4\left[1-e^{-17.96}(1+17.96+2(8.98)^2)\right]
\le K_b\le\frac14,
\]
whose lower bound exceeds $.20$. At $\mathrm{Amp}=.9$ the principal value
is below $-.047$; at $1.2$ it exceeds $.038$; its derivative is at least $.36$.
For sufficiently small $\lambda$ the exact polynomial has a unique root on
this interval. It is the positive branch
$(-b+\sqrt{b^2-4ac})/(2a)$, with its nonzero derivative already proved by
the bracket estimate. Implicit differentiation gives
$|\partial_\eta\mathrm{Amp}|\le C_{\rm pre}\lambda(1+\log(1/\lambda))$.
This constructs the smooth amplitude and proves both required moments, including
all terms in (A.19)--(A.20).

\section{Pressure datum and the cone margins on every outer stage}
\subsection{Exact pressure and its parameter analyticity}
Omit only the swirl bumps whose total $\int E^2dy$ was just proved to be zero.
The remaining scheduled field has the form $E=c(y)f(\eta)^{\vartheta(y)}$,
$0\le\vartheta\le1$, with $c,\vartheta$ and all interval lengths independent
of $\eta$ and $X_R$. On one simply connected complex neighborhood of $[-1,1]$
avoiding $\pm i$, choose the logarithm of $f$. The holomorphic functions
$\exp(\vartheta\log f)$ are uniformly bounded for $0\le\vartheta\le1$.
The squared integrand is bounded by $Ce^{y/5}$ as $y\to-\infty$ and
$Ce^{-(1+2h)y}$ as $y\to\infty$. Uniform convergence of the integral on
compact complex subsets therefore defines a holomorphic function
\[
\Pi_{\rm ax}(\eta)=-\tfrac12\int_{-\infty}^\infty E_{\rm sched}(y,\eta)^2dy.
\]
It is even, independent of $X_R$, and its derivative is an integral of
$\vartheta J_0'c^2f^{2\vartheta}$. This has the sign of $\eta$ and is
strict for $\eta\ne0$ because the reference interval has $\vartheta=1$.
The reference interval contributes exactly
$-\tfrac12P_*^2f^2\int_{-\infty}^0e^{y/5}dy=-\tfrac52P_*^2f^2$.
Restoring the pressure-neutral bumps leaves the total unchanged. Hence the
actual pressure is exactly
\[
\Pi(X,\eta)=-\tfrac12\int_{\log(X/X_R)}^\infty E(v,\eta)^2dv.
\]
This proves all assertions of [NS, Lemma A.5]. It also proves the precise
preservation map: an edit of zero total $\int\delta(E^2)/(2X)dX$ keeps
forward pressure unchanged before and after its support; it need not keep
pressure unchanged inside that support.

\subsection{The exact quantities used for the cone test}
Retain
\[
W=1-2D\eta M/X-dM_\eta/X,\quad H_c=D\eta+dU,
\quad a=2-2l,\quad b_s=2XU_X/E,
\]
\[
Q_s=-W+\frac{(1-h)I-D\eta I_\eta-dJ_\eta+2(h-D)\eta J}{XH},
\]
\[
N_s=-WU+\frac{D(M-\eta M_\eta)+4h\eta S-dS_\eta}{X}
+4A\eta\Pi-d\Pi_\eta.
\]
These formulas follow by integrating the source equations
\[
Q_s'+(1+l)Q_s=-Wl-h(1-2\eta U)-H_c(\log E)_\eta,
\]
\[
N_s'+N_s=-WU'-A(1-2\eta U)U-H_cU_\eta
-d\Pi_\eta+4A\eta\Pi+2\eta X\Pi_X.
\]
For example $(XW)_X=1-2D\eta U-dU_\eta$; integration by parts of
$-WXH_X$ produces $-XWH$ and the displayed $I,J$ terms. In the pressure
terms use $\int_0^X\Pi,dx=X\Pi-\int_0^XE^2/2,dx$.
These operations produce the exact $S$ and $M$ terms above, and establish the
map from the five original moments to stress coefficients in [NS, (4.16)].

Set $w=N_s/(EQ_s)$ where $Q_s>0$, and $p_{s,1}=XQ_s/L$,
$p_{s,2}=wp_{s,1}$. To explain precisely the sufficient test used below, put
\[
c=1-b_sw/a,\quad j=w+b_s/a,\quad v_s=a+b_s^2/a,
\quad P_c=p_{s,1}c,\quad J_c=p_{s,1}j.
\]
Direct expansion gives
\[
2c^2-(v_s-2)j^2
=\left(1+\frac{b_s^2}{a^2}\right)
\left[2-2b_sw-\frac{b_s^2}{a}-(a-2)w^2\right].
\]
Thus on a compact set with $a>0$, $a-b_sw>0$, and the bracket positive,
there are constants $c_*>0$, $G_*>0$, and $B<\infty$ bounding these
expressions away from zero and bounding $v_s,cv_s$. For
$p_{s,1}\ge\max((B+2)/c_*,8B/G_*)$,
$P_c>\max(2,v_s)$ and
\[
\frac{2(P_c-v_s)^2-(v_s-2)J_c^2}{p_{s,1}^2}
\ge G_*-4B/p_{s,1}\ge G_*/2.
\]
For $v_s>2$ these are exactly the admissible cone inequalities; for $v_s\le2$
they imply the relaxed cone. Equivalence follows by solving the quadratic
$2(P_c-v)^2-(v-2)J_c^2=0$, whose smaller root is the function in (4.21),
between $2$ and $P_c$. This verifies the finite threshold used in Lemma 4.5.

\subsection{Reference and axial-decrease stages}
On the ideal reference interval $U=4\eta$, $W=1-4L$, $l=3/5$, so
\[
Q_s=\frac{(3/5)(4L-1)-h(1-8\eta^2)+(D+4d)\eta J_0'}{8/5}>c>0
\]
for the chosen small $h$. Here $a=4/5$, $b_s=0$. Its $N_s$ geometric
source is $O(|\eta|)$ and its pressure source is $O(P_*^2|\eta|)$;
these estimates follow by substituting the constant $U$ and the bounded
backward pressure integral above. The vanishing at zero follows from even
pressure and odd $U$, and one parameter derivative removes the $|\eta|$
factor. The $N_s'+N_s$ convolution has finite bounds on each prescribed
reference interval $x\ge x_->0$.

During the first unit transition, $b_s=0$, $.8\le a\le2$, and
the angular source is at least $-h$ since $W<0$, $l\ge0$, and
$H_cJ_0'\ge0$. The initial positive lower bound and variation of constants
retain $Q_s>c$ for small $h$. With $b_s=0$ and $a\le2$, the second sufficient
cone bracket is automatically at least $2$.

During axial decrease, $l=0$, $a=2$, $b_s=2k'\eta/E$. The angular
source is exactly
$-h(1-2k\eta^2)+(D+dk)\eta J_0'$.
Since $\eta J_0'=2\eta^2/(1+\eta^2)\asymp\eta^2$, convolution with
$e^{-y}$ and $h\ll e^{-T_d}$ gives
$Q_s\ge c(\eta^2+e^{-y})$ for $0\le y\le T_d$.
The geometric axial source is $O(|\eta|)$. The pressure source is
$O(E^2|\eta|)$ because subsequent scheduled swirl decays at least
$e^{-\Delta y/2}$ and $|\partial_\eta\log E|\le |J_0'|$; the later
pressure-neutral angular bumps may be omitted from the backward integral.
Solving the $N_s$ equation yields
\[
|N_s|/E^2\le C|\eta|(e^{T_d}/P_*^2+1+y)
\le C|\eta|(1+y).
\]
Consequently
$|b_sw|\le C|k'|\eta^2(1+y)/(\eta^2+e^{-y})\le Ce_d$ and
$b_s^2\le Ce_d^2e^{T_d}/P_*^2$.
Choose $M_d$ so $Ce_d$ is small, then $P_*>e^{T_d}$ sufficiently large.
Both required strict margins follow, uniformly in $\eta$ and this whole stage.

\subsection{Intermediate power law and the complete pulse}
Write $M/X=\bar k\eta$ during axial decrease. Then $\bar k'=k-\bar k$.
Because $0\le k\le4$ and its final eleven units are zero,
$L\bar k<1/2$ at the next transition. When $U=0$ this gives $W\ge1/2$.
Write $c_\eta=D\eta J_0'\asymp\eta^2$. The angular source is
$(-l)W-h+c_\eta$. During the unit transition $0\ge l\ge-\lambda$,
the prior positive lower bound persists; $w$ has a finite bound depending
only on earlier choices. Thus $\sqrt\lambda|w|$ is small. The cone is
admissible as soon as $l<0$, and relaxed at its initial endpoint.

On the constant-slope interval, $\bar k(y)=\bar k(0)e^{-y}$ and
\[
Q_s'+(1-\lambda)Q_s=\lambda-h+c_\eta-\lambda L\bar k(0)e^{-y}.
\]
The right side is at least $\lambda/4+c_\eta$ for small $h/\lambda$.
Its explicit convolution therefore gives
\[
Q_s\ge c_{\rm pre}(\eta^2+\lambda+e^{-(1-\lambda)y}),\qquad
Q_s-\frac{\lambda-h+c_\eta}{1-\lambda}
=O_{C^1}(C_{\rm pre}(1+y)e^{-(1-\lambda)y}).
\]
There is no geometric $N_s$ source when $U=0$; the pressure source is
$O(E^2|\eta|)$. Its convolution gives
$|N_s/E|\le C_{\rm pre}|\eta|(1+y)e^{-(1/2-\lambda)y}$.
Using $|\eta|/(\eta^2+e^{-(1-\lambda)y})
\le\tfrac12e^{(1-\lambda)y/2}$ proves
\[
\sqrt\lambda|w|\le C_{\rm pre}\sqrt\lambda(1+y)e^{\lambda y/2}=o(1)
\quad(0\le y\le60\log(1/\lambda)).
\]
Since $a=2+2\lambda$, $b_s=0$, this supplies the strict cone bracket.

During the pulse set $m=(M/X)/E$ and $\beta=1/2-\lambda$.
The exact identity $m'+\beta m=R_b$ follows by differentiating $M/X$ and $E$.
With the pulse extended flatly to the left, write its convolution as
$\int_0^\infty e^{-\beta v}R_b(\lambda(y-v))\,dv$ plus
$m(0)e^{-\beta y}$. Taylor's formula under this integral gives
\[
m=R_b/\beta-\lambda\partial_{\xi_b}R_b/\beta^2
+O(\lambda^2)+m(0)e^{-\beta y},
\]
because the integral of $v^2e^{-\beta v}$ is finite uniformly for
$\beta>.4$. The end bumps have exponentially small fixed $\xi_b$ derivatives.
The same estimate holds after one $\eta$ derivative.

Since $E\le C_{\rm pre}\lambda^{30}$, all terms from
$W=1-2D\eta Em-d(Em)_\eta$ and from $dU$ in the angular source have
size $O(C_{\rm pre}E)$, including one parameter derivative. The initial
constant-slope estimate and exponential convolution give
\[
Q_s=\frac{\lambda-h+c_\eta}{1-\lambda}+O(C_{\rm pre}\lambda^{28}).
\]
The backward pressure and its first derivative are $O(E^2)$.
For $S$, integration over $0\le\lambda y\le13$ bounds
$S/(XE^2)$ by $C_{\rm pre}e^{26}/\lambda$ in $C^1$.
The numerator accumulates at most a fixed integral in $\xi_b$ divided by
$\lambda$ and the denominator loses at most $e^{-26}$; this explains the
constant $e^{26}/\lambda$, not $e^{26/\lambda}$.
Substituting $M/X=Em$ in the exact moment formula for $N_s$ gives
\[
N_s/E=-R_b+(D+c_\eta)m-D\eta m_\eta+O(C_{\rm pre}\lambda^{27}).
\]
For example $D(M-\eta M_\eta)/(XE)
=D(m-\eta m_\eta+\eta J_0'm)$, giving exactly $(D+c_\eta)m$.
The $S$ and pressure remainders are at most $CEe^{26}/\lambda$ in this
quotient, which is smaller than the stated bound for sufficiently small
$\lambda$. The difference $W-1$ similarly costs $O(E)$.

The amplitude derivative and the initial-moment bounds show
$|m_\eta|\le C_{\rm pre}\lambda(1+\log(1/\lambda))$.
Since $|\eta|/(\lambda+c_\eta)\le C/\sqrt\lambda$,
the quotient of $\eta m_\eta$ by $Q_s$ is $o(1)$.
The preceding formulas therefore imply, uniformly throughout the pulse,
\[
w=2R_b-C_d\partial_{\xi_b}R_b+o(1),\qquad
C_d=\frac{\lambda(D+c_\eta)(1-\lambda)}
{\beta^2(\lambda-h+c_\eta)}\le2+o(1).
\]
The coefficient of $R_b$ before this last replacement is exactly
$(1-\lambda)/\beta$, since $D+c_\eta-\beta=\lambda-h+c_\eta$.
The quotient defining $C_d$ is maximal at $c_\eta=0$, because its derivative
in $c_\eta$ has sign $\lambda-h-D<0$; $h/\lambda\to0$ gives the bound.
Also $b_s=-(1+2\lambda)R_b+2\lambda\partial_{\xi_b}R_b=-R_b+O(\lambda)$.
Apart from exponentially small end bumps, $R_b\ge0$ and
$\partial_{\xi_b}R_b\le1.2$, since the declining cutoff has a nonpositive
contribution. The exact upper bounds for the two leading quadratics are
\[
\sup_{R\ge0}(-2R^2+2.41R)=2.41^2/8<.74,
\]
\[
\sup_{R\ge0}(-3.5R^2+4.82R)=4.82^2/14<1.68.
\]
The uniformly $o(1)$ remainders can be made smaller than the strict numerical
gaps. Hence $a-b_sw>1$ and
$2b_sw+b_s^2/a+(a-2)w^2<2$ on the complete pulse, including its descending
and correction parts. Here $v_s\ge a=2+2\lambda>2$.

\subsection{Interpolation and exterior transitions}
After the pulse $U=M=J=0$, so $W=1$.
The angular source during interpolation is $-l-h+\vartheta_fc_\eta$,
giving $Q_s\ge c\lambda$. The tiny angular correction preserves this bound.
Since the complete amplitude solve gives $S(\infty)=0$,
\[
S(X)=\tfrac12\int_X^\infty E(X')^2dX'.
\]
Before exterior release, $l\le-\lambda/2$ and the later release has the
$h$ uniform integral bound proved above. Therefore
$(|S|+|S_\eta|)/X\le CE^2/\lambda$ and
$|\Pi|+|\Pi_\eta|\le CE^2$. The exact moment formula yields
$|w|\le CE/\lambda^2\ll1$, because pulse end reduces $E$ exponentially
in $1/\lambda$. Here $b_s=0$ and $2<a\le2+2\lambda+.2+o(1)$,
so both cone margins remain strict.

On the parameter-independent release and tail, the same backward integrals
give the exact expression
\[
N_s=2\eta\int_y^\infty(he^{v-y}-A)E(v)^2dv.
\]
Since $l\le-3h/4$,
$E(v)^2\le E(y)^2e^{-(1+3h/2)(v-y)}$.
Thus the $he^{v-y}$ integral is at most $\tfrac23E(y)^2$ and the
other at most $E(y)^2/(1+3h/2)$. In particular $|N_s|\le CE^2$
with a constant independent of small $h$.
Before the steep hold, $Q_s\ge c\lambda$ and the amplitude is already
small relative to $\lambda$. During the hold it remains above this bound.
After the hold, $E\le Ce_{\rm rel}h^6$, while $Q_s\ge ch$ up to terminal
coordinate $y=1/2$. The latter follows from its explicit positive terminal
integral and its intervening monotone decay to $Q_p\asymp h$.
Consequently $|w|\le CE/Q_s\ll1$, even if $h\ll e_{\rm rel}$.
Here $b_s=0$, $2<a\le4$, so the sufficient cone test has uniform strict slack.

All these intervals are compact once the sequential parameters are fixed.
Under the exact moment scaling $X=X_Rx$, every $X_R$ factor cancels in
$Q_s,N_s$, while $p_{s,1}=X_RxQ_s/L$ grows linearly with $X_R$.
Thus one finite increase of $X_R$ meets the threshold proved above on the
whole designated interval. This completes the local reconstruction of
[NS, Proposition A.4], including each part of its cone verification.

\section{The heat factor, its exact equation, and the compensated replacement}
\subsection{A positive solution of the radial swirl heat equation}
Define for $Z\ge0$,
\[
\paHeat(Z)=\frac1{\Gamma(1+h)}\int_0^\infty e^{-v}v^h(1+Zv)^{-h}\,dv.
\]
Dominated differentiation gives for every nonnegative integer $m$
\[
\paHeat^{(m)}(Z)=\frac{(-1)^m(h)_m}{\Gamma(1+h)}
\int_0^\infty e^{-v}v^{h+m}(1+Zv)^{-h-m}\,dv,
\quad \paHeat^{(m)}(0)=(-1)^m(h)_m(1+h)_m.
\]
For fixed $m$ the integrand after deleting the last factor is integrable
uniformly for $Z\ge0$. Hence $\paHeat$ is positive, smooth up to zero from the
right, and has bounded fixed derivatives on every bounded nonnegative interval
(in fact on the entire half-line for absolute derivatives).
For $0<h\le h_0<1/2$ and fixed $m\ge1$, the last factorial product is at
most $C_mh$. Taylor's formula therefore proves
\[
\paHeat(Z)=1-h(1+h)Z+O_h(Z^2),\qquad
|\paHeat(Z)-1|+|Z\paHeat'(Z)|\le C hZ
\]
on a fixed bounded nonnegative interval.

The integral solves the required ODE by an exact polynomial calculation.
Multiply
$Z^2\paHeat''+(1+2(1+h)Z)\paHeat'+h(1+h)\paHeat$
by $\Gamma(1+h)/h$, and use common integrand
$e^{-v}v^h(1+Zv)^{-h-2}$. Its remaining polynomial is
\[
(1+h)Z^2v^2-(1+2(1+h)Z)v(1+Zv)+(1+h)(1+Zv)^2
=(1+h)-v-Zv^2.
\]
This is precisely the polynomial from
$\partial_v[e^{-v}v^{1+h}(1+Zv)^{-1-h}]$ in the same common factor.
The boundary values at zero and infinity are both zero. Thus
\begin{equation}\label{pa:eq:heatode}\tag{PA.HEAT}
Z^2\paHeat''+(1+2(1+h)Z)\paHeat'+h(1+h)\paHeat=0.
\end{equation}
Set $K(r,t)=c_\infty s^{-A}\paHeat(2\tau/s)$, $s=r^2/2$, $Z=2\tau/s$.
Direct differentiation gives
\[
K_t=-2c_\infty s^{-A-1}\paHeat',
\]
\[
(\partial_{rr}+r^{-1}\partial_r-r^{-2})K
=2c_\infty s^{-A-1}
\{Z^2\paHeat''+(2A+1)Z\paHeat'+(A^2-1/4)\paHeat\}.
\]
Since $2A+1=2(1+h)$ and $A^2-1/4=h(1+h)$,
\eqref{pa:eq:heatode} proves the heat equation, with its swirl term $-r^{-2}K$.

Normalize the positive integrand defining $\paHeat$ to a probability density.
Then $-Z\paHeat'/\paHeat$ is its average of $hZv/(1+Zv)$, so
$0\le-Z\paHeat'/\paHeat<h$. It follows that
$rK_r/(2K)=-A-Z\paHeat'/\paHeat<-1/2$, in particular $K_r<0$.
At $Z=0$ this is still strict because $h>0$.
Finally $Z=2(1-\eta^2)/X$ gives only factors $-4\eta/X$ and $-4/X$
when differentiated in $\eta$, never a denominator $1-\eta^2$.
Consequently all fixed mixed $X,\eta$ derivatives have one-sided limits at
$\eta=\pm1$, and for each fixed logarithmic radial and parameter derivative,
\[
\frac{E_{\rm heat}}{E_{\rm pow}}
=\paHeat(2d/X)=1-\frac{2h(1+h)d}{X}+O_{h,j,m}(X^{-2}).
\]
This proves [NS, Lemma A.6] on its stated nonnegative domain. No analytic
continuation through negative $Z$ is used or required.

\subsection{The three heat discrepancies and their exact cancellation}
Let $y=\log(X/X_{\rm tail})$, $X_K=e^{.2}X_{\rm tail}$,
$X_b=e^3X_{\rm tail}$, and choose $\chi_K=0$ for $y\le.2$, $\chi_K=1$
for $y\ge.5$. Replace the reference swirl $E_{\rm cl}$ by
\[
E_{\rm cl}[1+\chi_K(y)(\paHeat(2d/X)-1)].
\]
For $x=X/X_K\ge1$ and its reference amplitude $e_K$, the preceding
derivative bounds give exactly
\[
|(X\partial_X)^j\partial_\eta^m(E-E_{\rm cl})|
\le C_{j,m}e_KX_K^{-1}x^{-A-1}.
\]
The profile is comparable to $e_Kx^{-A}$ with constants depending on the
fixed schedule. Thus the pressure, $S$, and subtracted angular discrepancies
obey, after $m$ parameter derivatives,
\[
|\Delta C_p|_m\le C_m\frac{e_K^2}{X_K}\int_1^\infty x^{-2A-2}dx,
\quad |\Delta S|_m\le C_me_K^2\int_1^\infty x^{-2A-1}dx,
\]
\[
\left|\Delta\int(H-H_{\rm pow})dX\right|_m
\le C_me_KX_K^{1/2}\int_1^\infty x^{-A-1/2}dx
=C_m e_KX_K^{1/2}/h.
\]
The powers arise respectively from $\delta(E^2)dX/(2X)$,
$-\delta(E^2)dX/2$, and $\sqrt{2X}\delta E,dX$.
The final $1/h$ is retained: $h$ has already been fixed before $X_R$ grows.

At the earlier second reserved patch write $x_*=X/X_*$,
$E=e_*f x_*^{-1/2-\lambda}$, $U=0$, and add
$\delta E=e_*\sum_{i=1}^3 c_i(\eta)\beta_i(x_*)$.
Divide the three rows by $e_*^2$, $X_*e_*^2$, and $X_*^{3/2}e_*$.
The exact linear weights are then
\[
f x_*^{-3/2-\lambda},\qquad -f x_*^{-1/2-\lambda},\qquad
\sqrt2 x_*^{1/2}.
\]
The first two rows have exact quadratic terms and the third is linear.
The moment determinant is nonzero by the moment-map calculation above; its inverse is uniform in
$\eta$. Since $X_*/X_K$ and $e_*/e_K$ are fixed, all three normalized
discrepancies are $O_m(X_K^{-1})$. The contraction gives the exact correction
with $\|c\|_{C^m}\le C_mX_K^{-1}$ for any prescribed finite $m$, after
choosing $X_R$ large. More generally its higher fixed derivative bounds follow
by implicit differentiation with the same small forcing: the Jacobian stays
invertible and every differentiated equation has only lower derivative products
and an $O_m(X_K^{-1})$ inhomogeneity. This justifies the stated $O_m$ estimates
without selecting infinitely many smallness thresholds at once.

Both edits have $U=0$, so $M,J$ are unchanged exactly. All three other total
moments are restored by this solve. The total pressure integral is therefore
unchanged, so the analytic axis pressure and every profile before the patch
remain unchanged. On the compact logarithmic interval up to $y=.5$,
the exact moment formulas transfer the $O(X_K^{-1})$ changes of profiles and
moments to $Q_s,N_s$, costing one more $\eta$ derivative.
Indeed substituting $I=X_R^{3/2}\widehat I$, $J=X_R^{3/2}\widehat J$,
$M=X_R\widehat M$, $S=X_R\widehat S$, $H=X_R^{1/2}\widehat H$
cancels every $X_R$ in those formulas. Positive lower bounds for $E,Q_s$
and the already strict cone margins persist by continuity, uniformly on this
fixed compact interval. This proves [NS, Proposition A.7] and its preservation
claims; the remaining endpoint is treated next rather than being inferred
from a compact-interior estimate.


% Source body: axis_exact.md
\section{Analytic axis construction and endpoint inequalities}\label{pa:sec:component-axis}

\subsection{1. Axis data and the complementary parameter regions}

Keep \(A=1/2+h\), \(D=1/2-h\), \(d=1-\eta^2\), \(L=1-2h\eta^2\), and the already fixed \(0<h\le10^{-2}\). The datum constructed in (PA.4) is holomorphic near \([-1,1]\), real and even on that interval, and has
\[
\Pi_0\le-\frac52P_*^2(1+\eta^2)^{-2},\qquad
\eta\Pi_0'(\eta)>0\quad(\eta\ne0).
\]
Choose \(0<j_0\le0.05\) and set exactly
\[
U_*=4\eta+j_0,\quad H_*=D\eta+dU_*,
\quad W_*=1-dU_*'-2D\eta U_*,
\]
\[
Z_*=-A(1-2\eta U_*)U_*-H_*U_*'-d\Pi_0'+4A\eta\Pi_0.
\]
For \(-1<\eta<1\),
\[
\frac{H_*}{d}=\frac{D\eta}{1-\eta^2}+4\eta+j_0,
\quad \partial_\eta(H_*/d)
=D\frac{1+\eta^2}{(1-\eta^2)^2}+4>0.
\]
Its endpoint limits are \(-\infty,+\infty\), so it has exactly one zero \(\eta_0\). Its value at zero is \(j_0>0\), hence \(\eta_0<0\). The root equation yields
\[
|\eta_0|=\frac{j_0}{4+D/(1-\eta_0^2)}.
\]
First this is at most \(j_0/4\le1/80\). Substitution back gives the useful explicit bounds \(j_0/5<|\eta_0|<j_0/4\). Further,
\[
-W_*=3-8h\eta^2+(1-2h)j_0\eta
\ge3-0.08-0.05=2.87>2.8.
\]
At the zero, \(H_*U_*'=0\) and \(-d\Pi_0'\ge0\). The term \(4A\eta_0\Pi_0\) is positive. Using \((1+\eta_0^2)^{-2}\ge1/4\), it is at least \((A/2)j_0P_*^2\). The remaining negative term has magnitude bounded by a fixed multiple of \(j_0\), since \(U_*=-D\eta_0/d\) at the zero. Therefore the already large fixed \(P_*\) gives \(Z_*(\eta_0)>0\), quantitatively at least a positive constant times \(j_0P_*^2\).

Continuity gives an open neighborhood \(V\) of \(\eta_0\) on which \(Z_*\) is bounded below by a positive number. Choose \(\delta_*>0\) smaller than half that lower bound. The compact set \(K_\delta=\{\eta:|Z_*(\eta)|\le\delta_*\}\) avoids \(V\), hence \(|H_*|\ge m>0\) there (if this set is empty, the following restriction is vacuous). Choose \(\sigma_*>0\) with \(\sigma_*^2<m^2/99\). Then
\[
\chi=\frac{H_*^2}{H_*^2+\sigma_*^2}>0.99\quad\hbox{on }K_\delta.
\]
Thus \(\chi\le0.99\) implies \(|Z_*|>\delta_*\), exactly the later endpoint alternative. Strict margins allow the same choices on a slightly larger compact interval \(I\supset[-1,1]\). There is a bounded simply connected complex neighborhood \(\Omega\) of \(I\) on whose closure \(L\) and \(H_*^2+\sigma_*^2\) are nonzero and \(\Pi_0\) is holomorphic: take a sufficiently thin neighborhood of the compact real interval, avoiding the closed zero sets.

Define
\[
\zeta_*=-\frac{LH_*}{H_*^2+\sigma_*^2},\qquad \xi_0=\Lambda\zeta_*,
\quad\phi_*(\eta)=\exp\left(\Lambda\int_0^\eta\zeta_*(w)\,dw\right).
\]
The primitive is single-valued on \(\Omega\), and \(\phi_*>0\) on \(I\). The exact identity needed below is \(H_*\zeta_*/L=-\chi\). After choosing \(\Lambda\), impose
\[
C\ge\sup_{\eta\in\Omega}|\phi_*(\eta)|.
\]
This supremum is finite after choosing \(\Omega\) with compact closure inside the holomorphic domain. Then \(g=\phi_*/C\) has a complex bound independent of both later large parameters. A real-interval bound alone would not establish the required parameter derivative bounds.

\subsection{2. The coefficient space, with all derivative losses accounted for}

Set \(Y=\Lambda X\). For coefficient sequences \(F=\sum_{\alpha\ge0}F_\alpha(\eta)Y^\alpha\), use precisely
\[
a_{\alpha\beta}=
\frac{20^{-\alpha}\rho^{-\beta}\beta!\binom{\alpha+\beta}{\beta}}
{(\alpha+1)^2(\beta+1)^2},
\qquad
\|F\|_\rho=\sup_{\alpha,\beta\ge0}\sup_{\eta\in I}
\frac{|\partial_\eta^\beta F_\alpha(\eta)|}{a_{\alpha\beta}}.
\]
Completeness follows coefficient by coefficient: a norm-Cauchy sequence converges uniformly for every derivative, and the fundamental theorem of calculus identifies each derivative limit with the derivative of its predecessor. The supremum bound passes to the limit, so this is a Banach space.

Let \(C_{\rm sq}=8\sum_{n\ge1}n^{-2}\). For \(0\le i\le N/2\), \((N+1)^2/(N-i+1)^2\le4\), while the other half has the analogous bound after \(i\mapsto N-i\). Therefore
\[
\sum_{i=0}^N\frac{(N+1)^2}{(i+1)^2(N-i+1)^2}\le C_{\rm sq}.
\]
In a coefficient product and its \(\beta\)-th derivative, the Leibniz factor \(\binom\beta k\) cancels the factorials in the two input weights. The remaining binomial quotient is at most one because
\[
\binom{i+k}{k}\binom{\alpha-i+\beta-k}{\beta-k}
\le\binom{\alpha+\beta}{\beta}.
\]
It counts subsets selecting specified numbers from the two fixed blocks. Applying the preceding convolution bound once in the radial and once in the parameter index proves
\(\|FG\|_\rho\le C_{\rm sq}^2\|F\|_\rho\|G\|_\rho\).

Define \(\mathcal I F(Y)=\int_0^YF(v)\,dv\), \(\mathcal A F=Y^{-1}\mathcal I F\), and let \(\mathcal J_\nu F\), \(\nu=1,2\), be the regular solution with zero value at \(Y=0\) of \(YG_{YY}+\nu G_Y=F\). Comparing coefficients gives
\[
(\mathcal J_\nu F)_{\alpha+1}=\frac{F_\alpha}{(\alpha+1)(\alpha+\nu)},
\quad (\mathcal J_\nu F)_0=0.
\]
The exact weight ratios are
\[
\frac{a_{\alpha\beta}}{a_{\alpha+1,\beta}}
=20\frac{(\alpha+2)^2}{(\alpha+1)^2}
\frac{\alpha+1}{\alpha+\beta+1}\le80,
\]
\[
\frac{a_{i,k+1}}{a_{i+1,k}}
=\frac{20}{\rho}(i+1)\frac{(i+2)^2}{(i+1)^2}
\frac{(k+1)^2}{(k+2)^2}\le\frac{80}{\rho}(i+1).
\]
These prove boundedness of multiplication by \(Y\), \(\mathcal I\), \(\mathcal J_\nu\), averaging (which divides coefficient \(\alpha\) by \(\alpha+1\)), and \(\partial_\eta\mathcal I\). The last assertion uses its integration divisor \(i+1\) to cancel the last factor in the second ratio.

For \(\mathcal J_\nu[(\partial_\eta F)(Y\partial_YG)]\), consider output degree \(N=\alpha+1\), split the input radial degrees as \(i\) and \(\alpha-i\), and allocate the added degree to the differentiated \(F\)-factor. Put \(I_1=i+1\), \(J_1=\alpha-i\). The weight comparison and radial inverse give, after dividing by the output weight, at most
\[
\frac{80}{\rho}\|F\|_\rho\|G\|_\rho
\sum_{i=0}^\alpha\sum_{k=0}^\beta
\frac{I_1J_1}{N(\alpha+\nu)}
\frac{(N+1)^2}{(I_1+1)^2(J_1+1)^2}
\frac{(\beta+1)^2}{(k+1)^2(\beta-k+1)^2}.
\]
The first fraction is at most one, including the zero endpoints; the two sums are bounded by \(C_{\rm sq}\). If the logarithmic derivative is absent, replace \(J_1\) by one, still giving a fraction at most one. If \(F\) is averaged, its divisor \(i+1\) cancels \(I_1\). If the parameter derivative is absent, use the first weight ratio instead. Inserting further undifferentiated factors adds the same convolution bounds with unchanged degree allocation. This proves every nonlinear operator estimate invoked below, including the case where both differentiated factors are copies of the same unknown.

A holomorphic coefficient \(b(\eta)\) bounded on a larger neighborhood belongs as degree zero to this space: Cauchy's estimate gives \(|b^{(\beta)}|\le M\beta!r^{-\beta}\), and choosing \(\rho<r\) bounds \((\beta+1)^2(\rho/r)^\beta\). This applies uniformly to \(g\), since its complex supremum is at most one.

\subsection{3. Exact rescaling and contraction}

Write the original fields as
\[
\phi=\phi_*\Phi,\quad U=U_*+\Lambda^{-1}u,
\quad \Pi=\Pi_0+\Lambda^{-1}p,
\quad p=\mathcal I(g^2\Phi^2),\quad \Phi(0)=1,\quad u(0)=0.
\]
Thus \(\Pi_X=g^2\Phi^2=\phi^2/C^2\) exactly. Set
\[
B=-2D\eta\mathcal A u-d\partial_\eta\mathcal A u,
\quad W=W_*+\Lambda^{-1}B,\quad H_c=H_*+\Lambda^{-1}du.
\]
Substitute these into the two stress-free equations in \texttt{core\_exact.md}. The unchanged equations are
\[
2(Y\Phi_{YY}+2\Phi_Y)=-\chi\Phi+\Lambda^{-1}R_1,
\qquad 2(Yu_{YY}+u_Y)=-Z_*/L+\Lambda^{-1}R_2,
\]
where
\[
LR_1=[W+h(1-2\eta U)+du\zeta_*]\Phi+WY\Phi_Y+H_c\Phi_\eta,
\]
\[
LR_2=[A(1-4\eta U_*)+dU_*']u-2A\eta\Lambda^{-1}u^2
+WYu_Y+H_*u_\eta+d\Lambda^{-1}uu_\eta
-4A\eta p+d\partial_\eta p-2\eta Yp_Y.
\]
Both identities have been checked as exact rational expressions by \texttt{exact\_checks.py}, before any estimate. In particular the pressure terms retain all three signs and coefficients. The angular \(-\chi\Phi\) comes from \(H_*\zeta_*/L=-\chi\); this term is order one and must be inverted, not treated as small.

The only products with both an unintegrated parameter derivative and a logarithmic radial derivative are \((\partial_\eta\mathcal A u)Y\Phi_Y\) and \((\partial_\eta\mathcal A u)Yu_Y\). They have the precise structure proved above. The term \(\partial_\eta p=\partial_\eta\mathcal I(g^2\Phi^2)\) is bounded, and \(Yp_Y=Yg^2\Phi^2\). Every other monomial is controlled by multiplication and a radial inverse, with at most one of the two derivative types. Subtracting two monomials by replacing their factors one at a time proves the corresponding local Lipschitz estimates. All constants are finite on a fixed ball, uniformly for \(\Lambda\ge1\) and \(C\) satisfying the complex bound.

Let \(M_\chi\) be the bounded multiplier norm of \(\chi\), and set \(T=\mathcal J_2\chi/2\), multiplication preceding integration. A sequence vanishing below degree \(b\) obeys
\[
\|\mathcal J_2F\|_\rho\le\frac{80}{(b+1)(b+2)}\|F\|_\rho.
\]
Multiplication by \(\chi\) preserves that vanishing order and \(\mathcal J_2\) increases it by one. Therefore
\[
\|T^k\|\le\frac{(40M_\chi)^k}{k!(k+1)!}.
\]
The absolutely convergent series \(\sum_{k\ge0}(-T)^k\) is a two-sided inverse of \(1+T\): multiply finite partial sums by \(1+T\) and use \(\|T^{k+1}\|\to0\). No smallness of \(\chi\) is required.

The center of the contraction is
\[
\Phi_0=f_0(Y\chi),\quad u_0=-\frac{YZ_*}{2L},
\qquad f_0(z)=\sum_{\alpha\ge0}\frac{(-z/2)^\alpha}{\alpha!(\alpha+1)!}.
\]
Its recurrence verifies \(2(zf_0''+2f_0')=-f_0\), \(f_0(0)=1\), so \(\Phi_0=(1+T)^{-1}1\). On the closed radius-one ball about \((\Phi_0,u_0)\), let \(M\) bound the pair
\(((1+T)^{-1}\mathcal J_2R_1/2,\mathcal J_1R_2/2)\), and let \(K\) be its Lipschitz constant. These are the finite constants supplied by the proved estimates. Choosing
\(\Lambda\ge\max\{1,2M,2K\}\) makes the map
\[
(\Phi,u)\longmapsto
\left(\Phi_0+\frac{(1+T)^{-1}\mathcal J_2R_1}{2\Lambda},
u_0+\frac{\mathcal J_1R_2}{2\Lambda}\right)
\]
invariant and a contraction with factor at most \(1/2\). Its unique fixed point has norm distance at most \(M/\Lambda\). This explicitly supplies the existential large-parameter threshold from finite operator constants.

To recover analytic functions, for \(R<20\) use
\[
\sum_{\alpha\ge0}\binom{\alpha+\beta}{\beta}(R/20)^\alpha
=(1-R/20)^{-\beta-1}.
\]
The coefficient norm therefore bounds parameter derivatives on \(|Y|\le R\) by a constant times \(\beta![\rho(1-R/20)]^{-\beta}\). It gives analytic continuation in a common positive parameter radius. Choose \(4.1<R_1<R_2<20\); Cauchy's formula on the larger radial disk and a slightly smaller parameter neighborhood gives, for every fixed \(r,s\),
\[
\sup_{[0,4.1]\times I}
\left|\partial_Y^r\partial_\eta^s(\Phi-f_0(Y\chi),u+YZ_*/(2L))\right|
\le C_{r,s}/\Lambda.
\]
The same parameter domain serves every fixed radial derivative because those derivatives use radial Cauchy estimates on the fixed larger disk. For original radial derivatives the factor is exactly \(\partial_X^r=\Lambda^r\partial_Y^r\), giving \(C_{r,s}\Lambda^{r-1}\). Real coefficients and uniqueness imply real profiles on the real rectangle.

\subsection{4. Positivity and endpoint alternative with the original constants}

For \(0\le z\le4.1\), put \(t=z/2\le2.05\). The magnitude ratio of consecutive terms of \(f_0\) is \(t/[(n+1)(n+2)]\), so the alternating tail after the cubic has decreasing magnitudes. Consequently
\[
f_0(z)\ge1-t/2+t^2/12-t^3/144
\ge305719/1152000>0.265.
\]
The cubic derivative is \(-1/2+t/6-t^2/48<0\): its discriminant is negative and its leading coefficient is negative. Its minimum is thus the displayed exact endpoint value. Grouping the alternating series after the constant term gives \(f_0\le1\). The proved \(O(\Lambda^{-1})\) bound gives \(\Phi\ge c_0>0\) when \(\Lambda\) is sufficiently large. In particular division by \(\phi\) in the original equation is justified. Uniform complex derivative bounds then give a smaller common complex domain with no zeros; derivatives of \(\log\Phi\) are bounded there.

Differentiate \(\chi=H_*^2/(H_*^2+\sigma_*^2)\):
\[
|H_*\chi'|=2|H_*'|\frac{H_*^2\sigma_*^2}{(H_*^2+\sigma_*^2)^2}
\le2\|H_*'\|_\infty\chi.
\]
Since \(f_0\) is positive and has bounded fixed derivatives on the stated compact range,
\[
|D_X\log\Phi|+|H_*\partial_\eta\log\Phi|
\le C\chi+C/\Lambda.
\]
Here \(D_X=Y\partial_Y\) and no division by \(\chi\) is used. Also
\[
-H_c\xi_0=\Lambda L\chi-du\zeta_*,\qquad
|du\zeta_*|\le C_{\sigma_*}\sqrt\chi,
\]
because \(|H_*|/(H_*^2+\sigma_*^2)\le\sigma_*^{-1}\sqrt\chi\). Thus the full source has the lower bound
\[
S_q\ge\Lambda L\chi-W_*-10h-C\chi-C_{\sigma_*}\sqrt\chi-C/\Lambda.
\]
For any fixed small \(\varepsilon>0\), Young's inequality in the exact form
\(K\sqrt\chi\le\varepsilon\Lambda L\chi+K^2/(4\varepsilon\Lambda L)\)
absorbs the square-root term. Increase \(\Lambda\) so the linear \(C\chi\) term uses the remaining part of a total \(0.05\Lambda L\chi\) allowance. Since \(-W_*>2.8\), \(10h\le0.1\), and the remaining constant error can be made at most \(0.2\), this gives
\[
S_q\ge2.5+0.95\Lambda L\chi.
\]

The integral formula
\[
Q_s(X,\eta)=\frac{\int_0^Xx\Phi(\Lambda x,\eta)S_q(x,\eta)\,dx}
{X^2\Phi(\Lambda X,\eta)}
\]
proves \(p_1/X=Q_s/L\ge c_1>0\). Stress freedom gives the exact identities
\[
p_1=\frac{XQ_s}L=a=-2Y\Phi_Y/\Phi,
\quad n_s=N_s/L=-2U_X=-2u_Y=Z_*/L+O(\Lambda^{-1}),
\quad p_2=Xn_s/E.
\]
All fixed parameter derivatives of \(\Phi,\log\Phi,u,p_1,n_s\) are bounded uniformly in subsequent \(C\), with \(\Lambda\) fixed; the radial scaling factor has already been accounted for.

At \(X_0=4/\Lambda\), first take \(\chi>0.99\), so \(t=2\chi\in(1.98,2]\). The series for \(f_0(z)+zf_0'(z)\) has terms \((-t)^n/(n!)^2\). Its alternating remainder beyond degree four is negative, giving
\[
f_0+zf_0'\le1-t+t^2/4-t^3/36+t^4/576.
\]
The derivative is \(-1+t/2-t^2/12+t^3/144\); on \([1.98,2]\) its first two terms are nonpositive and its last two have sum \(-t^2(12-t)/144<0\). Thus the polynomial decreases, and its value at \(1.98\) is exactly \(-75535511/400000000<-0.188<-0.18\). Since \(f_0\le1\),
\[
-2zf_0'/f_0=2-2(f_0+zf_0')/f_0>2.36.
\]
Taking \(\Lambda\) large gives \(p_1>2.3\).

On \(\chi\le0.99\), \(|Z_*|>\delta_*\). Since \(L\le1\) on \([-1,1]\), the preceding error estimate gives \(|n_s|\ge\delta_*/2\). At the same endpoint, exactly
\[
|p_2|=C\sqrt{2/\Lambda}\,\frac{|n_s|}{\phi_*\Phi}.
\]
For fixed \(\Lambda\), \(\phi_*\Phi\) has a finite positive upper bound independent of subsequent \(C\); \(p_1\) has a finite upper bound independent of \(C\). Thus choosing \(C\) large makes \(p_2^2/p_1>2.3\) uniformly on this complementary compact set. The two alternatives yield
\[
p_1+p_2^2/p_1>2+0.2
\quad\text{at }X_0=4/\Lambda,
\]
retaining the source's allowed \(c_{\rm ex}=0.2\). The proof uses distinct regions only after constructing their exact covering through \(\chi\) and \(Z_*\); no missing case at \(H_*=0\) remains.



% Source body: continuation_exact.md
\section{Continuation, exact moment repair, and assembly of the annular profile}\label{pa:sec:component-continuation}

This file gives the independent continuation derivation for source pp. 150-157 and checks the assembly on pp. 36-45. Original source parameters remain \(0<h<1/100\), \(X_0=X_a=4/\Lambda\), \(X_b^{\rm ref}=100\), \(X_i=110\). The notation \(X_b^{\rm ref}\) distinguishes the temporary number 100 in Appendix B from the final outer annular edge \(X_b=e^3X_{\rm tail}\) in Section 4; neither numerical value or role is changed. The core and axis sections above provide all field definitions and identities.

\subsection{1. The reference continuation and its estimates}

The smooth step is exactly
\[
\sigma(y)=\frac{e^{-1/y^2}}{e^{-1/y^2}+e^{-1/(1-y)^2}}\quad(0<y<1),
\qquad \sigma=0\ (y\le0),\quad\sigma=1\ (y\ge1).
\]
It is increasing: the ratio of its second denominator term to its first has logarithmic derivative \(-2/(1-y)^3-2/y^3<0\). It is smooth and flat at both endpoints by the corresponding exponential formulas. Choose fixed \(\bar t>0\) with \(4e^{2\bar t}<4.1\). Set \(y=\log(X/X_0)\). For \(0<t_1\le\bar t\), retain the analytic axis profile through \(y=t_1\), then prescribe on \((t_1,2t_1)\)
\[
\partial_y\log\phi_r=(1-\sigma((y-t_1)/t_1))D_X\log\phi_{\rm nat},
\quad \partial_yU_r=(1-\sigma((y-t_1)/t_1))D_XU_{\rm nat}.
\]
After \(2t_1\) keep \(\phi_r,U_r\) constant in \(y\). All endpoint jets match because the step is flat. Pressure and the integrated coefficients continue to be computed from the axis by the original integral definitions.

The axis approximation bounds \(U_{\rm nat}-U_*\) and its fixed parameter derivatives by \(O(\Lambda^{-1})\); its cutoff variation integrates over an interval of length \(t_1\) and is \(O(t_1/\Lambda)\). Thus \(U_r-U_*\) and its radial average remain small uniformly through \(X_i\): explicitly
\[
\|\mathcal A_XU_r-U_*\|_{C^k_\eta}
\le\sup_{0\le X\le X_i}\|U_r-U_*\|_{C^k_\eta}.
\]
No factor equal to the long logarithmic interval appears. Likewise integrate the bounded natural logarithmic slope to bound \(\log\phi_r-\log\phi_*\). Its leading parameter derivative remains \(\xi_0\), with the \(O(\chi)+O(\Lambda^{-1})\) weighted-gradient errors proved for the natural solution and an additional error tending to zero with \(t_1\) after \(\Lambda\) is fixed. In particular \(CE_r=\sqrt{2X}\phi_r\), its reciprocal on \([X_0,X_i]\), and all their fixed parameter derivatives have finite bounds independent of subsequent large \(C\) and sufficiently small \(t_1\).

Pressure satisfies the exact identity and estimate
\[
\Pi_r-\Pi_0=C^{-2}\int_0^X\phi_r(x,\eta)^2\,dx,
\qquad \sup_{X\le X_i}\|\Pi_r-\Pi_0\|_{C^k_\eta}\le K_kC^{-2}.
\]
The axis factor makes this integral finite at zero. The natural shear has \(a=p_1>0\); cutting its nonpositive \(D_X\log\phi\) to zero proves \(l_r\le1\). In the angular source the unchanged large term is \(-H_*\xi_0=\Lambda L\chi\). Replacing \(H_*\) by \(H_{c,r}\) costs \(K_{\sigma_*}\sqrt\chi+o_{t_1}(1)\); the other weighted gradients cost \(K\chi+K/\Lambda\). Young's inequality as in \texttt{axis\_exact.md}, now allocating a total \(0.06\Lambda L\chi\), and the strict \(-W_*>2.8\) margin give
\[
S_{q,r}\ge0.94\Lambda L\chi+2.4.
\]
Here one first enlarges \(\Lambda\), then \(C\), and then decreases \(t_1\); all residual constant errors have a strictly positive allowance. The axial source is
\[
S_{n,r}=Z_*+O(\Lambda^{-1})+o_{t_1}(1)+O(C^{-2})
\]
uniformly on the finite interval. Indeed all changes of \(U\) and of its parameter derivative are small, the logarithmic radial derivative is \(O(\Lambda^{-1})\) on the cutoff and zero later, and pressure is controlled by the preceding exact integral.

The defining primitive equations imply
\[
D_Xp_{1,r}=XS_{q,r}/L-l_rp_{1,r},\qquad
D_Xn_{s,r}+n_{s,r}=S_{n,r}/L.
\]
Positive regular initial data and \(l_r\le1\) imply \(p_{1,r}>0\). Multiplying the first inequality by the integrating factor \(e^y\), keeping the two positive source terms separately, gives
\[
p_{1,r}(y)\ge p_{1,r}(0)e^{-y}
+1.88\chi(e^y-e^{-y}),
\]
\[
p_{1,r}(X)\ge p_{1,r}(X_0)X_0/X
+1.2(X-X_0^2/X).
\]
For the first, \(X=X_0e^y=(4/\Lambda)e^y\) supplies the exact factor \(0.94\Lambda X_0/2=1.88\). For the second, use \(L\le1\), integrate \(2.4X\) against the same factor, and convert back to \(X\). On \(\chi>0.99\), the first lower bound remains above 2.3: replace its initial value by 2.3; its derivative is positive for \(y\ge0\), since \(1.88\chi>2.3/2\). On \(\chi\le0.99\), \(|Z_*|>\delta_*\), and the exact primitive for \(n_{s,r}\) retains its sign and a fixed positive lower bound when the displayed errors are small. Because \(p_{2,r}=Xn_{s,r}/E_r\), increasing \(C\) then makes \(p_{2,r}^2/p_{1,r}\) large uniformly. This proves
\[
v_r:=p_{1,r}+p_{2,r}^2/p_{1,r}>2+c
\]
on the entire reference interval, for a fixed \(c>0\). The second comparison yields \(p_{1,r}>3\) on \([100,110]\), since \(X_0\le4\). Upper bounds for \(p_{1,r},n_{s,r}\) and their fixed parameter derivatives follow directly from their integral representations and the finite coefficient bounds. They are independent of subsequent \(C\) and of the small reference widths.

\subsection{2. Activation of stress, with division by the flat factor proved}

Fix a sufficiently large finite \(C\). Choose \(t_*>0\) so the preceding bounds hold for every \(0<t_1\le t_*\). On this family, \(v_r\) has a finite common upper bound \(V_{\max}\); it may depend on \(C\), which has already been fixed. Choose \(0<\kappa_0<1/2\) later small enough relative to this bound. On \(0<y<t_1\) put
\[
e_a=(1-\kappa_0)\sigma(y/t_1),\quad\kappa=1-e_a,
\quad a=\kappa p_{1,r},\quad D_XU=-\kappa Xn_{s,r}/2,
\quad \partial_y\log\phi=-\kappa p_{1,r}/2.
\]
The natural reference is stress-free here. Therefore the exact differences obey
\[
\partial_y(\log\phi-\log\phi_r)=e_ap_{1,r}/2,
\quad \partial_y(U-U_r)=e_aXn_{s,r}/2.
\]
For a fixed parameter derivative, the integral from zero is bounded by a constant times \(ye_a(y)\), since \(e_a\) is increasing. The constants are uniform in smaller \(t_1,\kappa_0\) after \(\Lambda,C\) are fixed. The five moment differences are integrals of these value differences, with fixed smooth weights on the activation interval; their parameter derivatives have the same bound. The proved moment map in \texttt{core\_exact.md}, using one further parameter derivative, then gives
\[
p_s-p_{s,r}=O(ye_a),\qquad E_r/E=1+O(ye_a).
\]
The actual shear vector is exactly
\[
s_{\rm shear}=\kappa(p_{1,r},p_{2,r}E_r/E).
\]
Because \(\kappa\) cancels before forming its slope, there is no inverse power of \(\kappa_0\) in
\[
t_s=(p_{2,r}/p_{1,r})(E_r/E),\quad
v_s=\kappa v_r+O(ye_a),\quad P_c=v_r+O(ye_a),\quad J_c=O(ye_a).
\]
Consequently \(P_c-v_s=e_av_r+O(ye_a)\ge ce_a\) for small fixed \(t_1\), and \(P_c>2+c'\). The ratio
\((v_s-2)_+J_c^2/(P_c-v_s)^2\) is at most a fixed constant times \(y^2\), hence below 2 for small \(t_1\). The exact cone equivalence proves admissibility wherever \(v_s>2\) and the relaxed condition everywhere else. Since \(v_s(0)=v_r(0)>2+c\), continuity on the compact parameter interval supplies a first collar of positive width with \(v_s>2+c/2\).

The stronger endpoint assertion requires more than an \(O(ye_a)\) estimate. Here is the exact flat-integral calculation. For \(c>0\) and smooth \(b\), substitute \(u=y/\sqrt{1+y^2v}\) to obtain
\[
\int_0^y e^{-c/u^2}u^{-j}b(u,\eta)\,du
=\tfrac12e^{-c/y^2}y^{3-j}
\int_0^\infty e^{-cv}(1+y^2v)^{(j-3)/2}
b\left(\frac y{\sqrt{1+y^2v}},\eta\right)\,dv.
\]
For each fixed derivative the final integrand is bounded by \(e^{-cv}\) times a polynomial in \(v\), uniformly on the compact parameter set; denominators \(1+y^2v\) are at least one. Differentiation under the integral proves a smooth coefficient through \(y=0\), with limiting value \(b(0,\eta)/c\). This proves the needed result from Appendix A.9 directly.

Near zero, the specific step gives \(e_a=e^{-t_1^2/y^2}g_a(y)\), where
\[
g_a(y)=\frac{1-\kappa_0}{e^{-t_1^2/y^2}+e^{-1/(1-y/t_1)^2}}
\]
extends smoothly with \(g_a(0)=(1-\kappa_0)e>0\). Taking \(j=0\) in the flat-integral formula proves
\[
\log\phi-\log\phi_r,\ U-U_r\ \in e_ay^3C^\infty.
\]
This class is preserved by smooth compositions, products with smooth functions, and division by nonvanishing fields: for example \((e^{e_ay^3b}-1)/(e_ay^3)=b\int_0^1e^{ve_ay^3b}\,dv\) is smooth. Moment and pressure differences acquire a second integration and belong to \(e_ay^6C^\infty\). Substitution into the exact moment formulas therefore shows
\[
T_0=e_aB_0(y,\eta),\qquad
B_0\in C^\infty,\qquad
B_0(0,\eta)=F(X_0,\eta)p_{s,r}(X_0,\eta)\ne0.
\]
At the endpoint the latter vector is a positive multiple of the limiting shear. Its component along \((1,t_s)\) is positive, and that along \((-t_s,1)\) is zero. Both cone-direction margins consequently stay positive on a smaller collar. Every derivative of the exponential adds only finitely many powers of \(y^{-1}\), so for every fixed mixed derivative
\[
|T_0|\ge c e^{-t_1^2/y^2},\qquad
|\partial^\alpha T_0|\le C_\alpha e^{-t_1^2/y^2}y^{-N_\alpha}.
\]
Since \(X_0>0\), converting \(y\)-derivatives to \(X\)-derivatives changes only fixed smooth coefficients. Thus the vector direction is smooth through the zero-stress edge, not merely convergent along individual paths.

\subsection{3. Continuation to 110 and preservation of parameter analyticity}

After activation retain the same shear prescription with \(\kappa=\kappa_0\) through the reference cutoff and then through \(X=100\). The natural reference slopes are nonzero only in its initial interval of length \(2t_1\). Integrating the difference on that interval gives \(O(t_1)\). Beyond it, the reference slopes vanish, and the new slopes have size at most \(\kappa_0\) times the bounded functions \(p_{1,r}/2\), \(Xn_{s,r}/2\). Integration on the fixed interval gives \(O(\kappa_0)\). Thus field values, logarithms and the required parameter derivatives differ from reference by \(O(t_1+\kappa_0)\). Their integrated coefficients obey the same estimate with finite constants depending on the already fixed \(C,\Lambda\).

Choose \(\kappa_0\) small enough that \(\kappa_0V_{\max}<1/4\) and that its comparison errors use at most one quarter of the minimum \(v_r-2\). Then choose \(t_1\) small for the other error quarters and for the activation inequalities. The same exact slope formula yields \(P_c>2+c/2\), \(v_s<1\). The relaxed cone therefore holds regardless of \(J_c\). This order works for the entire family of reference widths, avoiding a circular choice of \(\kappa_0\) from a width that itself depends on \(\kappa_0\).

Starting at \(X=100\), multiply the axial prescription by a smooth \(\beta\) from one to zero on a short logarithmic interval. The integrated coefficients change by arbitrarily small errors when the interval is short, while
\[
P_c=p_{1,r}+\beta p_{2,r}^2/p_{1,r}+o(1)>2,
\quad v_s<1,
\]
using \(p_{1,r}>3\) there. Then hold \(D_XU=0\) and interpolate \(a\) convexly from \(\kappa_0p_{1,r}\) to 0.8, after arranging \(\kappa_0\sup p_{1,r}<0.8\). During this second interval \(v_s=a\le0.8\), \(P_c=p_1>2\) by continuity and its strictly positive initial margin. Both intervals can fit strictly between 100 and 110.

On the final interval, \(a=0.8\), \(l=0.6\), \(D_XU=0\). The unchanged leading gradient and the already small field changes give \(S_q>1\): its constant source contribution is \(0.6(-W_*)>1.68\), while the large \(\Lambda L\chi\) term absorbs the previously identified weighted errors. At a hypothetical first downward crossing \(p_1=2\),
\[
D_Xp_1=XS_q/L-0.6p_1>X-1.2>0,
\]
a contradiction. Hence \(p_1>2\) persists through \(X_i=110\).

On a small first collar, all reference coefficients arise from the analytic axis solution, its derivatives and forward integrals. The only cutoffs depend on \(y\), not on \(\eta\). Choose a compact complex neighborhood on which the natural \(\Phi\), \(L\), and the relevant denominators do not vanish. Integration in the real radial variable preserves holomorphy with uniform bounds on that neighborhood. Exponentiating the prescribed logarithm preserves nonvanishing. Radial differentiation changes the cutoff bounds but does not change the holomorphic parameter domain. A smaller fixed \(t_c\in(0,t_1)\), with \(4e^{t_c}<4.1\), therefore gives the common parameter domain for \(F,U,V_0/X,\Pi\) and every fixed radial derivative throughout \([0,X_0e^{t_c}]\). Subsequent changes have later support, so their forward primitives leave this rectangle exactly unchanged.

\subsection{4. Transition length fixed before amplitude}

Let \(\ell_i=\log(CE(X_i,\eta))\), \(G_i=U(X_i,\eta)\), and \(L_0=\log(X_i/X_0)\). Write \(M_k,N_k\) for the reference bounds on \(p_{1,r},n_{s,r}\), which are independent of subsequent \(C\). Before the final interpolation the logarithmic \(\phi\) slope is \(-\kappa p_{1,r}/2\); during it the slope is a convex interpolation with \(-0.4\); after it the slope is exactly \(-0.4\). Hence
\[
\|\log\phi(X_i)\|_{C^k_\eta}
\le\|\log\phi_*+\log\Phi(4,\cdot)\|_{C^k_\eta}
+\tfrac12(M_k+0.8)L_0.
\]
Adding \(\frac12\log(2X_i)\) proves \(\|\ell_i\|_{C^k_\eta}\le B_k\), where \(B_k\) is independent of subsequent \(C\) and smaller transition widths. The axial change on activation is at most \(X_0e^{\bar t}N_kt_1/2\); the remaining small-shear part contributes at most \(X_iN_k\kappa_0/2\). The natural-axis difference is \(O(\Lambda^{-1})\) with a constant independent of \(\Lambda\). Therefore
\[
\|G_i-U_*\|_{C^k_\eta}
\le C_k^{\rm nat}/\Lambda+C_k^{\rm join}(\Lambda)(t_1+\kappa_0+\omega_{\rm fin}),
\]
where \(\omega_{\rm fin}\) is the sum of the two final logarithmic widths. No inverse cutoff width occurs in these integrals.

With \(f=(1+\eta^2)^{-1}\), choose a finite positive \(T_{\rm sh}\) satisfying exactly the source lower bound
\[
T_{\rm sh}\ge20\|\sigma'\|_\infty(B_0+\|\log f\|_\infty).
\]
For \(y_i=\log(X/X_i)\), prescribe
\[
\log E=-\log C+y_i/10+(1-\sigma(y_i/T_{\rm sh}))\ell_i
+\sigma(y_i/T_{\rm sh})\log f,\qquad U=G_i.
\]
At the left endpoint this matches \(E(X_i)\) and its radial jet, because both original and new profiles have \(D_X\log E=0.1\) there and the step is flat. At the right endpoint it matches \(C^{-1}f(X/X_i)^{1/10}\) with its entire radial jet. Direct differentiation gives
\[
l=0.6+T_{\rm sh}^{-1}\sigma'(y_i/T_{\rm sh})(\log f-\ell_i)
\in[0.55,0.65],\quad a\in[0.7,0.9],\quad b_s=0.
\]
The old parameter gradient satisfies
\(-H_c\ell_i'\ge\Lambda L\chi-K\chi-K_{\sigma_*}\sqrt\chi-K/\Lambda-o(1)\).
The new one is
\[
-H_c(\log f)'=\frac{2\eta H_c}{1+\eta^2}
\ge-Kj_0^2-K/\Lambda-o(1).
\]
For the latter bound, the exact identity
\(\eta H_*=(D+4d)\eta^2+dj_0\eta\)
and completion of the square bound its negative part by \(Kj_0^2\); the \(U-U_*\) error is already controlled. These gradients are combined with nonnegative coefficients summing to one. The old large positive term absorbs its correspondingly weighted errors, and the new negative part is small. Since \(-W\) is close to \(-W_*>2.8\) and \(l\ge0.55\), the constant contribution exceeds 1.54 before the small \(h,j_0,\Lambda^{-1}\) errors. Thus \(S_q>1\). The same crossing argument at \(p_1=2\), now \(l\le0.65\), gives \(D_Xp_1>X-1.3>0\). Since \(b_s=0\), this proves \(P_c=p_1>2\), \(v_s=a<1\), hence the relaxed cone, along the full transition and the following ideal-angular interval.

\subsection{5. The five matching equations, including all scaling factors}

Set
\[
X_R=110(CP_*)^{10},\quad x=X/X_R,\quad
X_{\rm sep}=110e^{T_{\rm sh}},\quad
x_{\rm sep}=\frac{e^{T_{\rm sh}}}{(CP_*)^{10}}.
\]
Here \(X_{\rm sep}\) is fixed before \(C\). The ideal pair is \(U_0=4\eta\), \(E_0=P_*fx^{1/10}\). The exact normalized moments are
\[
\widehat M_0=4\eta x,\quad\widehat I_0=\frac{5\sqrt2}{8}P_*fx^{8/5},
\quad\widehat J_0=4\eta\widehat I_0,
\]
\[
\widehat S_0=16\eta^2x-\frac5{12}P_*^2f^2x^{6/5},
\quad C_{p,0}=\frac52P_*^2f^2x^{1/5}.
\]
These follow by direct integration, retaining the factors
\[
(M,I,J,S,C_p)=(X_R\widehat M,X_R^{3/2}\widehat I,
X_R^{3/2}\widehat J,X_R\widehat S,C_p).
\]

On \([0,X_{\rm sep}]\), fixed parameter derivatives of \(U\) and \(CE\) are bounded independently of later \(C\); near zero \(CE=O(\sqrt X)\). Thus
\[
\|\widehat M\|_{C^k}+\|\widehat S\|_{C^k}\le K_kx_{\rm sep},
\quad\|\widehat I\|_{C^k}+\|\widehat J\|_{C^k}
\le K_kC^{-1}x_{\rm sep}^{3/2},
\quad\|C_p\|_{C^k}\le K_kC^{-2}
\]
at \(x=x_{\rm sep}\). For pressure the proof uses the physical integral \(C^{-2}\int_0^{X_{\rm sep}}(CE)^2/(2X)\,dX\), finite because \(CE=O(\sqrt X)\). All constants are fixed before increasing \(C\). The ideal moments also tend to zero with the displayed exact powers. Their pressure increment is \(O(C^{-2})\), because \(x_{\rm sep}^{1/5}=e^{T_{\rm sh}/5}(CP_*)^{-2}\).

Beyond \(X_{\rm sep}\), the angular profile is exactly ideal, since
\(C^{-1}f(X/X_i)^{1/10}=P_*fx^{1/10}\). Choose \(C\) so \(x_{\rm sep}<e^{-8}\). Restore \(G_i\) to \(4\eta\) by a fixed smooth step on \(-8<\log x<-7\), retaining \(E_0\). The size of this change, and its fixed logarithmic radial and parameter derivatives, is bounded by fixed constants times \(\|G_i-4\eta\|\). The moment errors accumulated to the correction interval are bounded by
\[
K_k\|G_i-4\eta\|_{C^{k+1}_\eta}+o_C(1).
\]
Every relevant radial weight is integrable at zero: their actual exponents are \(0,3/5,1/5\) and \(-4/5\) for the pressure integrand, all greater than \(-1\). No divergent contribution is discarded.

On \(-6<\log x<-5\), choose two fixed ordered nonnegative bumps for \(U\), three for \(E\). At the ideal pair make the invertible row operations \(J\mapsto J-4\eta I\), \(S\mapsto S-8\eta M\), with \(4\eta\) held at its base value. The linearized axial block has weights \(1,fx^{3/5}\); the azimuthal block has weights \(x^{1/2},fx^{1/10},fx^{-9/10}\), retaining fixed nonzero factors \(P_*,\sqrt2\) in the matrix. The exponents in each block are distinct. The Chebyshev determinant proof and quadratic contraction in \texttt{core\_exact.md} therefore solve all five moment equations exactly, with small smooth coefficients. The row operations do not impose that the corrected \(U\) remain constant; the full nonlinear differences still include \(\delta U\delta E\), \((\delta U)^2\), and \((\delta E)^2\).

To verify one tolerance can precede \(C\), put \(R=[e^{-8},e^{-5}]\), \(Q_{\min}=\min Q_{s,0}>0\), \(e_* =\min E_0>0\), and \(w_*=1+\max|N_{s,0}/(E_0Q_{s,0})|\) on \(R\times[-1,1]\). The ideal source is positive directly: for this pair
\[
S_{q,0}=\frac95-h+\frac{16}{5}h\eta^2
+\frac{2\eta^2(D+4d)}{1+\eta^2}>0,
\]
as \(W=-3+8h\eta^2\), \(l=3/5\), \(H_c=(D+4d)\eta\). Its regular weighted primitive gives \(Q_{s,0}>0\). The compact minima thus exist without importing the positivity of a separate abbreviated formula. The exact moment map has no \(X_R\) in \(Q_s,N_s\), so a fixed small tolerance in the fields and normalized moments through one extra parameter derivative ensures
\[
Q_s\ge Q_{\min}/2,\quad E\ge e_*/2,\quad |a-0.8|\le0.1,
\quad|N_s/(EQ_s)|\le w_*,\quad |b_s|\le0.1/(1+w_*).
\]
Then, retaining the source numbers,
\[
v_s\le0.9+0.01/0.7<1,
\quad G:=Q_s-b_sN_s/(aE)\ge(1-1/7)Q_s\ge3Q_{\min}/7,
\quad P_c=X_RxG/L.
\]
Choosing \(X_R\) large now gives \(P_c>2\) without shrinking the tolerance. First choose \(j_0\), then \(\Lambda\), so \(j_0+C_k^{\rm nat}/\Lambda\) fits the tolerance. Next fix the \(B_k\) and \(T_{\rm sh}\), then enlarge \(C\) for the vanishing moment contributions. Finally decrease \(\kappa_0,t_1,\omega_{\rm fin}\) for the remaining \(G_i\) error. These choices prove the exact order
\[
(M_d,T_d,P_*,\lambda,h)\to\text{moment tolerance},j_0
\to\delta_*,\sigma_*,\Lambda\to(B_k),T_{\rm sh}
\to C,X_R\to\kappa_0,t_1,\text{final widths}.
\]
Only finitely many parameter derivative orders need to be small. Smooth implicit differentiation supplies all remaining fixed orders after the choices. At \(X_h=X_Re^{-5}\), all five moments and both fields agree with the ideal pair, with the same axis pressure \(\Pi_0\), so the exact exterior-preservation map from \texttt{core\_exact.md} applies.


\section{Exterior stress, exact flat factors, and the order-one edge estimate}\label{pa:sec:component-edge}
\subsection{Convergent moment subtraction and backward stress formulas}
Apply this calculation to the actual regular joined field constructed above.
The five-row join has proved all the required total moments. In particular, $M(\infty)=J(\infty)=S(\infty)=0$ and the compensated
subtracted angular moment gives
\[
I(X)=\frac{XH_{\rm pow}(X)}{1-h}-\int_X^\infty(H-H_{\rm pow})dx,
\qquad S(X)=\tfrac12\int_X^\infty E^2dx.
\]
The exact pressure is the backward integral already established. These
identities, substituted into the moment formulas, determine the same forward
stress everywhere in the tail.

To prove the backward physical stress representation rather than assume its
boundary constants vanish, put $K_{\rm pow}=c_\infty(r^2/2)^{-A}$.
For $r\to\infty$, uniformly on bounded $\tau\ge0$ intervals, the heat
expansion gives
\[
K-K_{\rm pow}=O_h(\tau r^{-3-2h}),\quad K_t=O_h(r^{-3-2h}),
\quad r^2K_r-rK=O_h(r^{-2h}).
\]
Thus the weighted difference and its time derivative are dominated at infinity
by $Cr^{-1-2h}$, whose integral from $R_*$ is $R_*^{-2h}/(2h)$.
At zero the subtracted pure power has weight $r^{1-2h}$, integrable for
$h<1$, while the regular field is integrable by axis smoothness.
The substitution $r=\sqrt qR$, $dX=R\,dR$ proves exactly
\[
\int_0^\infty r^2(u_\theta-K_{\rm pow})dr
=q^{3/2-A}\int_0^\infty(H-H_{\rm pow})dX=0.
\]
Time differentiation is justified by the displayed integrable majorant on
compact physical parameter neighborhoods, proving \eqref{pa:eq:submoment} and
the zero integrated time term.
The angular transport flux has zero axial integral because it is
$q^{3/2-2A}J(\infty)=0$; its radial boundary vanishes by regularity at
zero and $U=M=V_0=0$ in the exterior. The radial viscous boundary is zero
by the $r^{-2h}$ estimate above.

For the axial component $\int ru_zdr=q^{1-A}M(\infty)=0$.
Fubini with the absolutely convergent positive pressure integrand yields
\[
\int_0^\infty rp(r)dr
=-\int_0^\infty\frac{u_\theta(r')^2}{r'}
\left(\int_0^{r'}r\,dr\right)dr'
=-\tfrac12\int_0^\infty r'u_\theta(r')^2dr'.
\]
The pressure has decay $O(r^{-2-4h})$, hence $r^2p\to0$, and its contribution
to the integrated axial flux is exactly $q^{1-2A}S(\infty)=0$.
All axial radial transport and viscosity boundaries vanish as well.
Therefore both leading residuals have zero weighted totals, so the stresses
defined with a minus sign by integration from the axis are exactly
\begin{equation}\label{pa:eq:backward}\tag{PA.BACK}
T_\theta(r)=r^{-2}\int_r^\infty r'^2\paResidual_\theta^{(0)}(r')dr',\qquad
T_z(r)=r^{-1}\int_r^\infty r'\paResidual_z^{(0)}(r')dr'.
\end{equation}
For $X\ge X_b$ the field is $(u_r,u_\theta,u_z)=(0,K,0)$ and pressure
is its $z$ independent backward radial integral. The heat equation makes both
leading residuals zero there, proving zero exterior stress.

\subsection{An exact flat integral map}
For $c>0$, integer $j\ge0$, and a smooth $b(u,\eta)$ on a closed compact
parameter rectangle, substitute $u=\delta(1+\delta^2v)^{-1/2}$.
Then $du=-\tfrac12\delta^3(1+\delta^2v)^{-3/2}dv$ and
$c/u^2=c/\delta^2+cv$. Hence
\begin{equation}\label{pa:eq:flat}\tag{PA.FLAT}
\int_0^\delta e^{-c/u^2}u^{-j}b(u,\eta)du
=\tfrac12e^{-c/\delta^2}\delta^{3-j}B(\delta,\eta),
\end{equation}
where the coefficient is explicitly
\[
B(\delta,\eta)=\int_0^\infty e^{-cv}(1+\delta^2v)^{(j-3)/2}
b(\delta/\sqrt{1+\delta^2v},\eta)dv.
\]
Every fixed derivative of the integrand is bounded by $e^{-cv}$ times a
polynomial in $v$, uniformly in the compact parameters and $0\le\delta\le\delta_0$.
This follows directly from differentiating $\delta(1+\delta^2v)^{-1/2}$;
its denominators are at least one and all its numerators are polynomials in
$\delta,v$. Dominated differentiation proves smoothness of $B$ with all
fixed derivatives bounded, and $B(0,\eta)=b(0,\eta)/c$.
This proves the exact map in [NS, Lemma A.9], including the factor $1/2$.

\subsection{Positive angular stress and the limiting direction}
For terminal $y\ge.5$, the heat cutoff is one, so $u_\theta=Kf$, $f=f_o(y)$,
and $u_r=u_z=0$. At fixed $z,t$, $y_r=2/r$, while
$y_t=(qL)^{-1}$, $y_z=-2\eta/(q^DL)$. The leading angular residual,
which omits axial viscosity, is exactly
\[
\paResidual_\theta^{(0)}=Kf'/(qL)-K(f_{rr}+r^{-1}f_r)-2K_rf_r.
\]
Differentiate the physical pressure $p=-\int_r^\infty K(r')^2f(r')^2dr'/r'$.
Since $K_z=0$ and $dr'/r'=dy'/2$, this gives the sign and factor
\[
p_z(r)=\frac{2\eta}{q^DL}\int_y^3K(y')^2f(y')f'(y')dy'.
\]
Inserting the angular residual into \eqref{pa:eq:backward} and integrating its
viscous terms by parts gives
\begin{equation}\label{pa:eq:positive}\tag{PA.POS}
T_\theta=Kf_r+\frac1{r^2}\int_r^\infty(r'K-r'^2K_{r'})f_{r'}dr'
+\frac1{qLr^2}\int_r^\infty r'^2Kf'\,dr'.
\end{equation}
To check the integration sign, the term $-\int r'^2Kf_{r'r'}dr'$ contributes
\[
r^2Kf_r+\int(2r'K+r'^2K_{r'})f_{r'}dr'.
\]
The two remaining
viscous terms leave exactly $r'K-r'^2K_{r'}$ in this integral.
Every term in \eqref{pa:eq:positive} is nonnegative since $K>0$, $K_r<0$,
$f'\ge0$. The remaining interval has logarithmic length at most $2.5$;
radii and heat factors there have uniformly bounded ratios. Since
$\int_y^3f'\,dv=\rho_o\psi_o(y)$, its last term implies
\[
T_\theta\ge c\frac{rK}{qL}\rho_o\psi_o(y)>0\quad(y<3).
\]
The pressure formula and its next radial integral similarly give
\[
|p_z|\le Cq^{-D}K^2\rho_o\psi_o,
\quad |T_z|\le Crq^{-D}K^2\rho_o\psi_o,
\quad |T_z/T_\theta|\le Cq^{1-D}K=C E_{\rm pow}\paHeat(2d/X).
\]
The final equality uses $1-D=A$. The $h^6$ amplitude reduction makes this
ratio small. Furthermore, since $\partial_yZ=-Z$,
\[
a=2+2h+2Z\paHeat'/\paHeat-2f'/f.
\]
For sufficiently large $X_R$, $-Z\paHeat'/\paHeat<h/4$ uniformly in $\eta$;
also $f'/f<h/4$ by construction. Hence $2+h<a\le2+2h$ and $b_s=0$.
The exact stress identities are
$P_c-a=T_{0,\theta}/F>0$, $J_c=T_{0,z}/F$ with $F=E/\sqrt{2X}>0$.
Thus the cone is exactly $(a-2)(T_z/T_\theta)^2<2$, with uniform strict
slack throughout the terminal interval. This also covers the plateau $.5\le y\le1$:
the positive backward time integral persists there even though $f'(y)=0$.

Near $\delta=3-y=0$, the specified step has the exact factorization
\[
\psi_o=e^{-4/\delta^2}g(\delta),\qquad
g(\delta)=\frac1{e^{-(1-\delta/2)^{-2}}+e^{-4/\delta^2}},\qquad g(0)=e,
\]
\[
f'=\rho_oe^{-4/\delta^2}\delta^{-3}[8g(\delta)+\delta^3g'(\delta)].
\]
The normalized boundary term in \eqref{pa:eq:positive} is
\[
q^{A+1/2}Kf_r=\frac{2E_{\rm pow}(X)\paHeat(2d/X)}{\sqrt{2X}}f'.
\]
The two remaining angular integrals contain $f'$, so \eqref{pa:eq:flat} with
$c=4,j=3$ makes each $e^{-4/\delta^2}$ times a smooth coefficient.
Relative to the boundary factor $e^{-4/\delta^2}\delta^{-3}$ they vanish
to at least order $\delta^3$. Therefore
\[
T_{0,\theta}=e^{-4/\delta^2}\delta^{-3}b_\theta(\delta,\eta),\qquad
b_\theta(0,\eta)=\frac{16\rho_oE_{\rm pow}(X_b)}{\sqrt{2X_b}}
\paHeat(2d/X_b)g(0)>0.
\]
The heat factor has a positive minimum on $[0,2/X_b]$, so $b_\theta$
has a positive lower bound on a sufficiently short closed collar uniformly
for $|\eta|\le1$.
The pressure integral contains one $f'$ factor; the same substitution gives
$p_z=q^{-D-2A}e^{-4/\delta^2}$ times a smooth coefficient.
Its further radial integration for $T_z$ uses \eqref{pa:eq:flat} with $j=0$,
giving
\[
T_{0,z}=e^{-4/\delta^2}\delta^3b_z(\delta,\eta).
\]
The powers of $q$ cancel exactly:
$q^{A+1/2}q^{1/2}q^{-D-2A}=q^{1-D-A}=1$.
The remaining coefficients depend only on $X,\eta,\paHeat(2d/X),L^{-1}$,
with their already proved closed-parameter smoothness. Consequently
\[
T_{0,z}/T_{0,\theta}=\delta^6b_z/b_\theta\longrightarrow0
\]
smoothly, and the stress direction extends to $(1,0)$ at the outer edge.
Every fixed profile derivative costs only a finite inverse power of $\delta$,
and $|T_0|\ge ce^{-4/\delta^2}\delta^{-3}$ there.
This proves (A.48)--(A.51) including the precise leading coefficient and
the independence from the physical scale $q$.


\section{Admissible shear loop, radial modulation, and moment restoration}\label{pa:sec:component-loop}\label{pa:admissible-shear-loop-radial-modulation-and-moment-restoration}

\subsection{Scalar cone and its four gaps}\label{pa:scalar-cone-and-its-four-gaps}

At \(a>0\), set \[
t_s=-b_s/a,\quad v_s=a(1+t_s^2),\quad
P_c=p_{s,1}+t_sp_{s,2},\quad
J_c=p_{s,2}-t_sp_{s,1}.
\] The relaxed condition is \(P_c>2\), \(v_s<U(P_c,J_c)\), with \[
U(P,J)=P+\frac{J^2}{4}
-|J|\sqrt{\frac{P-2}{2}+\frac{J^2}{16}}.
\] The admissible condition also requires \(v_s>2\). The quadratic expression \[
G(v)=2(P-v)^2-(v-2)J^2
\] has roots \(U(P,J)\) and the same expression with a plus before the square root. For \(P>2\), \(U>2\): writing \(w=P-2>0\), the inequality \[
\left(w+\frac{J^2}{4}\right)^2
-J^2\left(\frac w2+\frac{J^2}{16}\right)=w^2>0
\] proves \(w+J^2/4>|J|\sqrt{w/2+J^2/16}\). Also \(U\le P\), because the square-root term is at least \(J^2/4\), with equality only for \(J=0\). When \(J=0\), \(U=P>2\). Therefore for \(v>2\), admissibility is exactly \[
P-v>0,\qquad 2(P-v)^2-(v-2)J^2>0.
\] The four continuous smooth gaps on \(a>0\) are \(a\), \(v_s-2\), \(P_c-v_s\), and \(2(P_c-v_s)^2-(v_s-2)J_c^2\). Positivity on a compact parameter set gives one positive lower margin for all four. The root \(U\) need only be continuous; no differentiation of its \(|J|\) term at zero is used.

\subsection{A smooth loop whose average is exactly the original shear}\label{pa:a-smooth-loop-whose-average-is-exactly-the-original-shear}

Use the actual joined data on the compact rectangle \(I\times[-1,1]\) chosen in Section~\ref{pa:sec:assembly-loop}; the preceding construction proves its relaxed condition and admissibility on both radial boundary collars. The positive minimum of \(P_c(t_s)-2\) permits a constant \[
0<d_0<\tfrac12\min(P_c(t_s)-2).
\] Here \(P_c(t)=p_{s,1}+p_{s,2}t\) and \(J_c(t)=p_{s,2}-p_{s,1}t\) for a variable slope \(t\). With normalized \(\theta'\)-average on \([0,2\pi]\), define \[
M_e(z)=\frac1{2\pi}\int_0^{2\pi}e^{z\sin\theta'}\,d\theta',
\quad
t(\theta';\mu)=t_s+
\frac{d_0}{p}\left(\frac{e^{\mu p\sin\theta'}}{M_e(\mu p)}-1\right),
\quad p=p_{s,2}.
\] For \(p=0\) the quotient is defined by continuity. Its numerator is smooth in \(p,\mu,\theta'\), zero at \(p=0\), and the identity \[
\frac{G(p)}p=\int_0^1G'(up)\,du
\] proves joint smoothness, including all parameters and derivatives. Since \(M_e(0)=1,M_e'(0)=0\), the value there is \(t=t_s+d_0\mu\sin\theta'\).

The definitions give \[
\langle t\rangle_{\theta'}=t_s,\qquad
P_c(t)=P_c(t_s)+d_0\left(\frac{e^{\mu p\sin\theta'}}{M_e(\mu p)}-1\right)
\ge P_c(t_s)-d_0>2.
\] The variance is \[
V(\mu,p)=\langle(t-t_s)^2\rangle
=\frac{d_0^2}{p^2}
\left(\frac{M_e(2\mu p)}{M_e(\mu p)^2}-1\right),
\] with value \(d_0^2\mu^2/2\) at \(p=0\). To show strict increase, let \(g(z)=\log M_e(z)\). Differentiating under the finite integral gives \[
g''(z)=
\langle(\sin\theta'-\langle\sin\theta'\rangle_z)^2\rangle_z>0,
\] where the tilted probability density is \(e^{z\sin\theta'}/(2\pi M_e(z))\). It is strictly positive on the circle and \(\sin\theta'\) is not constant, so the variance is strictly positive. For \(p\ne0,\mu>0\), \[
\partial_\mu(g(2\mu p)-2g(\mu p))
=2p(g'(2\mu p)-g'(\mu p))>0.
\] For \(p>0\) both factors on the right are positive; for \(p<0\) both are negative. Thus \(V\) strictly increases in \(\mu>0\), including its separately proved formula at \(p=0\).

Here is the required growth bound, without an unproved asymptotic invocation. For \(z\ge1\), write \(\theta'=\pi/2+x\) near the maximum. On a fixed small interval \(|x|\le a<\pi\), there are \(c_1,c_2>0\) with \(1-c_2x^2\le\cos x\le1-c_1x^2\). A lower bound integrates over \(|x|\le z^{-1/2}\min(a,1)\), giving \(ce^zz^{-1/2}\). An upper bound on the same interval uses the Gaussian integral and gives \(Ce^zz^{-1/2}\). On its complement \(\sin\theta'\le1-c_a\), so the contribution is at most \(e^{(1-c_a)z}\), also bounded by \(Ce^zz^{-1/2}\). For \(z<0\), shift \(\theta'\) by \(\pi\) to obtain \(M_e(z)=M_e(-z)\). Hence \[
c\,e^{|z|}/\sqrt{|z|}
\le M_e(z)\le C\,e^{|z|}/\sqrt{|z|}\quad (|z|\ge1).
\] Taking the ratio yields positive constants bounding \(M_e(2z)/M_e(z)^2\) above and below by multiples of \(\sqrt{|z|}\). Therefore \(V(\mu,p)\to\infty\) for each fixed \(p\ne0\); for \(p=0\) its explicit quadratic formula proves the same.

The compact range of \(p\) admits a uniform finite \(\mu_{\max}\) such that \(V(\mu_{\max},p)>3/\min a\). To prove uniformity, for each \(p_0\) select a finite \(\mu_{p_0}\) giving the strict bound there; continuity gives a neighborhood in \(p\); a finite subcover gives finitely many such \(\mu\), and their maximum works by monotonicity. This step is valid at \(p=0\), so no divergent factor \(p^{-2}\) is being used near zero.

All slopes \(t(\theta';\mu)\) for \(0\le\mu\le\mu_{\max}\) form a compact smooth family with \(P_c(t)>2\). Continuity of \(U\) and the strict inequality \(U(P,J)>2\) established above give a positive minimum of \(U-2\). Choose \(0<\delta_L<1\) smaller than that minimum, and smaller than the positive minimum of \(v_s-2\) on chosen smaller radial boundary collars.

Choose a smooth \(0\le\zeta_L\le1\) as a function of \(v_s\), equal to one for \(v_s\le2+\delta_L/8\) and zero for \(v_s\ge2+\delta_L/4\). Put \[
v_*=2+\delta_L/2,\qquad
\rho=\zeta_L(v_s)^2(v_*-v_s),\qquad
v=v_s+\rho,
\] and define \(\rho=0\) where \(\zeta_L=0\). On its possible nonzero support, \(v_*-v_s\ge\delta_L/4\), so \[
\sqrt\rho=\zeta_L(v_s)\sqrt{v_*-v_s}
\] extends smoothly by zero. Since \(v_s>0\), \(0\le\rho<v_*<3\).

There is a unique \(0\le\mu<\mu_{\max}\) solving \(V(\mu,p)=\rho/a\), by the preceding strict monotonicity and endpoint bounds. Its smoothness at \(\rho=0\) needs a distinct verification. Expanding the finite integral in its convergent power series gives \(M_e(z)=1+z^2/4+O(z^4)\), and therefore \[
V(\mu,p)=\tfrac12d_0^2\mu^2+O(\mu^4p^2)
\] uniformly for bounded \(p\). Thus the signed square root is \(\mu(d_0/\sqrt2+O(\mu^2p^2))\), a smooth odd function of \(\mu\) near zero with positive derivative \(d_0/\sqrt2\). Applying the implicit-function theorem to this function and \(\sqrt{\rho/a}\), already proved smooth, gives smooth \(\mu\) there. For \(\rho>0\), strict positivity of \(\partial_\mu V\) gives the same conclusion. On the overlap the two local inverses agree by uniqueness.

The three cutoff regimes show \(v>2\) everywhere. Where \(\zeta_L\ne0\), \(v\le v_*<2+\delta_L<U(P_c(t),J_c(t))\). Where \(\zeta_L=0\), \(\mu=0,t=t_s,v=v_s>2\), and the original relaxed condition gives \(v<U\). Thus each proposed state lies in the admissible cone.

Define a lifted phase coordinate by \[
\phi(0)=0,\qquad
\frac{d\phi}{d\theta'}=\frac{a(1+t^2)}{2\pi v},
\quad
(a_L,-b_L)=\frac{v(1,t)}{1+t^2}.
\] Since \(\langle t\rangle=t_s\) and \(\langle(t-t_s)^2\rangle=V\), \[
a\langle1+t^2\rangle=a(1+t_s^2)+aV=v_s+\rho=v.
\] Therefore \(\phi(\theta'+2\pi)=\phi(\theta')+1\). Its derivative has a uniform positive lower bound on the compact parameter set; it defines an orientation-preserving circle diffeomorphism with a smooth parameter-dependent inverse. Changing variables proves both original mean components exactly: \[
\int_0^1a_L\,d\phi
=\frac a{2\pi}\int_0^{2\pi}1\,d\theta'=a,
\quad
\int_0^1(-b_L)\,d\phi
=\frac a{2\pi}\int_0^{2\pi}t\,d\theta'=at_s=-b_s.
\] Also \(-b_L/a_L=t\), \(a_L(1+t^2)=v\), so the verified scalar cone inequalities are the desired vector cone gaps. At boundary collars \(\zeta_L=0\), and the reparametrized shear is exactly the original constant-in-\(\phi\) shear. Compactness gives one \(\kappa_L>0\) for all four cone gaps. This completes the complete loop construction, including the mean, boundary, and smoothness claims.


\subsection{Exact signed square root at zero variance}
The following identity also proves the preceding local smoothness without an
asymptotic remainder. Define
\[
 R_e(z)=M_e(2z)/M_e(z)^2,\qquad
 H_e(z)=(R_e(z)-1)/z^2\ (z\ne0),\qquad H_e(0)=1/2.
\]
The convergent power series of the finite circle integral gives an even real
analytic $M_e$, with $M_e(z)=1+z^2/4+O(z^4)$. Its positive real values
give an even real analytic $R_e$, and multiplication of the power series gives
$R_e(z)=1+z^2/2+O(z^4)$. Thus the quotient $H_e$ extends real analytically
and evenly through zero. For $z\ne0$, Cauchy--Schwarz gives the strict inequality
\[
 \langle e^{z\sin\theta'}\rangle^2
 <\langle e^{2z\sin\theta'}\rangle\langle1\rangle=M_e(2z).
\]
Equality would require $e^{z\sin\theta'}$ to be constant almost everywhere;
continuity and the positive uniform measure make that impossible for $z\ne0$.
Consequently $H_e(z)>0$ on the entire real line. Exactly,
\[
 V(\mu,p)=d_0^2\mu^2H_e(\mu p),\qquad
 F_e(\mu,p)=d_0\mu\sqrt{H_e(\mu p)}.
\]
The latter function is a globally smooth signed square root. At zero its
$\mu$ derivative is $d_0/\sqrt2$ for every $p$. For $\mu>0$ it is positive,
and $2F_e\partial_\mu F_e=\partial_\mu V>0$, so the derivative stays
positive. The equation $F_e(\mu,p)=\sqrt\rho/\sqrt a$ has a smooth
implicit solution, including points where $p=\rho=0$, because
$\sqrt\rho=\zeta_L(v_s)\sqrt{v_*-v_s}$ is the smooth zero extension
already proved above. The strict variance monotonicity identifies all local
solutions with the same globally defined $\mu$. This proves every required
parameter derivative at zero variance as well as at $p=0$.

\subsection{Exact insertion into the unchanged original radial coordinate}\label{pa:exact-insertion-into-the-unchanged-original-radial-coordinate}

Let \(\mathcal A,\mathcal B\) be the zero-mean periodic antiderivatives \[
\partial_\phi\mathcal A=-\tfrac12(a_L-a),\qquad
\partial_\phi\mathcal B=\tfrac12E(b_L-b_s).
\] They exist because the right sides have zero \(\phi\)-mean. One explicit construction for any zero-mean \(g\) is \[
\mathcal I g(\phi)=\int_0^\phi g(w)\,dw-
\int_0^1\int_0^s g(w)\,dw\,ds.
\] It is periodic, zero mean, smooth in all parameters, and has derivative \(g\). It vanishes identically when \(g\equiv0\). Therefore these antiderivatives vanish on the boundary collars and extend smoothly by zero outside \(I\).

For an integer \(N>0\), keep the source's exact definitions \[
E_N=E\exp(\mathcal A(X,\eta,N\log X)/N),\qquad
U_N=U+\mathcal B(X,\eta,N\log X)/N.
\] The modified \(E_N\) is positive without approximation. The chain rule in the logarithmic coordinate is \(X\partial_X=D_X+N\partial_\phi\) after substitution. Hence \[
\begin{aligned}
a_N
&=1-2X\partial_X\log E_N\\
&=a-2D_X\mathcal A/N-2\partial_\phi\mathcal A
=a_L-2D_X\mathcal A/N,
\end{aligned}
\] \[
\begin{aligned}
b_N
&=\frac{2X\partial_XU_N}{E_N}\\
&=e^{-\mathcal A/N}
\left(\frac{2D_XU+2\partial_\phi\mathcal B}{E}
+\frac{2D_X\mathcal B}{NE}\right)\\
&=e^{-\mathcal A/N}
\left(b_L+\frac{2D_X\mathcal B}{NE}\right).
\end{aligned}
\] The axial exponential denominator is retained; replacing it by one would destroy this exact shear formula.

On the fixed compact interval from \(X_-\) through the first reserved patch, \(X,E,H=\sqrt{2X}E,L\) have positive lower bounds. The phase \(N\log X\) is independent of \(\eta\), so every fixed number \(m\) of \(\eta\)-derivatives of \(E_N-E,U_N-U,a_N-a_L,b_N-b_L\) is bounded by \(C_m/N\). For the exponential difference, use \[
e^{\mathcal A/N}-1
=\frac{\mathcal A}N\int_0^1e^{s\mathcal A/N}\,ds
\] and then differentiate under its fixed finite integral; all factors of \(\mathcal A\) and its parameter derivatives are uniformly bounded on the periodic compact domain. For radial derivatives, \(r\ge1\) differentiations produce at most \(N^r\) from the phase and retain the initial \(N^{-1}\), giving \(C_{r,m}N^{r-1}\). Thus no small-radial-derivative assertion is used for the modulation.

\subsection{All five moments and the dependence of the integrated stress}\label{pa:all-five-moments-and-the-dependence-of-the-integrated-stress}

Retain exactly \[
M=\int_0^XU\,dx,\quad I=\int_0^X\sqrt{2x}E\,dx,\quad
J=\int_0^XU\sqrt{2x}E\,dx,
\] \[
S=\int_0^X(U^2-E^2/2)\,dx,\qquad
C_p=\int_0^X E^2/(2x)\,dx,\qquad \Pi=\Pi_0+C_p.
\] On the support of the differences, all scalar weights are bounded, since it is a compact interval with positive left endpoint. For instance \[
U_N^2-U^2=(U_N-U)(U_N+U),\qquad
E_N^2-E^2=(E_N-E)(E_N+E).
\] These identities, the product rule, and integration of the uniform \(C_m/N\) profile bounds prove every fixed \(\eta\)-derivative of the full prefix moment difference is \(O_m(N^{-1})\), uniformly in the endpoint \(X\). The pressure datum is the same, so \(\Pi_N-\Pi=C_{p,N}-C_p\) has exactly that bound.

The formulas whose continuity matters are \[
W=1-\frac{2D\eta M+(1-\eta^2)M_\eta}{X},
\] \[
Q_s=-W+
\frac{(1-h)I-D\eta I_\eta-(1-\eta^2)J_\eta+2(h-D)\eta J}{XH},
\] \[
N_s=-WU+
\frac{D(M-\eta M_\eta)+4h\eta S-(1-\eta^2)S_\eta}{X}
+4A\eta\Pi-(1-\eta^2)\Pi_\eta,
\] \[
p_s=(XQ_s/L,\ XN_s/(LE)).
\] All denominators are bounded away from zero on the fixed interval. Differences of reciprocals are exactly \[
E_N^{-1}-E^{-1}=-(E_N-E)/(E_NE),\qquad
H_N^{-1}-H^{-1}=-(H_N-H)/(H_NH).
\] Applying the product rule to the displayed formulas proves the \(m\)-th parameter derivative bound for \(p_{s,N}-p_s\) from the profile and moment bounds through order \(m+1\), with no radial derivative of the differences. This proves the precise assertion needed for the cone stability.

The compact loop has positive four-gap margin. Its values and the original \(p_s\) range in a compact subset of \(a>0\); take a fixed small neighborhood still in \(a>0\), on which the four-gap map has a finite Lipschitz constant. The just-proved \(O(N^{-1})\) errors in \(a,b,p_s\) preserve each gap for sufficiently large finite \(N\). Between \(I\) and the first correction patch, \(E_N=E,U_N=U\), so the shears themselves equal the originals; the \(p_s\) error remains \(O(N^{-1})\). The original strict admissible margins on that compact intervening interval are preserved by the same argument.

\subsection{Nonlinear moment correction with an exact endpoint map}\label{pa:nonlinear-moment-correction-with-an-exact-endpoint-map}

On the first reserved patch, the original profile is \(U=0,E=K(\eta)X^{-1/2-\lambda}\), where fixed \(\lambda>0\) and positive \(K\) have been selected before \(N\). The earlier modulation is zero on this patch. Add two smooth compactly supported \(U\)-bumps and three \(E\)-bumps with five coefficient functions \(c(\eta)\). Each block's bumps are nonnegative, nonzero, and supported on ordered disjoint intervals.

Differentiate the five original moment functionals at zero correction. In row order \((M,J;I,S,C_p)\), the two \(U\)-columns have weights \(1,H=\sqrt{2X}E\), and the three \(E\)-columns have weights \(\sqrt{2X},-E,E/X\). The off-diagonal blocks are zero because \(U=0\) before the correction. Up to the displayed nonzero factors, the exponent lists are \[
(0,-\lambda),\qquad
(1/2,-1/2-\lambda,-3/2-\lambda).
\] All powers within each block are distinct for the fixed \(\lambda>0\). To establish invertibility directly, a nonzero linear combination of \(m\) distinct powers has at most \(m-1\) distinct positive zeros: divide by the least power, use Rolle's theorem on \(m\) putative zeros, and apply induction to the derivative, which has \(m-1\) distinct powers and nonzero coefficients wherever the original nonconstant coefficients were nonzero. Thus the evaluation determinant \(\det[x_j^{\alpha_i}]\) never vanishes for ordered \(0<x_1<\cdots<x_m\). Its sign is constant on that connected set. Multilinearity and Fubini give \[
\det\left[\int x^{\alpha_i}\beta_j(x)\,dx\right]
=\int\det[x_j^{\alpha_i}]\prod_j\beta_j(x_j)\,dx_1\cdots dx_m\ne0.
\] The integrand has a constant nonzero sign on a set of positive measure. This proves each block's rank and hence the full moment Jacobian's rank. Smooth dependence on \(\eta\), compactness of \([-1,1]\), and the derivative identity for the inverse give uniform finite bounds at every fixed parameter order. No uniform estimate as \(\lambda\to0\) is claimed or needed.

The exact moment increment map is \(B(\eta)c+Q_\eta(c,c)\), because \(M,I\) are linear, \(J\) is bilinear in \(U,E\), and \(S,C_p\) are quadratic. Let \(d_N(\eta)\) be the negative moment discrepancy accumulated before the patch; it obeys \(\|d_N\|_{C^m}\le C_m/N\). On the \(C^0\) ball of radius \(r=2\beta_0\|d_N\|_{C^0}\), where \(\beta_0=\|B^{-1}\|_{C^0\to C^0}\), the map \[
\mathcal T(c)=B^{-1}(d_N-Q(c,c))
\] maps the ball to itself and has Lipschitz constant at most \(2\beta_0\kappa_0r\le1/2\), provided \(8\beta_0^2\kappa_0\|d_N\|_{C^0}\le1\). Thus it has a unique fixed point selected by iteration from zero. The derivative \(B+D_cQ(c,c)\) is invertible by the same estimate, giving smooth parameter dependence including one-sided endpoint derivatives. In each preselected finite \(C^m\) norm the identical contraction proof works after increasing \(N\), with the relevant multiplication and bilinear norms, and yields \(C_m/N\) for the same branch.

The assertion that every other fixed derivative is \(O_m(N^{-1})\) can also be justified without demanding simultaneous smallness of infinitely many norms. Differentiate the pointwise equation \(m\) times. The highest coefficient derivative is multiplied by the uniformly invertible \(B+D_cQ(c,c)\). All other terms contain either a derivative of \(d_N\), \(O_m(N^{-1})\), or a lower derivative of \(c\); induction gives \(O_m(N^{-1})\) for them, with the fixed data derivatives included in the constants. The same inverse bound therefore proves the estimate for the highest derivative. Only the finite orders needed for positivity and the cone gaps determine the actual choice of \(N\).

The bump shapes and widths are fixed, so their radial derivatives multiply these small coefficients by fixed constants. The correction is therefore small in the radial derivatives entering \(a,b_s\), in the first parameter derivatives entering \(p_s\), and in all prefix moments. The compact input cone margin on the correction patch persists. Positivity of \(E\) there also persists, after one further finite lower bound on \(N\).

At the right endpoint \(X_{\rm rep}\), the five corrected moments equal the original moments exactly as functions of \(\eta\), while the corrected profiles equal the original profiles for every \(X\ge X_{\rm rep}\). Their derivatives as functions of \(\eta\) also agree. Each difference moment has zero derivative in \(X\) there, by the original moment integrands, so it stays zero at every later radius. The common axis pressure datum gives identical \(\Pi\). The displayed formulas for \(Q_s,N_s,p_s\) and the original incompressibility formula \[
V_0=\frac X L\left(2\eta U-2D\eta M/X-(1-\eta^2)\partial_\eta(M/X)\right)
\] then prove exact equality of these fields at every later radius. This is the exact morphism from the five-moment endpoint data to exterior pressure, velocity, and stress; it proves preservation of the terminal compensation, both unused reserved correction patches, and the entire exterior. The inner collar is unchanged because both modifications start later and pressure was integrated from the same axis datum.

\subsection{Edge direction and the common weight}\label{pa:edge-direction-and-the-common-weight}

Appendix C leaves both original collars unchanged. The input edge facts it uses are: at the inner edge, with \(y_a=\log(X/X_a)\), the stress is \[
T_0=e^{-t_1^2/y_a^2}g_a(y_a)B_0(X,\eta),\qquad
g_a(0)>0,\quad B_0(X_a,\eta)=F p_{s,r}\ne0;
\] its limiting direction is a positive multiple of the limiting shear \((a,-b_s)=a(1,t_s)\). At the outer edge, with \(y_b=\log(X_b/X)\), the component factorization has common exponential \(e^{-4/y_b^2}\), angular prefactor \(y_b^{-3}b_\theta\) with \(b_\theta(0,\eta)>0\), and the direction \[
n=\frac{(b_\theta,y_b^6b_z)}
{\sqrt{b_\theta^2+y_b^{12}b_z^2}}.
\] The inner continuation and outer endpoint calculations above prove these smooth factors and all fixed-order derivative bounds for the constructed field. Their exact use is as follows.

The positive factors cancel in \(n=T_0/|T_0|\), giving a smooth direction through both endpoints. At the inner edge \(n_z-t_sn_\theta=0\) and \(n_\theta+t_sn_z>0\). At the outer edge \(n=(1,0)\), \(t_s=0\), giving the same two conclusions. In the interior, \[
n_\theta+t_sn_z=\frac{P_c-v_s}{|p_s-(a,-b_s)|}>0,\qquad
n_z-t_sn_\theta=\frac{J_c}{|p_s-(a,-b_s)|}.
\] The admissible cone makes \[
2-(v_s-2)\frac{(n_z-t_sn_\theta)^2}{(n_\theta+t_sn_z)^2}>0.
\] Both the denominator and this gap extend continuously to the closed annulus, where the gap has value two at both radial edges. Their positive compact minima permit one \(0<\kappa<2\) with \[
n_\theta+t_sn_z\ge\kappa,\quad
(v_s-2)(n_z-t_sn_\theta)^2
\le(2-\kappa)(n_\theta+t_sn_z)^2.
\] The positive lower bounds on \(F,a,v_s-2\) in the interior and the original endpoint bounds likewise extend uniformly by compactness.

The original stress identity \(T_0=F(p_s-(a,-b_s))\) proves nonvanishing throughout the open shell: if it vanished, \(P_c=v_s\), contrary to admissibility. The outside stress remains zero by the exact endpoint restoration and unchanged inner region. Define the original proposed common edge weight \[
\zeta(X)=\exp(-t_1^2/y_a^2-4/y_b^2),\qquad
\delta=\min(1,y_a,y_b),
\] extended by zero outside the open shell. On a sufficiently short inner collar, the second exponential is smooth and bounded above and below by positive constants, while \(g_a,B_0\) have the required nonzero value bounds and fixed derivative bounds. The inner factorization therefore proves \(|T_0|\ge c\zeta\) and \(|\partial^IT_0|\le C_I\zeta\delta^{-m_I}\). On a short outer collar the first exponential is similarly a smooth positive factor; \(y_b^{-3}b_\theta\) has a lower positive bound and its derivatives have only finite inverse powers of \(y_b\). The supplied outer factor derivative bounds give the same conclusion. On the remaining compact interior both \(\zeta\) and \(|T_0|\) have positive minima and every fixed derivative is bounded, so the same pair of inequalities follows after increasing constants. These three regions cover the annulus. The exponential derivative calculation in the closed theorem proof proves smooth zero extension of all stress derivatives at both boundaries.

All remaining completion claims in C.3 are exact preservation statements: the axis Cartesian formulas remain unchanged on the analytic rectangle; the pressure at infinity and exterior moment values remain unchanged by the proved five-moment map; the exterior heat function is unchanged pointwise; and on \(X\ge X_v=X_{\rm end}\) the original \(U=M=0\) gives \(V_0=0\) directly by its formula. The third and fourth reserved intervals remain the exact original \(U=0,E=c_{\rm patch}(1+\eta^2)^{-1}X^{-1/2-\lambda}\) intervals, with their original order. Fixing this one finite \(N\) before the physical scale \(q\), all dyadic bands, and all later correction stages proves the claimed independence of the resulting derivative constants from those later parameters.


% --- included proof body: stage_inputs_body.tex ---
\newtheorem{siproposition}{Proposition}[section]

\section{Exact coordinate, evaluation, and covering maps}
\label{si:maps-audit:coordinates}

This section audits the coordinate and auxiliary maps on their original
domains.  For the physical and profile coordinate maps retain the base
construction's range $0<h<1/2$.  Its constants are
\[
 A=\tfrac12+h,\quad D=\tfrac12-h,
 \quad d=1-\eta^2,\quad L=1-2h\eta^2.
\]
For the auxiliary and oscillatory construction below retain its narrower
original range $0<h<1/100$.
The letter $t$ denotes physical time.  The symbol $\mathsf t$ below is an
extended differential operator, not a second time coordinate.  No statement
at the excluded point $(z,t)=(0,1)$ is inferred from the maps.

\subsection{The global q map and its boundary}
Define
\[
 \Psi:(0,\infty)\times(-1,1)\longrightarrow
 \mathbb R\times(-\infty,1),\qquad
 \Psi(q,\eta)=(z,t)=(q^D\eta,1-q(1-\eta^2)).
\]
This is a smooth bijection with smooth inverse.  Indeed, for given $z$ and
$\tau=1-t>0$, put $q_0=|z|^{1/D}$.  On $q>q_0$ the function
$f_z(q)=q-z^2q^{2h}$ has derivative
\[
 f_z'(q)=1-2hz^2q^{2h-1}\ge1-2h>0.
\]
It tends to zero as $q\downarrow q_0$ (also when $z=0$), and tends to
infinity as $q\to\infty$, since $2h<1$.  The intermediate value theorem
and strict monotonicity give exactly one root $f_z(q)=\tau$.  Setting
$\eta=zq^{-D}$ gives $|\eta|<1$ and both displayed equations for $\Psi$.
Conversely, any preimage must give that root and that value of $\eta$.
The exact derivative matrix and determinant are
\[
 \frac{\partial(z,t)}{\partial(q,\eta)}
 =\begin{pmatrix}Dq^{D-1}\eta&q^D\\-d&2q\eta\end{pmatrix},
 \qquad \det\frac{\partial(z,t)}{\partial(q,\eta)}=q^DL>0.
\]
Inverting this matrix proves smoothness of the inverse by the inverse
function theorem, and proves, with the other physical coordinate fixed,
\[
 q_t=-\frac1L,\qquad \eta_t=\frac{D\eta}{qL},\qquad
 q_z=\frac{2\eta q^{1-D}}L,\qquad \eta_z=\frac{d}{q^DL}.
\]
The same map on $(0,\infty)\times[-1,1]$ is a bijection onto
\[
 \{(z,t):t<1\}\ \cup\ \{(z,1):z\ne0\}.
\]
For $t=1$ the root is $q=|z|^{1/D}>0$ and
$\eta=\operatorname{sign}z$.  Since the determinant remains positive
there, the inverse has a smooth extension to a neighborhood of every such
boundary point and hence all one-sided derivatives.  There is no point
of this domain above $(0,1)$.

Retain $s=\tfrac12r^2$ and $X=s/q$ for $r\ge0$.  For $r>0$,
\[
 \partial_rX=\frac rq,\qquad
 X_t=\frac X{qL},\qquad X_z=-\frac{2\eta X}{q^DL}.
\]
The inverse radial relation is $r=\sqrt{2qX}$, so the original physical
radius is preserved.  For every $b\in\mathbb R$ and smooth $f(X,\eta)$,
the ordinary product and chain rules, with no change of $q,X,\eta$, give
\[
 \partial_t(q^bf)=q^{b-1}T_bf,\qquad
 T_bf=L^{-1}(-bf+D\eta f_\eta+Xf_X),
\]
\[
 \partial_z(q^bf)=q^{b-D}Z_bf,\qquad
 Z_bf=L^{-1}(2b\eta f+df_\eta-2\eta Xf_X).
\]
Consequently $\partial_{zz}(q^bf)=q^{b-2D}Z_{b-D}Z_bf$ with that exact
order and shifted subscript.  These formulas contain no division by
$d$ and remain valid at $\eta=\pm1$ through the boundary extension.

\subsection{Chart-to-profile maps with fixed Q}
For a fixed integer $\ell$, set
\[
 Q=2^{-\ell},\quad\varepsilon=Q^h,\quad S_*=\ell^2,
 \qquad R=\frac r{\sqrt Q},\quad Z=\frac z{Q^D},\quad
 T=\frac{1-t}Q.
\]
Use $s_Q=q/Q$ in the next formulas; it is different from
$s=\tfrac12r^2$.  The exact transition maps are
\[
 Z=s_Q^D\eta,\quad T=s_Q(1-\eta^2),\quad
 X=\frac{R^2}{2s_Q},\quad
 R_{\rm profile}=\frac r{\sqrt q}=\frac R{\sqrt{s_Q}}.
\]
In particular $T=s_Q-Z^2s_Q^{2h}$, whose derivative in $s_Q$ at fixed
$Z$ is $L$.  Direct inversion, or the physical formulas above, gives
\[
 (s_Q)_T=\frac1L,\quad(s_Q)_Z=\frac{2\eta s_Q^{1-D}}L,
 \quad\eta_T=-\frac{D\eta}{s_QL},\quad
 \eta_Z=\frac{d}{s_Q^DL},
\]
\[
 X_R=\frac R{s_Q},\qquad X_T=-\frac X{s_QL},\qquad
 X_Z=-\frac{2\eta X}{s_Q^DL}.
\]
The inverse map from $(X,s_Q,\eta)$ to $(R,Z,T)$ is explicitly
\[
 (R,Z,T)=(\sqrt{2Xs_Q},s_Q^D\eta,s_Qd).
\]
For any fixed $0<X_a<X_b<\infty$, all fixed derivatives of this map
and its inverse are bounded on
\[
 X_a\le X\le X_b,\qquad \tfrac12\le s_Q\le2,
 \qquad -1\le\eta\le1.
\]
To justify the assertion at every derivative order, the first derivative
formulas have denominators consisting only of positive powers of $s_Q$,
$L$, and $R$ for the inverse radial relation.  On this set these are
bounded below by $1/2$, $1-2h$, and $\sqrt{X_a}$, respectively.
Differentiating any finite product of these factors again produces a
finite sum with the same types of denominators and bounded numerator
factors.  Induction on the derivative order proves the assertion.  This
establishes the exact chart/profile derivative morphism including all
edge parameters; it never identifies $R$ with $R_{\rm profile}$.

\subsection{Physical evaluation as a differential algebra map}
From this subsection onward take $0<h<1/100$, as in the original
oscillatory construction.
Let $\mathbb T^2=\mathbb R^2/\mathbb Z^2$ with Haar measure of mass one,
let $\theta\in\mathbb R/(2\pi\mathbb Z)$, and let $r>0$, $z\in\mathbb R$,
$t<1$ or a retained boundary point with $q>0$.  Set
\[
 J_g=\begin{pmatrix}3&1\\1&5\end{pmatrix},\quad b_g=\sqrt2-1,
 \quad v_r=(1,-b_g),\quad v_t=(b_g,1),
\]
\[
 \Lambda_g=4-\sqrt2,\quad T_g=4+\sqrt2,
 \quad\rho_g=\frac{\log\Lambda_g}{\log T_g},\quad
 \kappa_s=10^{-5},\quad d_r=2((1+h)\rho_g-h\kappa_s).
\]
Matrix multiplication gives $J_gv_r=\Lambda_gv_r$ and
$J_gv_t=T_gv_t$.  Also
$v_r\cdot v_t=0$, $|v_r|^2=|v_t|^2=1+b_g^2$, and
$\det(v_r,v_t)=1+b_g^2>0$.  Since $\Lambda_g>2$ and $T_g<6$,
$\rho_g>\log2/\log6>1/3>\kappa_s$, so $d_r>0$.

For a smooth extended scalar $F(r,\theta,z,t,Y)$ define
\[
 \gamma(r,t)=v_rr^{d_r}+v_tt\pmod{\mathbb Z^2},\qquad
 (\mathcal EF)(r,\theta,z,t)=F(r,\theta,z,t,\gamma(r,t)).
\]
This map preserves addition, multiplication, and constants.  Its
intertwining identities are exactly
\[
 \partial_t\mathcal E=\mathcal E\mathsf t,\qquad
 \partial_r\mathcal E=\mathcal E\mathsf r,\qquad
 \partial_z\mathcal E=\mathcal E\partial_z,\qquad
 \partial_\theta\mathcal E=\mathcal E\partial_\theta,
\]
where
\[
 \mathsf t=\partial_t+v_t\cdot\partial_Y,\qquad
 \mathsf r=\partial_r+d_rr^{d_r-1}v_r\cdot\partial_Y.
\]
For example, differentiating $F(r,\theta,z,t,\gamma)$ in $r$ produces
$F_r+(d_rr^{d_r-1}v_r)\cdot F_Y$, evaluated at $Y=\gamma$; the other
three identities follow by the same explicit chain rule.  All four
unscaled coordinate derivations on the right commute: every possible
coefficient derivative in a commutator is zero, except for radial
differentiation of $d_rr^{d_r-1}$, which never appears in a commutator
with a different operator.  Iterating the identities proves them for
every ordered composition of the coordinate derivatives.

The use of $r^{d_r}$ is restricted to $r>0$.  If an extended field is
supported in $X_a\le r^2/(2q)\le X_b$ with $X_a>0$, its evaluation
vanishes in a neighborhood of $r=0$ at every fixed $(z,t)$ with $q>0$.
More precisely, on a neighborhood where $q\ge q_*>0$ it vanishes for
$r<\sqrt{2X_aq_*}$.  Its zero extension in Cartesian coordinates is
therefore smooth near that axis neighborhood without differentiating
$r^{d_r}$ at zero.

\subsection{Integer coverings, differential intertwining, and Haar measure}
For each integer $i\ge0$, define
\[
 \pi_i:\mathbb T^2\longrightarrow\mathbb T^2,
 \qquad \pi_i([Y])=[J_g^iY],\qquad P_if=f\circ\pi_i.
\]
Integer entries make the map well-defined.  A lift $y$ of any target
point has preimage $[J_g^{-i}y]$, so the map is surjective.  For $i=1$
the integer column swap followed by the row operation
$\mathrm{row}_2\mapsto\mathrm{row}_2-5\mathrm{row}_1$ and the column
operation $\mathrm{col}_2\mapsto\mathrm{col}_2-3\mathrm{col}_1$ reduces
$J_g$ to $\operatorname{diag}(1,-14)$.  All these operations are invertible
integer torus maps.  The diagonal map has fourteen kernel points, so
$|\ker\pi_1|=14$.  The map $\pi_1:\ker\pi_{i+1}\to\ker\pi_i$ is
surjective: lift each point of $\ker\pi_i$ through the surjective
$\pi_1$.  Every fiber is a coset of $\ker\pi_1$ and has fourteen
points.  Induction gives $|\ker\pi_i|=14^i$.

Write $L_i=v_r\cdot\partial_{Y_i}$ and
$N_i=v_t\cdot\partial_{Y_i}$.  The chain rule and the two eigenvector
identities prove
\[
 (v_r\cdot\partial_Y)P_if=\Lambda_g^iP_iL_if,\qquad
 (v_t\cdot\partial_Y)P_if=T_g^iP_iN_if.
\]
For $n\in\mathbb Z^2$, the character $e^{2\pi in\cdot Y_i}$ pulls
back to the character of frequency $(J_g^i)^Tn$.  Since $J_g^i$ is
invertible over $\mathbb R$, this frequency is zero precisely for
$n=0$.  Integrating the Fourier series of a smooth function, which
converges absolutely, thus gives
\[
 \int_{\mathbb T^2}P_if\,dY=\int_{\mathbb T^2}f\,dY_i.
\]
This proves preservation of the normalized Haar measure exactly.

The image of $P_i$ consists precisely of the smooth functions invariant
under translations by $\ker\pi_i$.  Necessity follows from
$\pi_i(Y+k)=\pi_i(Y)$ for $k$ in that kernel.  For sufficiency define
$f(y)=F(Y)$ for any preimage $Y$ of $y$; invariance makes it independent
of the preimage, and local inverses of the nonsingular linear cover
prove smoothness.  In Fourier coordinates the same condition is that
the support is contained in $(J_g^i)^T\mathbb Z^2$.  To prove the
nontrivial direction of this last equivalence, invariance implies that
any nonzero Fourier coefficient at $m$ satisfies
$m\cdot J_g^{-i}k\in\mathbb Z$ for every $k\in\mathbb Z^2$, since
$[J_g^{-i}k]\in\ker\pi_i$.  Thus $(J_g^{-i})^Tm\in\mathbb Z^2$,
which is exactly the asserted sublattice condition.

For completeness, the directional inverses used with these maps are
unambiguous on the zero-mean subspace.  If $n\ne0$ then
\[
 (v_r\cdot n)((n_1+n_2)+\sqrt2n_2)=(n_1+n_2)^2-2n_2^2,
\]
\[
 (v_t\cdot n)((n_2-n_1)-\sqrt2n_1)=(n_2-n_1)^2-2n_1^2.
\]
Each right side is a nonzero integer: a zero would force the indicated
integer ratio to equal $\sqrt2$, or force both coordinates to vanish.
The conjugate factors have absolute value at most
$C_g|n|$, $C_g=\sqrt{4+2\sqrt2}$, by the Euclidean Cauchy--Schwarz
inequality.  Hence $|v_r\cdot n|,|v_t\cdot n|\ge(C_g|n|)^{-1}$.
The inverse multiplier is exactly $(2\pi i(v\cdot n))^{-1}$ at
$n\ne0$ and zero at $n=0$.  It grows at most linearly in $|n|$,
so preserves rapid Fourier decay and gives
$\|(v\cdot\partial_Y)^{-1}f\|_{H^s}\le C_s\|f\|_{H^{s+1}}$.
Uniqueness follows because any smooth zero-mean function in the kernel
has every Fourier coefficient zero.  The exact intertwining for the
inverses is
\[
 L_{\rm abs}^{-1}P_i=\Lambda_g^{-i}P_iL_i^{-1},\qquad
 N_{\rm abs}^{-1}P_i=T_g^{-i}P_iN_i^{-1}.
\]
Translation, averaging, and these inverses preserve the Fourier
sublattice, hence preserve descent.  Integration in $r$ at fixed
$z,t,Y$ also preserves descent whenever the integrand and integration
limits are defined and the latter are independent of $Y$; deck
translation commutes pointwise with that integral.

\subsection{The original scaled operators and their threshold}
For $\ell\ge4$, $2^\ell\ge\ell^2$: it holds at four, and its induction
step follows from $2\ell^2\ge(\ell+1)^2$ for $\ell\ge4$.
Therefore $Q^{-1-h}/S_*>1$ and the integer
\[
 i(\ell)=\left\lfloor\log_{T_g}\frac{Q^{-1-h}}{S_*}\right\rfloor
\]
is nonnegative.  The floor inequality is
\[
 T_g^{-1}Q^{-1-h}/S_*<T_g^{i(\ell)}\le Q^{-1-h}/S_*.
\]
Combining the fixed-$Q$ coordinate map with the cover intertwining
gives, on the band representative,
\[
 \mathsf t_*=Q^{1+h}\mathsf t=-\varepsilon\partial_T+c_iN_i,
 \qquad c_i=T_g^iQ^{1+h},
\]
\[
 D_r=\sqrt Q\,\mathsf r
 =\partial_R+M_id_rR^{d_r-1}L_i,
 \qquad M_i=\Lambda_g^iQ^{d_r/2},
\]
\[
 D_z=\sqrt Q\partial_z=\varepsilon\partial_Z,
 \qquad D_\theta=R^{-1}\partial_\theta.
\]
Here the sign of the time term is fixed by $t=1-QT$.  The constants obey
\[
 T_g^{-1}S_*^{-1}<c_i\le S_*^{-1},\qquad
 T_g^{-\rho_g}\varepsilon^{-\kappa_s}S_*^{-\rho_g}
 <M_i\le\varepsilon^{-\kappa_s}S_*^{-\rho_g},
\]
because $\Lambda_g^i=(T_g^i)^{\rho_g}$ and
$d_r/2=(1+h)\rho_g-h\kappa_s$.  Thus each application of $D_r$ to
a coefficient incurs at most the retained factor
$\varepsilon^{-\kappa_s}$, with the power $S_*^{-\rho_g}$ still
present in its fast part.  Every derivative of $d_rR^{d_r-1}$ is
bounded on the fixed positive chart annulus.  Applying $D_z$ incurs
the factor $\varepsilon$, and the two parts of $\mathsf t_*$ are
exactly the displayed slow and fast parts.

The operators $D_r,D_z,\mathsf t_*,\partial_\theta$ commute pairwise.
The orthonormal angular derivative has instead the exact frame
commutator
\[
 [D_r,D_\theta]=-R^{-1}D_\theta.
\]
Indeed, $D_r(R^{-1})=-R^{-2}$ and $[D_r,\partial_\theta]=0$.
Accordingly, the cylindrical formulas retain their frame terms.  For
an extended vector in component order $(r,\theta,z)$ they are
\[
 \operatorname{div}_*u=(D_r+R^{-1})u_r+D_\theta u_\theta+D_zu_z,
\]
\[
 \operatorname{curl}_*u=
 \bigl(D_\theta u_z-D_zu_\theta,\;
       D_zu_r-D_ru_z,\;
       (D_r+R^{-1})u_\theta-D_\theta u_r\bigr).
\]
These follow by inserting the evaluation intertwining into the usual
cylindrical formulas, or directly differentiating the basis vectors
$\partial_\theta e_r=e_\theta$, $\partial_\theta e_\theta=-e_r$.
The scalar spatial Laplacian is
$D_r^2+R^{-1}D_r+D_z^2+R^{-2}\partial_\theta^2$.
For physical fields $u_{\rm phys}=Q^{-A}\mathcal Eu_*$ and
$p_{\rm phys}=Q^{-2A}\mathcal Ep_*$, multiplying the physical equation
by $Q^{2A+1/2}$ gives the original time coefficient $Q^{1+h}$,
first-spatial-derivative coefficient $\sqrt Q$, and viscosity-one
coefficient $Q^h=\varepsilon$.  These are the identities
$2A+1/2-A=1+h$ and $2A+1/2-A-1=h$; none of $D_r,D_z,\mathsf t_*$
is replaced by its ordinary slow-coordinate part.

\subsection{Common torus, local band range, and exact pulse path}
Fix a physical point with $q=q_0>0$ and restrict to a neighborhood $U$
such that $q(U)\subset(q_0/2,2q_0)$.  For bands with
$q/Q_\ell\in[1/2,2]$ somewhere on $U$, one has
$Q_\ell\in(q_0/4,4q_0)$.  Hence only finitely many occur and the
largest difference of their indices is at most four.  This restriction
on $U$ is necessary for a uniform bound; finiteness of an arbitrary
band set alone would not provide it.

Choose an integer $\ell_0>4$ below all active indices.  On
$x\ge\ell_0$ put
\[
 f(x)=\frac{(1+h)x\log2-2\log x}{\log T_g}.
\]
For $|\ell-\ell'|\le4$, the mean value theorem and
$|\lfloor a\rfloor-\lfloor b\rfloor|\le |a-b|+1$ give
\[
 |i(\ell)-i(\ell')|
 \le1+\frac{4(1+h)\log2+8/\ell_0}{\log T_g}.
\]
Take the ceiling of this expression as $\Delta_{\max}$.  Define
$i_0=\min_{\ell\in\mathcal L(U)}i(\ell)$ and the common auxiliary
coordinate $H=J_g^{i_0}Y$.  For a particular band set $\Delta=i-i_0$.
Its representative is pulled back through
\[
 P_{\Delta}:f(Y_i)\longmapsto f(J_g^\Delta H),
 \qquad 0\le\Delta\le\Delta_{\max}.
\]
The equality $P_{i_0}P_\Delta=P_i$ proves compatibility with the
absolute torus.  On an overlap of two neighborhoods, choose the lesser
of their two $i_0$ values and pull both representatives to that torus;
their pullbacks agree whenever they represent the same absolute field.
Injectivity of pullback by a surjective covering proves this agreement
as an equality of functions, not merely of averages.  Because the maps
are independent of the slow coordinates, slow derivatives preserve it.
For the same fixed $Q$, the exact common coefficients are
\[
 c_{i_0}=T_g^{-\Delta}c_i,\qquad
 M_{i_0}=\Lambda_g^{-\Delta}M_i.
\]
Thus the common and band formulas give the same absolute scaled
operators.  Factors coming from $J_g^\Delta$ are uniformly bounded
because $\Delta$ lies in the displayed finite set.

Consider an injective rectangle in the band torus with lift
\[
 Y_i=c_\gamma+k+\xi_gv_r+\eta_gv_t,
 \qquad |\xi_g|<r_0,\quad |\eta_g|<r_0.
\]
Its full preimage consists exactly of the copies
\[
 H=J_g^{-\Delta}(c_\gamma+k+\xi_gv_r+\eta_gv_t)
       \pmod{\mathbb Z^2},\qquad
 [k]\in\mathbb Z^2/J_g^\Delta\mathbb Z^2.
\]
Each copy maps diffeomorphically onto the rectangle.  Any preimage
gives one such integer $k$ by lifting its image.  If two copies meet,
the corresponding band points agree; injectivity of the band
rectangle forces their $\xi_g,\eta_g$ coordinates to agree, and
then the two $k$ values differ by $J_g^\Delta\mathbb Z^2$.
This proves completeness and pairwise disjointness of the copies.
Their number is $14^\Delta$ by the proved kernel count.

Set $v=(\eta_g+r_0)/c_i$ and $L_s=2r_0/c_i$.  With
$\lambda_t(w)=v_t\cdot w/(1+b_g^2)$ one has
$\lambda_t(v_t)=1$, $\lambda_t(v_r)=0$, and
$\lambda_tJ_g^\Delta=T_g^\Delta\lambda_t$ because $J_g$ is symmetric
and $v_t$ is its eigenvector.  In the band rectangle,
$L_i\eta_g=N_i\xi_g=0$ and $L_i\xi_g=N_i\eta_g=1$.
Consequently
\[
 D_rv=D_zv=0,\qquad\mathsf t_*v=1.
\]
For a fixed lift of a common rectangle, the path at fixed transverse
and slow coordinates from pulse value $v(H)$ to $w\in[0,L_s]$ is
exactly
\[
 H_w=H+T_g^{-\Delta}c_i(w-v(H))v_t.
\]
Its image in the band rectangle changes by
$c_i(w-v(H))v_t$, and hence has pulse value $w$ and unchanged
$\xi_g$.  Its derivatives are
\[
 D_Hv=\frac{T_g^\Delta}{c_i}\lambda_t,
 \qquad D_HH_w=I-v_t\lambda_t,
 \qquad \partial_wH_w=T_g^{-\Delta}c_iv_t.
\]
All derivatives of order at least two in $H,w$ vanish.  The first
displayed derivative is $O(S_*)$ and $D_HH_w$ is bounded independently
of the band.  Changing a lift or moving by a deck translation just
changes the integer $k$; the formula for the path changes by the same
deck translation.  Thus an initial-value problem whose coefficients
and initial datum are pullbacks of the same band data has equal
solutions on deck-related copies: subtract their pulled-back equations
to obtain the same initial-value problem with zero difference datum,
whose uniqueness gives zero difference.  This is the exact descent
statement for that solution.  The path itself takes values at
individual auxiliary points and performs no averaging.

\subsection{Separation for all products used in covariance assembly}
The slow labels are boxes $\beta=(\ell,a)$, with the two signs retained
as distinct labels $\gamma=(\ell,a,\sigma)$.  The squared partitions
are $\sum_\ell\chi_\ell(q)^2=1$ and, in each band chart,
$\sum_a\chi_{\ell,a}^2=1$.  Define
$\eta_\beta=\chi_\ell(q)\chi_{\ell,a}$.  Then
\[
 \sum_{\beta=(\ell,a)}\eta_\beta^2=1
\]
on their active domain; the signs do not duplicate this sum.  Products
of these nonnegative summands give the identity directly.  Choose a
representative in the intersection of the active shell and every
selected grid support.  The mesh in the three slow chart coordinates
is $S_*^{-3}=\ell^{-6}$, and each grid support and its fixed enlarged
version are contained in a box with side length at most a fixed
constant times that mesh.

Join two distinct labels whenever their enlarged slow supports meet
and their band indices differ by at most four, and join the two signs
of each box.  This countable graph has a finite upper bound $d_g$ on
its vertex degree.  Here is the geometric count: for bands differing
by at most four, each chart coordinate changes by a diagonal factor
bounded above and below by fixed powers of $2^4$.  Their mesh ratio
is $(\ell/\ell')^6$, also bounded above and below since
$\ell,\ell'\ge\ell_0>4$ and $|\ell-\ell'|\le4$.  Thus the image of
one enlarged box meets at most a fixed number of grid boxes in each
of three coordinates, and hence at most the product of those three
numbers.  There are at most nine band indices and two signs to count.
These bounds prove the finite degree assertion.  Enumerate the labels
and at each step choose the least unused color among
$1,\ldots,d_g+1$ on its already colored neighbors.  There are at most
$d_g$ such neighbors, so this yields a proper coloring of all labels.

For these finitely many colors choose centers $c_\mu\in\mathbb T^2$
such that
\[
 c_\mu\ne J_g^\Delta c_\nu\pmod{\mathbb Z^2},
 \qquad 0\le\Delta\le\Delta_{\max},
\]
apart from the unavoidable case $\Delta=0,\mu=\nu$.  This choice can
be made with rational coordinates.  To prove existence, regard the
centers as a point in the finite product of tori.  For distinct colors
each forbidden equality is a closed set with empty interior: holding
the other center fixed, arbitrary small changes of $c_\mu$ leave that
equality.  For equal colors and $\Delta>0$, the matrix
$J_g^\Delta-I$ is nonsingular, since its eigenvalues are
$\Lambda_g^\Delta-1,T_g^\Delta-1>0$.  Its torus kernel is finite:
preimages can be represented by integer vectors in the bounded image
of a unit fundamental square, and there are only finitely many such
integer vectors.  Thus this forbidden set also has empty interior.
The complement of each of the finitely many closed forbidden sets is
open and dense.  Repeatedly intersecting any nonempty open set with
these complements shows that their full intersection is nonempty and
open.  Rational tuples are dense, so one lies in that intersection.
For this fixed choice the finitely many forbidden differences have
a positive minimum distance $d_c$ from $\mathbb Z^2$.

Let
\[
 C_J=\max_{0\le\Delta\le\Delta_{\max}}\|J_g^\Delta\|,
 \qquad C_v=|v_r|+|v_t|.
\]
Choose $r_0>0$ so small that
\[
 4r_0C_v<1,\qquad 2r_0C_v(C_J+1)<d_c.
\]
The closed enlarged rectangles with
$|\xi_g|,|\eta_g|\le2r_0$ are injective, because the difference
between any two of their lifts has length at most $4r_0C_v<1$ and
so cannot be a nonzero integer vector.  Suppose the absolute
preimages of two such rectangles for joined labels, at levels
$i$ and $i+\Delta$, have a common point.  Writing the two rectangle
errors as $e_\nu,e_\mu$, each of length at most $2r_0C_v$, gives
\[
 c_\mu-J_g^\Delta c_\nu
   =J_g^\Delta e_\nu-e_\mu\pmod{\mathbb Z^2}.
\]
The left side is one of the forbidden differences: at $\Delta=0$
joined vertices have distinct colors, and at $\Delta>0$ all color
pairs were excluded.  Its distance from the lattice is at least
$d_c$.  The right side has distance at most
$2r_0C_v(C_J+1)<d_c$, a contradiction.  Hence the closed absolute
preimages are disjoint for every pair that can interact at a slow
point.  Passing to a common torus preserves their disjointness: any
intersection there would lift by its surjective map from the absolute
torus to an intersection already excluded.

For actual smooth coefficients supported in these products of slow
supports and auxiliary rectangles, every derivative has support
contained in the original closed support.  Outside that support the
function vanishes on a neighborhood, and therefore all derivatives
vanish there, which proves this support assertion.  Consequently every
product of coefficients or of any of their derivatives from distinct
active labels is identically zero on the extended domain.  Since
$\mathcal E$ preserves products, the same vanishing holds after
physical evaluation.  This proves the cross-label cancellations used
in covariance and differential covariance assembly, including the
two opposite signs of a single box and their support boundaries.

\subsection{Averaging, evaluation, and the preserved covariance target}
Let $\mathcal P F=\int_{\mathbb T^2}F\,dY$ and let
$\mathcal J$ embed an auxiliary-independent function as a constant
function of $Y$.  Then
\[
 \mathcal P\mathcal J=I,\qquad\mathcal E\mathcal J=I,
 \qquad \mathcal E F=\mathcal P F+
                  \mathcal E(F-\mathcal J\mathcal P F).
\]
The cover identities prove agreement of absolute, common, and band
averages.  Integrating a directional derivative of a periodic smooth
function gives zero, so
\[
 \mathcal P\mathsf r F=\partial_r\mathcal PF,\qquad
 \mathcal P\mathsf t F=\partial_t\mathcal PF,
\]
and the analogous $z,\theta$ identities hold.  This gives an exact
differential map to the auxiliary mean.

There is no operation on the evaluated field alone that recovers this
mean for all extended fields.  For any $n\in\mathbb Z^2\setminus\{0\}$,
take
\[
 F(r,\theta,z,t,Y)=e^{2\pi in\cdot Y}
                       -e^{2\pi in\cdot\gamma(r,t)}.
\]
It is smooth on the specified $r>0$ domain, $\mathcal EF=0$, and
$\mathcal PF=-e^{2\pi in\cdot\gamma(r,t)}\ne0$.
Any proposed recovery map would therefore assign to the zero evaluated
field both zero and this nonzero function, which is impossible.  Real
and imaginary parts give real counterexamples.  In particular an
auxiliary covariance identity refers to the extended construction;
it is not an identity obtained by averaging the physically evaluated
field over a missing auxiliary variable.

Finally, for a real $e$, an integrable coefficient with
$f_*=Q^{a_f}f_{\rm phys}$, and corresponding radial integration
intervals $I_r=\sqrt Q I_R$, the substitution $r=\sqrt Q R$ proves
\[
 \int_{I_R}R^e\langle f_*\rangle_Y\,dR
 =Q^{a_f-(e+1)/2}
  \int_{I_r}r^e\langle f_{\rm phys}\rangle_Y\,dr.
\]
This includes both the original power weight and the measure.
Agreement of the three auxiliary averages follows from the proved
cover morphisms.  The identity makes no claim at a nonintegrable
endpoint or at $q=0$.


\section{Phase, constrained evolution, and every fixed derivative}\label{si:phase-constrained-evolution-and-every-fixed-derivative}

This section derives the pulse equations and estimates of source pp.~73--81. The exact auxiliary maps and operators were proved above. The actual-base expansion proved in the initial-state section below gives, on the fixed enlarged annulus, the chart radial bound \(b=O(\varepsilon)\) and the tangential differences \(V-V_0,G-G_0=O(\varepsilon^2)\), at every fixed slow derivative. Here those exact chart data enter the pulse system, whose phase, constrained evolution, and derivative estimates are calculated explicitly.

\subsection{Cone coordinates and the phase gradient}\label{si:cone-coordinates-and-the-phase-gradient}

Write \(F=V/R\), \(g=(RF_R,G_R)\) in component order \((\theta,z)\). At the fixed representative of a slow box, the leading shear is \[
g_0=F_0(-a,b_s),\quad N=g_0/|g_0|,\quad K=(-N_z,N_\theta).
\] Let \(t_s=-b_s/a\) and \(v_s=a(1+t_s^2)\). Since \(a,F_0>0\), \[
N=-\frac{(1,t_s)}{\sqrt{1+t_s^2}},\quad
K=\frac{(t_s,-1)}{\sqrt{1+t_s^2}},\quad
|g_0|=aF_0\sqrt{1+t_s^2}.
\] Retain the source's \[
\lambda_0^2=-2F_0N_\theta(2F_0N_\theta+|g_0|),
\qquad c_0=\lambda_0/(2F_0N_\theta).
\] Substitution gives \[
\lambda_0^2
=2F_0^2\left(a-\frac2{1+t_s^2}\right)
=2aF_0^2(1-2/v_s)>0,\qquad
c_0^2=(v_s-2)/2,\quad c_0<0.
\] The positive lower bounds on \(F_0,a,v_s-2\) and compact upper bounds give positive lower bounds on \(\lambda_0,|c_0|\), and finite upper bounds.

At the same, unfrozen profile point, \(T_0=F(p_s-(a,-b_s))\), \(P_c=p_{s,1}+t_sp_{s,2}\), \(J_c=p_{s,2}-t_sp_{s,1}\). Consequently \[
T_0\cdot N=-\frac{F(P_c-v_s)}{\sqrt{1+t_s^2}},
\qquad
T_0\cdot K=-\frac{FJ_c}{\sqrt{1+t_s^2}}.
\] Multiplication of \(T_0\) by the positive chart factor \((Q/q)^{A+1/2}\) changes neither sign nor quotient. Thus the two stress cone inequalities are exactly \(T_{0,*}\cdot N<0\) and \[
\left|c_0\frac{T_{0,*}\cdot K}{T_{0,*}\cdot N}\right|<1.
\] The weighted direction at a vanishing edge is used, not a division of zero values. Its smooth extension and strict compact margin give a supremum \(\mu<1\). Choosing \(u_*>0\) with \(u_*/\sqrt{1+u_*^2}>\mu\) is possible because that function increases to one. All quantities now frozen at representatives are kept fixed under differentiation. Uniform continuity in the compact chart domain makes the same strict cone margin valid on every enlarged box of diameter \(CS_*^{-3}\), after one lower bound on \(\ell\).

Retain \(k=\lceil\varepsilon^{-1/2}\rceil\), so for \(0<\varepsilon\le1\), \[
1\le\varepsilon k^2\le(1+\sqrt\varepsilon)^2\le4,\qquad
k^{-1}\le\sqrt\varepsilon.
\] Define \[
B_s^2=\frac{\lambda_0}{\varepsilon k^2(1+u_*^2)^{3/2}},
\quad s(v)=\sigma(u_*/2+u_*v/L_s),\quad x_0=\sigma B_su_*/2.
\] The \(B_s\) have fixed positive lower and upper bounds. Set \[
(\widetilde p/R_0,p_z)
=B_s\left(K-\frac{\sigma u_*g_0}{L_s|g_0|^2}\right).
\] Choose \(kp\) as the specified nearest nonzero integer to \(k\widetilde p\). The error \(|p-\widetilde p|\le k^{-1}\) suffices, including when the nearest unrestricted integer would have been zero. The integer is constant with the label. Put \[
\Phi=p\theta+p_zZ/\varepsilon+x_0R-vH_\Phi,\qquad H_\Phi=pF+p_zG.
\] Because \(D_rv=D_zv=0\) and the base is auxiliary-independent, \[
n_\Phi=\nabla_*\Phi
=\bigl(x_0-v(H_\Phi)_R,\ p/R,\ p_z-\varepsilon v(H_\Phi)_Z\bigr).
\] At the representative, before rounding, \((\widetilde p/R_0,p_z)\cdot g_0=-\sigma B_su_*/L_s\). Variation across the slow box changes this by \(O(S_*^{-3})\); the background change contributes \(O(\varepsilon^2)\); rounding contributes \(O(k^{-1})\). Hence \[
|(H_\Phi)_R+\sigma B_su_*/L_s|
\le C(S_*^{-3}+\varepsilon^2+k^{-1}).
\] Multiplication by \(v\le L_s=O(S_*)\) gives the radial comparison with \(B_ss(v)\). The tangential comparison has contributions \(O(S_*^{-1})\) from the deliberate tilt, \(O(S_*^{-3}+\varepsilon^2+k^{-1})\) from freezing and rounding, and \(O(\varepsilon S_*)\) from the axial phase term. A single sufficiently large \(\ell_*\) satisfies \[
S_*^2(\varepsilon+\varepsilon^2+k^{-1})\le1
\] for every \(\ell\ge\ell_*\), since \(S_*=\ell^2\) and \(\varepsilon=2^{-h\ell}\). Therefore \[
|n_\Phi-B_s(s(v),K)|\le C/S_*.
\] Differentiation in \(v\) gives \(n_\Phi'=(-(H_\Phi)_R,0,-\varepsilon(H_\Phi)_Z)\), also \(O(S_*^{-1})\). Each fixed number of slow derivatives of these coefficients is bounded by a polynomial in \(S_*\): the only growing factors in these formulas are \(v\), \(\varepsilon^{-1}\) in the phase itself, and the fixed chart derivatives. The \(\varepsilon^{-1}\) term of \(\Phi\) disappears in \(n_\Phi\); it is not discarded from the full exponential when physical derivatives are later taken.

The exact transport defect follows term by term: \[
\mathsf t_*\Phi=-H_\Phi+\varepsilon v(H_\Phi)_T,\qquad
F\partial_\theta\Phi+GD_z\Phi=H_\Phi-\varepsilon vG(H_\Phi)_Z.
\] Thus \[
E_{\rm ik}:=(\mathsf t_*+bD_r+F\partial_\theta+GD_z)\Phi
=\varepsilon v((H_\Phi)_T-G(H_\Phi)_Z)+bn_{\Phi,r}.
\] The \(b\) term is \(O(\varepsilon)\) at every fixed slow derivative. Auxiliary coefficient derivatives of \(v\) are \(O(S_*)\), while derivatives of the fixed affine path of order two or greater vanish. Leibniz's rule therefore proves \(D^IE_{\rm ik}=O_I(\varepsilon S_*^{b_I})\) for every fixed \(I\), on the same domain. It imposes no smallness condition on the higher derivatives.

\subsection{Exact pressure elimination and plane isomorphism}\label{si:exact-pressure-elimination-and-plane-isomorphism}

Set \[
\mathcal K=
\begin{pmatrix}0&-2F&0\\2F+RF_R&0&0\\G_R&0&0\end{pmatrix},
\quad d=\varepsilon k^2|n_\Phi|^2,\quad
\mathcal A_\Phi=-\mathcal K+
\frac{n_\Phi(n_\Phi^T\mathcal K-(n_\Phi')^T)}{|n_\Phi|^2}.
\] For the original equation \[
t_m'+\mathcal Kt_m+m^2dt_m+ikm n_\Phi\pi_m=-f_m,
\quad n_\Phi\cdot t_m=0,
\] differentiate its constraint to get \(n_\Phi\cdot t_m'=-n_\Phi'\cdot t_m\). Taking its normal component then forces \[
\pi_m=\frac{i}{km}\,
\frac{n_\Phi\cdot\mathcal Kt_m-n_\Phi'\cdot t_m+n_\Phi\cdot f_m}{|n_\Phi|^2}.
\] The factor has sign \(+i/(km)\): multiplying it by \(ikm n_\Phi\) gives the negative of the required normal vector. Substituting gives the projected equation \[
t_m'=\mathcal A_\Phi t_m-m^2dt_m-\operatorname{proj}_{n_\Phi^\perp}f_m.
\] Conversely this projected equation and the displayed pressure imply the original equation provided the constraint holds. For any solution of the projected equation, \[
(n_\Phi\cdot t_m)'=-m^2d(n_\Phi\cdot t_m).
\] An initially zero constraint remains zero, proving equivalence, existence, and uniqueness once the finite-dimensional initial-value equation is solved.

Since \(n_{\Phi,\tan}\) stays close to \(B_sK\), its norm has a uniform positive lower bound. The nonzero \(p/R\) also prevents it from being literally zero. Define \[
K_a=n_{\Phi,\tan}/|n_{\Phi,\tan}|,\quad
N_a=((K_a)_z,-(K_a)_\theta),\quad
s_a=n_{\Phi,r}/|n_{\Phi,\tan}|.
\] Then \(N_a\perp K_a\), \(|N_a|=|K_a|=1\), and \[
U=[e_r-s_aK_a,N_a],\quad
M=\begin{pmatrix}1&1\\c_0\sqrt{1+s^2}&-c_0\sqrt{1+s^2}\end{pmatrix},
\quad B=UM
\] maps \(\mathbb C^2\) isomorphically onto the actual complex plane \(n_\Phi^\perp\). The exact left inverse is \[
B^\ell=M^{-1}\begin{pmatrix}e_r^T\\N_a^T\end{pmatrix}
\] because the last displayed row matrix times \(U\) is \(I_2\). The determinant of \(M\) is \(-2c_0\sqrt{1+s^2}\), bounded away from zero. All factors and their inverses are uniformly bounded in value and polynomially bounded in each fixed coefficient derivative.

For the reference normal \(n_{\rm ref}=B_s(s,K)\), use \(e=e_r-sK\) and \(N\) as an orthogonal plane basis, of lengths \(\sqrt{1+s^2}\) and \(1\). The actual source matrix \(\mathcal K_0\) gives \[
\mathcal K_0e=(2F_0sK_\theta,\ 2F_0e_\theta+g_0),\quad
\mathcal K_0N=(-2F_0N_\theta,0).
\] Consequently \[
e\cdot(-\mathcal K_0e)=sK\cdot g_0=0,\quad
e\cdot(-\mathcal K_0N)=2F_0N_\theta,
\] \[
N\cdot(-\mathcal K_0e)=-(2F_0N_\theta+|g_0|),\quad
N\cdot(-\mathcal K_0N)=0.
\] Orthogonal projection onto \(n_{\rm ref}^\perp\) leaves these scalar products unchanged. Dividing the first coordinate by \(1+s^2\) produces \[
\mathcal D=
\begin{pmatrix}
0&2F_0N_\theta/(1+s^2)\\
-(2F_0N_\theta+|g_0|)&0
\end{pmatrix}.
\] Its two eigenvectors are exactly the columns of \(M\), with eigenvalues \(\lambda=\lambda_0/\sqrt{1+s^2}\) and \(-\lambda\), since \(c_0=\lambda_0/(2F_0N_\theta)\) and \(\lambda_0^2=-2F_0N_\theta(2F_0N_\theta+|g_0|)\). This verifies the orientation and signs of the growing coordinate.

Replacing the background and normal by the actual ones changes the algebraic projected matrix by \(O(S_*^{-1})\). This follows from the explicit rational expression for the projection, the uniform positive denominators, the value bounds already proved, and the mean-value formula on their common compact range. Also \(K_a-K,s_a-s=O(S_*^{-1})\), while their \(v\)-derivatives and \(s'\) are \(O(S_*^{-1})\). Hence \(B'=O(S_*^{-1})\) and the normal-motion term in \(\mathcal A_\Phi\) is \(O(S_*^{-1})\). Therefore \[
B^\ell(\mathcal A_\Phi B-B')=
\operatorname{diag}(\lambda,-\lambda)+E,\qquad |E|\le C/S_*.
\] Every fixed derivative of \(E\) is polynomially bounded; no derivative-smallness claim is needed. The damping comparison is \[
d_{\rm ref}=\varepsilon k^2B_s^2(1+s^2),\qquad
|d-d_{\rm ref}|\le C/S_*,
\] because \(\varepsilon k^2\) is bounded and \(|n_\Phi|^2-B_s^2(1+s^2)\) is \(O(S_*^{-1})\).

\subsection{Envelope and uniform forward inverse}\label{si:envelope-and-uniform-forward-inverse}

Let \[
a_{\rm net}(y)=\frac{\lambda_0}{\sqrt{1+y^2}}-
\frac{\lambda_0(1+y^2)}{(1+u_*^2)^{3/2}},
\qquad y=|s(v)|=u_*/2+u_*v/L_s.
\] It is zero at \(v=L_s/2\), and direct differentiation gives \[
\frac{d}{dv}a_{\rm net}(|s(v)|)
=-\frac{u_*\lambda_0|s(v)|}{L_s}
\left((1+s(v)^2)^{-3/2}+2(1+u_*^2)^{-3/2}\right).
\] Since \(u_*/2\le|s|\le3u_*/2\), its value lies between \(-C/L_s\) and \(-c/L_s\), for fixed positive \(c,C\). Define \(P(v)=\exp(\int_{L_s/2}^v a_{\rm net}(|s(w)|)\,dw)\). Integrating the derivative inequality on either side of the midpoint, keeping the direction of each integral, gives \[
-C_1(v-L_s/2)^2/L_s\le\log P(v)\le-c_1(v-L_s/2)^2/L_s.
\] This proves the two Gaussian bounds and \(0<P\le1\) with the source's integral normalization.

Writing \(t_m=Bz_m\) gives \[
z_m'=A_mz_m+g_m,\quad
A_m=\operatorname{diag}(\lambda,-\lambda)+E-m^2dI_2,
\quad g_m=-B^\ell\operatorname{proj}_{n_\Phi^\perp}f_m.
\] For the complex Euclidean norm, taking the real part of \(z^*A_1z\) bounds the upper right derivative of \(|z|\) by \((\lambda-d_{\rm ref}+C/S_*)|z|\). Integration gives \[
\|V_1(v,w)\|\le
\exp(C(v-w)/S_*)\,P(v)/P(w)\le C'P(v)/P(w).
\] The factor \(C'\) is fixed because \(0\le v-w\le L_s\asymp S_*\). Since the difference of \(A_m\) from \(A_1\) is the scalar matrix \(-(m^2-1)dI_2\), differentiation directly verifies \[
V_m(v,w)=\exp\left(-(m^2-1)\int_w^v d(a)\,da\right)V_1(v,w).
\] For every nonzero integer \(m\), \(m^2-1\ge0\) and \(d>0\), so the same value bound is valid for all such harmonics. This is the exact mechanism preventing a new small-domain requirement when the finite harmonic set increases.

If the source obeys \[
|D^If_m|\le C_I\varepsilon^\alpha S_*^{b_I}
\sqrt\zeta\,\delta^{-a_I}P(v),
\] the zero-data Duhamel integral gives \[
z_m(v)=\int_0^vV_m(v,w)g_m(w)\,dw,\quad
|z_m(v)|\le C\varepsilon^\alpha S_*^{b+1}
\sqrt\zeta\,\delta^{-a}P(v).
\] The source's \(P(w)\) cancels the propagator's denominator exactly. The shell weights are constant along the integration path because \(X\) is constant there.

For completeness, parameter differentiation can be performed without differentiating a possibly exponentially large propagator in isolation. Hold \(v\) fixed and use slow variables and the initial transverse coordinate as parameters. They commute with \(\partial_v\); the zero initial data and all their parameter derivatives vanish. For each multi-index \(I\), \[
(\partial^Iz_m)'=A_m\partial^Iz_m+\partial^Ig_m+
\sum_{0<J\le I}\binom IJ(\partial^JA_m)\partial^{I-J}z_m.
\] At fixed finite harmonic bound \(M\), each \(\partial^JA_m\) is bounded by \(C_{J,M}S_*^{c_J}\). Induction on \(|I|\), using the already proved propagator bound, shows that every source term on the right has the same \(\varepsilon^\alpha\sqrt\zeta P\) weight with a larger inverse-distance exponent and polynomial degree. A further time integral contributes only \(L_s\). One may choose the degree at step \(I\) to exceed \[
1+\max\{b_I,\max_{0<J\le I}(c_J+b'_{I-J})\}.
\] Pulse derivatives follow from the differential equation, and mixed pulse/parameter derivatives follow by repeated differentiation of that equation. Returning to the current \(H\) uses \(D_Hv=O(S_*)\), \(D_HH_w=I-v_t\lambda_t\), with higher affine derivatives zero, and therefore introduces only additional polynomial powers. Multiplying by \(B\) yields the stated bound for \(t_m\). The explicit pressure has the extra factor \((km)^{-1}=O(\varepsilon^{1/2})\); its other factors have the proved derivative bounds. Thus the exponent for pressure is \(\alpha+1/2\).

The support argument also follows from this equation. If a fixed slow/transverse point is outside the source's closed support, the entire source path is zero; uniqueness gives a zero solution. At a boundary all source jets vanish by the prescribed smooth zero extension, so the differentiated equations inductively give zero solution jets. Across a shell edge the retained \(\sqrt\zeta\delta^{-a_I}\) bounds give the same conclusion. Deck translations preserve the path and the prescribed data, giving common-torus descent by uniqueness. Smooth extension to the rectangle enlargement follows from smooth finite-dimensional ODE dependence on its coefficients and data; no temporal zero extension is asserted before cutoff. On compact sets with \(q>0\), one-sided background/source derivatives at \(\tau=0\) pass through the same parameter equations and pressure formula.

Only the lower bounds for denominators and finitely many value comparisons determine \(q_*\). The induction above changes constants and polynomial degrees for each fixed \(M,I\); it does not impose further smallness inequalities. This proves the quantifier statement of Corollary 7.3.

\subsection{The real homogeneous pulse and energy transfer}\label{si:the-real-homogeneous-pulse-and-energy-transfer}

For \(m=1,f=0\), specify \(z_+(0)=P(0)>0,z_-(0)=0\). While \(z_+>0\), the ratio \(r=z_-/z_+\) obeys exactly \[
r'=E_{21}+(-2\lambda+E_{22}-E_{11})r-E_{12}r^2.
\] Let \(c_\lambda=\inf\lambda>0\). On \(r=\pm K_b/S_*\), the outward derivative is at most \[
\frac{-2c_\lambda K_b+C(1+K_b/S_*+K_b^2/S_*^2)}{S_*}.
\] Choose fixed \(K_b>C/c_\lambda\), then enlarge the fixed band threshold to make the numerator negative. The closed interval \([-K_b/S_*,K_b/S_*]\) is invariant until any putative first zero of \(z_+\). On this interval, \[
\left(\log(z_+/P)\right)'=-(d-d_{\rm ref})+E_{11}+E_{12}r=O(S_*^{-1}).
\] Since \(z_+(0)=P(0)\) and \(L_s/S_*\) is bounded, integration proves \(cP\le z_+\le CP\), ruling out that putative zero. For all \(v\), \[
x=z_++z_-=z_+(1+r),\quad
y=c_0\sqrt{1+s^2}(z_+-z_-),
\] so \[
cP\le x\le CP,\qquad
y/x=c_0\sqrt{1+s^2}+O(S_*^{-1}).
\] The same differentiated equation used for the inverse proves every fixed derivative bound, with the nonzero initial value retained: it is fixed with the label, so all slow/transverse derivatives of that datum vanish. The value solution contributes \(\|V_1(v,0)\|P(0)\le CP(v)\). Pulse derivatives are obtained from the equation. Pressure again has the extra \(\varepsilon^{1/2}\).

For the real homogeneous amplitude \(t\), \(n_\Phi\cdot t=0\) makes the normal component of \(\mathcal A_\Phi t\) perpendicular to \(t\). Direct multiplication, keeping both cylindrical \(2F\) terms, gives \[
t\cdot\mathcal Kt
=-2Ft_rt_\theta+(2F+RF_R)t_rt_\theta+G_Rt_rt_z
=g\cdot(t_r(t_\theta,t_z)).
\] Taking \(t\cdot t'\) therefore proves the exact identity \[
\frac12\frac{d}{dv}|t|^2
=-g\cdot(t_r(t_\theta,t_z))-\varepsilon k^2|n_\Phi|^2|t|^2.
\] At the representative the first term is \(-|g_0|T_N\) for the flux \(T=T_NN+T_KK\). It is positive for the negative \(N\)-component of the realized stress. The carrier physical wavelength is \(\sqrt Q/k\asymp Q^{1/2+h/2}\), while the leading physical wave amplitude has power \(Q^{-A}\sqrt\varepsilon=Q^{-1/2-h/2}\), up to polynomial factors in \(S_*\). Their product has power \(Q^0\); this agrees with the retained leading damping \(\varepsilon k^2\asymp1\).




% Standalone inclusion fragment; requires amsmath, amssymb, amsthm.
% The wrapper in this directory supplies the theorem environments.
\section{The complete radial obstruction and the five-equation correction}
\label{si:sec:rf-audit}

This fragment proves the radial and finite-dimensional maps used in the
initial mean correction. The input fields remain the actual fields: a
nonzero flat field is retained in the field and its products. In particular,
the proof does not identify Haar cancellation with exact characteristic
cancellation. The reference formulas checked here are
\texttt{mean\_corrections/mean\_corrections\_body.tex}, (MC1)--(MC10),
(MC13)--(MC15), and (MC22)--(MC30), together with the reserved-support
construction in \texttt{base\_heat/derivation.tex}.

\subsection{Coordinates, domains, and fixed parameters}

Fix the original parameters
\[
0<h<\tfrac1{100},\quad A=\tfrac12+h,\quad D=\tfrac12-h,
\quad\kappa_s=10^{-5},\quad Q=2^{-\ell},\quad
\epsilon=Q^h,\quad S_*=\ell^2,
\]
where the integer band index satisfies \(\ell\ge1\). The physical variables
are \(r>0,z\in\mathbb R,t<1\). Write \(\tau=1-t\). The equations
\[
z=q^D\eta,\qquad \tau=q(1-\eta^2),\qquad
X=\frac{r^2}{2q},\qquad L_q=1-2h\eta^2
\tag{RF1}
\]
specify \(q>0,|\eta|<1\) uniquely: the derivative of
\(q-z^2q^{2h}\), on \(q>|z|^{1/D}\), equals
\(1-2hz^2q^{2h-1}\ge1-2h>0\); the function starts at zero and
tends to infinity. Implicit differentiation gives
\[
q_t=-L_q^{-1},\quad q_z=2\eta q^{1-D}/L_q,
\quad\eta_t=D\eta/(qL_q),\quad
\eta_z=(1-\eta^2)/(q^DL_q).
\tag{RF2}
\]
These formulas also give one-sided derivatives at \(|\eta|=1\) when
\(q>0\). There is no division by \(1-\eta^2\).

In a fixed band put
\[
R=Q^{-1/2}r,\quad Z=Q^{-D}z,\quad T=Q^{-1}\tau,
\quad s=q/Q,\qquad X=R^2/(2s).
\]
The source band has \(s_-=1/2\le s\le s_+=2\); we retain the
symbols \(s_-,s_+\) when writing its fixed endpoint bounds.
They give a compact range for \(Z=s^D\eta,T=s(1-\eta^2)\),
including its one-sided boundary. Equations (RF2) give
\[
\partial_Zs=2\eta s^{1-D}/L_q,\quad
\partial_Ts=L_q^{-1},\quad
\partial_Z\eta=(1-\eta^2)/(s^DL_q),\quad
\partial_T\eta=-D\eta/(sL_q).
\tag{RF3}
\]
All repeated derivatives are bounded on that compact domain, because
every differentiated denominator is a power of \(s\) or \(L_q\), with
positive lower bounds. This proves the needed slow-coordinate bounds.
The active shell is
\[
X_a<X<X_b,\qquad
0<X_a<X_b<\infty,
\quad R_a(Z,T)=\sqrt{2sX_a},\quad R_b(Z,T)=\sqrt{2sX_b}.
\tag{RF4}
\]
Every field supported in its closed shell is extended smoothly by zero
across both edges. The resulting uniform radial interval stays away from
\(R=0\).

Retain the integer matrix and its exact eigenvectors:
\[
J_g=\begin{pmatrix}3&1\\1&5\end{pmatrix},\quad
\Lambda_g=4-\sqrt2,\quad T_g=4+\sqrt2,\quad
v_r=(1,1-\sqrt2),\quad v_t=(\sqrt2-1,1),
\]
\[
\rho_g=\frac{\log\Lambda_g}{\log T_g},\qquad
d_r=2((1+h)\rho_g-h\kappa_s)>0.
\]
The positivity follows, for example, from \(\Lambda_g>2\), \(T_g<8\),
so \(\rho_g>1/3\), while \(0<h<1/2\). Direct matrix
multiplication proves \(J_gv_r=\Lambda_gv_r\) and
\(J_gv_t=T_gv_t\). At a fixed common covering level \(i_0\),
\[
y=J_g^{i_0}Y\in\mathbb T^2=\mathbb R^2/\mathbb Z^2,
\quad M=\Lambda_g^{i_0}Q^{d_r/2},\quad c=T_g^{i_0}Q^{1+h},
\quad L=v_r\cdot\partial_y,\quad N=v_t\cdot\partial_y.
\]
The exact normalized differential operators are
\[
D_r=\partial_R+Md_rR^{d_r-1}L,\quad
D_z=\epsilon\partial_Z,\quad
t_*=-\epsilon\partial_T+cN,\qquad
\Delta_0=D_r^2+R^{-1}D_r+D_z^2.
\tag{RF5}
\]
Here \(Q,i_0,\epsilon,M,c\) are fixed under all coefficient
differentiation. The phase in physical variables is
\(Y=v_rr^{d_r}+v_tt\); the chain rule gives exactly (RF5).
The operators \(D_r,D_z,t_*\) commute pairwise, since the only
nonconstant coefficient in them depends just on \(R\).

For \(i(\ell)=\lfloor\log_{T_g}(Q^{-1-h}/S_*)\rfloor\),
\[
T_g^{-1}S_*^{-1}<T_g^{i(\ell)}Q^{1+h}\le S_*^{-1},
\quad
T_g^{-\rho_g}\epsilon^{-\kappa_s}S_*^{-\rho_g}
<\Lambda_g^{i(\ell)}Q^{d_r/2}
\le\epsilon^{-\kappa_s}S_*^{-\rho_g}.
\tag{RF6}
\]
The second line follows by raising the first floor inequality to
\(\rho_g\) and substituting \(d_r/2=(1+h)\rho_g-h\kappa_s\).
If \(|i_0-i(\ell)|\le B\) for one fixed integer \(B\), the bounds
only change by the fixed factors \(T_g^{\pm B}\) and
\(\Lambda_g^{\pm B}\). In particular
\(M^{-p}\le C_p\epsilon^{p\kappa_s}S_*^{p\rho_g}\).

Define
\[
d_a=\log(X/X_a),\quad d_b=\log(X_b/X),\quad
\delta=\min(1,d_a,d_b),\qquad
\zeta=\exp(-a_a/d_a^2-a_b/d_b^2),\quad a_a,a_b>0.
\]
The space \(\mathcal M_\alpha\) is the space of compatible smooth
torus coefficients with this shell support such that every fixed
coefficient derivative \(I\) satisfies
\[
|\partial^If|\le C_I\epsilon^\alpha S_*^{B_I}
\zeta\delta^{-C'_I}.
\tag{RF7}
\]
All constants may depend on the fixed stage, fixed profiles and derivative
order, but not on the band. The space \(\mathcal S_\alpha\) consists
of slow functions of \((Z,T)\) with the same estimate without
\(\zeta\delta^{-C'_I}\). Smooth coefficient derivatives preserve the
exponent. Formula (RF5) shows that \(D_r\) loses at most
\(\kappa_s\), \(D_z\) gains one, and \(t_*\) preserves the
exponent; on a torus-independent field \(t_*\) gains one.
Leibniz' rule and \(0<\zeta\le1\) prove
\(\mathcal M_a\mathcal M_b\subset\mathcal M_{a+b}\).

\subsection{The full radial exact sequence and its specified splitting}

At a fixed slow point let \(\mathcal C\) denote real
\(C^\infty\) functions of \((R,y)\in(0,\infty)\times\mathbb T^2\)
supported in the closed shell and smoothly zero across its edges. For
\(e\in\{0,1,2\}\) define \(\mathcal D_e=D_r+e/R\) and
\[
(\Gamma H)(R,w)=H(R,w+MR^{d_r}v_r).
\]
Its inverse uses the opposite sign. The chain rule gives
\[
\partial_R\Gamma(R^e\phi)=\Gamma(R^e\mathcal D_e\phi).
\tag{RF8}
\]
Define a continuous linear map from \(\mathcal C\) into smooth
functions on the entire torus by
\[
(K_ef)(w)=\int_0^\infty R^e f(R,w+MR^{d_r}v_r)\,dR.
\tag{RF9}
\]
Continuity in the smooth topologies follows by differentiating the fixed
compact integration interval; for each fixed \(M\) the required
coefficient bounds are finite.

\begin{siproposition}[All radial obstructions]
The sequence of real vector spaces
\[
0\longrightarrow\mathcal C\xrightarrow{\mathcal D_e}\mathcal C
\xrightarrow{K_e}C^\infty(\mathbb T^2;\mathbb R)
\longrightarrow0
\tag{RF10}
\]
is exact. A compactly supported solution of \(\mathcal D_e\phi=f\)
exists if and only if \(K_ef=0\) as a torus function, and is unique.
\end{siproposition}
\begin{proof}
Integrating (RF8) gives necessity, since its primitive is zero at both
radial endpoints. Conversely set
\[
\phi(R,y)=R^{-e}\int_0^R R'^e
f(R',y+M(R'^{d_r}-R^{d_r})v_r)\,dR'.
\tag{RF11}
\]
It vanishes below the shell. Above it the integral equals
\((K_ef)(y-MR^{d_r}v_r)\), which vanishes by the specified condition.
Inside and at the edges its derivatives exist and match its zero
extension, by differentiation of smooth compactly supported data and
(RF8). That identity proves \(\mathcal D_e\phi=f\). If
\(\mathcal D_e\phi=0\), then \(\Gamma(R^e\phi)\) is constant
in \(R\); its value below the shell is zero, proving injectivity and
uniqueness. Finally choose \(h_m\in C_c^\infty((R_a,R_b))\) of
integral one. For any prescribed \(F\in C^\infty(\mathbb T^2)\),
\(f(R,y)=R^{-e}h_m(R)F(y-MR^{d_r}v_r)\) satisfies
\(K_ef=F\), by direct substitution. This proves surjectivity.
\end{proof}

Choose the original slow radial step \(\chi_m=\chi_m(X)\), zero
near \(X_a\), one near \(X_b\), with \(\partial_R\chi_m\)
supported in a fixed interior subinterval. For \(g=R^ef\), put
\[
\begin{split}
Ig(R,y)&=\int_0^R g(R',y+M(R'^{d_r}-R^{d_r})v_r)\,dR',\\
Jg(R,y)&=\int_0^\infty g(R',y+M(R'^{d_r}-R^{d_r})v_r)\,dR',\\
T_ef&=R^{-e}(Ig-\chi_mJg),\qquad
A_ef=R^{-e}(\partial_R\chi_m)Jg.
\end{split}
\tag{RF12}
\]
In particular \(Jg(R,y)=(K_ef)(y-MR^{d_r}v_r)\).
Equation (RF8) gives \(D_rIg=g,D_rJg=0\), and hence
\[
\mathcal D_eT_e=1-A_e,\qquad K_eA_e=K_e,
\qquad A_e^2=A_e,\qquad T_e\mathcal D_e=1\quad\hbox{on }\mathcal C.
\tag{RF13}
\]
For the second identity, conjugate the integral for \(K_eA_ef\):
the integrand becomes \(\partial_R\chi_m\,K_ef(w)\), whose integral
equals \(K_ef(w)\). It then implies idempotence from (RF12).
For the last identity \(K_e\mathcal D_e=0\), so the compact primitive
\(T_e\mathcal D_e\phi\) equals \(\phi\) by uniqueness.
The same argument proves \(T_eA_e=0\), since the conjugated partial
integral of \(\partial_R\chi_mK_ef\) equals \(\chi_mK_ef\).
Thus the specified projector splits the exact sequence.

Let a bar mean Haar average on \(\mathbb T^2\), with total mass one.
Translation invariance and Fubini give
\[
\overline{K_ef}=\int_0^\infty R^e\bar f(R)\,dR.
\tag{RF14}
\]
This scalar condition is strictly weaker than \(K_ef=0\).
Indeed, for every nonzero \(k\in\mathbb Z^2\), the real function
\[
f_M(R,y)=e^{-M^2}R^{-e}h_m(R)
\cos(2\pi k\cdot(y-MR^{d_r}v_r))
\tag{RF15}
\]
has \(\bar f_M(R)=0\) at every radius, whereas
\(K_ef_M(w)=e^{-M^2}\cos(2\pi k\cdot w)\ne0\).
Taking a smoothly scaled fixed interior bump makes this a family on the
whole band domain. Every fixed derivative grows by at most a fixed
power of \(M\) before the factor \(e^{-M^2}\). Since (RF6) implies
\(M\to\infty\) exponentially in \(\ell\), that factor controls
every prescribed power of \(\epsilon\). Thus even an all-orders
flat, zero-Haar source need not admit exact compact radial inversion.

\subsection{Weighted primitive bounds with moving shell edges}

Set \(U=R^{d_r}\) and, on the positive radial interval, define
\[
F(U,Z,T,y)=\frac{g(R,Z,T,y)}{d_rR^{d_r-1}}.
\]
Extend it by zero to \(U\in\mathbb R\). This extension is smooth,
because its support stays in a compact subset of \(U>0\). The
change of variables and its inverse have uniformly bounded fixed
derivatives on the radial interval. The integrals (RF12) become
\[
A_-(U,y)=\int_{-\infty}^0F(U+u,y+Muv_r)\,du,
\quad
A_+(U,y)=\int_0^\infty F(U+u,y+Muv_r)\,du,
\tag{RF16}
\]
with \(Ig=A_-\), \(Jg=A_-+A_+\). Thus
\(Ig-\chi_mJg=(1-\chi_m)A_--\chi_m A_+\).
The argument shift \(Muv_r\) is independent of \(U,Z,T\).
Every coefficient derivative in (RF16) therefore differentiates just
\(F\); no extra factor of \(M\) is produced. The effective
integration length is uniformly bounded. Moving edges give no boundary
terms, because the integration is of the smooth zero extension on fixed
half-lines.

It remains to retain the prescribed weight. For any \(a>0,B\ge0\),
\[
\frac d{ds}\log(e^{-a/s^2}s^{-B})=2a/s^3-B/s>0
\tag{RF17}
\]
on a sufficiently small positive interval (on every positive interval
if \(B=0\)). Near the inner edge only \(A_-\) is present, and
its contributing radii satisfy \(0<d_a(R')\le d_a(R)\).
The other edge factor and all its fixed derivatives are bounded there.
Apply (RF17) to each differentiated input bound to obtain
\(C\zeta(R)\delta(R)^{-B'}\). Near the outer edge only
\(A_+\) occurs, and its contributing radii satisfy
\(0<d_b(R')\le d_b(R)\); the same calculation applies.
On the intervening closed interior region \(\zeta\) has a positive
lower bound, so the unweighted integral bound supplies the desired
estimate. Derivatives of \(\chi_m\) have this interior support.
The finite Leibniz and change-of-variable sums now prove
\[
T_e:\mathcal M_\alpha\longrightarrow\mathcal M_\alpha.
\tag{RF18}
\]
The operator cannot create a slow point outside the projection of the
source's closed support, since the entire source integral is zero on
every slow neighborhood where the source is zero.

\subsection{Every derivative of the projector is flat after Haar cancellation}

The directional inverses used here have exact small-divisor bounds. If
\(k=(k_1,k_2)\ne0\), then
\[
(v_r\cdot k)((k_1+k_2)+\sqrt2k_2)=(k_1+k_2)^2-2k_2^2
\]
is a nonzero integer. It cannot vanish unless both integer coordinates
are zero, by irrationality of \(\sqrt2\). The second factor is at
most \(C|k|\), hence \(|v_r\cdot k|^{-1}\le C(1+|k|)\).
The identical argument with \((k_2-k_1)+\sqrt2k_1=v_t\cdot k\)
proves the bound for \(N\).
On zero-mean torus functions define
\[
(L^{-p}H)\widehat{\ }(k)=
\begin{cases}(2\pi i v_r\cdot k)^{-p}\widehat H(k),&k\ne0,\\0,&k=0.
\end{cases}
\]
Integrate the Fourier coefficient \(m+p+3\) times in a coordinate
with \(|k_j|\ge |k|/\sqrt2\). After an output derivative of order
at most \(m\), each summand is bounded by
\(C\|H\|_{C^{m+p+3}}(1+|k|)^{-3}\). Since that sequence is
summable on \(\mathbb Z^2\),
\[
\|L^{-p}H\|_{C^m}\le C_{m,p}\|H\|_{C^{m+p+3}}.
\tag{RF19}
\]
The proof also justifies termwise differentiation, preservation of
reality, and commutation with coefficient derivatives outside the torus.
The same proof gives \(N^{-1}\), with a loss of four derivatives.

Suppose the exact identity
\(\int R^e\bar f\,dR=0\) holds at every slow point.
For \(F^\circ=F-\bar F\), integration of
\[
\frac d{du}L^{-1}F^\circ(U+u,y+Muv_r)
=\partial_UL^{-1}F^\circ(U+u,y+Muv_r)+MF^\circ(U+u,y+Muv_r)
\]
over the whole line has zero boundary terms. Repeating it \(p\) times
gives the exact identity, with its retained sign,
\[
Jg=(-M^{-1})^p\int_{\mathbb R}
\partial_U^pL^{-p}F^\circ(U+u,y+Muv_r)\,du.
\tag{RF20}
\]
The omitted zero Fourier mode and every one of its slow derivatives
integrate to zero by the specified moment identity. Applying (RF19) and
the fixed-shift differentiation argument gives, for every fixed
coefficient derivative \(I\) and every integer \(p\ge1\),
\[
|\partial^IA_ef|\le C_{I,p}\epsilon^{\alpha+p\kappa_s}
S_*^{B_{I,p}}\zeta\delta^{-C'_{I,p}},\qquad
\overline{A_ef}=0.
\tag{RF21}
\]
Only finitely many extra input derivatives are needed for each
\((I,p)\): at most \(p\) extra radial derivatives and
\(p+3\) extra torus derivatives in the displayed estimate. The output
is supported where \(\partial_R\chi_m\ne0\), so the weight
there is bounded below. If \(f\) is torus independent, then
\(Jg=\int R'^e f(R')\,dR'=0\) exactly; thus \(A_ef=0\).

Here is the quantifier statement needed for a fixed construction.
Choose the fixed parameters, band comparison bounds, covering-offset
bound and a fixed lower band index \(\ell_*\ge1\) once. Fix a
physical field normalization and an output Cartesian derivative order
\(m\). In the shell every graph derivative follows (RF2), (RF5),
and the derivatives of the cylindrical frame. The latter cost powers
of \(r^{-1}\), with \(r\asymp Q^{1/2}\). Consequently a
finite sum of coefficient derivatives controls the physical derivative,
with a loss \(Q^{-L_m}\ell^{B_m}\), for some finite \(L_m\)
independent of \(p\). For desired flat order \(N_0\ge0\),
choose an integer \(p\) with
\(h(\alpha+p\kappa_s)-L_m>N_0\). Formula (RF21) gives
\[
|\nabla_{x,t}^{\le m}(A_ef)_{\rm phys}|
\le C\,Q^{N_0+\beta}\ell^B\le C' q^{N_0},
\qquad\beta=h(\alpha+p\kappa_s)-L_m-N_0>0.
\tag{RF22}
\]
The last constant is finite on the same bands \(\ell\ge\ell_*\),
because \(\sup_{u\ge1}u^B2^{-\beta u}<\infty\), as follows by
differentiating \(B\log u-\beta(\log2)u\). If \(B<0\), the
bound is immediate. Thus \(p\) and the constants may depend on
\(m,N_0\), while the allowed physical domain does not shrink with
\(m,N_0\). This is all-orders flatness, proved on one fixed domain.

These maps also respect the covering maps. For an integer \(n\ge0\)
put \(P_nH(y)=H(J_g^ny)\), \(M'=\Lambda_g^nM\). Direct
substitution gives
\[
T_e^{M}P_n=P_nT_e^{M'},\quad A_e^{M}P_n=P_nA_e^{M'},
\quad K_e^{M}P_n=P_nK_e^{M'}.
\tag{RF23}
\]
The Haar mean is preserved: a Fourier mode \(k\) pulls back to
\((J_g^n)^Tk\), which is zero only if \(k=0\), since
\(\det J_g=14\ne0\). Thus noninvertibility of this torus covering
does not obstruct the exact compatibility of the radial maps.

\subsection{Pressure, temporal inversion, and the actual divergence-free field}

Choose the fixed physical bump
\(\rho_{\rm phys}=q^{-1/2}\widehat\rho(r/\sqrt q)\),
\(\int\widehat\rho(x)dx=1\), with support in a fixed proper
interior interval. Its chart representative
\(\rho=Q^{1/2}\rho_{\rm phys}=s^{-1/2}\widehat\rho(R/\sqrt s)\)
has \(\int\rho\,dR=1\). For a radial source \(g_r\), define
\[
P=\int\bar g_r\,dR,\qquad p_m=T_0(g_r-\rho P).
\]
Then the exact residual, including its flat part, is
\[
D_rp_m-g_r=-\rho P-A_0(g_r-\rho P),
\qquad\partial_R\bar p_m=\bar g_r-\rho P.
\tag{RF24}
\]
The first term is a specified slow bump residual. The second has zero
Haar mean and is flat by (RF21)--(RF22). Radial integration preserves
the exponent into \(\mathcal S_\alpha\), because every
\(\zeta\delta^{-B}\) is bounded and the interval has bounded
length. Thus (RF18) proves the pressure and pressure-difference estimates.

For a requested axial increment \(\gamma_d\) satisfying
\(\int R\bar\gamma_d\,dR=0\), define the actual potential and field
\[
\Psi_*=T_1\gamma_d,\qquad
\Delta\beta=-D_z\Psi_*=-\epsilon\partial_Z\Psi_*,\qquad
\Delta\gamma=\mathcal D_1\Psi_*
=\gamma_d-A_1\gamma_d.
\tag{RF25}
\]
Their divergence vanishes exactly:
\[
\mathcal D_1\Delta\beta+D_z\Delta\gamma
=-\epsilon\mathcal D_1\partial_Z\Psi_*
+\epsilon\partial_Z\mathcal D_1\Psi_*=0.
\]
This uses commutation of differential operators on the constructed
potential. It does not commute \(\partial_Z\) with \(T_1\),
whose cutoff and shell can depend on \(Z\). Equations (RF18) and
(RF21) give
\(\Psi_*,\Delta\gamma\in\mathcal M_\alpha\) and
\(\Delta\beta\in\mathcal M_{\alpha+1}\). Moreover
\(\overline{\Delta\gamma}=\bar\gamma_d\), so the weighted
axial moment remains exactly zero. When \(\gamma_d\) is slow,
\(A_1\gamma_d=0\) and this realization is pointwise exact.

For zero-Haar tangential errors \(E_\theta^\circ,E_z^\circ\), set
\[
\Delta v=-c^{-1}N^{-1}E_\theta^\circ,\qquad
\gamma_d=-c^{-1}N^{-1}E_z^\circ,\qquad a=-A_1\gamma_d.
\tag{RF26}
\]
The inverse Fourier estimate and \(c^{-1}\le CS_*\) give the
same coefficient exponents, and the inputs have zero Haar mean at
each radius. Realize the axial field by (RF25). Since torus
translations and the slow cutoff commute with \(N\),
\[
cN\Delta v=-E_\theta^\circ,\qquad
cN\Delta\gamma=-E_z^\circ+cNa.
\tag{RF27}
\]
The term \(cNa\) is flat; the terms
\(-\epsilon\partial_T\Delta v\) and
\(-\epsilon\partial_T\Delta\gamma\) are the retained slow-time
terms with one extra \(\epsilon\). The actual velocity contains
\(a\). In particular its products contain \(a^2\) and all cross
terms; zero Haar mean of \(a\) does not eliminate those products.
This gives the exact interpretation of the temporal mean step used
before the five-equation map.

\subsection{The original five rows and every profile-to-chart factor}

Let \(I_{\rm mean}\Subset(X_a,X_b)\) be the reserved open
power-law interval, and \(I_m=\{\sqrt{2X}:X\in I_{\rm mean}\}\).
The original leading profiles on this entire interval are
\[
U_0=0,\qquad E_0=c_{\rm patch}(1+\eta^2)^{-1}X^{-1/2-\lambda},
\qquad \lambda>0,\quad c_{\rm patch}>0.
\tag{RF28}
\]
Hence the \(q\)-profile tangential base, under the exact map
\(x=r/\sqrt q=\sqrt{2X}\), is
\[
G_q=q^AG_{\rm phys}=0,\quad
V_q=q^AV_{\rm phys}=a(\eta)x^{-1-2\lambda},\quad
a(\eta)=2^{1/2+\lambda}c_{\rm patch}(1+\eta^2)^{-1}.
\tag{RF29}
\]
In particular \(a(\eta)\ge a_0=2^{-1/2+\lambda}c_{\rm patch}>0\)
on \(|\eta|\le1\), and all derivatives of \(a^{-1}\) are
bounded there. Section below verifies that the positive-order base
construction preserves this exact law, not merely its leading term.

For the actual mean velocity \((\beta,v,\gamma)\) and fixed full
wave covariance \(W^{ab}\), the three slow defects are
\[
\begin{split}
P&=\int\bar g_r\,dR,\\
J_\theta&=\int R^2\overline{Gv+V\gamma+\gamma v+W^{z\theta}}\,dR,\\
J_z&=\int R\overline{2G\gamma+\gamma^2+W^{zz}}\,dR
-\frac12\int R^2\bar g_r\,dR.
\end{split}
\tag{RF30}
\]
The five requested correction rows in the fixed \(Q\) chart are
\[
\begin{aligned}
\int R^2\Delta v\,dR&=0,&\int R\gamma_d\,dR&=0,\\
\int(2V/R)\Delta v\,dR&=-P,\\
\int R^2(G\Delta v+V\gamma_d)\,dR&=-J_\theta,\\
\int(2RG\gamma_d-RV\Delta v)\,dR&=-J_z.
\end{aligned}
\tag{RF31}
\]
The source \(g_{r,*}=Q^{2A+1/2}g_{r,\rm phys}\) and
velocities \(u_*=Q^Au_{\rm phys}\) give
\[
P=Q^{2A}P_{\rm phys},\quad
J_\theta=Q^{2A-3/2}J_{\theta,\rm phys},\quad
J_z=Q^{2A-1}J_{z,\rm phys}.
\tag{RF32}
\]
Indeed \(R=Q^{-1/2}r,dR=Q^{-1/2}dr\): each product in
\(J_\theta\) has velocity factor \(Q^{2A}\) and measure
\(Q^{-3/2}\), while each term in \(J_z\) has factor
\(Q^{2A-1}\), including the source integral with its retained
\(-1/2\). Therefore the exact \(q\)-profile targets are
\[
P_q=s^{2A}P,\qquad
J_{\theta,q}=s^{2A-3/2}J_\theta,\qquad
J_{z,q}=s^{2A-1}J_z.
\tag{RF33}
\]
If \(u(x),d(x)\) are profile increments, their chart fields are
\[
\Delta v(R)=s^{-A}u(R/\sqrt s),\qquad
\gamma_d(R)=s^{-A}d(R/\sqrt s).
\tag{RF34}
\]
Substitution \(R=\sqrt s\,x\) into every row of (RF31) proves
that its left side is its \(x\)-profile left side multiplied,
respectively, by
\(s^{3/2-A},s^{1-A},s^{-2A},s^{3/2-2A},s^{1-2A}\).
Equations (RF33) are exactly the matching target factors.

Choose \(x_0\in I_m\) and a fixed \(d_0>0\) such that
\(x_0,x_0e^{d_0},x_0e^{2d_0}\) all lie in \(I_m\).
Choose a sufficiently narrow nonnegative nonzero
\(\eta_0\in C_c^\infty(I_m)\) around \(x_0\), and set
\[
\eta_j(x)=e^{-jd_0}\eta_0(e^{-jd_0}x),\quad j=0,1,2,
\qquad\mu_p=\int_0^\infty x^p\eta_0(x)dx>0.
\]
Their supports are disjoint, their closed convex hull is contained in
\(I_m\), and direct change of variables gives
\[
\int x^p\eta_j(x)dx=\mu_p e^{jd_0p}\qquad(p\in\mathbb R).
\tag{RF35}
\]
The positivity and finite separation of the three centers prove that
such a common support width exists.

Use \(u(x)=\sum_{j=0}^2u_j\eta_j(x)\),
\(d(x)=\sum_{j=0}^1s_j\eta_j(x)\). Set
\[
p_0=2,\quad p_1=-2-2\lambda,\quad p_2=-2\lambda,
\quad t_i=e^{d_0p_i},\quad C(s_0')=u_0+u_1s_0'+u_2(s_0')^2.
\]
The three angular rows are precisely
\[
C(t_0)=0,\qquad C(t_1)=-\frac{P_q}{2a\mu_{p_1}}=:b_1,
\qquad C(t_2)=\frac{J_{z,q}}{a\mu_{p_2}}=:b_2.
\]
Their unique solution is the polynomial
\[
C(s_0')=b_1\frac{(s_0'-t_0)(s_0'-t_2)}{(t_1-t_0)(t_1-t_2)}
+b_2\frac{(s_0'-t_0)(s_0'-t_1)}{(t_2-t_0)(t_2-t_1)}.
\tag{RF36}
\]
Every denominator is nonzero, since \(p_1<p_2<p_0\).
Evaluation at each of the three nodes proves the three equations.
Uniqueness follows because a polynomial of degree at most two cannot
have three distinct roots: division by their three linear factors
would force degree at least three. The signed determinant in row
order (angular moment, pressure, axial-flux contribution) is
\[
-2a^2\mu_2\mu_{-2-2\lambda}\mu_{-2\lambda}
(t_1-t_0)(t_2-t_0)(t_2-t_1)\ne0.
\tag{RF37}
\]

For the two axial rows set \(t_b=e^{d_0}\),
\(t_a=e^{d_0(1-2\lambda)}\). The exact solution is
\[
s_1=-\frac{J_{\theta,q}}{a\mu_{1-2\lambda}(t_a-t_b)},
\qquad s_0=-t_bs_1.
\tag{RF38}
\]
Indeed \(\mu_1(s_0+t_bs_1)=0\) and
\(a\mu_{1-2\lambda}(s_0+t_as_1)=-J_{\theta,q}\).
The signed determinant is
\(a\mu_1\mu_{1-2\lambda}(t_a-t_b)\ne0\), since
\(\lambda>0\). Equations (RF34)--(RF38) prove all five original
rows, including the minus sign of \(-RV\Delta v\).

This is a linear map from the three arbitrary target functions into
the fixed five-dimensional correction space; two of its five prescribed
outputs are zero. Its uniqueness concerns that chosen space. For every
fixed \(\lambda>0,c_{\rm patch}>0\), its inverse coefficients and
all fixed slow derivatives are bounded on \(|\eta|\le1\),
because all node differences and \(\mu_p\) are fixed nonzero
constants and \(a\) has the explicit positive lower bound. Equations
(RF3), (RF33), (RF34), and Leibniz' rule prove the map
\[
(\mathcal S_\alpha)^3\longrightarrow
\mathcal M_\alpha\times\mathcal M_\alpha.
\tag{RF39}
\]
The output weight has a positive lower bound on its fixed interior
profile support. The linear solve requires no small target size and
no new small-\(q\) threshold. Constants are not uniform as
\(\lambda\downarrow0\): at \(\lambda=0\),
\(x^2V_q=ax\), so \(\int xd\,dx=0\) implies
\(\int x^2V_qd\,dx=0\). A target \(J_{\theta,q}\ne0\)
is then impossible for any compact smooth axial correction, not only
for the selected bumps. This is the exact degeneracy excluded by the
original fixed positive parameter.

The slow input \(\gamma_d\) has zero weighted moment. Equation
(RF25) therefore gives \(\Delta\gamma=\gamma_d\) exactly.
Its potential is supported in the closed interval between the first
and last axial bump: below the first bump the integral is zero; above
the last it is its zero total. It need not vanish between bumps.
Thus the base law and coefficient estimates must hold throughout that
interval. Its induced radial component is
\(\Delta\beta\in\mathcal M_{\alpha+1}\).

\subsection{Why the Borel base retains the exact reserved power law}

The original leading law (RF28) follows from the actual schedule.
Retain the fixed step
\[
\sigma(v)=\frac{e^{-1/v^2}}{e^{-1/v^2}+e^{-1/(1-v)^2}}
\ (0<v<1),\qquad \sigma=0\ (v\le0),\quad \sigma=1\ (v\ge1).
\]
Its extensions are smooth: every fixed derivative of the vanishing
exponential is that exponential times a polynomial in the inverse
edge distance, and hence tends to zero at that edge.
The parameter order is \(M_d\), then \(T_d=e^{M_d}+10\) and
\(P_*>e^{T_d}\), then small positive \(\lambda\), then small positive
\(h\) relative to \(\lambda,e^{-T_d}\), then large \(X_R\).
The finite restriction \(0<\lambda<e^{-25/60}\) ensures
\(T_w=60\log(1/\lambda)>25\). Put \(X_w=X_Re^{T_d+2}\).
The four reserved open intervals are
\[
\begin{aligned}
I_1&=X_w\exp((T_w-25,T_w-20)),&
I_2&=X_w\exp((T_w-20,T_w-15)),\\
I_3&=X_w\exp((T_w-14,T_w-9)),&
I_4&=X_w\exp((T_w-8,T_w-3)).
\end{aligned}
\tag{RF40a}
\]
Their assignments are cone repair, heat compensation, positive-order
moments and mean correction: \(I_{\rm pos}=I_3\),
\(I_{\rm mean}=I_4\). They are proper subintervals of the power-law
stage; the gap between \(I_3,I_4\) has logarithmic length one.
The initial reference has \(U=4\eta,E=P_*f(X/X_R)^{1/10}\) for
\(X\le X_R\), with \(f=(1+\eta^2)^{-1}\). With the logarithmic
coordinate reset to zero at each stage, the next three prescriptions
are
\[
\begin{array}{c|c|c}
\text{length}&U&l:=\partial_y\log(\sqrt{2X}E)\\ \hline
1&4\eta&\frac35(1-\sigma(y))\\
T_d&k(y)\eta&0\\
1&0&-\lambda\sigma(y)
\end{array}
,\qquad
k(y)=4\left[1-\sigma\left(\frac{\log(1+y)}{M_d}\right)\right].
\tag{RF40b}
\]
The next stage holds \(U=0,l=-\lambda\) for length \(T_w\).
Integrating \(\partial_y\log E=l-1/2\) with the original endpoint
amplitudes defines the positive constants
\[
\begin{split}
P_1&=P_*\exp\!\left(\int_0^1
[\tfrac35(1-\sigma(v))-\tfrac12]\,dv\right),\\
e_w&=P_1e^{-T_d/2}
\exp\!\left(\int_0^1[-\lambda\sigma(v)-\tfrac12]\,dv\right),\qquad
c_{\rm patch}=e_wX_w^{1/2+\lambda}.
\end{split}
\tag{RF40c}
\]
Thus throughout the power-law stage, including \(I_4\),
\(E=e_wf(X/X_w)^{-1/2-\lambda}
=c_{\rm patch}fX^{-1/2-\lambda}\) and \(U=0\).

The subsequent modifications preserve that exact statement on the full
interval. The regular-axis join ends before \(X_h=X_Re^{-5}<X_R\);
cone modulation is followed by repair in \(I_1\); heat compensation
is in \(I_2\) and changes only \(E\); and the later pulse and
exterior heat edit start after \(I_4\). None changes \(E,U\) on
\(I_4\). For the radial field write
\(\mathfrak M(X,\eta)=\int_0^XU(\xi,\eta)\,d\xi\).
The axis and cone moment repairs restore \(\mathfrak M\) exactly,
while the heat edits have axial increment zero. Above a repair
endpoint \(X_{\rm rep}\), equality of the two \(U\)'s gives the
exact propagation map
\[
\mathfrak M_{\rm new}(X,\eta)-\mathfrak M_{\rm old}(X,\eta)
=\mathfrak M_{\rm new}(X_{\rm rep},\eta)
-\mathfrak M_{\rm old}(X_{\rm rep},\eta)=0.
\tag{RF40d}
\]
It also proves equality of every parameter derivative.

The exact restoration just used is enforced by the actual repair
equations. At base \(U_0=0,E_0=K(\eta)X^{-1/2-\lambda}\),
the increments \(u_c,e_c\) change
\((\mathfrak M,J,I,S,C_p)\), respectively, by
\[
\int u_c,\quad
\int(\sqrt{2X}E_0u_c+\sqrt{2X}u_ce_c),\quad
\int\sqrt{2X}e_c,\quad
\int(u_c^2-E_0e_c-e_c^2/2),\quad
\int(E_0e_c/X+e_c^2/(2X)),
\tag{RF40e}
\]
where the measure is \(dX\). With two axial and three angular
bumps this is precisely \(B(\eta)c+Q_\eta(c,c)\).
The linear blocks have powers \((0,-\lambda)\) and
\((1/2,-1/2-\lambda,-3/2-\lambda)\), with fixed nonzero amplitude
factors. Their matrices are invertible on ordered nonnegative bumps:
a nonzero sum of \(m\) distinct real powers has at most \(m-1\)
positive zeros. Divide by its least power, apply Rolle's theorem,
and induct on the derivative's fewer powers to prove this assertion.
The evaluation determinant is therefore nonzero with a fixed sign on
the connected set of ordered nodes. By multilinearity, the moment
determinant is its integral against the product of the ordered
nonnegative bumps and cannot vanish.
Write \(\beta=\sup\|B^{-1}\|\) and
\(\|Q(c,c')\|\le\kappa\|c\|\|c'\|\). For discrepancy \(d_*\),
the map \(c\mapsto B^{-1}(d_*-Q(c,c))\) maps
\(\|c\|\le2\beta\|d_*\|\) into itself and contracts by at most
\(1/2\) if \(8\beta^2\kappa\|d_*\|\le1\). Iteration from zero
converges by a geometric-series bound on successive differences,
and its limit solves the exact moment equation. Its differentiated
matrix is invertible by the same bound, so differentiating successively
supplies every fixed parameter derivative after this one choice.
For cone repair \(\|d_*\|\le D_0/N\), and the construction fixes
\(N\ge8\beta^2\kappa D_0\) before the physical scale or any band.
For the earlier axis repair the fixed invertible row operations
\(J\mapsto J-4\eta I\), \(S\mapsto S-8\eta\mathfrak M\)
give distinct powers \((0,3/5)\) and
\((1/2,1/10,-9/10)\); the same determinant argument and quadratic
iteration implement the construction's chosen finite join tolerance.
Thus (RF40d) uses exact solved equations, not estimates on residual
moments.

The reference cumulative axial moment can now be calculated. Put
\[
\mathfrak m=eX_R\left(4+\int_0^{T_d}k(v)e^v\,dv\right)>0.
\tag{RF40f}
\]
At \(X=eX_R\) the reference moment is \(4\eta eX_R\).
On the next stage \(X=eX_Re^v,dX=eX_Re^v\,dv\), so direct
integration gives \(\mathfrak M=\mathfrak m\eta\). Every following
stage through \(I_4\) has \(U=0\), and (RF40d) proves the same
moment for the final leading field. Its exact radial reconstruction
then gives
\[
V_0=\frac{2\eta XU_0-2D\eta\mathfrak M-(1-\eta^2)\mathfrak M_\eta}
{L_q}
=-\mathfrak m\,\frac{2D\eta^2+1-\eta^2}{1-2h\eta^2}
=-\mathfrak m.
\tag{RF40g}
\]
Here \(2D=1-2h\) makes the numerator factor exactly \(L_q\).
Consequently the full leading radial field on \(I_4\) is
\(u_{r,0}=-\mathfrak m/r\).

Here the source construction can be checked at the level of its actual
coefficients. Its positive-order angular and axial profiles \(E_n,U_n\)
are cut off before a fixed positive-order patch \(I_{\rm pos}\),
then corrected by three angular and two axial bumps in that patch.
This lies strictly before the separate mean patch. The imposed axial
moment is
\[
\int_0^\infty RU_n(R,\eta)\,dR
=\int_0^\infty U_n(X,\eta)\,dX=0,
\quad R=\sqrt{2X},\quad dX=R\,dR.
\tag{RF40}
\]
Beyond those velocity bump supports, \(E_n=U_n=0\) and
\(F_n(X,\eta)=\int_0^XU_n(x,\eta)dx=0\). Their radial profile
is the exact formula
\[
V_n=\frac{2\eta XU_n-2\eta(D+2nh)F_n-(1-\eta^2)\partial_\eta F_n}
{L_q}=0.
\tag{RF41}
\]
The zero function \(F_n\) is zero for every parameter on this open
interval, so its parameter derivatives vanish as well. Notice the
coefficient \(D+2nh\), which is retained even though the evaluated
terms vanish.

The actual Borel cutoff is applied to the divergence-free potential.
Its exact radial and axial components are
\[
\begin{split}
ru_{r,n}^{\rm cut}
&=q^{2nh}\left[\chi(c_nq)V_n-
\frac{2\eta}{L_q}(c_nq)\chi'(c_nq)F_n\right],\\
u_{z,n}^{\rm cut}&=q^{-A+2nh}\chi(c_nq)U_n,
\quad u_{\theta,n}^{\rm cut}=q^{-A+2nh}\chi(c_nq)E_n.
\end{split}
\tag{RF42}
\]
The minus sign comes from \(ru_r=-\partial_z(q^{1-A+2nh}F_n)\).
On the entire mean interval all three components in (RF42) vanish
identically, including the derivative-of-cutoff term. Summing any
number of these zero functions, or the locally finite Borel sum, changes
none of the leading velocities there. Positive-order pressure is not
asserted zero on this interval, and is not used in the five matrix rows.
This proves the exact passage from the leading law (RF28) to the summed
base law (RF29). It also explains why asymptotic closeness alone would
have been insufficient for the advertised exact matrix.

The just-proved radial identity and the vanishing of every
positive-order cutoff contribution give the exact chart law
\[
b=Q^Ab_{\rm phys}
=-\epsilon\,\frac{\mathfrak m}{R},\qquad
V=a(\eta)s^{\lambda-h}R^{-1-2\lambda},\qquad G=0.
\tag{RF43}
\]
The swirl formula follows from the exact coordinate composition,
since \(-A+(1+2\lambda)/2=\lambda-h\).
In particular
\(\partial_R^jb=(-1)^{j+1}\epsilon\mathfrak m\,j!R^{-j-1}\);
every coefficient derivative of \(b\) involving \(Z,T\) or the torus
is zero. The positive lower radial bound proves
\(\partial^Ib=O_I(\epsilon)\) there, while (RF3), (RF43) give
\(\partial^IV,\partial^IG=O_I(1)\). These estimates hold on the
whole induced-potential interval, including its gaps. The constants
depend on the fixed profile and requested derivative order. The
construction parameters require no choice depending on the band,
derivative order or correction stage.

\subsection{Exact nonlinear recomputation after the five rows}

Keep the base and the actual wave \(w\), including its complete
covariance \(W\), fixed in this step. The full radial source is
\[
\begin{split}
g_r={}&-t_*\beta-(D_r+1/R)(2b\beta+\beta^2+W^{rr})\\
&-D_z(b\gamma+G\beta+\beta\gamma+W^{zr})
+(2Vv+v^2+W^{\theta\theta})/R
+\epsilon(\Delta_0-R^{-2})\beta.
\end{split}
\tag{RF44}
\]
Write \(B=\Delta\beta,u=\Delta v,d=\Delta\gamma=\gamma_d\)
for the actual increments just constructed. Expanding each product
with no deletion gives
\[
g_{r,\rm new}-g_r=(2V/R)u+R_g,
\tag{RF45}
\]
where
\[
\begin{split}
R_g={}&-t_*B-(D_r+1/R)(2bB+2\beta B+B^2)\\
&-D_z(bd+GB+\beta d+\gamma B+Bd)
+(2vu+u^2)/R+\epsilon(\Delta_0-R^{-2})B.
\end{split}
\tag{RF46}
\]
For example \((\beta+B)^2-\beta^2=2\beta B+B^2\),
\((\beta+B)(\gamma+d)-\beta\gamma=\beta d+\gamma B+Bd\),
and \((v+u)^2-v^2=2vu+u^2\). These account for every nonlinear
term of (RF44); \(W\) is unchanged. Integrating (RF45), and
expanding the products in (RF30), the five exact rows imply
\[
\begin{split}
P_{\rm new}&=\int\overline{R_g}\,dR,\\
(J_\theta)_{\rm new}&=\int R^2\overline{\gamma u+vd+du}\,dR,\\
(J_z)_{\rm new}&=\int R\overline{2\gamma d+d^2}\,dR
-\frac12\int R^2\overline{R_g}\,dR.
\end{split}
\tag{RF47}
\]
In the final row the linear pressure moment is exactly
\(-\tfrac12\int R^2(2V/R)u\,dR=-\int RVu\,dR\),
which explains the retained sign in the fifth matrix row. These formulas
cancel the base-linear terms; they do not assert the new defects are zero.

For the Section 9 budgets
\(v,\gamma\in\mathcal M_{0.9}\),
\(\beta\in\mathcal M_{1.9}\), and target order \(\alpha\ge0.9\),
(RF39) and (RF25) give
\(u,d\in\mathcal M_\alpha\),
\(B\in\mathcal M_{\alpha+1}\). Every increment is torus
independent. With the exact local base bounds just proved, the
individual guaranteed exponents in (RF46) are
\[
\begin{array}{c|c@{\qquad}c|c}
\text{term}&\text{exponent}&\text{term}&\text{exponent}\\\hline
-t_*B&\alpha+2&-(D_r+1/R)(2bB)&\alpha+2-\kappa_s\\
-(D_r+1/R)(2\beta B)&\alpha+2.9-\kappa_s&
-(D_r+1/R)B^2&2\alpha+2-\kappa_s\\
-D_z(bd)&\alpha+2&-D_z(GB)&\alpha+2\\
-D_z(\beta d)&\alpha+2.9&-D_z(\gamma B)&\alpha+2.9\\
-D_z(Bd)&2\alpha+2&2vu/R&\alpha+0.9\\
u^2/R&2\alpha&\epsilon D_r^2B&\alpha+2-2\kappa_s\\
\epsilon R^{-1}D_rB&\alpha+2-\kappa_s&\epsilon D_z^2B&\alpha+4\\
-\epsilon R^{-2}B&\alpha+2&&
\end{array}
\tag{RF48}
\]
The first row uses \(t_*B=-\epsilon\partial_TB\), so its
coefficient \(cN\) is exactly absent. Every listed exponent exceeds
or equals \(\alpha+0.9-2\kappa_s\): for \(2\alpha\) the
difference is \(\alpha-0.9+2\kappa_s\ge2\kappa_s>0\),
and subtraction proves the other entries directly. Thus
\[
R_g\in\mathcal M_{\alpha+0.9-2\kappa_s},\qquad
(P_{\rm new},(J_\theta)_{\rm new},(J_z)_{\rm new})
\in(\mathcal S_{\alpha+0.9-2\kappa_s})^3.
\tag{RF49}
\]
For the two product expressions in (RF47), the exponents are
\(\alpha+0.9\) and \(2\alpha\), before the source moment is
added. This proves the stated bound for each defect and all fixed
derivatives. The first two matrix rows and the exact solenoidal
realization preserve both original zero moments of the velocity.

\subsection{Use in the initial state and the fixed-domain conclusion}

The maps proved here apply to the actual state after temporal correction.
They require its computed defects as inputs; those defects include every
product with the flat axial remnant \(a\). The fixed prescription is
to recompute (RF44), then \(P\) and (RF24), then (RF30), using
the updated velocity at every step. The five-equation step preserves
the exact two zero moments and divergence, and gives the defect gain
(RF49). At the initial target order \(H_0=1-\kappa_s\), the
gain furnished by this calculation is
\[
H_0+0.9-2\kappa_s=1.9-3\kappa_s=1.89997>1.2=C_0.
\tag{RF50}
\]
The increments also satisfy
\(H_0=0.99999>0.9\), \(H_0+1=1.99999>1.9\),
so they preserve the retained velocity class budgets. These exact
statements concern the mean map. The estimates supplying the initial
wave covariance and the tangential mean residual are separate
calculations and are not inferred from (RF50).

For fixed source profile parameters \(h,\lambda,c_{\rm patch}\),
fixed shell and bump supports, fixed band comparison and covering
bounds, all conclusions above hold on the single inherited positive-
\(q\) construction domain. No inverse in this fragment imposes a
derivative-order dependent threshold or a uniform \(\lambda\to0\)
limit. All-orders flatness is the estimate (RF22) on that same domain.
The exact obstruction (RF15) proves why the flat projector remainder
cannot be replaced by zero, and the rank loss at \(\lambda=0\)
proves why the fixed positive power-law parameter must be retained.


\subsection{Signed covariance and original positive-amplitude audit}


The covariance matrix, positivity mechanism, exact signed differential inverse, shell weights, and displayed increment exponents are correct. The physical assembly statement at source line 403 needs its domain stated: the physical stress must have admissible local mean-class representatives with the prescribed active-shell support and edge vanishing. An arbitrary auxiliary-independent physical stress is too broad a domain. The precise counterexample and corrected wording appear below.

\subsubsection{1. Original objects and input provenance}

The source uses

\[
A=\tfrac12+h,\quad Q=2^{-\ell},\quad \varepsilon=Q^h,
\quad S_*=\ell^2,\quad 0<h<1/100,\quad \kappa_s=10^{-5}.
\]

The normalized angular average has measure \(d\theta/(2\pi)\); the auxiliary torus has normalized Haar measure. For real extended vector fields,

\[
\mathcal C(w)=\langle\langle w_rw_{\tan}\rangle_\theta\rangle_Y,
\qquad
\mathcal B(w,v)=\langle\langle w_rv_{\tan}+v_rw_{\tan}\rangle_\theta\rangle_Y,
\quad w_{\tan}=(w_\theta,w_z).
\]

Every entry of these maps is a product followed by the two specified linear integrations. Expanding each product gives

\[
\mathcal C(w+v)=\mathcal C(w)+\mathcal B(w,v)+\mathcal C(v),
\qquad D\mathcal C(w)[v]=\mathcal B(w,v).
\]

In particular, the derivative contains both cross terms and therefore contains a factor of two on a scalar multiple of one wave.

The following inputs are explicitly present outside the assigned interval:

\begin{itemize}

\item
  Source lines 55--90 construct disjoint auxiliary preimages for every pair of simultaneously active distinct labels, including the two signs, and prove preservation of normalized Haar measure under the covering maps. This is the support and measure input used below.
\item
  Source lines 121--143 define the original frozen vectors \(N,K\), the original negative \(c_0\), and the uniform strict target-direction cone margin. Source lines 673--695 identify the profile endpoint-direction and profile-weight construction imported from the profiles lane. This audit does not independently construct those profile factors.
\item
  Source lines 233--256 and 278--292 supply the two-sided Gaussian bounds for \(P\), the positive radial homogeneous amplitude \(x\), and their derivative bounds. The calculation below gives the complete covariance consequence of these inputs.
\item
  Source lines 94--108 define the shell weight, derivative variables, and coefficient spaces, with all derivative bounds required pointwise. Source lines 417--453 give the exact curl remainder used in the final increment calculation.
\end{itemize}

In the original tangential component order,

\[
N=(N_\theta,N_z),\qquad K=(-N_z,N_\theta),\qquad |N|=1.
\]

Thus \(N\cdot K=0\), \(|K|=1\), and

\[
\det[N\mid K]=N_\theta^2+N_z^2=1.
\]

Consequently, expressing the covariance columns in the ordered basis \((N,K)\) preserves the determinant sign in the original \((\theta,z)\) order. It is a coordinate expression of the same columns, not a replacement or renormalization of them.

The coefficient derivatives used throughout are, for a five-component multi-index \(I\),

\[
D^I=\partial_R^{I_1}\partial_Z^{I_2}\partial_T^{I_3}
\partial_{H_1}^{I_4}\partial_{H_2}^{I_5},
\]

where \((H_1,H_2)\) are coordinates on the common auxiliary torus; the covariance matrix \(H\) defined below is a different object. Every label, band scale, and discrete harmonic is fixed under these derivatives. For the original radial shell variable \(X\), define

\[
\zeta=\exp\left(-a_a/\log^2(X/X_a)-a_b/\log^2(X_b/X)\right),
\qquad \delta=\min\{1,\log(X/X_a),\log(X_b/X)\},
\]

on \(X_a<X<X_b\), with \(a_a,a_b>0\) the original constants, and extend \(\zeta\) by zero. The mean class \(\mathcal M_\alpha\) consists of compatible angular-independent coefficients with smooth zero extension outside this shell and, for every \(I\),

\[
|D^If|\le C_I\varepsilon^\alpha S_*^{b_I}\zeta\delta^{-d_I}.
\]

It has no auxiliary-rectangle or pulse-envelope requirement and may depend on the auxiliary torus. Its auxiliary-independent subclass is denoted \(\mathcal M_\alpha^{\rm independent}\). The wave class \(\mathcal W_\alpha\) consists of the nonzero-harmonic coefficients \(w_{\gamma,m}\) in finite labelled sums \(w_{\gamma,m}e^{ikm\Phi_\gamma}\), with the source's prescribed closed slow, transverse-rectangle, and shell supports, satisfying

\[
|D^Iw_{\gamma,m}|\le C_I\varepsilon^\alpha S_*^{b_I}
\sqrt\zeta\delta^{-d_I}P(v).
\]

The actual cutoff coefficients additionally have compact pulse support and smooth zero extension at that support. Equal label-harmonic contributions are collected before testing these bounds. The finite constants and nonnegative exponents \(C_I,b_I,d_I\) may depend on the derivative order and the fixed finite collection under consideration, but are uniform in the band. Compatibility means equality after the specified torus covering pullbacks on chart overlaps. These are coefficient estimates with the exponential factored out; every coefficient derivative is required to satisfy the stated pointwise envelope bound. No conclusion below differentiates a bare inequality for \(P\).

\subsubsection{2. Exact column integral and its constants}

For each sign of one slow box, retain

\[
b_\sigma=\chi_g(\xi_g)\psi(v)t_\sigma^h\cos(k\Phi_\sigma),
\qquad H=[\mathcal C(b_+)\mid\mathcal C(b_-)].
\]

The integer \(kp\) is nonzero. For every non-angular phase offset \(c\),

\[
\frac1{2\pi}\int_0^{2\pi}\cos^2(kp\theta+c)\,d\theta
=\frac12+
\left[\frac{\sin(2kp\theta+2c)}{8\pi kp}\right]_0^{2\pi}
=\frac12.
\]

The rectangle map is

\[
(\xi_g,\eta_g)\longmapsto c_\gamma+\xi_gv_r+\eta_gv_t,
\quad v_r=(1,-b_g),\quad v_t=(b_g,1),\quad b_g=\sqrt2-1.
\]

Its absolute Jacobian is \(1+b_g^2\); \(\eta_g=c_iv-r_0\) gives \(d\eta_g=c_i\,dv\). Haar preservation means no additional covering degree multiplies or divides the integral. The homogeneous amplitude is independent of \(\xi_g\). Writing its radial component as \(x\), direct substitution therefore gives

\[
H_\sigma=D_gc_i\int_0^{L_s}\psi(v)^2x(v)t_{\sigma,\tan}^h(v)\,dv,
\] \[
h_\sigma=D_gc_i\int_0^{L_s}\psi(v)^2x(v)^2\,dv,
\qquad
D_g=\frac{1+b_g^2}{2}\int_{\mathbb R}\chi_g(\xi)^2\,d\xi>0,
\qquad L_s=\frac{2r_0}{c_i}.
\]

All four original factors---the angular half, rectangle Jacobian, transverse integral, and \(c_i\,dv\)---are retained.

Write the proved pulse bounds with positive uniform constants as

\[
c_xP(v)\le x(v)\le C_xP(v),\qquad
e^{-C_P(v-L_s/2)^2/L_s}\le P(v)
\le e^{-c_P(v-L_s/2)^2/L_s}.
\]

Put \(d=v-L_s/2\). On the central interval \(|d|\le\sqrt{L_s}\), the cutoff equals one when \(L_s\ge25\), since the source cutoff is one for \(|d|<L_s/5\). Thus

\[
I_0:=\int_0^{L_s}\psi^2x^2\,dv
\ge 2c_x^2e^{-2C_P}\sqrt{L_s}=:C_{0,-}\sqrt{L_s}.
\]

Using \(0\le\psi\le1\) and integrating the upper Gaussian over the whole real line gives

\[
I_0\le C_x^2\sqrt{\frac{\pi}{2c_P}}\sqrt{L_s}
=:C_{0,+}\sqrt{L_s}.
\]

The first absolute moment is bounded by

\[
I_1:=\int_0^{L_s}\frac{|d|}{L_s}\psi^2x^2\,dv
\le\frac{C_x^2}{L_s}\int_{\mathbb R}|d|e^{-2c_Pd^2/L_s}\,dd
=\frac{C_x^2}{2c_P}.
\]

The source floor estimate is

\[
\frac{1}{T_gS_*}<c_i\le\frac1{S_*},\qquad T_g=4+\sqrt2.
\]

Hence \(L_s\ge2r_0S_*\), so \(S_*\ge25/(2r_0)\) suffices for the central-interval bound. Also

\[
c_i\sqrt{L_s}=\sqrt{2r_0c_i}.
\]

Therefore the positive original \(h_\sigma\) obey

\[
h_-^{\rm bd}S_*^{-1/2}\le h_\sigma\le h_+^{\rm bd}S_*^{-1/2},
\] \[
h_-^{\rm bd}=D_gC_{0,-}\sqrt{2r_0/T_g}>0,
\qquad h_+^{\rm bd}=D_gC_{0,+}\sqrt{2r_0}.
\]

The strict floor inequality also implies this non-strict lower bound. The symbols \(h_-^{\rm bd},h_+^{\rm bd}\) denote bounds, not a replacement for the two different original numbers \(h_+,h_-\).

\subsubsection{3. Column directions and the exact determinant}

The original frozen constants satisfy \(c_0<0\), \(u_*>0\), and

\[
A_c=-c_0\sqrt{1+u_*^2}>0.
\]

The homogeneous polarization satisfies the uniform error bound stated in source line 334 relative to

\[
f(s)=-sK+c_0\sqrt{1+s^2}N,
\qquad s(v)=\sigma(u_*/2+u_*v/L_s).
\]

Here

\[
f'(s)=-K+c_0\frac{s}{\sqrt{1+s^2}}N,
\qquad |f'(s)|\le\sqrt{1+c_0^2},
\qquad |s(v)-\sigma u_*|=u_*|d|/L_s.
\]

The mean-value integral for \(f(s(v))-f(\sigma u_*)\), followed by the probability measure \(\psi^2x^2\,dv/I_0\), bounds its average by

\[
\sqrt{1+c_0^2}\,u_*\frac{I_1}{I_0}
\le \sqrt{1+c_0^2}\,u_*\frac{C_x^2}{2c_PC_{0,-}\sqrt{L_s}}.
\]

Adding the source's uniform \(O(S_*^{-1})\) polarization error gives the asserted uniform \(O(S_*^{-1/2})\) column error without discarding any sign:

\[
H_\sigma=h_\sigma(-A_cN-\sigma u_*K+e_\sigma),
\qquad |e_\sigma|\le C_eS_*^{-1/2}.
\]

Set \(e_{\sigma,N}=e_\sigma\cdot N\), \(e_{\sigma,K}=e_\sigma\cdot K\). In \((N,K)\) coordinates the full, perturbed original matrix is

\[
H=\begin{pmatrix}
(-A_c+e_{+,N})h_+&(-A_c+e_{-,N})h_-\\
(-u_*+e_{+,K})h_+&(u_*+e_{-,K})h_-
\end{pmatrix}.
\]

Its exact determinant is

\[
\det H=h_+h_-\Delta,
\] \[
\begin{aligned}
\Delta={}&(-A_c+e_{+,N})(u_*+e_{-,K})
-(-A_c+e_{-,N})(-u_*+e_{+,K})\\
={}&-2A_cu_*+A_c(e_{+,K}-e_{-,K})
+u_*(e_{+,N}+e_{-,N})\\
&\hspace{1cm}+e_{+,N}e_{-,K}-e_{-,N}e_{+,K}.
\end{aligned}
\]

For a target \(T=T_NN+T_KK\), the exact inverse is

\[
(H^{-1}T)_+=\frac{(u_*+e_{-,K})T_N+(A_c-e_{-,N})T_K}{h_+\Delta},
\] \[
(H^{-1}T)_-=\frac{(u_*-e_{+,K})T_N+(-A_c+e_{+,N})T_K}{h_-\Delta}.
\]

At zero error these formulas become exactly

\[
h_+(H_{\rm ref}^{-1}T)_+=\tfrac12(-T_N/A_c-T_K/u_*),
\quad
h_-(H_{\rm ref}^{-1}T)_-=\tfrac12(-T_N/A_c+T_K/u_*).
\]

In particular, \(\det H_{\rm ref}=-2A_cu_*h_+h_-\), with the same sign in the original tangential coordinates.

\subsubsection{4. A uniform threshold that proves positivity of the original amplitudes}

The strict compact direction gap in source lines 141--143 supplies a number \(\mu<1\). With the original choice

\[
\frac{u_*}{\sqrt{1+u_*^2}}>\mu,
\]

the unfrozen direction ratio has a strict margin below \(u_*/\sqrt{1+u_*^2}\). Uniform continuity of the extended normalized target direction, and the fixed mesh diameter \(O(S_*^{-3})\), preserves a strictly smaller margin for the frozen box basis. Hence one fixed \(\eta_c>0\) satisfies on every enlarged box after the stated initial band threshold

\[
T_N<0,\qquad |T_K|/u_*\le(1-\eta_c)(-T_N)/A_c,
\qquad T=T_{0,*}.
\]

No division by \(T\) at a zero edge is used: the compact direction is extended first, and the inequalities for nonzero \(T\) are applied in the open shell.

To make the perturbation threshold fully quantitative, let \(A_{\min}\le A_c\le A_{\max}\) be the positive uniform bounds from source line 135 and define

\[
H_0=\begin{pmatrix}-A_c&-A_c\\-u_*&u_*\end{pmatrix},
\qquad \mathcal E=[e_+\mid e_-]_{(N,K)},
\qquad D_h=\operatorname{diag}(h_+,h_-).
\]

These identities retain \(H=(H_0+\mathcal E)D_h\) and the original \(H\); they only factor its two already defined scalar column weights. Direct multiplication gives

\[
H_0H_0^T=\operatorname{diag}(2A_c^2,2u_*^2),
\qquad
\|H_0^{-1}\|\le M:=\frac1{\sqrt2\min(A_{\min},u_*)}.
\]

Choose \(C_E\) with \(\|\mathcal E\|\le C_ES_*^{-1/2}\); for example \(C_E=\sqrt2C_e\) is valid in the Euclidean operator norm. For every vector \(z\),

\[
|(I+H_0^{-1}\mathcal E)z|\ge(1-MC_ES_*^{-1/2})|z|.
\]

Thus \(MC_ES_*^{-1/2}\le1/2\) proves invertibility and

\[
\|(H_0+\mathcal E)^{-1}\|\le 2M.
\]

This proves the needed estimate directly, with no unproved appeal to a ``finite Neumann series.'' Along the real path \(I+tH_0^{-1}\mathcal E\), \(0\le t\le1\), the same lower norm bound prevents a zero determinant, so its determinant remains positive. Both singular values are at least \(1/2\). Consequently

\[
\Delta<0,\qquad |\Delta|\ge \frac{|\det H_0|}{4}=\frac{A_cu_*}{2},
\] \[
|\det H|\ge\frac{A_{\min}u_*}{2}(h_-^{\rm bd})^2S_*^{-1},
\qquad
\|H^{-1}\|\le\frac{2M}{h_-^{\rm bd}}\sqrt{S_*}.
\]

For coefficient positivity retain \(y=H^{-1}T_{0,*}\) and introduce only the comparison coordinates \(c=D_hy\), \(c^{\rm ref}=H_0^{-1}T_{0,*}\). Put \(p=-T_N/A_c>0\), \(q=T_K/u_*\). Then \(|q|\le(1-\eta_c)p\), and

\[
c_+^{\rm ref}=\tfrac12(p-q),\quad c_-^{\rm ref}=\tfrac12(p+q),
\quad c_\sigma^{\rm ref}\ge\tfrac{\eta_c}{2}p.
\]

Moreover

\[
|T_{0,*}|^2=A_c^2p^2+u_*^2q^2
\le [A_{\max}^2+(1-\eta_c)^2u_*^2]p^2.
\]

Hence, with

\[
m_0=\frac{\eta_c}{2\sqrt{A_{\max}^2+(1-\eta_c)^2u_*^2}}>0,
\]

we have \(c_\sigma^{\rm ref}\ge m_0|T_{0,*}|\). The exact inverse difference identity gives

\[
c-c^{\rm ref}=-(H_0+\mathcal E)^{-1}\mathcal E H_0^{-1}T_{0,*},
\quad |c-c^{\rm ref}|\le 2M^2C_ES_*^{-1/2}|T_{0,*}|.
\]

Therefore one fixed threshold

\[
\sqrt{S_*}\ge
\max\left\{2MC_E,\frac{4M^2C_E}{m_0}\right\}
\]

in addition to the already described pulse and cone thresholds proves

\[
\frac{m_0}{2h_+^{\rm bd}}\sqrt{S_*}|T_{0,*}|
\le y_\sigma
\le\frac{2M}{h_-^{\rm bd}}\sqrt{S_*}|T_{0,*}|.
\]

The original amplitudes are exactly

\[
a_\sigma=\sqrt{y_\sigma}>0
\]

on the open shell. They have not been divided by their size, replaced with unit vectors, or signed to absorb a correction. At each shell edge the target and extended amplitudes are zero. The threshold uses only fixed value bounds and direction margins, not the later derivative order or signed correction.

The original local wave at which the covariance is differentiated is
\[
W_0=\sqrt\varepsilon\sum_{\sigma=\pm}a_\sigma b_\sigma.
\]

\subsubsection{5. Derivatives and smooth zero extension with the original shell weight}

Retain the source weight

\[
\zeta(X)=\exp\left(-\frac{a_a}{\log^2(X/X_a)}-
\frac{a_b}{\log^2(X_b/X)}\right),
\quad \delta=\min\{1,\log(X/X_a),\log(X_b/X)\},
\]

with its original positive constants. The profile-specific source choice is \(a_a=t_1^2\), \(a_b=4\); the argument applies to precisely those constants. The profile bounds used here are

\[
|T_{0,*}|\ge c_T\zeta,
\qquad |D^IT_{0,*}|\le C_I\zeta\delta^{-d_I}.
\]

The source's chart/profile map has bounded derivatives on the dyadic shell and a positive bounded chart multiplier, so those physical profile bounds have the displayed chart form. All differentiations here keep the label and \(Q,\varepsilon,S_*,c_i,L_s\) fixed. They therefore do not differentiate a discrete frequency or a moving endpoint.

Each derivative of \(H\) is polynomially bounded in \(S_*\) by differentiation of the fixed finite pulse integral. To see that this gives every derivative of its inverse with no new smallness condition, differentiate \(HH^{-1}=I\). For a nonzero multi-index \(I\),

\[
D^IH^{-1}=-H^{-1}\sum_{0<J\le I}\binom IJ(D^JH)(D^{I-J}H^{-1}).
\]

Induction on \(|I|\), the already established value bound on \(H^{-1}\), and the finite sum prove a polynomial bound at every fixed order. Applying the finite Leibniz expansion to \(H^{-1}T_{0,*}\) gives

\[
y_\sigma\ge c\sqrt{S_*}\zeta,
\qquad |D^Iy_\sigma|\le C_IS_*^{b_I}\zeta\delta^{-d_I}.
\]

For completeness, for either exponent \(\nu=1/2\) or \(\nu=-1/2\), the exact multivariate chain-rule formula on the open shell is

\[
D^I(y_\sigma^\nu)=I!\sum_{\substack{(k_J)_{0<J\le I}\\
\sum_J k_JJ=I}}
(\nu)_{\underline n}\,y_\sigma^{\nu-n}
\prod_{0<J\le I}\frac1{k_J!}
\left(\frac{D^Jy_\sigma}{J!}\right)^{k_J},
\qquad n=\sum_Jk_J,
\]

where \((\nu)_{\underline n}=\nu(\nu-1)\cdots(\nu-n+1)\) and the empty product for \(I=0\) is one. This formula follows by applying the one-variable Taylor expansion to \(y(x+h)-y(x)\) and extracting the \(h^I\) coefficient; only degrees at most \(|I|\) can contribute. It therefore has finitely many terms at each fixed order.

For \(\nu=1/2\) and \(|I|>0\), \(n\ge1\); the factor \(y^{1/2-n}\) has a negative exponent and is bounded by the established lower bound on \(y\). The product of the \(n\) derivative factors contributes \(\zeta^n\), while that inverse power contributes \(\zeta^{1/2-n}\). The result is exactly \(\sqrt\zeta\). For \(I=0\), the upper value bound on \(y\) supplies the same weight, with permitted inverse-distance powers. Thus

\[
|D^Ia_\sigma|\le C_IS_*^{b'_I}\sqrt\zeta\,\delta^{-d'_I}.
\]

For \(\nu=-1/2\), the same calculation, including \(I=0\), gives

\[
|D^I(a_\sigma^{-1})|=|D^I(y_\sigma^{-1/2})|
\le C_IS_*^{b''_I}\zeta^{-1/2}\delta^{-d''_I}.
\]

The inverse amplitude itself is not claimed to extend smoothly; the product with an admissible target will extend.

The flatness statement is rigorous: at either edge, \(\sqrt\zeta\,\delta^{-N}\) tends to zero faster than every positive power of the edge distance, because its logarithm has a negative term of order \(-c/\delta^2\), which dominates every \(N\log(1/\delta)\). The same holds for all the displayed derivatives. In local smooth coordinates consisting of edge distance and tangential variables, the zero-extended value is continuous. Its difference quotient normal to the edge tends to zero by the faster-than-linear value bound; its tangential derivative equals zero on the edge. Repeating this argument for each already identified derivative proves by induction that the zero extension is smooth and all its jets vanish. Smooth chart changes preserve this statement. This proves the shell extension claimed for \(a_\sigma\) and for the signed coefficients below.

\subsubsection{6. Exact signed differential inverse, including its full quadratic remainder}

The input is exactly the source's real, angular- and auxiliary-independent two-vector \(\Sigma(R,Z,T)\) in \(\mathcal M_\alpha\), with

\[
|D^I\Sigma|\le C_I\varepsilon^\alpha S_*^{b_I}\zeta\delta^{-d_I}.
\]

Its values may have arbitrary signs and arbitrary size within that class. Keep the already constructed original \(a_\sigma>0\) fixed and define

\[
d_\Sigma=H^{-1}(\Sigma/\varepsilon),\qquad
\delta a_\sigma=\frac{(d_\Sigma)_\sigma}{2a_\sigma},\qquad
\mathcal L\Sigma=\sqrt\varepsilon\sum_{\sigma=\pm}\delta a_\sigma b_\sigma.
\]

The inverse-matrix derivative bounds and Leibniz rule give

\[
|D^Id_\Sigma|\le C_I\varepsilon^{\alpha-1}S_*^{b_I}\zeta\delta^{-d_I}.
\]

Multiplication with the proved \(a_\sigma^{-1}\) estimates gives

\[
|D^I\delta a_\sigma|\le C_I\varepsilon^{\alpha-1}S_*^{b'_I}
\sqrt\zeta\delta^{-d'_I}.
\]

The inverse power of the original amplitude cancels exactly one half-power of the target shell weight. Multiplying by \(\sqrt\varepsilon\) and the original pulse coefficient, whose derivatives have polynomial factors times \(P(v)\), proves

\[
\mathcal L:\mathcal M_\alpha^{\rm independent}\longrightarrow
\mathcal W_{\alpha-1/2},
\]

with the prescribed smooth shell extension, transverse support, pulse cutoff, and later slow cutoff. Both maps in this statement refer to the fixed source chart and labels. Auxiliary independence is required to take \(\delta a_\sigma\) outside the torus integral; angular independence is required for the same angular calculation.

Because the two sign supports are disjoint,

\[
\mathcal C(W_0)=\varepsilon\sum_\sigma a_\sigma^2H_\sigma
=\varepsilon H(y_+,y_-)^T=\varepsilon T_{0,*},
\] \[
\mathcal B(W_0,\mathcal L\Sigma)
=\varepsilon\sum_\sigma2a_\sigma\delta a_\sigma H_\sigma
=\varepsilon H d_\Sigma=\Sigma.
\]

At shell edges both sides have their proved zero extension, so the identity persists there. The map is linear because \(H,a_\sigma,b_\sigma\) are fixed before \(\Sigma\) is chosen. No positivity or smallness of \(\Sigma\) was used in any equality.

The precise nonlinear remainder can also be written without an unspecified symbol. Set \(z=H^{-1}\Sigma\). Then

\[
\mathcal C(\mathcal L\Sigma)
=\frac1{4\varepsilon}H
\begin{pmatrix}z_+^2/y_+\\z_-^2/y_-\end{pmatrix}.
\]

Thus

\[
\mathcal C(W_0+\mathcal L\Sigma)
=\varepsilon T_{0,*}+\Sigma+
\frac1{4\varepsilon}H
\begin{pmatrix}(H^{-1}\Sigma)_+^2/y_+\\(H^{-1}\Sigma)_-^2/y_-\end{pmatrix}.
\]

The remainder has \(\varepsilon\) exponent \(2\alpha-1\); the product of the two wave weights is \(\zeta P^2\), and \(0\le P\le1\). After averaging it therefore belongs to \(\mathcal M_{2\alpha-1}\) with all derivatives, by the finite Leibniz rule. The exact signed inverse is an inverse of the derivative at the fixed original wave; it does not remove this nonzero quadratic term.

\subsubsection{7. Physical assembly, chart maps, and the missing domain qualification}

The original chart relation is

\[
T_{0,*}=(Q/q)^{A+1/2}T_0.
\]

Keeping every physical scaling factor gives

\[
\begin{aligned}
Q^{-2A}\varepsilon T_{0,*}
&=Q^{-2A+h}(Q/q)^{A+1/2}T_0\\
&=Q^{-A+h+1/2}q^{-A-1/2}T_0\\
&=q^{-A-1/2}T_0,
\end{aligned}
\]

because \(-A+h+1/2=0\). No equality of distinct \(Q_\beta\) is used.

For a physical stress with admissible local representatives, retain the assembled fields exactly as in the source:

\[
W_{0,\mathrm{phys}}^{\rm as}
=\sum_\beta Q_\beta^{-A}\eta_\beta W_{0,\beta},
\] \[
\mathcal L_{\mathrm{phys}}^{\rm as}\sigma
=\sum_\beta Q_\beta^{-A}\eta_\beta
\mathcal L_\beta(Q_\beta^{2A}\sigma).
\]

At a fixed physical slow point there are finitely many active labels with uniform bounded overlap. Products from distinct active labels vanish by the already proved common-torus support disjointness. The slow \(\eta_\beta\) and \(Q_\beta\) factors are independent of both averaging variables. Therefore the exact contribution of one label is

\[
Q_\beta^{-2A}\eta_\beta^2
\mathcal B(W_{0,\beta},\mathcal L_\beta(Q_\beta^{2A}\sigma))
=Q_\beta^{-2A}\eta_\beta^2Q_\beta^{2A}\sigma
=\eta_\beta^2\sigma.
\]

On the region covered by the squared partition, \(\sum_\beta\eta_\beta^2=1\). Outside the active shell an admissible stress is zero, as are the extended wave fields. These two facts establish the global identity on its proper domain:

\[
\mathcal B(W_{0,\mathrm{phys}}^{\rm as},
\mathcal L_{\mathrm{phys}}^{\rm as}\sigma)=\sigma.
\]

In a reference chart, the exact representative maps are

\[
W_0^{\rm as}=Q^AW_{0,\mathrm{phys}}^{\rm as},
\quad
\mathcal L^{\rm as}\Sigma
=Q^A\mathcal L_{\mathrm{phys}}^{\rm as}(Q^{-2A}\Sigma).
\]

Bilinearity gives

\[
\mathcal B(W_0^{\rm as},\mathcal L^{\rm as}\Sigma)
=Q^{2A}(Q^{-2A}\Sigma)=\Sigma.
\]

For overlapping bands the chart maps are exactly

\[
R_\beta=(Q/Q_\beta)^{1/2}R,\quad
Z_\beta=(Q/Q_\beta)^DZ,\quad
T_\beta=(Q/Q_\beta)T,
\quad \varepsilon_\beta=(Q_\beta/Q)^h\varepsilon.
\]

Their constants and inverses are uniformly bounded on the allowed band range, as are \(S_\beta/S_*\) and covering multiplicities. The shell variable \(X\) and its weights describe the same physical point. Each fixed derivative of these affine scale changes introduces bounded factors, and derivatives of the fixed mesh cutoffs add polynomial powers of \(S_*\). Finite sums and prior collection of equal label-harmonic contributions therefore preserve every stated coefficient exponent.

\textbf{Necessary qualification at source line 403.}

The phrase ``For a physical auxiliary-independent stress sigma(r,z,t)'' must retain the preceding local input class and support. Taken literally for every such physical stress, it is false. Choose any nonzero smooth two-vector stress with compact support in an open set outside the closed active shell. It is auxiliary-independent. Every assembled wave in the displayed formulas vanishes there, so its bilinear covariance is zero there, while the chosen stress is nonzero. Thus the asserted right-inverse identity cannot hold on that enlarged domain. Furthermore an arbitrary nonvanishing stress approaching a shell edge need not cancel \(1/a_\sigma\) and need not give any smooth extension.

A precise source replacement is:

\begin{quote}
For a real, auxiliary-independent physical stress supported in the prescribed active shell whose band representatives \(Q_\beta^{2A}\sigma\) satisfy the local \(\mathcal M_\alpha\) bounds and smooth zero extension used above, use the assembled maps below.
\end{quote}

This repairs the domain while retaining all original formulas. The stated local inverse at line 383 and actual increment at line 474 already include the mean-class restriction. No additional sign or smallness requirement should be inserted.

\subsubsection{8. Actual divergence-free covariance increment and every exponent}

Retain the source hypotheses and fields:

\[
U=W_0^{\rm as}+E,\qquad U\in\mathcal W_{1/2},\quad E\in\mathcal W_\rho,
\qquad \Sigma\in\mathcal M_\alpha^{\rm independent}.
\]

The exact curl construction applied to the original \(\mathcal L^{\rm as}\Sigma\) produces

\[
V=\mathcal L^{\rm as}\Sigma+R,
\qquad R\in\mathcal W_{\alpha-\kappa_s},
\]

since the input exponent is \(\alpha-1/2\) and the curl correction gains \(1/2-\kappa_s\). With \(L=\mathcal L^{\rm as}\Sigma\) only as a temporary notation for this unchanged field, expansion of all products gives

\[
\begin{aligned}
\mathcal C(U+V)-\mathcal C(U)
&=\mathcal B(U,V)+\mathcal C(V)\\
&=\mathcal B(W_0^{\rm as}+E,L)+\mathcal B(U,R)+\mathcal C(V)\\
&=\Sigma+\mathcal B(E,L)+\mathcal B(U,R)+\mathcal C(V).
\end{aligned}
\]

The first remainder includes precisely the interaction of the already present error with the uncurled signed inverse. The second retains its interactions with both \(W_0^{\rm as}\) and \(E\), because \(U\) is not replaced. The third retains all three quadratic pieces

\[
\mathcal C(V)=\mathcal C(L)+\mathcal B(L,R)+\mathcal C(R).
\]

Leibniz differentiation and the product weight \(\sqrt\zeta P\sqrt\zeta P=\zeta P^2\le\zeta\) give

\[
\mathcal B(E,L)\in\mathcal M_{\rho+\alpha-1/2},
\quad
\mathcal B(U,R)\in\mathcal M_{\alpha+1/2-\kappa_s}.
\]

The three quadratic pieces have exponents, respectively,

\[
2\alpha-1,\qquad 2\alpha-\tfrac12-\kappa_s,
\qquad 2\alpha-2\kappa_s.
\]

Because \(0<\varepsilon\le1\) and \(\kappa_s<1/2\), the latter two exponents are at least \(2\alpha-1\); hence their estimates imply the same weaker class \(\mathcal M_{2\alpha-1}\). Equivalently, \(R\in\mathcal W_{\alpha-\kappa_s}\) implies \(R\in\mathcal W_{\alpha-1/2}\), so \(V\) is in the latter class. This proves exactly the source's third displayed exponent without suppressing a quadratic term.

The original tangential radial-divergence components are \((D_r+2/R)T_\theta\) and \((D_r+1/R)T_z\). On the fixed positive radial interval the frame multipliers and all their derivatives are bounded. The established coefficient derivative map is \(D_r:\mathcal M_\beta\to\mathcal M_{\beta-\kappa_s}\). Consequently the three respective exponents after this conservative radial-divergence estimate are

\[
\rho+\alpha-\tfrac12-\kappa_s,
\qquad \alpha+\tfrac12-2\kappa_s,
\qquad 2\alpha-1-\kappa_s.
\]

For these particular fully averaged remainders one can also observe that they are auxiliary-independent; then the torus component of \(D_r\) annihilates them and the ordinary chart derivative need not lose \(\kappa_s\). The source's weaker estimate remains valid and is the one recorded for stage comparison. The local calculation itself imposes no extra inequality involving \(\alpha\) and \(\rho\); a later residual-improvement claim has to substitute its actual exponents into these proved expressions.

\subsubsection{9. Audit disposition}

The exact determinant, orientation, reference coefficients, signed inverse, covariance factors, physical chart factors, and all three remainder exponents have been checked. The positivity argument can be made fully explicit by the single value threshold in Section 4, and all fixed derivative orders follow without increasing that threshold. Division by the original positive amplitudes is valid on the prescribed weighted domain and gives smooth zero-extended signed corrections.

The substantive omission is the overly broad physical input wording at line 403; Section 7 provides both its counterexample and the exact corrected domain. The phrase ``finite Neumann series estimate'' at line 353 is imprecise; Section 4 supplies a direct norm proof and exact inverse-difference identity. The background endpoint factors and pulse evolution remain identifiable imported construction inputs, with their precise source locations recorded in Section 1; this covariance audit does not certify those separate constructions or the later numerical choice of iteration exponents.



\section{The actual initial finite state}

This calculation uses the objects just constructed, with their original physical
definitions.  An overbar is normalized auxiliary Haar average.  Angular average
is denoted by $\langle\cdot\rangle_\theta$ and always averages cylindrical
components.  The two operations are not identified with physical evaluation.
Set $\kappa=10^{-5}$ and $H=1-\kappa$.  All the estimates below retain
$\varepsilon=Q^h$, $S_*=\ell^2$, and the shell weights in the definitions of
$\mathcal W_\alpha,\mathcal M_\alpha,\mathcal S_\alpha$.

\subsection{The base expansion supplies the precise chart input}

Write the original diagonal sum of the positive-order coefficient potentials as
$\sum_{n\ge1}\chi(c_nq)\mathcal A_n$, where
$\mathcal A_n=q^{1-A+2nh}F_n e_\theta/r$.  Its radial and axial components
are exactly
\[
ru_{r,n}^{\rm cut}=q^{2nh}\left\{\chi(c_nq)V_n-
 \frac{2\eta}{L}(c_nq)\chi'(c_nq)F_n\right\},\qquad
u_{z,n}^{\rm cut}=q^{-A+2nh}\chi(c_nq)U_n.
\]
The first equality follows from $ru_r=-\partial_z(r\mathcal A_\theta)$,
$q_z=2\eta q^{1-D}/L$, and the original coefficient divergence identity.
In particular, no derivative of a cutoff is omitted from the background.
The diagonal choice is made for all normalized coefficient seminorms through
order $n$, with each corresponding term at most $2^{-n}q^{nh}$.
Such a choice is simultaneous on the original compact profile rectangle:
at step $n$ there are finitely many seminorms, each bounded by a fixed
constant times $q^{2nh}$ and a fixed power of $|\log q|$ after the
scale derivatives of $\chi(c_nq)$ have been included.  On their support
$c_nq\le1$, and on a nonzero differentiated cutoff $1/2\le c_nq\le1$;
hence $(c_nq)^j\chi^{(j)}(c_nq)$ is bounded independently of $c_n$.
For each of the finitely many constants, $q^{nh}(1+|\log q|)^b$ tends
to zero as $q\downarrow0$. Choose $c_n\ge2c_{n-1}$ large enough
that all the finitely many required bounds hold for $q\le c_n^{-1}$.
This fixes one entire sequence, with no subsequent choices depending on a
requested output derivative. The original finite coefficients are unchanged
near $q=0$, where each individual cutoff equals one.
For a fixed derivative order $m$, the finitely many terms with $n\le m+1$
retain their original powers $q^{2nh}$; the tail starting at
$n\ge\max(2,m+2)$ is bounded by
\[
\sum_{n\ge\max(2,m+2)}2^{-n}q^{nh}\le 2q^{2h},\qquad 0<q\le1.
\]
This proves the all-fixed-derivative estimate $O(q^{2h})$ for the tangential
background difference.  Multiplication by the original common radial factor
$q^h$ proves $b=O(\varepsilon S_*^B)$ and its positive-order difference
$O(\varepsilon^3 S_*^B)$ in a $Q$ chart, because $q/Q$ stays in a fixed
compact positive interval.  The slow coordinate transition proved above
passes every fixed derivative through this estimate.

For the stress, use the same diagonal selection on $T_n/\zeta$ for $n\ge2$.
All these quotients are smooth with fixed-order bounds on the compact shell;
the single $n=1$ edge term has the explicitly retained bounds
$|\partial^IT_1|\le C_I\zeta\delta^{-b_I}$.  Separating that single term
and the finitely many early terms from the tail proves
\[
|\partial^I(\widehat T-T_0)|
 \le C_Iq^{2h}\zeta\delta^{-b_I}.
\]
Here $\widehat T$ is the actual summed profile stress, not a truncated stress.
The original physical stress and its chart value are
\[
\Sigma_{\rm phys}=q^{-A-1/2}\widehat T,\qquad
\Sigma=Q^{2A}q^{-A-1/2}\widehat T,
\]
so the leading value $\Sigma^{(0)}=\varepsilon T_{0,*}$ satisfies
\[
\Sigma-\Sigma^{(0)}\in\mathcal M_3.
\tag{SI1}
\]
Indeed $Q^{2A}q^{-A-1/2}=\varepsilon(Q/q)^{A+1/2}$ because
$A=1/2+h$, and $q^{2h}=\varepsilon^2(q/Q)^{2h}$.  No factor in this
identity is replaced by one.  The weaker $\mathcal M_2$ bound would also
suffice below.

\subsection{The complete primary curl and its residual}

Let $w_0^{\rm tan}$ be the assembled real transverse primary field and let
$w=\operatorname{curl}_*\mathcal A_0$ be its actual cutoff curl.  For each
label and nonzero harmonic, its potential coefficient is
\[
C_m=\frac{i n_\Phi\times t_m}{km|n_\Phi|^2},\qquad
n_\Phi=\nabla_*\Phi,
\]
with $n_\Phi\cdot t_m=0$ and the cutoff included in $t_m$.  The vector
triple-product identity and the original right-handed cylindrical orientation
give
\[
\operatorname{curl}_*(C_me^{ikm\Phi})=(t_m+r_m)e^{ikm\Phi},\qquad
r_m=(-D_zC_{m,\theta},\ D_zC_{m,r}-D_rC_{m,z},\
 (D_r+R^{-1})C_{m,\theta}).
\]
The multiplier $n_\Phi/|n_\Phi|^2$ has polynomial $S_*$ derivative bounds,
$k^{-1}\le\sqrt\varepsilon$, and $D_r$ costs at most $\kappa$.
Consequently
\[
w_0^{\rm tan}\in\mathcal W_{1/2},\qquad
w-w_0^{\rm tan}\in\mathcal W_{1-\kappa},\qquad
w\in\mathcal W_{1/2}.
\tag{SI2}
\]
The full $w$ is divergence free.  Expansion of its divergence, for the
complete amplitude $a_m=t_m+r_m$, gives the exact longitudinal identity
\[
ikm(n_\Phi\cdot a_m)
=-\{(D_r+R^{-1})a_{m,r}+D_za_{m,z}\}.
\tag{SI3}
\]
Thus $n_\Phi\cdot a_m=n_\Phi\cdot r_m$ has exponent $1-\kappa$.
For two same-label harmonics, the apparently leading transport term is
\[
ikm'(a_m\cdot n_\Phi)a_{m'}
=-\frac{m'}m\{(D_r+R^{-1})a_{m,r}+D_za_{m,z}\}a_{m'}.
\]
It has exponent $1-\kappa$ for two primary waves.  Ordinary amplitude
transport and cylindrical connection terms have that exponent or a larger
one.  Distinct-label derivative products vanish by the support separation
proved above.  The product weight is $\zeta P^2$, bounded both by $\zeta$
and by $\sqrt\zeta P$.  This proves the stated exponent for both the zero
and nonzero angular harmonics, at every fixed coefficient derivative order.

For completeness the primary linear error has the following term bounds.
The operator on the transverse amplitude is the homogeneous constrained pulse
operator; its temporal cutoff error is kept in the flat residual.  The
remaining groups and their exponents, for a wave of exponent $\alpha$, are
\[
\begin{array}{c|c}
(-\varepsilon\partial_T+bD_r+GD_z)a&\alpha+1-\kappa\\
ikmE_{\rm ik}a&\alpha+1/2\\
a_r b_R e_r+a_zD_z(b,V,G)+(b/R)a_\theta e_\theta&\alpha+1\\
(D_r\pi,0,D_z\pi)&\alpha+1/2-\kappa\\
-\varepsilon(D_r^2+R^{-1}D_r+D_z^2)a&\alpha+1-2\kappa\\
-2i\varepsilon km(n_rD_r+n_zD_z)a&\alpha+1/2-\kappa\\
\text{phase divergence and cylindrical connection terms}&\alpha+1/2
\end{array}
\]
In the last row $n_\Phi$ is auxiliary independent in the radial direction
of the pulse: $D_rv=0$; its radial coefficient derivatives therefore have
polynomial $S_*$ bounds without an extra $\varepsilon^{-\kappa}$.
The phase defect is exactly
$E_{\rm ik}=\varepsilon v((H_\Phi)_T-G(H_\Phi)_Z)+bn_{\Phi,r}$.
These facts prove every entry by multiplication with the displayed differential
operators.  Applying the principal pulse velocity operator to $r_m$ preserves
its exponent: $cN$ and the coefficient matrix preserve exponents, while
$\varepsilon k^2m^2|n_\Phi|^2$ is bounded at each fixed $m$.
Taking $\alpha=1/2$ gives the safe common primary nonzero residual bound
\[
G_{\rm wave}^{\rm pre}\in\mathcal W_{1-3\kappa}.
\tag{SI4}
\]

\subsection{Initial full mean equations and the improved Haar average}

Before mean updates put $v=\gamma=\beta=0$ and retain
$W^{ab}=\langle w_aw_b\rangle_\theta$.  The actual radial source and
pressure are
\[
g_r=-(D_r+R^{-1})W^{rr}-D_zW^{zr}+W^{\theta\theta}/R,
\quad P=\int\bar g_r\,dR,\quad p_m=\mathcal T_0(g_r-\rho P).
\]
The product rules give $g_r,p_m\in\mathcal M_H$ and
$P\in\mathcal S_H$.  The two residuals are
\[
E_\theta=(D_r+2/R)(W^{r\theta}-\Sigma_\theta)+D_zW^{z\theta},
\]
\[
E_z=(D_r+1/R)(W^{rz}-\Sigma_z)+D_z(W^{zz}+p_m).
\tag{SI5}
\]
Hence $E_\theta,E_z\in\mathcal M_H$.  The integrated definitions
\[
J_\theta=\int R^2\overline{W^{z\theta}}\,dR,\quad
J_z=\int R\overline{W^{zz}}\,dR-\tfrac12\int R^2\bar g_r\,dR
\]
give $(P,J_\theta,J_z)\in\mathcal S_H$.

The signed covariance construction at the original positive amplitudes gives
$\overline{\langle(w_0^{\rm tan})_r(w_0^{\rm tan})_{\rm tan}
\rangle_\theta}=\Sigma^{(0)}$ exactly.  Expanding the complete curl
$w=w_0^{\rm tan}+r$ retains both cross terms and $\langle r_rr_{\rm tan}
\rangle_\theta$.  Their exponents are $3/2-\kappa$ and $2-2\kappa$.
Therefore
\[
\overline{W^{r\theta},W^{rz}}-\Sigma^{(0)}
 \in\mathcal M_{3/2-\kappa}.
\]
For Haar averages $\overline{D_rf}=\partial_R\bar f$ by integration of
the torus derivative.  Differentiating the preceding estimate therefore
preserves its exponent.  Equations (SI1) and (SI5) give
\[
\bar E_\theta,\bar E_z\in\mathcal M_{3/2-\kappa}
 \subset\mathcal M_{1.49}.
\tag{SI6}
\]
The axial flux terms have exponents $2$ and $H+1=2-\kappa$;
they do not reduce this bound.

\subsection{Temporal cancellation with the entire reconstructed field}

Use the fixed common-torus inverse to define
\[
u=-c^{-1}N^{-1}(E_\theta-\bar E_\theta),\quad
\gamma_d=-c^{-1}N^{-1}(E_z-\bar E_z),\quad
\Psi=\mathcal T_1\gamma_d,
\]
\[
B=-\varepsilon\partial_Z\Psi,\quad
d=(D_r+R^{-1})\Psi=\gamma_d+a,\quad a=-\mathcal A_1\gamma_d.
\tag{SI7}
\]
These are the actual increments.  In particular, $a$ is retained in $d$ in
every product.  They satisfy
$u,d,\Psi\in\mathcal M_H$, $B\in\mathcal M_{H+1}$,
$cNu=-E_\theta+\bar E_\theta$, and
$cNd=-E_z+\bar E_z+cNa$.
The projector proof gives flatness of $a$ and all its physical derivatives.
It also gives $\bar u=\bar d=\bar\Psi=\bar B=0$ pointwise in radius.
Consequently both original weighted mean moments remain zero exactly.

Recompute the radial source after this update.  Direct product expansion
gives
\[
\Delta g_r=(2V/R)u+R_g^{\rm time},
\]
\begin{align*}
R_g^{\rm time}={}&-t_*B-(D_r+R^{-1})(2bB+B^2)\\
&-D_z(bd+GB+Bd)+u^2/R
 +\varepsilon(\Delta_0-R^{-2})B.
\end{align*}
Here the exponents, in the order written after splitting the finite sums,
are
\[
H+1,\quad H+2-\kappa,\quad2H+2-\kappa,\quad
H+2,\quad H+2,\quad2H+2,\quad2H,\quad H+2-2\kappa.
\]
The term $t_*B$ includes $cNB$; unlike the later slow five-dimensional
correction, this term cannot be assigned exponent $H+2$ in its entirety.
Nevertheless $\Delta g_r\in\mathcal M_H$.  The fixed linear pressure
map gives $\Delta p_m\in\mathcal M_H$.

Subtracting the complete old and new mean equations now gives exactly
\begin{align*}
E_\theta^{\rm time}-\bar E_\theta={}&-\varepsilon\partial_Tu
-\varepsilon(\Delta_0-R^{-2})u\\
&+(D_r+2/R)(bu+VB+Bu)+D_z(Gu+Vd+du),\\
E_z^{\rm time}-\bar E_z={}&-\varepsilon\partial_Td
-\varepsilon\Delta_0d+cNa\\
&+(D_r+1/R)(bd+GB+Bd)+D_z(2Gd+d^2+\Delta p_m).
\end{align*}
The slow-time, viscous, linear radial flux, and linear axial flux groups
have exponents at least $H+1$, $H+1-2\kappa$, $H+1-\kappa$,
and $H+1$, respectively.  The radial quadratic flux has exponent
$2H+1-\kappa$ and the axial quadratic flux exponent $2H+1$.
All of these are at least $H+1-2\kappa=2-3\kappa$.
After the explicitly displayed $cNa$ is placed in the flat residual,
\[
E_\theta^{\rm time},E_z^{\rm time}\in\mathcal M_{1.49},\quad
(P,J_\theta,J_z)^{\rm time}\in\mathcal S_H.
\tag{SI8}
\]
This assertion does not say that the Haar averages are unchanged.
Their linear changes vanish because the base is auxiliary independent and
$\bar u=\bar d=\bar B=0$.  Their quadratic changes need not vanish;
the displayed equations retain them.  They have sufficient exponents to
preserve (SI8).

\subsection{The five exact moment rows and the resulting budgets}

Apply the proved fixed five-dimensional inverse on the original reserved
power-law interval to the recomputed $(P,J_\theta,J_z)^{\rm time}$.
Denote its slow azimuthal, axial, and radial increments by $u_2,d_2,B_2$.
They have exponents $H,H,H+1$, preserve the first two moments exactly,
and have identically zero axial projector remainder.  The last three rows
cancel the linear contributions in the original fixed base, with their
physical and chart factors retained.  Put $v=u$, $\gamma=d$, $\beta=B$
only in the following full difference identity:
\[
g_r^{\rm new}-g_r=(2V/R)u_2+R_g^{\rm slow},
\]
\begin{align*}
R_g^{\rm slow}={}&-t_*B_2
 -(D_r+R^{-1})(2bB_2+2\beta B_2+B_2^2)\\
&-D_z(bd_2+GB_2+\beta d_2+\gamma B_2+B_2d_2)\\
&+(2vu_2+u_2^2)/R+\varepsilon(\Delta_0-R^{-2})B_2.
\end{align*}
Every entry is the difference of the corresponding term in the original
radial equation.  No wave covariance changes in this update.  Since $B_2$
is auxiliary independent, its time term is $+\varepsilon\partial_TB_2$.
Using the deliberately retained bounds
$v,\gamma\in\mathcal M_{0.9}$, $\beta\in\mathcal M_{1.9}$, the
respective exponents of the displayed terms are
\[
\begin{gathered}
H+2;\quad H+2-\kappa, H+2.9-\kappa, 2H+2-\kappa;\\
H+2, H+2, H+2.9, H+2.9, 2H+2;\\
H+0.9, 2H;\quad H+2-2\kappa.
\end{gathered}
\]
The last common viscous value bounds all four entries of its displayed
operator.  Subtraction shows that every listed exponent is at least
$H+0.9-2\kappa=1.9-3\kappa$, since $H\ge0.9$.
The exact new defects are
\[
P^{\rm new}=\int\overline{R_g^{\rm slow}}\,dR,
\]
\[
J_\theta^{\rm new}=\int R^2\overline{\gamma u_2+vd_2+d_2u_2}\,dR,
\]
\[
J_z^{\rm new}=\int R\overline{2\gamma d_2+d_2^2}\,dR
-\tfrac12\int R^2\overline{R_g^{\rm slow}}\,dR.
\tag{SI9}
\]
In the last equation the radial-source term is indispensable: the linear
part $-\tfrac12\int R^2(2V/R)u_2=-\int RVu_2$ is canceled by the
fifth row, while the displayed nonlinear radial-source moment remains.
Thus all three new defects lie in $\mathcal S_{1.9-3\kappa}$.

For the tangential residuals the difference equations have exactly the same
flux patterns as the temporal difference equations, with $(b,V,G)$ replaced
in the products by $(b+\beta,V+v,G+\gamma)$ and $(B,u,d)$ replaced by
$(B_2,u_2,d_2)$.  The slow time has no $cN$ part.  Linear radial and
axial terms have exponents $H+1-\kappa$ and $H+1$, viscosity
$H+1-2\kappa$, and the additional existing-mean products have exponent
at least $H+1.9-\kappa$ in the radial flux and $H+1.9$ in the axial
flux.  The pressure change has exponent $H$, hence its $D_z$ term
$H+1$.  All changes lie in $\mathcal M_{2-3\kappa}$, so (SI8)
persists after the five-dimensional correction.

Neither mean correction changes $w$.  Its interactions with the wave in
the full residual do change.  For a tangential mean of exponent $\mu$
and radial mean of exponent $\mu+1$, the largest phase term has exponent
$1/2+\mu-1/2=\mu$.  Indeed $n_\Phi$ has bounded polynomial coefficients,
$k\le2\varepsilon^{-1/2}$, and the radial mean component has one
extra $\varepsilon$.  The amplitude derivative terms have exponent at
least $1/2+\mu-\kappa$, and the cylindrical connection terms have
exponent $1/2+\mu$.  Taking $\mu=H$ for both constructed updates and
combining with (SI4) yields $G_{\rm wave}^{[0]}\in\mathcal W_{1-3\kappa}$.

Every claim in the desired initial-state bounds now follows by subtraction
of explicit numbers:
\[
1-3\kappa=0.99997>0.7,
\quad1.49>1.2,
\quad1.9-3\kappa=1.89997>1.2,
\]
\[
1-\kappa=0.99999>0.68,
\quad H>0.9,
\quad H+1>1.9.
\]
Since $0<\varepsilon\le1$, a larger exponent gives the smaller class.
The actual state therefore has, with the original budgets unchanged,
\[
G_{\rm wave}^{[0]}\in\mathcal W_{0.7},\quad
E_\theta^{[0]},E_z^{[0]}\in\mathcal M_{1.2},\quad
(P,J_\theta,J_z)^{[0]}\in\mathcal S_{1.2},
\]
\[
w\in\mathcal W_{0.5},\quad w-w_0^{\rm tan}\in\mathcal W_{0.68},\quad
v,\gamma,p_m\in\mathcal M_{0.9},\quad\beta\in\mathcal M_{1.9},
\]
\[
\int R^2\bar v\,dR=0,\qquad\int R\bar\gamma\,dR=0.
\tag{SI10}
\]
The residual equations in (SI10) use the wave/mean representative obtained
by putting all the explicitly constructed flat terms in the additive flat
residual.  The velocity and pressure remain the complete constructed fields.

\section{One fixed domain and the actual residual source class}

Choose the background parameters and original profiles first.  Fix the
positive compact shell, reserved mean interval, strict cone margins,
color centers, auxiliary rectangles, and cutoff functions.  Only then choose
one integer $\ell_0$ large enough for the finite list of value comparisons:
nonnegative covering indices; bounded neighboring covering differences;
phase length and plane-isomorphism lower bounds; pulse growth perturbation
and invariant-interval bounds; the covariance determinant and positive
coefficient margins; and the elementary exponential-versus-polynomial bounds
used for those value comparisons.  The actual mean moment inverse uses
the exact power-law shell and a fixed invertible matrix, so adds no
source-dependent restriction.  Fix
\[
0<q_{\rm big}\le\min(q_{\rm base},2^{-(\ell_0+1)},1/2).
\tag{SI11}
\]
An active dyadic band has $Q\le2q$, so every band on this domain has
$\ell\ge\ell_0$.  At a point with $q_0>0$, take a slow neighborhood
$U$ with $q(U)\subset(q_0/2,2q_0)$.  Its active dyadic scales lie in
$(q_0/4,4q_0)$ and consequently have span at most four integers.
This is the needed local restriction for the uniformly bounded common
covering changes.  Mere finiteness of an arbitrary larger collection of
bands would not give that bound.

For clarity, increasing the stage or requested derivative order does not
change (SI11).  In the pulse equation the exact propagators obey
\[
V_m(v,w)=\exp\left(-(m^2-1)\int_w^v d(s)\,ds\right)V_1(v,w),
\quad d>0,\quad |m|\ge1.
\]
Thus the value bound $\|V_m(v,w)\|\le C P(v)/P(w)$ holds on the
same domain for every nonzero integer $m$.  Parameter differentiation of
the ODE has the exact triangular form
\[
(\partial^I z)'=A_m\partial^Iz+\partial^Ig+
\sum_{0<J\le I}\binom IJ(\partial^JA_m)\partial^{I-J}z.
\]
At each fixed derivative and finite harmonic set, the already bounded
lower derivatives give a new polynomial $S_*$ bound by the same
Duhamel integral.  No new denominator or smallness comparison is introduced.
The torus inverses have finite derivative loss for each requested order;
their value definitions are unchanged.  The radial inverse's repeated
integration by parts estimates the same fixed operator and does not change
it with the number of integrations.  The signed covariance inverse divides
only by the original $2a_\sigma$ and uses the original $H^{-1}$.
The five-dimensional map, pressure map, and support bumps are fixed linear
maps.  Finite sums, products, and derivatives exist throughout the fixed
domain.  Induction over the finite directed graph of these operations proves
that every finite state exists on (SI11), with constants depending on the
stage and derivative order.  There is no intersection of a stage-indexed
sequence of decreasing positive thresholds.

Here is the precise source invariant for that induction.  The nonzero
harmonics are finite sums
\[
\sum_{\gamma}\sum_{m\in H_j}
 f^{[j]}_{\gamma,m}e^{ik_\gamma m\Phi_\gamma},\qquad
H_j\Subset\mathbb Z\setminus\{0\},
\]
whose equal-label, equal-harmonic coefficients are collected.  Each
coefficient is defined on the original enlarged pulse rectangle, has the
original closed slow and transverse supports, is smooth by zero extension
at their boundaries and both radial edges, and satisfies
\[
|\partial^If^{[j]}_{\gamma,m}|
\le C_{j,I}\varepsilon^\alpha S_*^{b_{j,I}}
 \sqrt\zeta\delta^{-d_{j,I}}P(v).
\tag{SI12}
\]
Its forward-pulse source representative need not have zero data at the
terminal end; the actual correction is multiplied by the fixed terminal
cutoff before taking its curl.  The initial sources have (SI12) by the
primary amplitude, curl and residual calculations above.  A coefficient
derivative preserves supports, and shifts only the derivative index.
$D_r,D_z,t_*$ have the exact costs already proved.  A mean coefficient
can have full auxiliary-torus support but its product with a wave keeps
the wave support and $P(v)$.  A same-label wave product has $P(v)^2\le P(v)$;
distinct labels have zero products; multiplication adds angular harmonic
integers and differentiation preserves them.  These facts prove closure
under each local polynomial residual operation, not just an abstract
membership assertion.

The pulse Duhamel map preserves (SI12) because its integrand has
$P(v)/P(w)$ times the source factor $P(w)$; the radial and slow weights
are constant along that exact integration path.  Its zero-data support
property follows from ODE uniqueness, and each boundary jet follows from
the differentiated zero-data equations.  The signed stress map supplies
the exact factor $\varepsilon^{\alpha-1}\sqrt\zeta$ on division by
the original positive amplitude and hence a wave of exponent
$\alpha-1/2$ with the original envelope.  The weighted radial primitives,
interior moment subtractions, temporal inverse, and five-dimensional
corrections yield smooth shell-supported mean coefficients.  Radial
integration fixes $(z,t)$ and hence $q(z,t)$, so it never gathers an
additional dyadic band.  A radial line meets $O(S_*^3)$ boxes, which adds
only a polynomial to the estimates.  Subsequent products with waves again
preserve (SI12).

Finally define the flat class by requiring every fixed physical derivative
to be $O(q^N)$ for every $N$.  It is an ideal under multiplication by
the present finite-stage fields: if one factor has derivative bound
$Cq^{-K}(1+|\log q|)^b$, use the flat estimate with exponent
$N+K+1$ and the boundedness of $q(1+|\log q|)^b$.  Leibniz's rule has
finitely many terms at each order.  The physical derivatives of the fields
have such polynomial bounds by the exact coordinate and phase maps.
The temporal cutoff errors retain the terms $(1-\psi)f$ and $\psi't$;
on their support the Gaussian pulse has $P\le e^{-cS_*}$.  At any fixed
physical derivative order their bound has the form
\[
C Q^{-M}S_*^b e^{-c\ell^2}.
\]
For every $N$, multiplication by $q^{-N}$ is bounded because
$-c\ell^2+(M+N)\ell\log2+2b\log\ell\to-\infty$.
The radial projector errors are flat by the full-obstruction proof.
Their actual fields remain in every product; the resulting flat terms
are assigned to $F^{[j]}$, not used as forward-pulse sources.  Together
with the base's flat residual, these calculations prove on the single fixed
domain the actual decomposition, whose nonzero $G^{[j]}$ harmonics satisfy
(SI12):
\[
\mathcal R(u^{[j]},p^{[j]})=G^{[j]}+F^{[j]},\qquad
|F^{[j]}|_m\le C_{j,m,N}q^N
\]



% --- included proof body: summation_realization_body.tex ---
\section{Shrinking-cutoff summation and local realization: exact source audit}\label{shrinking-cutoff-summation-and-local-realization-exact-source-audit}

\subsection{Result and coverage}\label{result-and-coverage}

The shrinking-cutoff argument in Lemma 5.4 is valid under its stated hypotheses. The calculation below verifies its scale selection, every mixed derivative loss, local finiteness, preservation of divergence, the nonlinear differential-polynomial comparison, and its use of one fixed-stage flat remainder. No summability across stages of the flat remainders is required. The specialization to the physical Navier--Stokes residual gives the exact passage from the finite-stage bounds (9.17)--(9.18) to (9.20) and the representation (9.21).

This bounded audit found no concrete defect in that passage. It is not an independent construction of the preceding wave/mean increments, a proof of all earlier support and inverse estimates, a certification of every conclusion of Proposition 9.9, or a Lean/kernel acceptance result. In particular, validating the summation implication does not validate the input estimate (9.18) by itself.

The source is the current 166-page PDF:

\texttt{source/navier-stokes-166.pdf}

SHA256: \nolinkurl{0e779481c4da40bd28d1e642e1d8ca57447d129610df28dfa5a11e9af8ae228f}.

Its preserved Poppler layout text is \texttt{downloaded\_public\_release/current\_pdf\_extracted.txt}, under the same programme root. The following physical pages and exact text-line ranges were read:

{\def\LTcaptype{none} % do not increment counter
\begin{longtable}[]{@{}
  >{\raggedright\arraybackslash}p{(\linewidth - 6\tabcolsep) * \real{0.2500}}
  >{\raggedright\arraybackslash}p{(\linewidth - 6\tabcolsep) * \real{0.2500}}
  >{\raggedright\arraybackslash}p{(\linewidth - 6\tabcolsep) * \real{0.2500}}
  >{\raggedright\arraybackslash}p{(\linewidth - 6\tabcolsep) * \real{0.2500}}@{}}
\toprule\noalign{}
\begin{minipage}[b]{\linewidth}\raggedright
Source
\end{minipage} & \begin{minipage}[b]{\linewidth}\raggedright
Physical pages
\end{minipage} & \begin{minipage}[b]{\linewidth}\raggedright
Layout-text lines
\end{minipage} & \begin{minipage}[b]{\linewidth}\raggedright
Coverage
\end{minipage} \\
\midrule\noalign{}
\endhead
\bottomrule\noalign{}
\endlastfoot
(4.1), Lemma 4.1 and proof & {}24--25 & {}1251--1281 & Exact similarity map and derivatives used for (5.28). \\
(5.23), (5.27) & {}55--56 & {}2906--2910; 2981--3002 & Pointwise derivative seminorm, potential convention and tuple types. \\
Lemma 5.4, full hypotheses, statement and proof & {}57--59 & {}3005--3158 & Complete analytic summation argument, including (5.28)--(5.40). \\
(6.1)--(6.6), (6.11)--(6.12) & {}62--64, 66 & {}3295--3374; 3451--3463 & Physical graph derivative factors and pulse-coordinate derivatives. \\
(7.2)--(7.4) & {}74--75 & {}3899--3920 & Fixed label frequencies and exact phase used to check the derivative cost. \\
Definition 9.4 & {}106 & {}5535--5557 & Exact stage gains, residual orders and representation convention. \\
Lemma 9.7 and proof & {}111--112 & {}5804--5851 & Stated common-domain construction and finite derivative dependence. Earlier inverse proofs are not independently reconstructed here. \\
(9.15)--(9.19), Lemma 9.8 and proof & {}112--114 & {}5852--5937 & Physical increment tuples, derivative-loss arithmetic and finite-stage residual interface. \\
Proposition 9.9, statement and Steps 1--2 & {}114--115 & {}5938--5999 & Actual use of Lemma 5.4 and the full potential/direct-field/pressure assembly (9.21). \\
Proposition 9.9, Steps 3--5 & {}115--116 & {}6000--6055 & Read for the scope of its remaining endpoint/exterior/growth claims; these earlier dependent constructions are not certified by the summation calculation alone. \\
\end{longtable}
}

Rendered pages 57, 58, 59, 63, 113, 114 and 115 were also inspected. This resolves two potentially misleading layout extractions: the source has

\[
d_r=2\big((1+h)\rho_g-h\kappa_s\big),\qquad
s_x=\frac12+\frac h2,\qquad s_t=1+\frac{3h}{2},
\]

with \(\kappa_s=10^{-5}\). The factor 2 in \(d_r\) multiplies both terms. No source equation is changed in this audit.

\subsection{1. Original objects, types and derivative convention}\label{original-objects-types-and-derivative-convention}

Write \(x=(x_1,x_2,z)\), \(I=(\alpha_1,\alpha_2,\alpha_3,b)\in\mathbb N_0^4\), \(D^I=\partial_{x_1}^{\alpha_1}\partial_{x_2}^{\alpha_2}\partial_z^{\alpha_3}\partial_t^b\), and \(|I|=|\alpha|+b\). The source seminorm (5.23) is the pointwise quantity

\[
|V|_m=\max_{|I|\le m}|D^IV|.
\]

The lemma's domain is \(\Omega\), with \(0<q(z,t)<q_0\le1\), \(q\) smooth and bounded below on compact subsets. Its exact derivative hypothesis is

\[
|\partial_z^a\partial_t^bq|\le C_{a,b}q^{1-aD-b},\qquad 0<D<1.
\tag{A1}
\]

The base is \(U_0=(u_0,p_0,T_0)\), with \(\operatorname{div}u_0=0\). The optional stress entries may be omitted. Each increment is represented by

\[
Z_j=(A_j,B_j,p_j,T_j),\qquad B_j=b_j(r,z,t)e_\theta,
\] \[
U_j=(\operatorname{curl}A_j+B_j,p_j,T_j),\qquad
U^{[J]}=U_0+\sum_{j=1}^J U_j.
\tag{A2}
\]

These are actual physical fields on the same domain, with the Cartesian and zero-extension regularity specified before (5.30). The original bounds are

\[
|Z_j|_m\le C_{j,m}q^{g_j-\ell_m}\Lambda_{\log}(q)^{P_{j,m}},\qquad
\Lambda_{\log}(q)=1+|\log q|,
\tag{A3}
\]

where \(0<g_1\le g_2\le\cdots\to\infty\), and the nonnegative, nondecreasing \(\ell_m\) are independent of \(j\). Constants and logarithmic powers may depend on \(j\). No uniform bound for them in \(j\) is used below.

The realization map on one representative tuple is

\[
\mathcal C_a(A,B,p,T)
=\big(\operatorname{curl}(\chi(aq)A)+\chi(aq)B,
       \chi(aq)p,\chi(aq)T\big).
\tag{A4}
\]

It maps smooth potential/direct-azimuthal/scalar tuples to physical velocity/pressure/stress tuples on \(\Omega\). It is linear in the tuple. Its velocity is not defined by merely multiplying \(\operatorname{curl}A+B\) by \(\chi(aq)\); the exact extra term is retained in Section 3 below.

\subsection{2. Exact similarity map proves the required mixed derivatives of q}\label{exact-similarity-map-proves-the-required-mixed-derivatives-of-q}

For the application, the source fixes

\[
\tau=1-t,\quad A=\frac12+h,\quad D=\frac12-h,\quad
z=q^D\eta,\quad t=1-q(1-\eta^2),\quad
L=1-2h\eta^2,\quad d=1-\eta^2.
\tag{A5}
\]

The map \((q,\eta)\mapsto(z,t)\), from \((0,\infty)\times(-1,1)\) to \(\mathbb R\times(-\infty,1)\), has Jacobian

\[
\begin{pmatrix}
Dq^{D-1}\eta&q^D\\
-(1-\eta^2)&2q\eta
\end{pmatrix},\qquad
\det=q^D\big(2D\eta^2+1-\eta^2\big)=q^DL>0.
\tag{A6}
\]

For fixed \(z\) and \(\tau>0\), its inverse is the unique solution of \(q-z^2q^{2h}=\tau\) with \(q>|z|^{1/D}\), and \(\eta=zq^{-D}\). The left side is zero at the stated lower endpoint, tends to infinity, and has positive derivative \(1-2hz^2q^{2h-1}\ge1-2h\) there and thereafter. Thus existence and uniqueness follow without choosing a different branch. At \(z=0\), the inverse is \(q=\tau,\eta=0\).

Inverting (A6) gives exactly

\[
q_z=\frac{2\eta q^{1-D}}L,\quad q_t=-\frac1L,\quad
\eta_z=\frac{d}{q^DL},\quad \eta_t=\frac{D\eta}{qL}.
\tag{A7}
\]

For any real \(\beta\) and smooth \(F(\eta)\), the two derivative identities are

\[
\partial_z(q^\beta F)
=q^{\beta-D}\frac{2\beta\eta F+dF'}L,
\qquad
\partial_t(q^\beta F)
=q^{\beta-1}\frac{-\beta F+D\eta F'}L.
\tag{A8}
\]

Starting with \(F_{0,0}=1\), repeated application of these exact operators gives

\[
\partial_z^a\partial_t^bq=q^{1-aD-b}F_{a,b}(\eta).
\tag{A9}
\]

One concrete recursive construction is first to apply the time operator \(b\) times to \(q\), and then the axial operator \(a\) times, changing \(\beta\) by \(-1\) and \(-D\) respectively at each step. Commutation of the actual smooth derivatives identifies the same mixed derivative for any other order. Every resulting \(F_{a,b}\) is a finite expression obtained from polynomials and powers of \(L^{-1}\) by differentiation and multiplication. Since \(L\ge1-2h>0\) on \([-1,1]\), all these functions are smooth there and \(C_{a,b}:=\max_{[-1,1]}|F_{a,b}|<\infty\). This proves (A1) with the original exponent \(1-aD-b\).

Also \(q=\tau+q\eta^2\ge\tau\). A compact subset of the open time domain \(t<1\) has a positive minimum of \(\tau\), so it has a positive lower bound for \(q\). The stronger one-sided bounds at \(\tau=0,q\ge c>0\) used later in Proposition 9.9 are separate statements about the actual input representatives, not consequences of this compact-subset observation.

\subsection{3. Cutoff derivatives, all curl terms and an explicit fixed loss}\label{cutoff-derivatives-all-curl-terms-and-an-explicit-fixed-loss}

A fixed smooth cutoff with the stated plateaux exists explicitly. Define \(\eta_0(s)=0\) for \(s\le0\) and \(\eta_0(s)=e^{-1/s}\) for \(s>0\), and set

\[
\chi(s)=\frac{\eta_0(1-s)}{\eta_0(1-s)+\eta_0(s-1/2)}.
\tag{A10}
\]

The denominator is positive for every real \(s\). Each positive-side derivative of \(\eta_0\) is \(e^{-1/s}\) times a polynomial in \(s^{-1}\); the exponential series shows its limit is zero as \(s\downarrow0\). Hence \(\eta_0\), and then \(\chi\), are smooth. This cutoff equals one for \(s\le1/2\), is zero for \(s\ge1\), and lies between zero and one. Every positive-order derivative is supported in \([1/2,1]\).

For a positive-order mixed derivative \(\partial_z^{a'}\partial_t^{b'}\chi(aq)\), the repeated chain rule is a finite sum of terms

\[
c_{\mathcal P}\,a^k\chi^{(k)}(aq)
\prod_{\nu=1}^k\partial_z^{a_\nu}\partial_t^{b_\nu}q,
\quad a_\nu+b_\nu\ge1,
\quad\sum_\nu a_\nu=a',\quad\sum_\nu b_\nu=b',
\tag{A11}
\]

where \(c_{\mathcal P}\) is the finite integer multiplicity of the derivative partition. Applying (A1) bounds this term by

\[
|c_{\mathcal P}|\,\|\chi^{(k)}\|_\infty
\prod_{\nu=1}^k C_{a_\nu,b_\nu}\,
(aq)^kq^{-a'D-b'}.
\tag{A12}
\]

Its support has \(1/2\le aq\le1\); consequently \((aq)^k\le1\). Summing the finite list proves

\[
|\partial_z^{a'}\partial_t^{b'}\chi(aq)|
\le C^\chi_{a',b'}q^{-a'D-b'},
\tag{A13}
\]

with constants independent of the scale \(a\). Zeroth order uses \(|\chi|\le1\). Derivatives in \(x_1,x_2\) of the cutoff vanish because \(q=q(z,t)\). Thus there is no unrecorded factor of \(a\) left in the estimates.

Using the right-handed Cartesian curl and Levi-Civita symbol \(\epsilon_{ikl}\), the exact component expansion is

\[
\big(\operatorname{curl}(\chi(aq)A_j)\big)_i
=\sum_{k,l=1}^3\epsilon_{ikl}
\big(\chi(aq)\partial_k(A_j)_l
     +a\chi'(aq)(\partial_kq)(A_j)_l\big).
\tag{A14}
\]

Equivalently the additional velocity is \(a\chi'(aq)\nabla q\times A_j\), with its original sign. A derivative \(D^I\) of (A14) can instead be written directly as

\[
\sum_{k,l}\epsilon_{ikl}
\sum_{J\le I+e_k}\binom{I+e_k}{J}
 D^J\chi(aq)\,D^{I+e_k-J}(A_j)_l,
\tag{A15}
\]

where \(e_k\) is the spatial unit multi-index. Every nonzero cutoff derivative in this formula is purely in \((z,t)\). If it uses \(a'\) axial and \(b'\) time derivatives, while the derivative of the representative has total order \(d'\), then \(d'+a'+b'\le m+1\) for \(|I|\le m\), and the total loss is \(\ell_{d'}+a'D+b'\). The pressure, direct-field and stress entries have total derivative order at most \(m\), so the same upper list also covers them.

A concrete common choice is therefore

\[
\ell'_m:=\max_{d',a',b'\in\mathbb N_0,\ d'+a'+b'\le m+1}
                 (\ell_{d'}+a'D+b').
\tag{A16}
\]

It is nonnegative, nondecreasing and independent of \(j\) and \(a\). In particular \(\ell'_m\le\ell_{m+1}+m+1\), since \(0<D<1\). The larger explicit choice \(\ell_{m+1}+m+1\) is also valid if desired; no exact identity with (A16) is claimed.

Applying the finite product-rule sums to the original bounds (A3), and retaining their constants and logarithmic powers in their finite maximum/sum, yields finite numbers \(\widehat C_{j,m}\), \(\widehat P_{j,m}\), independent of \(a\), such that

\[
|\mathcal C_a Z_j|_m
\le \widehat C_{j,m}q^{g_j-\ell'_m}
          \Lambda_{\log}(q)^{\widehat P_{j,m}}.
\tag{A17}
\]

This calculation includes mixed derivatives striking the potential, the cutoff, the extra curl term, and the direct/scalar fields. It does not estimate only the differentiated velocity before applying the cutoff.

\subsection{4. Diagonal scales and exact tail estimate}\label{diagonal-scales-and-exact-tail-estimate}

For each \(j\), there are only \(j+1\) requirements in (5.37), one for every integer \(0\le m\le j\). Each has positive power \(g_j/2\):

\[
\widehat C_{j,m}\Lambda_{\log}(q)^{\widehat P_{j,m}}q^{g_j/2}\le2^{-j}
\quad(0<q\le a_j^{-1}).
\tag{A18}
\]

These inequalities can be imposed recursively together with \(a_j\ge2a_{j-1}\) and \(a_j^{-1}<q_0\). To make the selection explicit, take \(a_j\) to be the least positive integer satisfying these conditions and, for every \(m\le j\),

\[
1+\log a_j\ge\frac{2\max(\widehat P_{j,m},0)}{g_j},\qquad
\widehat C_{j,m}a_j^{-g_j/2}(1+\log a_j)^{\widehat P_{j,m}}\le2^{-j}.
\tag{A19}
\]

For \(j=1\), omit the preceding-scale condition and also require \(a_1>q_0^{-1}\). The defining set is nonempty. Indeed, for any fixed positive \(\varepsilon\) and finite real \(P\), \(a^{-\varepsilon}(1+\log a)^P\to0\). One direct proof puts \(y=\log a\), chooses an integer \(n>\max(P,0)\), and uses \(e^{\varepsilon y/2}\ge(\varepsilon y/2)^n/n!\) for \(y>0\); the remaining exponential and negative power tend to zero. Every condition is eventually true and the list at stage \(j\) is finite.

The endpoint test (A19) implies the full interval test (A18). For \(0<q<1\), direct differentiation gives

\[
\frac{q\,\frac{d}{dq}[q^{g_j/2}\Lambda_{\log}(q)^P]}
     {q^{g_j/2}\Lambda_{\log}(q)^P}
=\frac{g_j}{2}-\frac{P}{\Lambda_{\log}(q)}.
\tag{A20}
\]

This is nonnegative when \(\Lambda_{\log}(q)\ge2\max(P,0)/g_j\). Hence its maximum on \(0<q\le a_j^{-1}\) is attained at the upper endpoint once (A19) holds. All original logarithmic powers remain present in this check.

Combining (A17) and (A18) gives

\[
|\mathcal C_{a_j}Z_j|_m\le2^{-j}q^{g_j/2-\ell'_m}\qquad(m\le j).
\tag{A21}
\]

Outside the cutoff support, the summand and all its derivatives are zero, so this bound also applies there. For a compact \(K\Subset\Omega\), let \(\delta_K=\min_Kq>0\). Since \(a_j\ge2^{j-1}a_1\), all summands with \(a_j>\delta_K^{-1}\) vanish on a neighborhood of \(K\). Thus the sum is locally finite; all differentiation on \(K\) is differentiation of a finite sum. No convergence assertion at \(q=0\) is needed to obtain smoothness on \(\Omega\).

If \(J\ge\max(1,m)\) and \(0<q<(2a_J)^{-1}\), then \(a_jq<1/2\) for every \(j\le J\). On a neighborhood of such a point those cutoffs are identically one, including all their derivatives. Therefore

\[
U-U^{[J]}=\sum_{j>J}\mathcal C_{a_j}Z_j.
\]

Because \(g_j\ge g_{J+1}\) for \(j>J\), \(0<q<1\), and \(\sum_{j>J}2^{-j}=2^{-J}\),

\[
|U-U^{[J]}|_m
\le\sum_{j>J}2^{-j}q^{g_j/2-\ell'_m}
\le2^{-J}q^{g_{J+1}/2-\ell'_m}.
\tag{A22}
\]

This is exactly the exponent and constant in (5.35). It compares with the uncut finite prefix, not with a differently cut finite prefix.

All added terms vanish on the common collar \(a_1^{-1}<q<q_0\), because every \(a_j\ge a_1\). They therefore extend smoothly by zero across \(q=q_0\). At a lateral support boundary with \(q>0\), local finiteness reduces smooth zero extension to the original finite list of smooth representatives.

\subsection{5. Divergence and the coordinate-level representative map}\label{divergence-and-the-coordinate-level-representative-map}

For every smooth Cartesian vector field \(A\),

\[
\operatorname{div}(\operatorname{curl}A)
=\sum_{i,k,l}\epsilon_{ikl}\partial_i\partial_k A_l=0.
\tag{A23}
\]

The terms for \((i,k)\) and \((k,i)\) cancel because their signs are opposite and mixed derivatives commute. This applies to \(\chi(a_jq)A_j\) itself; it therefore includes every cutoff-curl contribution in (A14).

Away from the axis, the direct representative has Cartesian expression

\[
B_j=\left(-\frac{b_jx_2}{r},\frac{b_jx_1}{r},0\right),\quad
r=(x_1^2+x_2^2)^{1/2}.
\]

Its divergence is

\[
-x_2\frac{x_1}{r}\partial_r(b_j/r)
+x_1\frac{x_2}{r}\partial_r(b_j/r)=0.
\tag{A24}
\]

Moreover \(\nabla\chi(a_jq)=(0,0,a_j\chi'(a_jq)q_z)\), whose dot product with \(B_j\) is zero. Thus \(\operatorname{div}(\chi(a_jq)B_j)=0\). The smooth Cartesian extension makes the same identity hold at the axis by continuity. Adding the divergence-free base and using local finiteness proves \(\operatorname{div}u=0\).

For the mean-potential convention actually used in (9.15), the exact right-handed cylindrical map is

\[
\Psi_je_\theta\longmapsto
\operatorname{curl}(\Psi_je_\theta)
=(-\partial_z\Psi_j,0,(\partial_r+r^{-1})\Psi_j).
\tag{A25}
\]

Multiplying the potential by a function of \(q(z,t)\) gives the additional radial component \(-\partial_z\chi(a_jq)\Psi_j\), consistent with \(\nabla\chi\times\Psi_je_\theta\), since \(e_z\times e_\theta=-e_r\). The direct azimuthal term remains a different entry of the same realization map (A4), with its divergence calculation already proved.

\subsection{6. Full differential-polynomial residual comparison}\label{full-differential-polynomial-residual-comparison}

Retain the source polynomial, including its field-independent term:

\[
\mathcal F(U)^a=c_{a0}(x,t)
+\sum_{\nu=1}^{M_a}c_{a\nu}(x,t)
   \prod_{r=1}^{d_{a\nu}}D^{\beta_{a\nu r}}U^{k_{a\nu r}},
\]

with the fixed order/degree bounds \(|\beta_{a\nu r}|\le s\), \(1\le d_{a\nu}\le d\). The coefficient derivative bounds are exactly (5.32), on the fixed support regions specified there. No coefficient, support region, order or degree is allowed to depend on the truncation index.

For \(V=U^{[J]}\), \(e=U-V\), and one monomial, let \(a_r=D^{\beta_r}(V^{k_r}+e^{k_r})\) and \(b_r=D^{\beta_r}V^{k_r}\). The exact finite product identity is

\[
\prod_{r=1}^{d_\nu}a_r-\prod_{r=1}^{d_\nu}b_r
=\sum_{r=1}^{d_\nu}
  \left(\prod_{l<r}a_l\right)(a_r-b_r)
  \left(\prod_{l>r}b_l\right).
\tag{A26}
\]

Each term contains one differentiated entry of \(e\). Apply \(D^I\), \(|I|\le m\), distributing its derivatives among the coefficient and all factors with the full multi-index product-rule coefficients. Each field derivative then has order at most \(m+s\). Since the number of coefficients, monomials and factors is fixed, all sums have bounds independent of \(J\), apart from the displayed values of \(V\) and \(e\).

For this fixed \(m\), set \(H_m=1+\max(0,\max_{|I|\le m}K_I)\). The finitely many coefficient logarithms in (5.32) are bounded by a constant times \(q^{-1}\) on \(0<q<q_0\), by the exponential-versus-power estimate used above. It follows that their derivatives through order \(m\) are bounded by a fixed constant times \(q^{-H_m}\). The resulting comparison is

\[
|\mathcal F(V+e)-\mathcal F(V)|_m
\le C_mq^{-H_m}|e|_{m+s}
       (1+|V|_{m+s}+|e|_{m+s})^{d-1}.
\tag{A27}
\]

Outside the stated coefficient-support region a differentiated monomial vanishes, because the smooth field factors and their derivatives have the original supported extensions. Inside it the coefficient bound applies. Thus no additional support-dependent bound was used. The entire field-independent term \(c_{a0}\) cancels from this difference. It is still included in the finite-stage residual hypothesis and is not discarded from \(\mathcal F(V)\).

There is a \(J\)-independent power bound for each finite prefix. One may take \(K_m=\max(K_m^{\mathrm{base}},\ell_{m+1})\), enlarging a negative base exponent to zero if necessary. Indeed, taking the curl of each uncut potential costs at most one derivative, while \(q^{g_j}\le1\). For fixed \(J\), the finite sum of original constants and the maximum of the finite list of original logarithmic exponents give

\[
|U^{[J]}|_m\le C_{J,m}q^{-K_m}\Lambda_{\log}(q)^{P_{J,m}}.
\tag{A28}
\]

This does not bound \(C_{J,m}\) or \(P_{J,m}\) uniformly in \(J\).

Now fix derivative order \(m\) and decay order \(N\). Choose one integer \(J\ge\max(1,m+s)\) so large that

\[
\frac{g_{J+1}}2-\ell'_{m+s}
\ge N+H_m+(d-1)(K_{m+s}+1),
\qquad
\rho_J-K_m^{\mathcal F}\ge N+1,
\tag{A29}
\]

and the first left side is nonnegative. This is possible because both original gain sequences tend to infinity. After fixing \(J\), take \(q\) small enough that:

\begin{itemize}
\tightlist
\item
  \(q<(2a_J)^{-1}\), so (A22) compares with the actual uncut finite prefix;
\item
  the fixed constants/logarithms in (A28) are bounded by \(q^{-1}\);
\item
  the fixed constants/logarithms in the primary term of (5.33) are bounded by \(q^{-1}\).
\end{itemize}

The nonnegative exponent in (A22) gives \(|e|_{m+s}\le2^{-J}\le1\). Equation (A28) then implies \(1+|V|_{m+s}+|e|_{m+s}\le3q^{-(K_{m+s}+1)}\). Substitution into (A27) gives

\[
|\mathcal F(U)-\mathcal F(U^{[J]})|_m
\le C_m3^{d-1}2^{-J}
 q^{g_{J+1}/2-\ell'_{m+s}-H_m-(d-1)(K_{m+s}+1)}
\le C_m3^{d-1}2^{-J}q^N.
\tag{A30}
\]

The same fixed-stage primary residual is at most \(q^N\) by (A29). Its original flat term satisfies \(E_{J,m}\le C_{J,m,N}q^N\). Therefore

\[
|\mathcal F(U)|_m
\le\big(1+C_{J,m,N}+C_m3^{d-1}2^{-J}\big)q^N.
\tag{A31}
\]

The order of choices is exactly \((m,N)\), then \(J\), then the neighborhood and constants. The only flat remainder in (A31) is \(E_{J,m}\) for this one fixed \(J\). Neither \(\sum_jE_{j,m}\) nor a bound uniform in \(j\) for their constants occurs. This closes the residual part of Lemma 5.4.

\subsection{7. Check of the physical derivative losses in (9.17)}\label{check-of-the-physical-derivative-losses-in-9.17}

The source chart holds \(Q=2^{-\ell}\), \(\varepsilon=Q^h\), and \(S_*=\ell^2\) fixed while differentiating. Its coordinates are \((R,Z,T)=(r/Q^{1/2},z/Q^D,\tau/Q)\). On the perturbation annulus \(r\asymp Q^{1/2}\). The original auxiliary covering is determined by

\[
J_g=\begin{pmatrix}3&1\\1&5\end{pmatrix},\quad
\Lambda_g=4-\sqrt2,\quad T_g=4+\sqrt2,\quad
\rho_g=\frac{\log\Lambda_g}{\log T_g},
\] \[
\kappa_s=10^{-5},\quad d_r=2((1+h)\rho_g-h\kappa_s),\quad
i=\left\lfloor\log_{T_g}\frac{Q^{-1-h}}{S_*}\right\rfloor.
\]

From \(T_g^{-1}Q^{-1-h}/S_*<T_g^i\le Q^{-1-h}/S_*\) and \(\Lambda_g^i=(T_g^i)^{\rho_g}\),

\[
\Lambda_g^iQ^{d_r/2}\asymp
Q^{-(1+h)\rho_g}S_*^{-\rho_g}
Q^{(1+h)\rho_g-h\kappa_s}
=Q^{-h\kappa_s}S_*^{-\rho_g}.
\tag{A32}
\]

Thus a physical radial derivative of an evaluated coefficient has loss \(1/2+h\kappa_s\); an axial derivative has loss \(D=1/2-h\); a time derivative has loss \(1+h\). Cartesian differentiation of the cylindrical frame costs \(1/2\). Derivatives striking their smooth radial coefficients contribute further factors \(Q^{-1/2}\); the larger radial cost already covers each such differentiation. Powers of \(S_*\) enter the logarithmic factors and may depend on the fixed derivative order and stage.

For the label phase the exact source expression is

\[
\Phi=p\theta+p_zZ/\varepsilon+x_0R-v(pF+p_zG),\qquad
k=\lceil\varepsilon^{-1/2}\rceil.
\tag{A33}
\]

The quantities \(p,p_z,x_0,k\) are fixed for the label, including the source's integer rounding rule. The pulse coordinate obeys \(\nabla_xv=0\) and \(\partial_tv=Q^{-1-h}\), with higher physical derivatives of this first derivative zero within the fixed chart. These follow from (6.12), not from neglecting the auxiliary phase evaluation. The axial linear phase term has derivative factor \(Q^{-D}\varepsilon^{-1}=Q^{-1/2}\). Repeated differentiation of the remaining slow phase factors and the angular coordinate on \(r\asymp Q^{1/2}\) gives the phase bound quoted in Lemma 9.8,

\[
|\nabla_x^a\partial_t^b\Phi|
\le C_{a,b}Q^{-a/2-b(1+h)}S_*^{P_{a,b}},\qquad a+b\ge1,
\tag{A34}
\]

using the stated fixed-order slow-coefficient bounds. Verification of those earlier coefficient estimates is an upstream input; the derivative exponents and auxiliary-coordinate contribution here are checked explicitly.

To avoid confusing the derivative order with the source's harmonic integer in \(e^{ikm\Phi}\), write that integer as \(m_{\rm harm}\) in the next calculation; it is the same integer, not a changed frequency. A derivative partition with \(k'\) nonzero phase-derivative factors gives

\[
(ikm_{\rm harm})^{k'}e^{ikm_{\rm harm}\Phi}
\prod_{\nu=1}^{k'}\nabla_x^{a_\nu}\partial_t^{b_\nu}\Phi,
\quad k'\le a+b,\quad\sum a_\nu=a,\quad\sum b_\nu=b.
\]

Since \(0<Q\le1\), \(k\le2Q^{-h/2}\). With (A34), its power of \(Q\) is bounded by

\[
Q^{-a/2-b(1+h)-k'h/2}
\le Q^{-a(1/2+h/2)-b(1+3h/2)}.
\tag{A35}
\]

The factor \(2^{k'}|m_{\rm harm}|^{k'}\), derivative-partition coefficients and logarithmic powers are retained in the fixed-stage derivative constant. The source's harmonic set is finite at each fixed stage, independently of band; this permits dependence on the stage in that constant. It does not produce an additional stage-dependent power of \(Q\).

The resulting costs are \(s_x=1/2+h/2\), \(s_t=1+3h/2\). They dominate all four rows of (9.19): \(\kappa_s<1/2\), \(D=1/2-h\), \(1+h\le1+3h/2\), and \(1/2\le s_x\). Also \(s_x<s_t\). Physical rescaling has the four factors

\[
Q^{-A},\qquad Q^{-2A},\qquad Q^{1/2-A},\qquad Q^{1/2-A},
\]

for velocity, pressure, wave potential and mean potential. Each is bounded by \(Q^{-2A}\) for \(0<Q\le1\) and the original \(A=1/2+h\). The source keeps the factor \((km_{\rm harm})^{-1}\) inside the wave potential before this rescaling. One additional spatial derivative recovers velocity from a potential.

Hence the stated common choice

\[
\ell_m=2A+(m+1)\left(1+\frac{3h}{2}\right)
\tag{A36}
\]

covers all derivatives of order at most \(m\), including that extra curl derivative. Every retained stage-\(j\) increment has the source's class exponent at least \(j/10\), so its gain is \(g_j=hj/10\). This verifies the loss/gain arithmetic in (9.17). Passing from \(Q\) to comparable \(q\) may change constants and logarithmic factors, including constants depending on \(j\), but does not change the fixed exponent loss.

The residual conversion uses the fixed factor \(Q^{-2A-1/2}\). Together with (9.8), \(\sigma_j=1/5+j/10\), it gives gain \(\rho_j=h\sigma_j\) with a derivative loss independent of \(j\), as claimed in (9.18). The pointwise validity and all-label uniformity of the finite-stage remainder in (9.18) are supplied by the earlier residual decomposition, not by this exponent calculation. No such uniformity is inferred merely from a coefficient bound at one label.

\subsection{8. Exact Navier--Stokes specialization and truncation choice}\label{exact-navierstokes-specialization-and-truncation-choice}

For (9.20), the tuple is \(U=(u,p)\), without the optional stress variable. The viscosity in this local construction is exactly one. In Cartesian components,

\[
\mathcal R_i(u,p)=\partial_tu_i+\sum_{k=1}^3u_k\partial_ku_i
-\sum_{k=1}^3\partial_k^2u_i+\partial_ip.
\tag{A37}
\]

This is a fixed differential polynomial of order \(s=2\) and degree \(d=2\), with constant coefficients. No singular cylindrical coefficients occur in (A37); thus the comparison can take \(H_m=0\).

Put \(V=(v,p)=U^{[J]}\), \(e=(w,\pi)=U-U^{[J]}\). For every \(I\), the complete mixed-derivative difference is

\[
\begin{aligned}
D^I[\mathcal R_i(v+w,p+\pi)-\mathcal R_i(v,p)]
={}&\partial_tD^Iw_i-\sum_{k=1}^3\partial_k^2D^Iw_i+\partial_iD^I\pi\\
&+\sum_{k=1}^3\sum_{J'\le I}\binom I{J'}
\big[
(D^{J'}v_k)(D^{I-J'+e_k}w_i)
+(D^{J'}w_k)(D^{I-J'+e_k}v_i)\\
&\hspace{46mm}+(D^{J'}w_k)(D^{I-J'+e_k}w_i)
\big].
\end{aligned}
\tag{A38}
\]

Here \(J'\) is a derivative multi-index, not the truncation stage. Every linear, cross and quadratic term is present with its sign. For Euclidean tuple norms, each scalar component is bounded by the tuple norm and the three-component output is at most \(\sqrt3\) times its maximum component. Since \(\sum_{J'\le I}\binom I{J'}=2^{|I|}\), (A38) gives the explicit sufficient bound

\[
|\mathcal R(V+e)-\mathcal R(V)|_m
\le\sqrt3\,[5+6\,2^m|V|_{m+2}+3\,2^m|e|_{m+2}]\,|e|_{m+2}
\] \[
\le C_m(1+|V|_{m+2}+|e|_{m+2})|e|_{m+2},\qquad
C_m=\sqrt3(5+9\,2^m).
\tag{A39}
\]

If the source's absolute value convention is taken componentwise instead, the same upper constant remains valid. No norm identification is needed for this inequality.

For prescribed nonnegative integer \(m,N\), a concrete sufficient stage is any integer \(J\) satisfying

\[
J\ge\max\left\{1,m+2,
\left\lceil\frac{20}{h}(N+K_{m+2}+1+\ell'_{m+2})-1\right\rceil,
\left\lceil\frac{10}{h}(N+1+K_m^{\mathcal R})-2\right\rceil
\right\}.
\tag{A40}
\]

Indeed, the two gain inequalities become exactly

\[
\frac{h(J+1)}{20}-\ell'_{m+2}\ge N+K_{m+2}+1,
\qquad
\frac h5+\frac{hJ}{10}-K_m^{\mathcal R}\ge N+1.
\tag{A41}
\]

Keep this \(J\) fixed. On a sufficiently small common neighborhood, (A22), (A28), (A39) and the fixed-stage remainder estimate give

\[
|\mathcal R(U)|_m
\le\big(1+C_{J,m,N}+3C_m2^{-J}\big)q^N.
\tag{A42}
\]

This proves the exact quantified flatness conclusion (9.20) from the source's stated finite-state estimates. It does not assert that the finite-stage flat pieces were summed, nor that their constants are uniform in the stage.

\subsection{9. The actual representation (9.21), initialization and remaining scope}\label{the-actual-representation-9.21-initialization-and-remaining-scope}

Using the representatives in (9.15), define precisely

\[
\mathcal A=A_0+\sum_{j\ge1}\chi(a_jq)A_j,\qquad
\mathcal Be_\theta=B_0e_\theta+\sum_{j\ge1}\chi(a_jq)B_j,
\] \[
p_{\rm loc}=p_0+\sum_{j\ge1}\chi(a_jq)p_j,\qquad
u_{\rm loc}=\operatorname{curl}\mathcal A+\mathcal Be_\theta.
\tag{A43}
\]

Every sum is locally finite on \(q>0\). Differentiation therefore commutes with these sums locally, and (A43) is exactly (A4) summed over the original representative tuples. It provides the full map to the physical fields, not an identification of potentials with velocities.

The near-axis base potential in (5.27) has Cartesian form \(A_n=(S_n/r^2)(-x_2,x_1,0)\), where \(S_n/r^2\) is smooth in \((r^2,z,t)\) for \(q>0\). The base angular field similarly has the form \(b(r^2,z,t)(-x_2,x_1,0)\). The annular representatives vanish near the axis at each fixed \(t<1\). Their sum has a common such neighborhood locally because only finitely many summands are active there. Thus the Cartesian regularity asserted for the fields in (9.21) is preserved by this summation.

A finite initialization block with a nonpositive gain can be included in the base after one cutoff on its potentials, direct angular fields and pressures. On a fixed neighborhood of \(q=0\), this cutoff is identically one, so every original finite-state residual is unchanged there. Errors supported where that fixed cutoff varies are bounded away from \(q=0\). Once their fixed-stage uniform boundedness on that support is supplied by the stated field bounds, they satisfy a bound \(C_Nq^N\) there for every fixed \(N\), with \(C_N\) allowed to depend on the support's positive lower bound for \(q\). This operation creates no infinite family of errors that would need to be summed.

The source also states that all finite corrections agree with the exact exterior heat field outside one fixed interval of \(X\), so the finite residual is zero there. That is the input which extends the finite-stage estimate from a bounded profile interval to the whole local domain in Proposition 9.9 Step 1. The summation does not establish that exterior agreement independently. Likewise, the one-sided smoothness of the actual finite representatives at \(\tau=0,q\ge c\), fixed discrete label choices, and inner leading-coefficient growth are the further inputs used in Steps 3--5; local finiteness preserves their already established finite-term bounds, but does not replace their proofs.

The completed bounded result is therefore the analytic realization mechanism and the exact matching of its gains, losses, residual polynomial and representative map to (9.17)--(9.21). There is no extra summation gap caused by nonuniform fixed-stage flat-remainder constants. Verification of the global source construction and of the formal endpoint remains outside this result.

\subsection{10. Bounded formal-source correspondence and the initial cutoff}\label{bounded-formal-source-correspondence-and-the-initial-cutoff}

All formal paths in this section are under

\texttt{formal-source/NavierStokes/}.

The following declarations were located by a bounded read-only search. The summation and cutoff definitions/signatures listed in the first five rows were additionally read directly in this audit. No Lean/Lake/Elan process or proof replay was run.

{\def\LTcaptype{none} % do not increment counter
\begin{longtable}[]{@{}
  >{\raggedright\arraybackslash}p{(\linewidth - 4\tabcolsep) * \real{0.3333}}
  >{\raggedright\arraybackslash}p{(\linewidth - 4\tabcolsep) * \real{0.3333}}
  >{\raggedright\arraybackslash}p{(\linewidth - 4\tabcolsep) * \real{0.3333}}@{}}
\toprule\noalign{}
\begin{minipage}[b]{\linewidth}\raggedright
File and declaration
\end{minipage} & \begin{minipage}[b]{\linewidth}\raggedright
Exact source lines
\end{minipage} & \begin{minipage}[b]{\linewidth}\raggedright
Source scope
\end{minipage} \\
\midrule\noalign{}
\endhead
\bottomrule\noalign{}
\endlastfoot
\nolinkurl{SmoothCutoffs.lean}: \nolinkurl{cutoffBump}, \nolinkurl{cutoff}, \nolinkurl{scaledCutoff} & {}27--44; 134--146; 161--182 & Actual smooth bump with inner radius 1/2, outer radius 1, and \texttt{scaledCutoff\ a\ q\ =\ cutoff\ (a*q)}; finite families equal one on a common neighborhood of zero. \\
\nolinkurl{SolenoidalDiagonal.lean}: \nolinkurl{cutStage}, \nolinkurl{potentialSum}, \nolinkurl{eventually_zero_tail}, \nolinkurl{potentialSum_eventuallyEq_partial} & {}32--38; 46--67 & Literal cut-potential \nolinkurl{tsum}, local vanishing of a common tail and equality to a finite prefix. Its index set includes zero. \\
\nolinkurl{CutStageEstimates.lean}: \nolinkurl{RawStageBounds}, \nolinkurl{cutLoss}, \nolinkurl{exists_diagonal_cut_bounds}, \nolinkurl{exists_finite_diagonal_cut_bounds} & {}199--203; 192; 349--368; 379--417 & Raw stage estimates for positive stages; one doubling integer schedule; output gain is exactly \texttt{g\ j\ /\ 2}. One schedule covers a finite family of component spaces. \\
\nolinkurl{DiagonalJetBounds.lean}: \nolinkurl{CutStageBounds}, \nolinkurl{potential_tail_jet_bound}, \nolinkurl{exists_uncut_tail_order} & {}196--220; 275--303 & Positive stage j controls derivatives through j+2; tail after indices 0,\ldots,J is bounded by \texttt{2\^{}(-J)\ q\^{}(g(J+1)-L\ m)}; a single neighborhood compares all requested jets to the uncut prefix. \\
\nolinkurl{MixedDiagonalResidual.lean}: \nolinkurl{velocity}, \nolinkurl{pressure}, \nolinkurl{residual_eq_originalResidual}, \nolinkurl{exists_physical_schedule_residual_zero} & {}26--50; 200--232 & Curl of the cut potential sum plus direct cut-field sum, with pressure summed separately; literal nonlinear residual; one physical schedule with gain \texttt{g\ j\ /\ 2} and vanishing joint residual jets. \\
\nolinkurl{ActualIterationLedger.lean}: \nolinkurl{gain}, \nolinkurl{sigma_formula}, \nolinkurl{residualLoss}, \nolinkurl{residualRate_eq} & {}29; 36; 286; 292 & Literal gain \texttt{h*j/10}, sigma \texttt{1/5+j/10}, and a stage-independent physical residual loss. \\
\nolinkurl{ActualCycleResidualBounds.lean}: \nolinkurl{fixedLoss}, \nolinkurl{Invariant.residual_jetRate}, \nolinkurl{finite_residual_rates}; \nolinkurl{ActualStageEstimates.lean}: \nolinkurl{stageEstimates_of_representations} & {}1145, 1158, 1190; 347--403 & Binding of actual finite-cycle invariants and physical representations to the stage-estimate interface; upstream hypotheses were not traversed. \\
\nolinkurl{ActualCandidateAssembly.lean}: \nolinkurl{estimates}, \nolinkurl{Witness}, \nolinkurl{witness}, \nolinkurl{selected_witness} & {}1090; 1121; 1153; 1177 & Actual mixed candidate assembly together with later localization/activation. This locator does not identify the final globally modified object with the manuscript's local field. \\
\end{longtable}
}

The numerical \nolinkurl{DiagonalScale.lean} comments refer to ``Lemma 11.3,'' rather than the present ``Lemma 5.4.'' Numbering is not used as correspondence evidence.

\subsubsection{10.1 Exact half-gain accounting}\label{exact-half-gain-accounting}

The formal raw input has gain \(g_j\) in \nolinkurl{RawStageBounds}. The output of \nolinkurl{exists_diagonal_cut_bounds} is explicitly \texttt{CutStageBounds\ ...\ (fun\ j\ =\textgreater{}\ g\ j\ /\ 2)\ ...}. The tail theorem's parameter called \nolinkurl{g} is therefore the already reduced gain

\[
\widetilde g_j:=g_j/2.
\]

Its displayed output exponent is \(\widetilde g_{J+1}-L_m\), not \(\widetilde g_{J+1}/2-L_m\). Substituting the literal actual-stage gain gives

\[
\widetilde g_{J+1}=\frac{g_{J+1}}2
=\frac12\frac{h(J+1)}{10}=\frac{h(J+1)}{20}.
\tag{A44}
\]

Thus the formal sequence spends half the raw gain once, exactly as in (A18)--(A22). There is no further factor of 1/2 to turn the denominator into 40. The common loss in the formal theorem is \texttt{cutLoss\ L\ m\ =\ finiteBound\ L\ m\ +\ m}; it is a specified sufficient loss, and is not asserted to equal the sharper mixed-coordinate maximum (A16).

The formal derivative budget \(m\le j+2\) is at least as large as the paper's budget \(m\le j\). For a tail index \(j>J\), the formal condition \(m\le J+3\) implies \(m\le j+2\), explaining the tail theorem's offset. Recovering velocity still uses one additional spatial derivative of the potential, and the nonlinear residual still uses the needed further physical derivatives. Those operations are not paid for by another halving of the gain; they are handled by the fixed derivative-order losses and the finite derivative ceiling, as in Sections 3, 7 and 8.

\subsubsection{10.2 Exact map between cut stage zero and an uncut base}\label{exact-map-between-cut-stage-zero-and-an-uncut-base}

The manuscript's Lemma 5.4 leaves \(U_0\) uncut; formal \nolinkurl{potentialSum} cuts every natural-number index, including zero. The precise comparison uses the same original representative sequence and the same schedule. Let the formal zero entries be

\[
A(0)=A_0,\quad B(0)=B_0e_\theta,\quad P(0)=p_0,
\]

and let its entries at every \(j\ge1\) be the manuscript's \(A_j,B_j,p_j\), with no index shift. Then \texttt{uncutPrefix\ A\ (J+1)} is literally \(A_0+\sum_{j=1}^JA_j\), because the range \(0,\ldots,J\) is the disjoint union of \(\{0\}\) and \(\{1,\ldots,J\}\). The same identity holds for the direct field and pressure. Thus their uncut finite-state objects agree under this representative dictionary.

Use the formal cutoff \(\chi_f\), which has the same required plateaux on positive arguments. It need not equal the explicit example (A10); the proof of Lemma 5.4 uses only the stated smoothness/support/plateau properties, so applies to \(\chi_f\) itself. Let \((\mathcal A_{\rm base},B_{\rm base},p_{\rm base})\) denote the manuscript-style sum with its zero entries uncut, but these same positive-stage cutoffs. Let \((\mathcal A_f,B_f,p_f)\) be the all-indices formal-style sum. Put \(c_0=\chi_f(a_0q)\). Their exact differences on \(q>0\) are

\[
\mathcal A_f-\mathcal A_{\rm base}=(c_0-1)A_0,\quad
B_f-B_{\rm base}=(c_0-1)B_0e_\theta,\quad
p_f-p_{\rm base}=(c_0-1)p_0.
\tag{A45}
\]

For the physical velocities, with \(u_0=\operatorname{curl}A_0+B_0e_\theta\), this gives

\[
u_f-u_{\rm base}
=\operatorname{curl}((c_0-1)A_0)+(c_0-1)B_0e_\theta
=(c_0-1)u_0+\nabla c_0\times A_0.
\tag{A46}
\]

No cutoff commutator is omitted. The full residual difference is (A38) with \(w=(c_0-1)u_0+\nabla c_0\times A_0\) and \(\pi=(c_0-1)p_0\). On the open region

\[
0<q<\frac1{2a_0},
\tag{A47}
\]

the positive formal schedule has \(a_0>0\), so \(c_0\) is identically one on a neighborhood of every point. Equations (A45)--(A46) and every one of their mixed derivatives vanish there. Therefore the two representative sums, velocities, pressures and full nonlinear residuals agree exactly there, at every derivative order. Away from this region, (A45)--(A46), rather than equality, are the precise relation.

This proves the local initial-term-cutoff correspondence on the same representatives and schedule. It does not identify two independently selected schedules, different representative sequences, or the later localized/activated candidate. The formal \nolinkurl{exists_physical_schedule_residual_zero} conclusion is phrased in terms of joint residual jets at the endpoint; no equivalence of all its filter/domain quantifiers with every clause of the manuscript is claimed by this bounded locator.

\subsection{Transcription repair record}\label{transcription-repair-record}

The derivative-partition display in Section 7 contained the literal text \nolinkurl{quad} between \texttt{k\textquotesingle{}\textbackslash{}le\ a+b,} and \texttt{\textbackslash{}sum\ a\_\textbackslash{}nu=a}. It has been corrected to the TeX spacing command \texttt{\textbackslash{}quad}. This repairs transcription only; every mathematical symbol, coefficient, exponent, inequality and claim is retained. The accompanying \nolinkurl{summation_realization_audit_receipt.json} records the final note hash and the completed audit scope.


% --- included proof body: terminal_trace_body.tex ---
\section{The exact quotient of smooth terminal force extensions}

This note extends the explicit terminal-force construction in the accompanying
Navier--Stokes reader. It proves the topology of the actual trace quotient and
the impossibility of a continuous linear choice of all extensions. It makes no
additional claim that the preceding singular fluid construction has completed
independent verification. Spatial coordinates, vector components and the
right-hand time coordinate $\sigma=t-1$ are retained.

Fix the compact spatial set $K\subset\mathbb R^3$ used for localization,
with nonempty interior. For a smooth vector field $g:\mathbb R^3\to\mathbb R^3$
set $\|g\|_{C^m}=\max_{|\alpha|\le m}\sup_x|\partial_x^\alpha g(x)|_2$.
Let $V_K$ consist of such fields with support in $K$, and put
\[
 \J=\prod_{j=0}^\infty V_K,\qquad
 Q_m(F)=\max_{0\le j\le m}\|F_j\|_{C^m}.
\]
The topology defined by the increasing seminorms $Q_m$ is exactly the product
of the usual smooth topologies: every $Q_m$ uses finitely many coordinate
seminorms, while a finite list of coordinate derivative seminorms is bounded
by $Q_m$ for $m$ at least all of their indices and derivative orders.

Let $\E$ be the vector space of smooth fields
$H:\mathbb R^3\times[0,\infty)\to\mathbb R^3$, with all mixed derivatives
continuous from the right at zero, and with support contained in
$K\times[0,1]$. Use its increasing seminorms
\[
 P_m(H)=\max_{|\alpha|+r\le m}
       \sup_{x\in\mathbb R^3,\,\sigma\ge0}
                  |\partial_x^\alpha\partial_\sigma^rH(x,\sigma)|_2.
\]
Define the actual linear trace map and its kernel by
\[
 \mathcal T(H)=(\partial_\sigma^jH(\cdot,0))_{j\ge0},\qquad
 \N=\{H\in\E:\mathcal T(H)=0\}.
\]
The derivatives are evaluated in the original three vector components.
Their support is contained in $K$. Evaluation gives
\[
 Q_m(\mathcal T H)\le P_{2m}(H).
 \tag{T.1}
\]
Thus $\mathcal T$ is a continuous linear map and $\N$ is closed.

Define $\rho(s)=e^{-1/s}$ for $s>0$ and $\rho(s)=0$ for $s\le0$, and fix
\[
 \chi(s)=\frac{\rho(1-s)}{\rho(1-s)+\rho(s-3/4)}.
\]
The denominator is positive at every real $s$: its first term is positive
for $s<1$, and its second for $s>3/4$. Every right derivative of $\rho$
is a finite polynomial in $1/s$ times $e^{-1/s}$ and tends to zero as
$s\downarrow0$; induction and pasting with zero therefore prove its smoothness.
Thus $\chi$ is smooth, equals one for $s\le3/4$, is zero for $s\ge1$,
and satisfies $0\le\chi\le1$. Every positive-order derivative is supported
in $[3/4,1]$, so all the derivative supremums used below are finite.
For $b\ge1$ write
\[
 L_{j,b}(F_j)(x,\sigma)
       =\chi(b\sigma)\frac{\sigma^j}{j!}F_j(x).
\]
For each $r$ the full temporal derivative is
\[
 \partial_\sigma^r\left(\chi(b\sigma)\frac{\sigma^j}{j!}\right)
 =\sum_{\substack{0\le a\le r\\r-a\le j}}
   \binom{r}{a}b^a\chi^{(a)}(b\sigma)
                    \frac{\sigma^{j-r+a}}{(j-r+a)!}.
 \tag{T.2}
\]
Every nonzero term has $0\le\sigma\le1/b$. Consequently, with the finite
constant
\[
 C_{j,r}=\sum_{\substack{0\le a\le r\\r-a\le j}}
      \binom{r}{a}\frac{\|\chi^{(a)}\|_\infty}{(j-r+a)!},
\]
one has, for all $j,r,b$ and spatial multi-indices $\alpha$,
\[
 \|\partial_x^\alpha\partial_\sigma^rL_{j,b}(F_j)\|_\infty
 \le C_{j,r}b^{r-j}\|F_j\|_{C^{|\alpha|}}.
 \tag{T.3}
\]
All constants and factorials in this estimate are retained. Its exponent is
negative whenever $r<j$; no bound uniform in the index $j$ is asserted.

\begin{theorem}
The trace $\mathcal T$ is surjective. The induced map
\[
 \overline{\mathcal T}:\E/\N\longrightarrow\J,\qquad
 [H]\longmapsto\mathcal T H
\]
is a linear homeomorphism for the quotient topology and the specified product
topology. Nevertheless there is no continuous linear map
$B:\J\to\E$ satisfying $\mathcal T B=\operatorname{id}_{\J}$.
\end{theorem}

\begin{proof}
We first construct representatives with an explicit estimate on each quotient
seminorm. Fix $F\in\J$, an integer $m\ge0$, and a real $\epsilon>0$.
Keep the first $m+1$ entries in the finite field
\[
 H_{\le m}=\sum_{j=0}^m L_{j,1}(F_j).
\]
Equation (T.3) gives the finite constant and estimate
\[
 A_m=\sum_{j=0}^m\max_{0\le r\le m}C_{j,r},\qquad
 P_m(H_{\le m})\le A_mQ_m(F).
 \tag{T.4}
\]
Set $b_j=1$ for $0\le j\le m$. For every $j>m$, choose an integer $b_j\ge1$, increasing at least by a factor
two from its predecessor, such that both of the following finite lists hold:
\[
 C_{j,r}b_j^{r-j}\|F_j\|_{C^{|\alpha|}}
       \le\epsilon 2^{-j}\quad (|\alpha|+r\le m),
 \tag{T.5}
\]
\[
 C_{j,r}b_j^{r-j}\|F_j\|_{C^{|\alpha|}}
       \le2^{-j}\quad (|\alpha|+r\le\lfloor j/2\rfloor).
 \tag{T.6}
\]
There are finitely many $(\alpha,r)$ in either list, all with $r<j$.
Thus every nonzero left side tends to zero as $b_j$ increases, and every
zero left side already satisfies its inequality. The set of admissible
integers is nonempty; choosing its least member above the previous lower
bound specifies the sequence. This construction uses the actual input
family $F$, $m$, and $\epsilon$.

For each fixed derivative order $k$, all sufficiently large $j$ have
$k\le\lfloor j/2\rfloor$. Equations (T.3) and (T.6) show uniform
convergence of every differentiated tail in
\[
 H=H_{\le m}+\sum_{j>m}L_{j,b_j}(F_j).
 \tag{T.7}
\]
The finitely many remaining terms have finite smooth norms. To identify the
sum of first derivatives with the derivative of the sum, integrate each
partial sum along a coordinate segment and pass to the limit using uniform
convergence of the function and of that derivative. Repeat at every order,
using one-sided time segments at $\sigma=0$. This proves smoothness with all
the prescribed mixed right derivatives. All terms vanish off $K$ and for
$\sigma\ge1$; uniform convergence of their derivatives preserves these facts
and smoothness on both support boundaries. Hence $H\in\E$.

At $\sigma=0$, each fixed cutoff is constant on a neighborhood. The $r$th
time derivative of its monomial has value zero unless $j=r$, and value one
if $j=r$. Uniform convergence of the differentiated series therefore gives
$\partial_\sigma^rH(\cdot,0)=F_r$ for every $r$. In particular $\mathcal T$
is surjective. Summing (T.5) additionally gives
\[
 P_m(H)\le A_mQ_m(F)+\epsilon\sum_{j>m}2^{-j}
                    \le A_mQ_m(F)+\epsilon.
 \tag{T.8}
\]

For a class $[H]\in\E/\N$, write
\[
 \overline P_m([H])=\inf_{N\in\N}P_m(H+N).
\]
These seminorms generate the quotient topology. Indeed the $P_m$ are
increasing, so their single-seminorm balls form a neighborhood basis in
$\E$. The quotient image of $\{P_m<\delta\}$ is exactly
$\{\overline P_m<\delta\}$: either membership is witnessed by a
representative with $P_m<\delta$. The quotient projection is open, since
the full inverse image of the image of an open set $O$ is
$O+\N=\bigcup_{N\in\N}(O+N)$, which is open. This proves the claimed
quotient description directly.

The kernel definition makes $\overline{\mathcal T}$ well-defined and
injective; surjectivity was proved by (T.7). Its linear inverse is
well-defined by $F\mapsto[H]$, independently of the chosen realization:
two such realizations differ by an element of $\N$. Taking the infimum
of (T.1) over a class and then using (T.8), with arbitrary $\epsilon>0$,
proves both estimates
\[
 Q_m(\overline{\mathcal T}[H])\le\overline P_{2m}([H]),\qquad
 \overline P_m(\overline{\mathcal T}^{-1}F)\le A_mQ_m(F).
 \tag{T.9}
\]
They prove continuity of the map and its inverse with the stated topologies.
In particular the quotient is Hausdorff: if all its seminorms vanish,
(T.9) gives all $Q_m(F)=0$, hence $F=0$ and the class is zero.

Suppose finally that a continuous linear right inverse $B$ existed.
Continuity at zero, the increasing family $Q_m$, and the target norm $P_0$
give an integer $M$ and $\delta>0$ such that
\[
 Q_M(F)<\delta\quad\Longrightarrow\quad P_0(BF)<1.
\]
If $Q_M(F)=0$, this applies to $cF$ for every real $c$. Linearity and
homogeneity give $|c|P_0(BF)<1$ for all $c$, whence $P_0(BF)=0$.
Since $P_0$ is the supremum of the vector length, $BF$ is identically zero.
Choose a nonzero $g\in V_K$: a nonzero smooth scalar bump inside a ball
whose closure is contained in the interior of $K$, multiplied by a fixed
nonzero coordinate vector, supplies one. Take $F_{M+1}=g$ and all other
entries zero. Then $Q_M(F)=0$, whereas the right-inverse identity forces
$\partial_\sigma^{M+1}BF(\cdot,0)=g\ne0$, a contradiction.
\end{proof}

For the actual analytical and formal force constructions this identifies the
precise relation of two extensions after their full coordinate and vector
component maps have been applied. Equal terminal jets imply equality in
$\E/\N$; their pointwise difference is a smooth field flat at the terminal
slice. This relation does not identify their positive-time values. If both
are pasted to the same already specified past force at $\sigma\le0$, the
difference extends by zero there and is jointly smooth, since every mixed
trace vanishes. Its closed support may meet $\sigma=0$: flatness does not
imply a positive support gap. The fluid and all forcing before the terminal
time remain identical in this map of extensions.

Fixing $m=0$, $\epsilon=1$, $b_0=1$ in the least-integer construction
defines one specific right inverse $\mathcal B:F\mapsto H$, with the
remaining cutoff scales depending on the jets. The theorem explains why the
homeomorphism of the quotient is not a continuous linear choice of
representatives. No missing fluid estimate or singularity claim is inserted
as an assumption in this statement about the explicitly defined trace spaces.

\end{document}
